\documentclass[11pt]{article}
\usepackage[margin=1.02in]{geometry}
\usepackage{amsmath,amssymb,amsthm,mathtools}
\usepackage{microtype}
\usepackage{enumitem}
\usepackage{booktabs,longtable,array}
\usepackage{xcolor}
\usepackage[hidelinks]{hyperref}
\usepackage{lmodern}

\definecolor{deepblue}{RGB}{31,63,104}
\definecolor{softgray}{RGB}{245,247,250}
\newtheorem{theorem}{Theorem}[section]
\newtheorem{lemma}[theorem]{Lemma}
\newtheorem{proposition}[theorem]{Proposition}
\newtheorem{corollary}[theorem]{Corollary}
\newtheorem{conjecture}[theorem]{Conjecture}
\theoremstyle{remark}
\newtheorem{remark}[theorem]{Remark}
\newcommand{\norm}[1]{\left\lVert #1\right\rVert}
\newcommand{\ip}[2]{\left\langle #1,#2\right\rangle}


\usepackage{mathrsfs}
\hypersetup{pdftitle={Fixed-Space Prime Compatibility in Burnol's Sonine Quotients -- v1.25},pdfauthor={Alexander Eastwood}}
\newtheorem{assumption}[theorem]{Assumption}
\newtheorem{conditionalresult}[theorem]{Conditional result}
\newcommand{\Q}{\mathcal Q}
\newcommand{\Jc}{J^{\mathrm c}}
\newcommand{\Ju}{J^{\mathrm u}}
\newcommand{\Jr}{\widetilde J}
\newcommand{\E}{\mathscr E}
\newcommand{\Z}{\mathcal Z}
\newcommand{\dd}{\mathrm d}
\setlength{\emergencystretch}{2em}
\setcounter{tocdepth}{1}
\title{\textbf{Fixed-Space Prime Compatibility in Burnol's Sonine Quotients}\\[5pt]
\large Complete working manuscript, version 1.25\\
\normalsize Finite sampling, endpoint recovery, and the remaining graph defect}
\author{Alexander Eastwood}
\date{September 21, 2026}
\begin{document}
\maketitle
\begin{abstract}
We study fixed-space prime actions on Burnol's Sonine quotients and
finite Weil forms. Source covariance, compact Mellin division and
Fredholm reduction identify arithmetic boundary and range defects.
For the repaired fixed-order prolate source, tail, endpoint, correlation
and domain estimates establish weak G1. Its polynomial Fourier
projection recovers the physical endpoint and has an exponentially
small ordinary operator residual. These estimates do not identify
the lowest eigenspace or control a growing source block.

The centered Hardy sampler with an endpoint shell preserves ordinary
Gram and Weil-form limits while retaining an explicit graph defect.
Reflected tests and sharp cuts expose endpoint and complex-strip
obstructions. A finite-core metric criterion bypasses the global
arithmetic operator realization. Full-space projection preserves
strip-product bounds and Weil transfer, but its signed commutator
remains uncontrolled.

For the canonical semilocal Weil operator, signed Schur certificates
prove complete positivity and ground ordering at $\lambda=3$.
At $\lambda=4$ the entire tail $|n|>16$ is bounded below by its
positive archimedean energy. A structured inverse estimate retains
this certified tail metric in the low-mode correction, with complete
remote-row and dimension-aware error bounds. Finite-support trials
certify a zero head-error term at this window. Parity signs sharpen the complete inverse
iteration, with finite-prefix bounds retaining all remote rows.
Complete mixed-term estimates rule out the first inverse degree
with the saved head certificate, including its positive update.
Retaining that update gives a monotone family of complete inverse
upper bounds for subsequent degrees. Energy-scaled frozen trials now certify positivity of the entire even
form at this window, with all residual cross-Grams and every omitted
row retained. The odd sector remains open. The present scalar far
majorant requires exponentially growing cutoffs as the window grows.
Strong cancellation persists in the exact
source complement. Every complete eigenfunction has zero physical
endpoint. A cofinal complete-Weil lower error tending to zero would
imply RH directly, but remains unproved.
\end{abstract}

\paragraph{Status of this version.}
Version 1.25 proves positivity of the complete even Weil form at
$\lambda=4$, using a simultaneous $17$-column certificate and all
infinite-tail corrections. Its certified generalized relative margin
exceeds $0.62629$. This is a fixed-window sign, not an ordinary
spectral-gap number or a cofinal estimate. The odd-sector test remains
inconclusive. The growing-window analysis proves an exponential
cutoff cost for the current scalar diagonal far bound; it is not
an impossibility theorem for other methods. Weak G1, all 72 historical
dispositions, the physical endpoint and explicit sampler graph defect
are retained. G2 and RH are not proved.

\paragraph{Reading conventions.}
The Hermitian product is linear in the first variable:
$\ip{x}{cy}=\overline c\,\ip{x}{y}$.
In the Sonine development the multiplicative cutoff is $R>1$.
In the finite Weil development the multiplicative cutoff is $\lambda>1$
and its logarithmic length is $L=2\log\lambda$.
The source graph functional $\eta_{\mathrm B}$, the equal-coefficient
boundary functional $\eta_\partial$, and physical evaluation
$\delta=L^{-1/2}\eta_\partial$ are different objects.
The smoothed Cauchy value is $\mathscr E(F)=\sigma F(\sigma)$.
An off-line root space is used only under a contradiction hypothesis;
no off-line zero is asserted to exist. Independence from a zero list is
a property of a construction, not a substitute for proving its estimates.

\clearpage
\tableofcontents
\clearpage
\section{Objective and correction}
Fix a prime number $p$ and write
\[
  \lambda_\rho=p^{1/2-\rho}.
\]
The target is a fixed positive Hilbert-space action that retains this multiplier on the arithmetic modes attached to the nontrivial zeros, with positivity or contractivity proved independently of their locations.

A convention correction is essential. Burnol's $Z_w$ is defined by a complex bilinear pairing, whereas the project uses a positive Hermitian product linear in the first variable. If $k_w$ denotes the Hermitian Riesz vector, then
\[
  k_w=Z_{\bar w}.
\]
The project mode is therefore
\[
  v_\rho^R=k_{\bar\rho}^R=Z_\rho^R,
\]
not $Z_{\bar\rho}^R$. With this convention the original compression has the correct critical-line phase. The historical reciprocal-phase rejection resulted from a second, erroneous conjugation.

\section{Burnol framework}
Work in
\[
 K=L^2((0,\infty),dx),
\]
with Hermitian inner product linear in the first slot. Let
\[
 If(x)=x^{-1}f(1/x),\qquad G=I\mathcal F_+I,
\]
where
\[
 (\mathcal F_+f)(x)=2\int_0^\infty \cos(2\pi xt)f(t)\,dt
\]
is the unitary cosine transform, initially on integrable square-integrable
functions and then by its $L^2$ extension. The Mellin conventions are
\[
 M_0f(s)=\int_0^\infty f(x)x^{s-1}\,dx,\qquad
 c_0(s)=\pi^{-s/2}\Gamma(s/2),\qquad M f(s)=c_0(s)M_0f(s).
\]
These formulas are first used where the integrals converge and elsewhere
by the Sonine continuation. In particular $M(E\psi)=c_0\zeta M_0\psi$.
The evaluation and division properties of the completed Sonine space are
the framework inputs from \cite{Burnol2001,Burnol2002Sonine,BurnolCoPoisson2004}.
They do not include the additional quantitative assumptions below. For $R>1$ define
\[
 H_R=\{f:\operatorname{supp}f,\operatorname{supp}Gf\subset(0,R]\},
\]
\[
 V_R=C_c^\infty((1/R,R)),\qquad
 \mathcal D=x\frac{d^2}{dx^2}x,
\]
and
\[
 E(\psi)(x)=\sum_{n\ge1}\psi(nx)-\frac{1}{x}\int_0^\infty\psi(t)\,dt.
\]
The arithmetic co-Poisson space and quotient are
\[
 W_R=\overline{\{E(\mathcal D\phi):\phi\in V_R\}},\qquad
 Q_R=H_R\ominus W_R.
\]
Let $S_R$ and $P_R$ be the orthogonal projections onto $H_R$ and $Q_R$ respectively.

For a nontrivial zero $\rho$, Burnol's construction gives a nonzero vector
\[
 v_\rho^R=Z_\rho^R=k_{\bar\rho}^R\in Q_R.
\]
Put
\[
 v=v_\rho^R,\qquad v_+=v_\rho^{pR},\qquad b=v_+-v.
\]
Because the ambient Sonine spaces are nested and the two vectors represent the same evaluation on $H_R$,
\begin{equation}\label{eq:nested}
 S_Rv_+=v,\qquad b\perp H_R,\qquad
 \norm b^2=\norm{v_+}^2-\norm v^2.
\end{equation}

The unitary dilation is
\[
 (D_pf)(x)=p^{-1/2}f(x/p),
\]
and the completed Mellin transform obeys
\[
 M(D_pf)(s)=p^{s-1/2}M(f)(s).
\]

\subsection*{Quantitative evaluator inputs and compact division}
\begin{assumption}[Evaluator estimates used in the asymptotic sections]\label{ass:evaluators}
For each fixed critical-line parameter $w$, the chosen completed-Mellin
Riesz vectors obey $\|v_w^R\|^2=C_w\log R+O_w(1)$, $C_w>0$.
For fixed distinct critical parameters their cross inner product is
$O_{w,z}(1)$. For a fixed off-line parameter with
$d=|\Re w-1/2|$, $\|v_w^R\|\asymp_w R^d$.
For each fixed jet order $m$, the corresponding ambient jet Gram matrix
$V_{w,R}$ satisfies $\det V_{w,R}\ll_{w,m}R^{2md}$.
Where raw representatives are used below, their local expansion and
Sonine projection error are assumed to have the stated uniform bounds,
including the $O_w(R^{-1})$ projection correction.
\end{assumption}
These quantitative statements were attributed to earlier calculations in
the previous draft without a complete derivation in the supplied files.
They are not consequences of completeness or minimality alone. Sections
using their rates are conditional on this assumption. Exact compression,
covariance, and projection identities do not require it.
Proposition~\ref{prop:v110-normalized-shell} below proves the off-line
\emph{upper} norm bound directly and uses only nested nonzero evaluation
for a lower bound. It does not prove the matching asymptotic, the
critical-line estimates, or the projection-error assertions above.

\begin{lemma}[Compact Mellin division]\label{lem:compact-division}
Let $\theta\in C_c^\infty((1/R,R))$, let $\Theta=M_0\theta$, and let
$w\notin\{0,1\}$. There exists a source $\eta$ in the same compact
support interval with $M_0\eta(s)=\Theta(s)/(s-w)^r$ if and only if
$\Theta$ vanishes to order at least $r$ at $w$. If
$\Theta(0)=\Theta(1)=0$, the divided source has the same two moments.
Moreover $\mathcal D V_R$ is exactly the compact smooth source class with
these two moments zero.
\end{lemma}
\begin{proof}
For one division choose $b$ below the support and set
\[
 \eta(x)=-x^{-w}\int_b^x t^{w-1}\theta(t)\,dt.
\]
The function vanishes below the support. It vanishes above it precisely
when $\Theta(w)=0$, and then is smooth and compactly supported in the
same interval. It solves $-x\eta'-w\eta=\theta$. Integration by parts
gives $(s-w)M_0\eta=\Theta$. Repeat the argument for higher order; the
successive removability conditions are exactly the $r$ vanishing jets.
Necessity follows because a compact source has an entire Mellin transform.
Evaluation at 0 and 1 proves moment preservation when $w\notin\{0,1\}$.
Finally $M_0(\mathcal D\phi)=s(s-1)M_0\phi$. Conversely divide first by
$s$ and then by $s-1$ using the two zero moments and the same compact
construction. The resulting $\phi$ satisfies $\mathcal D\phi=\theta$.
\end{proof}

\begin{lemma}[Finite-jet realization and passage to the closed kernel]\label{lem:jet-density}
Fix $w\notin\{0,1\}$ and a finite order $m$. Admissible compact sources
realize arbitrary jets $(\Psi(w),\ldots,\Psi^{(m-1)}(w))$.
If their associated functionals extend continuously to $W_R$, the
intersection of their kernels in the source core is dense in the
intersection of their kernels in $W_R$.
\end{lemma}
\begin{proof}
For $\psi=\mathcal D\phi$, multiplication by $s(s-1)$ is an invertible
triangular map on the jet at $w$. The distributions
$x^{w-1}(\log x)^j$, $0\le j<m$, are linearly independent on every open
source interval. Hence the jet map from compact smooth $\phi$ is onto:
otherwise a nonzero linear combination of those distributions would
annihilate all test functions. Choose core vectors $u_1,\ldots,u_m$
whose joint jets are the coordinate vectors. If $f_n$ are core
approximants to a vector $f$ with zero jet, replace them by
$f_n-\sum_j\ell_j(f_n)u_j$. Continuity makes the correction tend to zero
and each corrected approximant has zero jet.
\end{proof}


\section{The corrected fixed compression}
Define
\[
 T_{R,p}=P_RD_p\big|_{Q_R}.
\]
It is a contraction independently of RH. The corrected exact boundary identity is
\begin{equation}\label{eq:boundary}
 (T_{R,p}-\lambda_\rho I)v_\rho^R
 =-P_RD_p(v_\rho^{pR}-v_\rho^R).
\end{equation}
Indeed the first-slot Riesz convention gives
\[
 S_RD_pv_\rho^{pR}=\lambda_\rho v_\rho^R,
\]
and \eqref{eq:boundary} follows from \eqref{eq:nested}.

Under Assumption~\ref{ass:evaluators}, the critical-line evaluator norm law is
\[
 \norm{v_\rho^R}^2=C_\rho\log R+O_\rho(1),\qquad C_\rho>0.
\]
Hence \eqref{eq:nested} gives $\norm b=O_{\rho,p}(1)$ and
\begin{equation}\label{eq:criticalrate}
 \frac{\norm{(T_{R,p}-\lambda_\rho I)v_\rho^R}}{\norm{v_\rho^R}}
 =O_{\rho,p}((\log R)^{-1/2}).
\end{equation}
Thus every fixed critical-line mode is asymptotically retained with the correct multiplier. The all-zero version of this estimate remains unproved and is RH-strength.

\section{New source covariance at adjacent scales}
The next lemma is a project deduction from the exact co-Poisson identities. It does not use a zero location.

For a source function define
\[
 (\tau_p\phi)(x)=\phi(x/p).
\]
A direct calculation gives
\begin{equation}\label{eq:covariance}
 \mathcal D\tau_p=\tau_p\mathcal D,
 \qquad
 D_pE(\psi)=p^{-1/2}E(\tau_p\psi).
\end{equation}
Consequently
\begin{equation}\label{eq:DWinclusion}
 D_pW_R\subset W_{pR}.
\end{equation}
Also $V_R\subset V_{pR}$, so $W_R\subset W_{pR}$. Burnol's identity $GE=EI$, together with commutation of $I$ and $\mathcal D$ on the compact source domain, gives
\[
 GW_{pR}=W_{pR}.
\]

\begin{theorem}[Three-piece arithmetic source decomposition]\label{thm:threepiece}
If $R^2>p$, then
\begin{equation}\label{eq:Wdecomp}
 W_{pR}=\overline{W_R+D_pW_R+GD_pW_R}.
\end{equation}
Since $GW_R=W_R$ and $GD_p=D_{1/p}G$, the last summand may equivalently be written $D_{1/p}W_R$.
\end{theorem}

\begin{proof}
The inclusion from right to left follows from \eqref{eq:DWinclusion}, $W_R\subset W_{pR}$, and $G$-invariance of $W_{pR}$.

For the reverse inclusion, the three source intervals
\[
 [1/(pR),R/p],\qquad [1/R,R],\qquad [p/R,pR]
\]
cover $[1/(pR),pR]$ when $R^2>p$. A smooth partition of unity therefore splits every $\phi\in V_{pR}$ as
\[
 \phi=\phi_-+\phi_0+\phi_+,
\]
where $\phi_0\in V_R$, $\phi_+=\tau_p\psi_+$ for some $\psi_+\in V_R$, and $\phi_-=I\tau_p\psi_-$ for some $\psi_-\in V_R$. Applying $E\mathcal D$, using \eqref{eq:covariance} and $GE=EI$, places the three pieces in $W_R$, $D_pW_R$, and $GD_pW_R$ respectively. Taking closures proves \eqref{eq:Wdecomp}.
\end{proof}

\begin{corollary}[Arithmetic projection disappears in two adjoint channels]\label{cor:projectiondrop}
Assume $R^2>p$ and let $y\in Q_{pR}$. Then
\begin{equation}\label{eq:drop}
 P_Ry=S_Ry,\qquad
 P_RD_py=S_RD_py,\qquad
 P_RD_{1/p}y=S_RD_{1/p}y.
\end{equation}
\end{corollary}

\begin{proof}
Since $Q_{pR}\perp W_{pR}$, Theorem~\ref{thm:threepiece} gives
$y\perp W_R$, $y\perp D_pW_R$, and $y\perp D_{1/p}W_R$. These imply respectively that $S_Ry$, $S_RD_{1/p}y$, and $S_RD_py$ are orthogonal to $W_R$. Each therefore lies in $Q_R$, yielding \eqref{eq:drop}.
\end{proof}

\section{A quotient-native prime operator pencil}
Define two contractions from the old quotient into the enlarged quotient,
\[
 J_{R,p}=P_{pR}\big|_{Q_R}:Q_R\to Q_{pR},
\]
\[
 B_{R,p}=P_{pR}D_{1/p}\big|_{Q_R}:Q_R\to Q_{pR}.
\]
Their adjoints are
\[
 J_{R,p}^*=P_R\big|_{Q_{pR}},\qquad
 B_{R,p}^*=P_RD_p\big|_{Q_{pR}}.
\]
For the zero vector $v_+=v_\rho^{pR}$, nested evaluation and the corrected boundary calculation give the exact identities
\begin{equation}\label{eq:pencildual}
 J_{R,p}^*v_+=v,
 \qquad
 B_{R,p}^*v_+=\lambda_\rho v.
\end{equation}
Equivalently,
\begin{equation}\label{eq:pencil}
 (B_{R,p}^*-\lambda_\rho J_{R,p}^*)v_\rho^{pR}=0.
\end{equation}
Thus the arithmetic prime multiplier already appears as an exact generalized eigenvalue of a fixed, zero-independent operator pencil. The unresolved issue is converting this generalized eigenrelation into a single positive fixed-space action without inserting the zero location by hand.

\section{Exact defect kernel and a sufficient fixed-space theorem}
Define on $Q_{pR}$
\begin{equation}\label{eq:Delta}
 \Delta_{R,p}=J_{R,p}J_{R,p}^*-B_{R,p}B_{R,p}^*.
\end{equation}
This is self-adjoint. From \eqref{eq:pencildual}, for every nontrivial zero,
\begin{equation}\label{eq:Deltadiag}
 \ip{\Delta_{R,p}v_\rho^{pR}}{v_\rho^{pR}}
 =(1-|\lambda_\rho|^2)\norm{v_\rho^R}^2.
\end{equation}
More generally, for zeros $\rho,\sigma$,
\begin{equation}\label{eq:Pick}
 \ip{\Delta_{R,p}v_\sigma^{pR}}{v_\rho^{pR}}
 =\bigl(1-\lambda_\sigma\overline{\lambda_\rho}\bigr)
   \ip{v_\sigma^R}{v_\rho^R}.
\end{equation}
This is a Pick-type difference-of-Gram kernel on the zero-evaluation system.

\begin{proposition}[Finite-cutoff span positivity would imply RH]\label{prop:positivity}
If, for one fixed $p>1$ and admissible cutoffs, $\Delta_{R,p}$ were nonnegative on the span of the zero-evaluation vectors $v_\rho^{pR}$, then every nontrivial zero would satisfy $|\lambda_\rho|\le1$. Functional-equation symmetry would then imply RH.
\end{proposition}

\begin{proof}
Apply nonnegativity to a single vector in \eqref{eq:Deltadiag}. Since $v_\rho^R\ne0$,
$1-|\lambda_\rho|^2\ge0$. Hence $p^{1/2-\Re\rho}\le1$, or $\Re\rho\ge1/2$. The functional equation supplies the reflected zero $1-\rho$, forcing equality.
\end{proof}

The span-positivity hypothesis is substantially stronger than RH and should not be treated as the main target. Indeed, assume RH for the moment. Then every diagonal quantity in \eqref{eq:Deltadiag} is zero. If $\Delta_{R,p}$ were positive semidefinite on the whole zero span, positivity would force each $v_\rho^{pR}$ to lie in the nullspace of the restricted form. Hence for all critical zeros $\rho,\sigma$,
\begin{equation}\label{eq:crossvanish}
 \bigl(1-\lambda_\sigma\overline{\lambda_\rho}\bigr)
 \ip{v_\sigma^R}{v_\rho^R}=0.
\end{equation}
Thus distinct prime phases would require the corresponding evaluation vectors to be orthogonal. Burnol's zero systems are studied as complete/minimal systems, not supplied here as an orthogonal basis. Therefore RH by itself does not furnish the extra cross-term vanishing in \eqref{eq:crossvanish}. The safe target is the \emph{diagonal} sign on each zero evaluator, not positivity of the full Pick matrix.

For the same reason, the stronger global operator inequality
\begin{equation}\label{eq:globalpositive}
 B_{R,p}B_{R,p}^*\le J_{R,p}J_{R,p}^*
\end{equation}
should be regarded only as an over-strong sufficient condition. Were it true, the Douglas factorization lemma would produce a contraction $C_{R,p}:Q_R\to Q_R$ with $B_{R,p}=J_{R,p}C_{R,p}$ and hence $C_{R,p}^*v_\rho^R=\lambda_\rho v_\rho^R$. Nothing in the present work proves the cross-term conditions required for such a global factorization, even under RH.

\section{The defect becomes a pure Sonine comparison}
Theorem~\ref{thm:threepiece} removes the arithmetic projection from the defect form. For $y\in Q_{pR}$ and $R^2>p$, Corollary~\ref{cor:projectiondrop} gives
\[
 J_{R,p}^*y=S_Ry,
 \qquad
 B_{R,p}^*y=S_RD_py.
\]
Therefore
\begin{equation}\label{eq:puresonine}
 \ip{\Delta_{R,p}y}{y}
 =\norm{S_Ry}^2-\norm{S_RD_py}^2.
\end{equation}
The diagonal sign problem has thus been reduced to the following concrete question:
\begin{quote}
For each zero evaluator $y=v_\rho^{pR}$, can an independent arithmetic estimate prove $\|S_RD_py\|\le\|S_Ry\|$?
\end{quote}
A proof of
\[
 \norm{S_RD_py}\le\norm{S_Ry}
\]
on each zero evaluator would prove RH by \eqref{eq:Deltadiag}. Positivity on the whole span is a stronger sufficient condition, as explained in Proposition~\ref{prop:positivity}. A proof on all of $Q_{pR}$ would give the stronger Douglas factorization above.

This is the current fixed-space compatibility target. It is more explicit than the earlier boundary-distance statement because the co-Poisson projection has disappeared from the quadratic form exactly. No sign for \eqref{eq:puresonine} has yet been proved.


\section{Exact Fredholm reduction to a boundary-shell imbalance}
We now make \eqref{eq:puresonine} explicit using Burnol's finite-interval Fredholm projection. Put
\[
 a=R^{-1},\qquad q=Iy\in K_{a/p},
\]
where $I f(x)=x^{-1}f(1/x)$ and $K_{a/p}$ is the cosine-Sonine space with gap $a/p$. Let $\mathcal F_+$ be the unitary cosine transform, let $P_a$ denote restriction/projection to $(0,a)$, and set
\[
 F_a=P_a\mathcal F_+P_a,\qquad D_a=F_a^2.
\]
Under inversion, $S_R$ becomes the orthogonal projection $\pi_a$ onto $K_a$, while $ID_p=D_{1/p}I$. Hence
\begin{equation}\label{eq:fredholm-start}
 \norm{S_Ry}^2-\norm{S_RD_py}^2
 =\norm{\pi_a q}^2-\norm{\pi_aD_{1/p}q}^2.
\end{equation}

For a general $f\in L^2(0,\infty)$ define
\[
 u=P_af,\qquad v=P_a\mathcal F_+f.
\]
Burnol's projection theorem (equivalently the $2\times2$ Fredholm inverse in Lemma 2 of \cite{Burnol2002Sonine}) gives the exact lost energy
\begin{align}
 \mathcal E_a(f)&:=\norm{f-\pi_af}^2\\
 &=\frac12\ip{u+v}{(I+F_a)^{-1}(u+v)}
 +\frac12\ip{u-v}{(I-F_a)^{-1}(u-v)}.\label{eq:FredholmEnergy}
\end{align}
In particular $\mathcal E_a(f)\ge0$ and
\[
 \norm{\pi_af}^2=\norm f^2-\mathcal E_a(f).
\]
Since dilation is unitary, \eqref{eq:fredholm-start} becomes
\begin{equation}\label{eq:fredholm-difference}
 \norm{S_Ry}^2-\norm{S_RD_py}^2
 =\mathcal E_a(D_{1/p}q)-\mathcal E_a(q).
\end{equation}

The Sonine condition $q\in K_{a/p}$ now produces a one-channel collapse after dilation. Indeed $\mathcal F_+q$ vanishes on $(0,a/p)$ and
\[
 \mathcal F_+D_{1/p}=D_p\mathcal F_+,
\]
so
\[
 P_a\mathcal F_+D_{1/p}q=0.
\]
Writing
\[
 u=P_aq,\qquad v=P_a\mathcal F_+q,
 \qquad u'=P_aD_{1/p}q,
\]
we obtain exactly
\begin{equation}\label{eq:dilated-energy}
 \mathcal E_a(D_{1/p}q)=\ip{u'}{(I-D_a)^{-1}u'}.
\end{equation}
Equations \eqref{eq:FredholmEnergy}--\eqref{eq:dilated-energy} are the desired finite-$R$ Fredholm formula for the sign problem.

There is also a transparent shell approximation. The Hilbert--Schmidt estimate gives $\norm{F_a}\le2a$. Consequently, for $a<1/2$,
\begin{align}
 \left|\mathcal E_a(q)-(\norm u^2+\norm v^2)\right|
 &\le \frac{2a}{1-2a}(\norm u^2+\norm v^2),\label{eq:Eapprox1}\\
 \left|\mathcal E_a(D_{1/p}q)-\norm{u'}^2\right|
 &\le \frac{4a^2}{1-4a^2}\norm{u'}^2.\label{eq:Eapprox2}
\end{align}
Because $q$ and $\mathcal F_+q$ both vanish below $a/p$,
\[
 \norm{u'}^2-\norm u^2
 =\int_a^{pa}|q(t)|^2\,dt,
 \qquad
 \norm v^2=\int_{a/p}^{a}|\mathcal F_+q(t)|^2\,dt.
\]
Thus
\begin{equation}\label{eq:shell-balance}
 \boxed{
 \norm{S_Ry}^2-\norm{S_RD_py}^2
 =\int_a^{pa}|q(t)|^2dt
 -\int_{a/p}^{a}|\mathcal F_+q(t)|^2dt
 +R_{a,p}(q),}
\end{equation}
where, explicitly,
\begin{align*}
 |R_{a,p}(q)|
 &\le \frac{4a^2}{1-4a^2}\int_{a/p}^{pa}|q(t)|^2dt\\
 &\quad+\frac{2a}{1-2a}
 \left(\int_{a/p}^{a}|q(t)|^2dt+
 \int_{a/p}^{a}|\mathcal F_+q(t)|^2dt\right).
\end{align*}
The sign problem has therefore been reduced, up to a quantitatively small Fredholm correction, to a signed imbalance between two adjacent multiplicative shells: a physical outer shell and a Fourier inner shell.

The arithmetic decomposition from Theorem~\ref{thm:threepiece} implies more than mere Sonine membership. For $y\in Q_{pR}$,
\[
 S_Ry,\quad S_RD_py,\quad S_RD_{1/p}y\in Q_R.
\]
Equivalently, the shell pair in \eqref{eq:shell-balance} comes from a vector orthogonal to the old, prime-expanded, and inverse-prime arithmetic source directions simultaneously. This is the precise additional arithmetic structure available for a future sign theorem.

\paragraph{What this does not prove.}
There is no generic Sonine inequality forcing the first shell integral in \eqref{eq:shell-balance} to dominate the second. On a zero evaluator $y=v_\rho^{pR}$ the exact dual identity already gives
\[
 \norm{S_Ry}^2-\norm{S_RD_py}^2
 =(1-|\lambda_\rho|^2)\norm{v_\rho^R}^2.
\]
Hence positivity on all zero evaluators is itself equivalent to the required one-sided zero exclusion. The Fredholm reduction is useful because it identifies the missing mechanism as a sign-sensitive arithmetic shell theorem; it does not manufacture that theorem from generic uncertainty geometry.


\subsection{Signed moments and an unconditional relative error budget}
The $O(a)$ correction in \eqref{eq:shell-balance} contains a signed
cross term. It can be retained explicitly, and even resummed with all
constant-kernel contributions, before estimating the error.

\begin{lemma}[Signed Fredholm formula and rank-one resummation]\label{lem:v110-signed}
Let $0<a<1/2$ and use $u,v,u'$ as above. Put
$A_a=(I-F_a^2)^{-1}$ and
\[
 H_a=\norm{u'}^2+\norm u^2+\norm v^2,\qquad
 \mathcal B_{a,p}(q)=\int_a^{pa}|q|^2-\int_{a/p}^a|\mathcal F_+q|^2.
\]
For $\mathfrak d_a(q)=\mathcal E_a(D_{1/p}q)-\mathcal E_a(q)$,
\begin{align}
 \mathfrak d_a(q)
 &=\ip{u'}{A_au'}-\ip{u}{A_au}-\ip{v}{A_av}
       +2\Re\ip{u}{F_aA_av}.\label{eq:v110-exact-signed}
\end{align}
Define $m=\int_0^a u$, $n=\int_0^a v$, and $m'=\int_0^a u'$.
Then
\begin{align}
 \mathfrak d_a(q)&=\mathcal B_{a,p}(q)+\mathcal C_a(q)
                      +\mathcal R_a(q),\label{eq:v110-rankone}\\
 \mathcal C_a(q)&=
 \frac{4\Re(m\overline n)+4a(|m'|^2-|m|^2-|n|^2)}{1-4a^2},
 \qquad
 |\mathcal R_a(q)|\le
 \frac{4\pi^2a^5}{5(1-2a)^2}H_a.\label{eq:v110-rankone-error}
\end{align}
In particular, $m'=p^{-1/2}\int_{a/p}^{pa}q(t)\,dt$.
These formulas require Sonine membership, not arithmetic orthogonality.
\end{lemma}
\begin{proof}
The identities $(I\pm F_a)^{-1}=A_a(I\mp F_a)$ in
\eqref{eq:FredholmEnergy} give \eqref{eq:v110-exact-signed}.
Self-adjointness and commutation of $F_a,A_a$ make the two mixed
terms conjugate; they enter the \emph{difference} with a plus sign.

Let $T_af=2\int_0^a f(t)\,dt$, the constant-kernel operator on $(0,a)$.
The actual cosine kernel gives
\[
 |2\cos(2\pi xt)-2|\le4\pi^2x^2t^2,\qquad
 \norm{F_a-T_a}\le\norm{F_a-T_a}_{\mathrm{HS}}
 \le\frac{4\pi^2}{5}a^5.
\]
Both $\norm{F_a}$ and $\norm{T_a}$ are at most $2a$.
On $L^2(0,a)\oplus L^2(0,a)$ set
\[
 K_F=\begin{pmatrix}I&F_a\\F_a&I\end{pmatrix},\qquad
 K_T=\begin{pmatrix}I&T_a\\T_a&I\end{pmatrix}.
\]
The resolvent identity yields
\[
 \norm{K_F^{-1}-K_T^{-1}}
 \le\frac{4\pi^2a^5}{5(1-2a)^2}.
\]
Apply this bound separately to $(u,v)$ and $(u',0)$ in the lost-energy
quadratic forms. With $e=1_{(0,a)}$ and
$e\otimes e:f\mapsto\ip{f}{e}e$, one has
\[
 (I-T_a^2)^{-1}=I+\frac{4a}{1-4a^2}e\otimes e,
 \qquad T_a(I-T_a^2)^{-1}=\frac{2}{1-4a^2}e\otimes e.
\]
Substitution proves the stated moment formula and error. The final
identity for $m'$ follows from $D_{1/p}q(x)=\sqrt p\,q(px)$ and
the gap below $a/p$.
\end{proof}

\begin{proposition}[Normalized shell reduction without evaluator asymptotics]\label{prop:v110-normalized-shell}
Fix a nontrivial off-line zero $\rho$, put
$\beta=\Re\rho$, $d=|\beta-1/2|\in(0,1/2)$, and fix a prime $p$.
For $y_R=v_\rho^{pR}$, $q_R=Iy_R$, $a=R^{-1}$ and
$N_R=\norm{v_\rho^R}^2$, the following estimates do not use
Assumption~\ref{ass:evaluators}:
\begin{align}
 \frac{|\mathcal R_a(q_R)|}{N_R}&=O_{\rho,p}(R^{2d-5}),
 &\frac{|\mathcal C_a(q_R)|}{N_R}&=O_{\rho,p}(R^{2d-1}),
 \label{eq:v110-normalized-errors}\\
 \frac{\mathcal B_{1/R,p}(q_R)}{N_R}
 &=1-p^{1-2\beta}+O_{\rho,p}(R^{2d-1}).
 \label{eq:v110-shell-limit}
\end{align}
Constants may also depend on a fixed initial cutoff with nonzero
evaluation. In particular all displayed error terms tend to zero.
\end{proposition}
\begin{proof}
For $\Re s=b>1/2$ and $f\in H_R$, support in $(0,R]$ and
Cauchy--Schwarz give
\[
 |M f(s)|\le |c_0(s)|\left(\frac{R^{2b-1}}{2b-1}\right)^{1/2}\norm f.
\]
The Mellin cosine multiplier, with the left-Mellin convention of
Section 2, gives
\[
 M_0(Gf)(s)=\frac{c_0(1-s)}{c_0(s)}M_0f(1-s),
 \qquad M(Gf)(s)=M f(1-s).
\]
These identities hold first on the Mellin line and extend by the
Sonine continuation; see \cite[\S1, (7)]{Burnol2004}, converting its
right-Mellin convention to the present one. Since $G$ is a unitary
involution preserving $H_R$, the same bound on the reflected parameter
applies when $b<1/2$. The norm of the evaluator at $\bar\rho$ is thus
\[
 \norm{v_\rho^R}^2\le C_\rho R^{2d},\qquad
 C_\rho=
 \begin{cases}
 |c_0(\rho)|^2/(2\beta-1),&\beta>1/2,\\
 |c_0(1-\rho)|^2/(1-2\beta),&\beta<1/2.
 \end{cases}
\]
Nested evaluation gives $N_R\ge N_{R_0}>0$ for $R\ge R_0$;
this uses only the nonzero evaluator and orthogonal projection, not
a lower growth asymptotic. Also
$H_a\le3\norm{q_R}^2=3\norm{v_\rho^{pR}}^2\ll_{\rho,p}R^{2d}$.
The remainder estimate follows. By Cauchy--Schwarz,
$|m|^2\le a\norm u^2$, and likewise for $n,m'$, whence
\[
 |\mathcal C_a(q_R)|\le
 \frac{2a(\norm u^2+\norm v^2)+4a^2H_a}{1-4a^2}.
\]
This proves the second error estimate. Finally use the exact arithmetic
diagonal identity \eqref{eq:Deltadiag} in \eqref{eq:v110-rankone}.
\end{proof}

\paragraph{The remaining arithmetic obligation, now with no Fredholm ambiguity.}
It suffices to fix $p=2$ and prove independently, for each hypothetical
zero with $\Re\rho<1/2$, that
\begin{equation}\label{eq:v110-sign-target}
 \liminf_{R\to\infty}
 \frac{\displaystyle\int_{1/R}^{2/R}|Iv_\rho^{2R}|^2
       -\displaystyle\int_{1/(2R)}^{1/R}|\mathcal F_+Iv_\rho^{2R}|^2}
      {\norm{v_\rho^R}^2}\ge0.
\end{equation}
Indeed \eqref{eq:v110-shell-limit} would instead give a strictly negative
limit, and reflection then excludes right-hand zeros too. Multiplicity
does not alter this criterion, since the value evaluator already exists
at every zero. No uniformity in $\rho$, no full-span positivity, and no
matching evaluator lower asymptotic is required. The rates are not
uniform as the parameter approaches the critical line or strip edges.
The three arithmetic source orthogonalities remain available, but neither
they nor the signed moment formula prove \eqref{eq:v110-sign-target}.
In particular the explicit cross term is asymptotically negligible on
each fixed off-line evaluator; controlling it alone cannot settle the
leading shell sign. This closes a normalization prerequisite of this
route, not the signed G2 estimate or RH.

\section{Projection-angle form of the same obstruction}
There is a second exact reformulation that removes the Fredholm coordinates altogether. For any closed subspace $M\subset K$ and unitary $U$, the projection onto $U^{-1}M$ is $U^{-1}P_MU$. Hence
\[
 D_{1/p}S_RD_p
\]
is the orthogonal projection onto $D_{1/p}H_R$. If $y\in Q_{pR}$, Theorem~\ref{thm:threepiece} gives $y\perp W_R$ and $y\perp D_{1/p}W_R$. Therefore
\[
 S_Ry\in Q_R,
 \qquad
 D_{1/p}S_RD_py\in D_{1/p}Q_R.
\]
The two vectors are in fact the orthogonal projections of $y$ onto the indicated arithmetic quotient spaces. Thus
\begin{equation}\label{eq:angleform}
 \boxed{
 \norm{S_Ry}^2-\norm{S_RD_py}^2
 =\norm{P_{Q_R}y}^2-\norm{P_{D_{1/p}Q_R}y}^2.}
\end{equation}
Here the second $P$ denotes the ordinary Hilbert projection onto the closed subspace $D_{1/p}Q_R$.

Equation~\eqref{eq:angleform} is useful conceptually. The two comparison spaces $Q_R$ and $D_{1/p}Q_R$ are unitarily congruent; neither is generically contained in the other. The missing sign is therefore not a consequence of dimension, positivity, or ordinary projection monotonicity. It asks for a genuinely prime-oriented arithmetic preference inside $Q_{pR}$.

For a zero evaluator $y=v_\rho^{pR}$ the two projections are explicit:
\[
 P_{Q_R}y=v_\rho^R,
 \qquad
 P_{D_{1/p}Q_R}y
 =\lambda_\rho D_{1/p}v_\rho^R,
\]
so \eqref{eq:angleform} reduces again to
\[
 (1-|\lambda_\rho|^2)\norm{v_\rho^R}^2.
\]
Thus any proof based only on the abstract fact that the two spaces are unitary images of one another cannot decide the sign. A successful argument must exploit how the co-Poisson arithmetic range sits relative to the dilation direction, not merely the ambient Sonine geometry.

\subsection*{The next concrete arithmetic object}
Theorem~\ref{thm:threepiece} identifies the enlarged arithmetic space with the closure of the range of the synthesis map
\[
 \mathcal A_{R,p}:W_R\oplus W_R\oplus W_R\longrightarrow H_{pR},
 \qquad
 \mathcal A_{R,p}(w_0,w_+,w_-)
 =w_0+D_pw_++D_{1/p}w_-.
\]
Thus $W_{pR}=\overline{\operatorname{Ran}\mathcal A_{R,p}}$; closed range is not asserted. Its Gram operator is the $3\times3$ Toeplitz-type block matrix
\[
 \mathcal A_{R,p}^*\mathcal A_{R,p}
 =\begin{pmatrix}
 I & C_p & C_{1/p}\\
 C_p^* & I & C_{1/p^2}\\
 C_{1/p}^* & C_{1/p^2}^* & I
 \end{pmatrix},
\]
where, for example,
\[
 C_p=P_{W_R}D_p\big|_{W_R}.
\]
The exact placement of adjoints depends only on the chosen first-slot convention. Positivity of this Gram matrix is automatic because it is a Gram operator; what is not automatic is the Schur-complement inequality needed to orient \eqref{eq:angleform}. This identifies a new, genuinely arithmetic target: obtain sign-sensitive control of the cross-scale correlation operators $C_p$ and $C_{p^2}$ from the co-Poisson source formula. Generic Hilbert-space positivity alone cannot provide it.


\subsection{A convergent arithmetic adjoint for the shell constraints}
To use the source ranges in a signed estimate, their infinite sums must
first be paired legitimately with arbitrary $L^2$ vectors. The formal
sum $\sum_{n\ge1}y(t/n)/n$ need not converge. The following cutoff
supplies an explicit renormalization in the conventions of Section 2.
The integer $M$ here is a co-Poisson summation cutoff, not a Fourier cut
or the Hardy smoothing order used in later sections.

\begin{lemma}[Integral-corrected source cutoff]\label{lem:v111-cutoff}
Let $h\in C_c^\infty((\alpha,b))$, $0<\alpha<b$, put
$V(h)=\int|h'|$, and define
\[
 E_Mh(x)=\sum_{n=1}^M h(nx)-\frac1x\int_0^{Mx}h(t)\,dt.
\]
Then $E_Mh=Eh$ for $x\ge b/M$, and
\begin{equation}\label{eq:v111-source-error}
 \norm{E_Mh-Eh}_2\le V(h)\sqrt{b/M}.
\end{equation}
If $h\ne0$ and $\int h=0$, put $H(t)=\int_0^t h$.
The uncorrected sum $U_Mh=\sum_{n\le M}h(nx)$ instead satisfies
\begin{equation}\label{eq:v111-raw-divergence}
 \norm{U_Mh-Eh}_2
 =\sqrt M\left(\int_\alpha^b\frac{|H(t)|^2}{t^2}\,dt\right)^{1/2}
       +O_h(M^{-1/2}).
\end{equation}
The leading constant is strictly positive.
\end{lemma}
\begin{proof}
Subtract the defining co-Poisson counterterm to obtain
\[
 E_Mh-Eh=-\sum_{n>M}h(nx)+\frac1x\int_{Mx}^\infty h(t)\,dt.
\]
On each interval $((n-1)x,nx)$, the error between the right endpoint
value and the integral average is at most the variation of $h$ on that
interval. Thus the displayed difference is bounded by
$\int_{Mx}^\infty|h'|\le V(h)$ and vanishes above $b/M$.
This proves \eqref{eq:v111-source-error}. The same quadrature bound
shows $Eh$ bounded near zero; above $b$ it is $-x^{-1}\int h$, so
$Eh\in L^2$. If $\int h=0$, then $H$ is supported in $[\alpha,b]$ and
$U_Mh-Eh=(E_Mh-Eh)+H(Mx)/x$.
Changing variables gives
$\norm{H(Mx)/x}_2=\sqrt M\norm{H(t)/t}_2$.
The triangle inequality proves \eqref{eq:v111-raw-divergence}; its
constant cannot vanish unless $h=H'=0$.
\end{proof}

\begin{proposition}[Regularized adjoint and arithmetic source equation]\label{prop:v111-adjoint}
For $y\in L^2(0,\infty)$ set
\begin{equation}\label{eq:v111-adjoint-finite}
 (\mathscr A_My)(t)=\sum_{n=1}^M\frac{y(t/n)}n
                       -\int_{t/M}^\infty\frac{y(x)}x\,dx.
\end{equation}
There is a distribution $\mathscr Ay\in H^{-1}_{\rm loc}(0,\infty)$
such that $\mathscr A_My\to\mathscr Ay$ locally in $H^{-1}$, and
\begin{align}
 \ip{E_Mh}{y}&=\int h(t)\overline{\mathscr A_My(t)}\,dt,
 &\ip{Eh}{y}&=\ip{h}{\mathscr Ay},\label{eq:v111-duality}\\
 |\ip{h}{(\mathscr A_M-\mathscr A)y}|
 &\le V(h)\sqrt{b/M}\norm y_2.\label{eq:v111-adjoint-error}
\end{align}
The distribution pairing in the second slot extends the first-slot-linear
Hermitian convention. For $y\in H_R$, one has the exact characterization
\begin{equation}\label{eq:v111-arithmetic-equation}
 y\in Q_R\quad\Longleftrightarrow\quad
 \mathscr Ay(t)=A_y+B_y/t\quad\hbox{on }(1/R,R)
\end{equation}
as distributions, for some constants $A_y,B_y$.
For every dilation factor $c>0$, $\mathscr A D_c=D_c\mathscr A$.
\end{proposition}
\begin{proof}
For finite $M$ all sums may be integrated term by term. In the integral
term, Fubini is justified by
\[
 \int_{\alpha/M}^\infty\frac{|y(x)|}{x}\,dx
 \le\sqrt{M/\alpha}\norm y_2.
\]
Changing $t=nx$ in each summand and reversing the counterterm integral
gives \eqref{eq:v111-duality}. Lemma~\ref{lem:v111-cutoff} gives
\eqref{eq:v111-adjoint-error}. Since
$V(h)\le\sqrt{b-\alpha}\norm{h'}_2$, the differences are Cauchy in
the dual of $H_0^1(\alpha,b)$, using its derivative norm. These local
limits agree on overlaps and define $\mathscr Ay$.

By the definition of $W_R$, arithmetic orthogonality is exactly
$\ip{\mathcal D\phi}{\mathscr Ay}=0$ for all
$\phi\in C_c^\infty((1/R,R))$. The operator
$\mathcal D\phi=(t^2\phi')'$ is formally self-adjoint. Hence this is
$(t^2(\mathscr Ay)')'=0$ in distributions. Integrating twice on the
connected interval gives precisely $A_y+B_y/t$. The converse follows
by integration by parts and passage to the closure defining $W_R$.
Finally a change of variables in \eqref{eq:v111-adjoint-finite} gives
$\mathscr A_MD_c=D_c\mathscr A_M$ exactly. Distributional passage
to the limit proves the covariance assertion.
\end{proof}

For the three-scale problem this becomes
$\mathscr Ay=A_y+B_y/t$ on $(1/(pR),pR)$, in addition to
$y\in H_{pR}$. This makes all three source constraints explicit;
it does not add three independent scalar conditions to that one local
equation. In particular the two coefficients need not vanish.
No pointwise or local $L^2$ convergence of $\mathscr A_My$ is claimed
for arbitrary $y$.
There is also an explicit growing-interval budget. With
$I_R=(1/(pR),pR)$ and the derivative norm on $H_0^1(I_R)$,
\eqref{eq:v111-adjoint-error} gives
\[
 \norm{(\mathscr A_M-\mathscr A)y}_{(H_0^1(I_R))'}
 \le\sqrt{pR(pR-1/(pR))/M}\norm y_2.
\]
For the fixed off-line evaluator $y=v_\rho^{pR}$,
Proposition~\ref{prop:v110-normalized-shell} therefore bounds this error
divided by $\norm{v_\rho^R}$ by $O_{\rho,p}(R^{1+d}/\sqrt M)$.
Choosing $M=\lceil R^4\rceil$ makes it tend to zero. This controls
source testing in the stated dual norm, not shell energies or a source
Gram inverse; it does not change any later Fourier cutoff.

\begin{corollary}[Controlled three-channel source Gram limits]\label{cor:v111-gram}
Fix finitely many smooth sources $h_i$ supported in $(\alpha,b)$, and
write $w_i=Eh_i$, $w_i^{(M)}=E_Mh_i$ and
$\epsilon_i=V(h_i)\sqrt{b/M}$. For
$c,d\in\{1,p,1/p\}$,
\begin{align*}
 &\left|\ip{D_cw_i^{(M)}}{D_dw_j^{(M)}}-
                  \ip{D_cw_i}{D_dw_j}\right|\\
 &\hspace{15mm}\le
 \epsilon_i\norm{w_j}+\epsilon_j\norm{w_i}+\epsilon_i\epsilon_j.
\end{align*}
Thus every fixed finite source-Gram matrix has a convergent arithmetic
quadrature representation. No convergence of the associated inverse or
uniform lower singular-value bound is implied.
\end{corollary}
\begin{proof}
Expand the difference into its two single-error terms and its double-error
term, use unitarity of both dilations and apply Cauchy--Schwarz.
\end{proof}

\paragraph{A check against a circular sign argument.}
For $1/2<\Re w<1$ and $y(x)=x^{w-1}1_{(0,B)}(x)$, which is in $L^2$,
the finite formula on $0<t<B$ gives
\[
 \mathscr A_My(t)=t^{w-1}
 \left(\sum_{n\le M}n^{-w}-\frac{M^{1-w}}{1-w}\right)
                  +\frac{B^{w-1}}{1-w}.
\]
The elementary Euler-summation continuation of $\zeta$ therefore gives
\begin{equation}\label{eq:v111-power-response}
 \mathscr Ay(t)=\zeta(w)t^{w-1}+\frac{B^{w-1}}{1-w}.
\end{equation}
Indeed, comparison of $n^{-w}$ with its integral over $(n,n+1)$
gives a locally uniform tail $O_w(M^{-\Re w})$ on $\Re w>0$,
$w\ne1$; the limiting analytic function agrees with $\zeta$ on
$\Re w>1$ and hence throughout that domain.
This is consistent with the Mellin co-Poisson identities in
\cite{Burnol2004,BurnolCoPoisson2004}. At a hypothetical zero the
nonconstant power disappears, but this only re-encodes the zero condition.
The raw power is not asserted to satisfy the second Sonine support
condition. Thus neither this scalar calculation nor positivity of the
source Gram matrix proves the signed shell target
\eqref{eq:v110-sign-target}. The new adjoint supplies a justified
arithmetic constraint for a future quadratic estimate, not that estimate.

\section{A local-prime sieve identity inside the three-scale range}
The three-piece decomposition contains more arithmetic information than support covariance alone. Let $h=\mathcal D\phi$ with $\phi\in V_R$. Then $\int h=0$, so the co-Poisson counterterm vanishes and
\[
 E(h)(x)=\sum_{n\ge1}h(nx).
\]
Since
\[
 p^{-1/2}D_{1/p}f(x)=f(px),
\]
we obtain the exact identity
\begin{equation}\label{eq:psieve}
 \boxed{
 \left(I-p^{-1/2}D_{1/p}\right)E(h)(x)
 =\sum_{\substack{n\ge1\\ p\nmid n}}h(nx).}
\end{equation}
Thus the difference between the old arithmetic direction and its inverse-prime dilation is precisely the co-Poisson sum with the $p$-divisible lattice points removed.

On Mellin transforms, \eqref{eq:psieve} is the familiar local Euler-factor cancellation
\begin{equation}\label{eq:eulerremove}
 \widehat{E(h)}(s)\longmapsto
 (1-p^{-s})\widehat{E(h)}(s)
 =\zeta(s)(1-p^{-s})\widehat h(s),
\end{equation}
within the strip where the co-Poisson/Mellin identity is interpreted. The factor $\zeta(s)(1-p^{-s})$ is the zeta Euler product with the local factor at $p$ removed.

Applying the functional-equation involution $G$ gives the companion direction
\[
 G\left(I-p^{-1/2}D_{1/p}\right)G
 =I-p^{-1/2}D_p,
\]
whose Mellin factor is $1-p^{s-1}$. On the critical line these two local factors are complex conjugates and have equal modulus; off the line their radial behavior differs. This is a genuine prime-local structure absent from a generic pair of Sonine subspaces.

The identity does not by itself prove the shell sign: $y\in Q_{pR}$ is already orthogonal to both $W_R$ and $D_{1/p}W_R$, hence also to their $p$-sieved difference. Its value is that it identifies a concrete arithmetic correlation operator to study. If
\[
 C_c=P_{W_R}D_c\big|_{W_R},
\]
then norms of the $p$-sieved vectors determine the real part of $C_{1/p}$ through
\[
 \left\|w-p^{-1/2}D_{1/p}w\right\|^2
 =\left(1+p^{-1}\right)\norm w^2
 -2p^{-1/2}\Re\ip{w}{D_{1/p}w}.
\]
Consequently a quantitative lower bound for the $p$-sieved co-Poisson norm, uniform over the relevant source class and strong enough to survive the three-space Schur complement, would give new control of the cross-scale Gram operator appearing above. This is a narrower arithmetic target than an unrestricted positivity claim.


\section{Why a finite Euler-product sieve fails in the Hilbert norm}
The one-prime identity \eqref{eq:psieve} suggests iterating over all primes up to a cutoff $z$. Put
\[
 \mathfrak P(z)=\prod_{p\le z}p,
 \qquad
 \mathcal S_z=\prod_{p\le z}\left(I-p^{-1/2}D_{1/p}\right).
\]
For an admissible compact source $h=\mathcal D\phi$ one obtains exactly
\begin{equation}\label{eq:finite-sieve}
 \mathcal S_zE(h)(x)
 =\sum_{(n,\mathfrak P(z))=1}h(nx).
\end{equation}
For each fixed $x>0$ the right side eventually equals $h(x)$, because only finitely many integers occur and every fixed $n>1$ eventually has a prime divisor at most $z$. Thus the finite Euler-product sieve converges pointwise to the bare source. The Hilbert norm behaves very differently.

Let
\[
 \Psi(s)=\int_0^\infty h(t)t^{s-1}\,dt.
\]
Admissibility gives $\Psi(1)=0$ (and also $\Psi(0)=0$). If $h\ne0$, analyticity implies that there is a least integer $m\ge1$ such that $\Psi^{(m)}(1)\ne0$. Define
\[
 R_z(x)=\mathcal S_zE(h)(x)-h(x).
\]
Choose a compact interval $U\Subset(0,\inf\operatorname{supp}h)$. For $x=u/z$, $u\in U$, and $z$ sufficiently large, every integer contributing to $R_z(u/z)$ lies between $z$ and $C_Uz<z^2$. Hence an integer coprime to $\mathfrak P(z)$ in this range is necessarily prime. The prime number theorem with a smooth weight then gives, uniformly for $u\in U$,
\begin{equation}\label{eq:prime-shell-asymp}
 R_z(u/z)
 =\frac{(-1)^m\Psi^{(m)}(1)}{u}\,
   \frac{z}{(\log z)^{m+1}}
 +O_{h,U}\!\left(\frac{z}{(\log z)^{m+2}}\right).
\end{equation}
Indeed, the main term is
\[
 \frac zu\int \frac{h(t)}{\log(zt/u)}\,dt,
\]
and expansion of the denominator in powers of $1/\log z$ kills precisely the first $m$ logarithmic moments. Consequently
\begin{equation}\label{eq:sieve-R2-divergence}
 \int_{U/z}|R_z(x)|^2\,dx
 \sim
 |\Psi^{(m)}(1)|^2
 \left(\int_U\frac{du}{u^2}\right)
 \frac{z}{(\log z)^{2m+2}},
\end{equation}
which tends to $+\infty$.

Thus the sharp finite Euler-product inversion has the paradoxical behavior
\[
 \mathcal S_zE(h)(x)\to h(x)\quad\text{for every fixed }x>0,
 \qquad
 \norm{\mathcal S_zE(h)-h}_2\not\to0;
\]
in fact a fixed nonzero admissible source has a diverging prime-shell contribution. This is a scoped no-go theorem for the raw finite-prime sieve, not for all smoothed or nonlocal arithmetic inverses. It also explains why imposing any \emph{fixed} finite number of additional logarithmic moments cannot repair the mechanism: powers of $\log z$ cannot beat the factor $z^{1/2}$ in the $L^2$ norm. A successful Euler-factor construction would need a genuinely nonlocal tail or a number of matched conditions growing with the scale.

\section{The arithmetic-reflection defect as a dual zero channel}
Let $w$ be a hypothetical off-line zero and put $w^\sharp=1-\overline w$. The ambient de Branges reflection multiplier
\[
 B_w(s)=\frac{s-w^\sharp}{s-w}
\]
has modulus one on the critical line. The compact-division argument below shows that, at an off-line zero of multiplicity $m$, the finite-cutoff leakage
\[
 K_{w,R}=P_{Q_R}B_w\big|_{W_R}
\]
has rank one and factors as
\[
 K_{w,R}f=\ell_{w,R}(f)q_{w,R},
 \qquad
 \ell_{w,R}(E(\psi))=\Psi(w).
\]
At a simple zero the vector $q_{w,R}$ is characterized by vanishing completed Mellin values at all other zero constraints and a fixed nonzero value at $w$. This is the finite-cutoff analogue of the biorthogonal zero direction appearing in Burnol's complete/minimal systems; at the critical cutoff $a=1$ the dual system is represented, up to triangular changes within a multiple zero, by the functions $\zeta(s)/(s-\rho)^k$ \cite{Burnol2004}.

The rank-one channel has an exact prime transport law. Dilation covariance of the source gives
\[
 \ell_{w,pR}(D_pf)=p^{w-1/2}\ell_{w,R}(f),
 \qquad
 \ell_{w,pR}(D_{1/p}f)=p^{1/2-w}\ell_{w,R}(f).
\]
Since $B_w$ commutes with dilation and $D_{\pm1/p}W_R\subset W_{pR}$, comparison of the two rank-one factorizations yields
\begin{equation}\label{eq:dual-prime-transport}
 \boxed{P_{Q_{pR}}D_pq_{w,R}=p^{w-1/2}q_{w,pR}},
 \qquad
 \boxed{P_{Q_{pR}}D_{1/p}q_{w,R}=p^{1/2-w}q_{w,pR}}.
\end{equation}
Writing $d=|\Re w-1/2|$, contractivity gives
\begin{equation}\label{eq:q-upper}
 p^d\norm{q_{w,pR}}\le\norm{q_{w,R}}.
\end{equation}
For a simple zero, $M(q_{w,R})(w)$ is the fixed nonzero number
\[
 (w-w^\sharp)c_0(w)\zeta'(w),
\]
when the normalizing source has $\Psi(w)=1$. Cauchy--Schwarz against the zero-evaluation Riesz vector, together with the off-line evaluator growth in Assumption~\ref{ass:evaluators} $\norm{v_w^R}\asymp_w R^d$, therefore gives
\begin{equation}\label{eq:q-lower}
 \norm{q_{w,R}}\gg_w R^{-d}.
\end{equation}
Iterating \eqref{eq:q-upper} supplies the matching upper estimate along geometric cutoffs $R=p^nR_0$. Hence
\begin{equation}\label{eq:q-asymp}
 \norm{q_{w,p^nR_0}}\asymp_{w,p,R_0}(p^nR_0)^{-d}.
\end{equation}

Let $r_{w,R}\in W_R$ be the Riesz vector of the source-moment functional $\ell_{w,R}$. The same two dilation inclusions imply
\[
 \norm{r_{w,pR}}\ge p^d\norm{r_{w,R}}.
\]
Since $B_w$ is unitary on its ambient domain,
\[
 \norm{K_{w,R}}=\norm{r_{w,R}}\norm{q_{w,R}}\le1.
\]
Combining these estimates with \eqref{eq:q-lower} yields, for a hypothetical simple off-line zero,
\begin{equation}\label{eq:rankone-persistent}
 0<c_{w,p,R_0}
 \le \norm{K_{w,p^nR_0}}\le1
 \qquad(n\ge0).
\end{equation}
Thus the arithmetic-reflection defect does not disappear in operator norm even though every fixed source can be repaired asymptotically in the enlarging arithmetic spaces. The source-moment Riesz vector grows like $R^d$, the dual zero channel shrinks like $R^{-d}$, and the product retains order-one strength. This gives a precise moving-boundary explanation for why fixed-input convergence is insufficient.

\subsection{Full multiplicity channel: exact jet transport and persistent leakage volume}
The rank-one calculation above uses one division by $s-w$ and therefore exposes only the first source moment, even when the zero $w$ has higher multiplicity.  To treat a zero of multiplicity $m$ without discarding its jet structure, use the $m$-fold ambient reflection
\[
 \mathcal B_w^{(m)}(s)=\left(\frac{s-w^\sharp}{s-w}\right)^m,
 \qquad w^\sharp=1-\overline w,
\]
and define
\[
 \mathcal K_{w,R}^{(m)}=P_{Q_R}\mathcal B_w^{(m)}\big|_{W_R}.
\]
For $f=E(\psi)$ write $\Psi=M_0\psi$ and introduce the factorially normalized source jet
\[
 \mathcal J_{w,R}f=
 \left(\Psi(w),\Psi'(w),\ldots,\frac{\Psi^{(m-1)}(w)}{(m-1)!}\right)^T\in\mathbb C^m.
\]
Because $\zeta$ has order exactly $m$ at $w$, these functionals extend continuously from the source core to $W_R$: they are fixed triangular linear combinations of the completed-Mellin output derivatives of orders $m,\ldots,2m-1$.

\begin{theorem}[Exact multiplicity-$m$ leakage factorization]\label{thm:mleak}
Conditionally suppose that $w$ is a nontrivial zero of multiplicity $m$ with $\Re w\ne1/2$.  Then there is an injective map
\[
 \mathcal Q_{w,R}:\mathbb C^m\longrightarrow Q_R
\]
such that
\begin{equation}\label{eq:mfactor}
 \mathcal K_{w,R}^{(m)}=\mathcal Q_{w,R}\mathcal J_{w,R},
 \qquad
 \operatorname{rank}\mathcal K_{w,R}^{(m)}=m.
\end{equation}
If $h=\log p$ and $N$ is the nilpotent shift on factorial jet coordinates, set
\[
 M_p=p^{w-1/2}e^{hN}.
\]
Then the leakage basis obeys the exact prime transport laws
\begin{equation}\label{eq:mtransport}
 P_{Q_{pR}}D_p\mathcal Q_{w,R}=\mathcal Q_{w,pR}M_p,
 \qquad
 P_{Q_{pR}}D_{1/p}\mathcal Q_{w,R}=\mathcal Q_{w,pR}M_p^{-1}.
\end{equation}
\end{theorem}

\begin{proof}
Write locally
\[
 \zeta(s)=(s-w)^m a_w(s),\qquad a_w(w)\ne0.
\]
For a core vector $f=E(\psi)$,
\[
 M(\mathcal B_w^{(m)}f)(s)
 =c_0(s)a_w(s)(s-w^\sharp)^m\Psi(s).
\]
The first $m$ completed-Mellin derivatives at $w$ are therefore an invertible lower-triangular transformation of the source jet $\mathcal J_{w,R}f$. Every element of $W_R$ has those first $m$ derivatives equal to zero.
This alone does not prove factorization. If the source jet vanishes,
Lemma~\ref{lem:compact-division} shows that
$\Psi(s)/(s-w)^m$ is again the transform of an admissible compact source.
Multiplication by $(s-w^\sharp)^m$ preserves that class, so
$\mathcal B_w^{(m)}f\in W_R$. The kernel inclusion extends to $W_R$ by
Lemma~\ref{lem:jet-density}, using continuity of the jet and of the ambient
unitary multiplier. Therefore the leakage factors through the source jet.
The jet map is onto by Lemma~\ref{lem:jet-density}. The displayed triangular
output map is invertible because $w\ne w^\sharp$, so its factor into $Q_R$
is injective and the rank is exactly $m$. On the critical line the
multiplier is the identity and the leakage is zero; that case is excluded.

Dilation gives
\[
 \mathcal J_{w,pR}(D_pf)=p^{w-1/2}e^{hN}\mathcal J_{w,R}f=M_p\mathcal J_{w,R}f.
\]
The multiplier $\mathcal B_w^{(m)}$ commutes with dilations, and $D_pW_R,D_{1/p}W_R\subset W_{pR}$. Projecting the identity $D_{\pm p}\mathcal B_w^{(m)}f=\mathcal B_w^{(m)}D_{\pm p}f$ to $Q_{pR}$ gives \eqref{eq:mtransport}.  Notice that $\det e^{hN}=1$, hence
\[
 |\det M_p|=p^{m(\Re w-1/2)}.
\]
\end{proof}

The theorem has a basis-independent consequence.  Let
\[
 G_R=\mathcal Q_{w,R}^*\mathcal Q_{w,R},
 \qquad
 H_R=\mathcal J_{w,R}\mathcal J_{w,R}^*.
\]
Thus $G_R$ is the Gram matrix of the dual leakage directions and $H_R$ is the Gram matrix of the Riesz representers of the source-jet functionals.  Contractivity of the two projected dilation maps in \eqref{eq:mtransport} gives
\begin{align}
 M_p^*G_{pR}M_p&\preceq G_R,
 &M_p^{-*}G_{pR}M_p^{-1}&\preceq G_R,\label{eq:Gineq}\\
 H_{pR}&\succeq M_pH_RM_p^*,
 &H_{pR}&\succeq M_p^{-1}H_RM_p^{-*}.\label{eq:Hineq}
\end{align}
Writing
\[
 d=\left|\Re w-\frac12\right|,
\]
taking determinants yields
\begin{equation}\label{eq:det-one-sided}
 \det G_{pR}\le p^{-2md}\det G_R,
 \qquad
 \det H_{pR}\ge p^{2md}\det H_R.
\end{equation}

The remaining quantitative leakage conclusions in this section assume Assumption~\ref{ass:evaluators}. Only a one-sided ambient jet bound is needed.  Let $V_{w,R}$ be the Gram matrix of the first $m$ completed-Mellin evaluation jets at $w$ on the ambient Sonine space $H_R$.  For the raw representatives specified in Assumption~\ref{ass:evaluators}, the fixed finite raw jet family has the exact cutoff scaling
\[
 R^{1/2-w}e^{(\log R)N}
 \bigl(1_{t\ge1}t^{w-1},\,1_{t\ge1}t^{w-1}\log t,\ldots,\,
 1_{t\ge1}t^{w-1}(\log t)^{m-1}\bigr)
\]
for $\Re w<1/2$; for $\Re w>1/2$ the Fourier involution gives the reflected model.  The nilpotent factor has determinant one, and the fixed raw Gram matrix is finite.  Orthogonal Sonine projection can only decrease the Gram matrix in Loewner order.  Therefore
\begin{equation}\label{eq:jetgram-upper}
 \boxed{\det V_{w,R}\ll_{w,m}R^{2md}.}
\end{equation}
No relative $o(1)$ estimate for the whole projected jet family is used here.

Let $A_w$ be the fixed invertible triangular matrix taking the source jet to the first $m$ output jets in the proof of Theorem~\ref{thm:mleak}.  The Gram determinant inequality applied to
\[
 \Gamma_{w,R}\mathcal Q_{w,R}=A_w
\]
gives
\[
 |\det A_w|^2\le \det V_{w,R}\det G_R.
\]
Combining this with \eqref{eq:jetgram-upper} and the upper bound obtained by iterating \eqref{eq:det-one-sided}, along every geometric sequence $R=p^nR_0$,
\begin{equation}\label{eq:Gdet-asymp}
 \boxed{\det G_R\asymp_{w,m,p,R_0}R^{-2md}.}
\end{equation}

The nonzero singular values of $\mathcal K_{w,R}^{(m)}=\mathcal Q_{w,R}\mathcal J_{w,R}$ satisfy
\begin{equation}\label{eq:singprod}
 \prod_{j=1}^m\sigma_j(\mathcal K_{w,R}^{(m)})^2
 =\det G_R\det H_R.
\end{equation}
Since $\mathcal B_w^{(m)}$ is unitary on the ambient critical-line space and $P_{Q_R}$ is contractive, every $\sigma_j\le1$.  Equations \eqref{eq:Hineq} and \eqref{eq:Gdet-asymp} imply
\[
 \det H_R\gg R^{2md},
\]
whereas \eqref{eq:singprod}, $\sigma_j\le1$, and the lower bound for $\det G_R$ give
\[
 \det H_R\ll R^{2md}.
\]
Thus
\begin{equation}\label{eq:Hdet-asymp}
 \boxed{\det H_R\asymp_{w,m,p,R_0}R^{2md}}
\end{equation}
and, most importantly,
\begin{equation}\label{eq:singular-persistent}
 \boxed{
 0<c_{w,m,p,R_0}
 \le \prod_{j=1}^m\sigma_j(\mathcal K_{w,R}^{(m)})
 \le1
 }
 \qquad(R=p^nR_0).
\end{equation}
Because each singular value is at most one, \eqref{eq:singular-persistent} also bounds the smallest active singular value away from zero.  Hence the entire multiplicity-$m$ arithmetic-reflection leakage channel remains uniformly nondegenerate along geometric cutoffs.

\begin{remark}
This is the multiplicity analogue of the simple-zero compensation law.  The dual leakage volume decays like $R^{-md}$, the source-jet volume grows like $R^{md}$, and the nilpotent logarithmic factors cancel from the determinant.  Higher multiplicity therefore does not make the off-line defect disappear; it converts the scalar renormalization into a Jordan-coupled $m$-channel renormalization.
\end{remark}

\begin{remark}[What this still does not prove]
The persistent leakage theorem is conditional on the existence of an off-line zero and describes what the arithmetic reflection channel would then do.  It is not itself a contradiction: a unitary ambient reflection can have a uniformly nontrivial finite-dimensional leakage into $Q_R$.  To deduce RH one still needs an independent arithmetic or geometric principle forcing this leakage channel to vanish, become asymptotically singular, or conflict with a positive fixed-space prime action.  No such principle is proved here.
\end{remark}


\section{Bridge to the prime-compression defect}
The persistent reflection leakage is not merely an auxiliary finite-dimensional construction.  It is itself the defect of a canonical compression of the unitary zero-reflection multiplier, and it intertwines exactly with the cross-scale prime compression.

Let $R_R$ denote orthogonal projection onto $W_R$, and for a fixed off-line parameter $w$ define
\[
 A_{w,R}=R_RB_w\big|_{W_R},
 \qquad
 K_{w,R}=P_RB_w\big|_{W_R}.
\]
At a zeta zero $w$, every $f\in W_R$ has completed Mellin value zero at $w$, so Burnol's division theorem places $B_wf$ back in $H_R=W_R\oplus Q_R$.  Since $B_w$ is unitary on the ambient critical-line $L^2$ space,
\begin{equation}\label{eq:reflection-defect}
 \boxed{I_{W_R}-A_{w,R}^*A_{w,R}=K_{w,R}^*K_{w,R}.}
\end{equation}
Thus the rank-one leakage for one reflection, and the multiplicity-$m$ leakage for $B_w^m$, are canonical positive defect operators of the arithmetic reflection compression.

Now define the cross-scale prime compression
\[
 C_{R,p}=P_{pR}D_p\big|_{Q_R}:Q_R\longrightarrow Q_{pR}.
\]
Because $D_pW_R\subset W_{pR}$ and $B_wD_p=D_pB_w$, for every $f\in W_R$ one has
\begin{equation}\label{eq:leak-intertwine}
 \boxed{K_{w,pR}D_pf=C_{R,p}K_{w,R}f.}
\end{equation}
The same identity holds for the multiplicity-$m$ leakage operator $\mathcal K_{w,R}^{(m)}$.  In the jet basis of Theorem~\ref{thm:mleak},
\[
 C_{R,p}\mathcal Q_{w,R}=\mathcal Q_{w,pR}M_p,
 \qquad M_p=p^{w-1/2}e^{(\log p)N}.
\]
Consequently the cross-scale defect restricted to the leakage channel has the exact Gram form
\begin{equation}\label{eq:cross-defect-gram}
 \boxed{
 \mathcal Q_{w,R}^*(I-C_{R,p}^*C_{R,p})\mathcal Q_{w,R}
 =G_R-M_p^*G_{pR}M_p\succeq0.
 }
\end{equation}
This is the first exact positive operator identity in the project in which the persistent arithmetic-reflection channel and the prime compression appear in the same formula.

The fixed-space compression is nevertheless different.  Since $D_pQ_R\subset H_{pR}$, write
\[
 F_{R,p}=R_{pR}D_p\big|_{Q_R}.
\]
Then
\begin{equation}\label{eq:cross-defect-factor}
 \boxed{I-C_{R,p}^*C_{R,p}=F_{R,p}^*F_{R,p}.}
\end{equation}
Let
\[
 \mathsf R_{R,p}=P_R\big|_{Q_{pR}}:Q_{pR}\to Q_R,
 \qquad
 E_{R,p}=P_RR_{pR}D_p\big|_{Q_R}:Q_R\to Q_R.
\]
Decomposing $D_pq$ into its $Q_{pR}$ and $W_{pR}$ components and then projecting to $Q_R$ gives the exact return formula
\begin{equation}\label{eq:return-alias}
 \boxed{T_{R,p}=\mathsf R_{R,p}C_{R,p}+E_{R,p}.}
\end{equation}
Moreover $E_{R,p}=P_RF_{R,p}$, hence
\begin{equation}\label{eq:alias-bound}
 \boxed{E_{R,p}^*E_{R,p}\preceq I-C_{R,p}^*C_{R,p}.}
\end{equation}
The term $E_{R,p}$ is the arithmetic \emph{aliasing} created by projecting the newly available relations in $W_{pR}$ back into the older quotient $Q_R$.  It vanishes if the quotients are nested in the required direction, but no such nesting is available here.

Equations \eqref{eq:cross-defect-gram}--\eqref{eq:alias-bound} isolate the precise missing bridge.  Reflection leakage controls the cross-scale prime defect, while the RH-sensitive fixed-space operator also depends on the angle map $\mathsf R_{R,p}$ and the aliasing $E_{R,p}$.  Hilbert-space positivity alone does not compare the two defects.

\begin{proposition}[No generic domination between cross-scale and fixed-space defects]\label{prop:no-domination}
Even with nested ambient spaces $H_R\subset H_{pR}$, nested arithmetic spaces $W_R\subset W_{pR}$, and a unitary $U$ satisfying $UW_R\subset W_{pR}$, there is no universal operator inequality in either direction between the defects of
\[
 P_{Q_{pR}}U\big|_{Q_R}
 \quad\text{and}\quad
 P_{Q_R}U\big|_{Q_R}.
\]
\end{proposition}
\begin{proof}
Take an orthonormal basis $e_1,e_2,e_3$, set
\[
 H_R=\operatorname{span}\{e_1,e_2\},\quad W_R=\operatorname{span}\{e_2\},\quad Q_R=\operatorname{span}\{e_1\},
\]
and
\[
 H_{pR}=\mathbb C^3,\qquad
 W_{pR}=\operatorname{span}\{e_2,\cos\theta\,e_1+\sin\theta\,e_3\}.
\]
Then
\[
 Q_{pR}=\operatorname{span}\{-\sin\theta\,e_1+\cos\theta\,e_3\}.
\]
For $U=I$, one has $UW_R\subset W_{pR}$, the fixed-space compression on $Q_R$ is the identity, while the cross-scale compression has norm $|\sin\theta|$.  Thus the fixed-space defect is zero while the cross-scale defect is $\cos^2\theta$.

Conversely choose a unitary $U$ fixing $e_2$ and sending $e_1$ to the unit generator of $Q_{pR}$.  Again $UW_R\subset W_{pR}$, but now the cross-scale compression is isometric while the fixed-space compression has norm $|\sin\theta|$.  The two defects therefore reverse their sizes.  This rules out a comparison based only on positivity, nesting, and unitary transport.
\end{proof}

The proposition does not use the additional arithmetic structure of Burnol's spaces, so it is not a no-go theorem for the RH program.  Its role is diagnostic: any successful passage from \eqref{eq:cross-defect-gram} to the fixed-space defect must exploit a genuinely arithmetic estimate on $\mathsf R_{R,p}$ or $E_{R,p}$, not a generic projection inequality.


\section{Infinitesimal dilation and the exact source-range defect}
The finite-prime picture has an infinitesimal counterpart which unifies the source-moment obstruction, the reflection leakage, and the Jordan transport law.  Let
\[
 (\mathsf A f)(x)=-\frac12 f(x)-x f'(x)
\]
be the skew-adjoint generator of the ambient unitary dilation group $D_{e^t}=e^{t\mathsf A}$.  On the uncompleted Mellin transform,
\begin{equation}\label{eq:generator-mellin}
 M_0(\mathsf A f)(s)=(s-\tfrac12)M_0f(s).
\end{equation}
Let $\mathcal A_R=\mathcal D V_R$ be the admissible compact source class and let
\[
 \mathfrak C_R=E(\mathcal A_R),
\]
which is dense in $W_R$ by definition.  Dilation covariance implies
\[
 \mathsf A E(\psi)=E(\mathsf A\psi),
 \qquad
 \mathsf A\mathcal D=\mathcal D\mathsf A.
\]
Moreover, if $\Psi(0)=\Psi(1)=0$, then $(s-\tfrac12)\Psi(s)$ has the same two zero moments.  Hence
\begin{equation}\label{eq:generator-core-invariance}
 \boxed{\mathsf A\mathfrak C_R\subset\mathfrak C_R.}
\end{equation}
The restriction $\mathsf A_{W,R}=\mathsf A|_{\mathfrak C_R}$ is therefore a densely defined skew-symmetric operator in $W_R$.

Fix $w$ in the open strip, put $\lambda_w=w-\tfrac12$, and let $r\ge1$.  For $f=E(\theta)$ with $\Theta=M_0\theta$, equation \eqref{eq:generator-mellin} gives
\[
 (\mathsf A_{W,R}-\lambda_w)^r E(\eta)=E(\theta)
 \quad\Longleftrightarrow\quad
 (s-w)^r H(s)=\Theta(s),
\]
where $H=M_0\eta$.  Lemma~\ref{lem:compact-division} shows that $H=\Theta/(s-w)^r$ is again the Mellin transform of an admissible source in the same cutoff exactly when
\[
 \Theta(w)=\Theta'(w)=\cdots=\Theta^{(r-1)}(w)=0.
\]
Thus the source-core range is exactly
\begin{equation}\label{eq:generator-range-core}
 \operatorname{Ran}(\mathsf A_{W,R}-\lambda_w)^r
 =\left\{E(\theta):\theta\in\mathcal A_R,\ \Theta^{(j)}(w)=0\ (0\le j<r)\right\}.
\end{equation}

At a zeta zero $w$ of multiplicity $m$, the source-jet functionals
\[
 \ell_j(E(\theta))=\frac{\Theta^{(j)}(w)}{j!},\qquad 0\le j<m,
\]
extend continuously to $W_R$.  Indeed they are triangular linear combinations of completed-Mellin output derivatives of orders $m,\ldots,2m-1$, with nonzero leading coefficient determined by $\zeta^{(m)}(w)$.  They are linearly independent: choose one interior source $\eta$ with $H(w)\ne0$ and several sufficiently small dilations $D_{a_k}\eta$; the jet matrix is a nonzero triangular factor times a Vandermonde matrix in $\log a_k$.  Therefore
\begin{theorem}[Exact generator range defect]\label{thm:generator-defect}
Conditionally suppose that $w$ is a zeta zero of multiplicity $m$.  Then
\begin{equation}\label{eq:generator-codim}
 \boxed{
 \overline{\operatorname{Ran}(\mathsf A_{W,R}-\lambda_w)^m}
 =\bigcap_{j=0}^{m-1}\ker\ell_j,
 \qquad
 \operatorname{codim}_{W_R}=m.
 }
\end{equation}
The codimension is exact at every finite cutoff for which the source class is nontrivial.
\end{theorem}
\begin{proof}
The range on the left is the range of the algebraic power on the invariant
source core, not an asserted power of a closed operator. Compact division
gives \eqref{eq:generator-range-core}. The continuous independent source
jets have a surjective finite-dimensional joint map. Lemma~\ref{lem:jet-density}
identifies the closure of its core kernel with its full kernel, which has
codimension $m$.
\end{proof}

This theorem identifies the multiplicity-$m$ leakage channel as a higher-order range defect rather than an ad hoc correction.  The basic deficiency index is one; multiplicity $m$ appears through a length-$m$ generalized deficiency chain, as made precise below.  Indeed on the Mellin line
\begin{equation}\label{eq:cayley-reflection}
 B_w(s)=\frac{s-w^\sharp}{s-w}
 =1+\frac{w-w^\sharp}{s-w},
 \qquad w^\sharp=1-\overline w,
\end{equation}
so $B_w$ is the Cayley--Blaschke transform of the dilation generator at the spectral parameter $\lambda_w$.  The condition $\ell_0=0$ is precisely the condition under which the resolvent term in \eqref{eq:cayley-reflection} remains inside the same compact arithmetic source space.  At a zeta zero, the output factor $\zeta(w)=0$ guarantees that $B_wW_R\subset H_R$, and the unresolved source coordinate is therefore forced to appear internally as the leakage map $K_{w,R}:W_R\to Q_R$.  Repeating the division $m$ times produces the full jet channel of Theorem~\ref{thm:mleak}.

The prime transport matrix from that theorem is exactly the exponential of the jet generator:
\begin{equation}\label{eq:jet-exponential}
 M_p
 =p^{w-1/2}e^{(\log p)N}
 =\exp\!\left((\log p)\bigl((w-\tfrac12)I+N\bigr)\right).
\end{equation}
Thus the persistent finite-cutoff leakage, its Jordan factors, and its prime scaling are the finite-time form of the same infinitesimal range defect.

This also rules out a tempting shortcut.  The ambient generator $\mathsf A$ is skew-adjoint, but its arithmetic restriction on $W_R$ is not automatically maximal skew-symmetric: equation \eqref{eq:generator-codim} exhibits a finite-codimensional resolvent obstruction.  Merely invoking the ambient unitary dilation group therefore cannot yield a Hilbert--Polya proof.  One would have to show that the relevant deficiency channel is eliminated by an additional arithmetic boundary condition while retaining every zero mode.  That is precisely the compatibility theorem still missing from the fixed-space prime construction.


\section{Boundary-triplet audit and an exact zero-incidence criterion}
The infinitesimal formulation permits a sharper operator-theoretic audit.  Let
\[
 \mathsf S_R=\overline{\mathsf A_{W,R}}
\]
be the closure in $W_R$ of the skew-symmetric arithmetic dilation generator on the compact source core.  The first conclusion is that the basic deficiency dimension is one, not the multiplicity of a zeta zero.

For $w$ with $0<\Re w<1$ and $\zeta(w)\ne0$, define the source-evaluation functional on $W_R$ by
\[
 \ell_w(f)=\frac{\mathcal M f(w)}{c_0(w)\zeta(w)}.
\]
It is continuous because completed Mellin evaluation is continuous on $H_R$.  At a zeta zero the same functional extends by the derivative formula already used in the source-moment audit.  On the source core,
\[
 \operatorname{Ran}(\mathsf A_{W,R}-(w-\tfrac12))
 =\ker\ell_w\cap\mathfrak C_R.
\]
Since one may correct a core approximant by a fixed core vector with source moment one, closure gives
\begin{equation}\label{eq:first-range-kernel}
 \overline{\operatorname{Ran}(\mathsf S_R-(w-\tfrac12))}
 =\ker\ell_w.
\end{equation}
Choose $0<a<1/2$ and the real points $w_\pm=1/2\pm a$, where $\zeta(w_\pm)\ne0$.  Because $\mathsf S_R$ is skew-symmetric,
\[
 \norm{(\mathsf S_R\mp a)f}^2=\norm{\mathsf S_Rf}^2+a^2\norm f^2,
\]
so these ranges are closed.  Equation \eqref{eq:first-range-kernel} therefore implies
\[
 \dim\ker(\mathsf S_R^*-a)=\dim\ker(\mathsf S_R^*+a)=1.
\]
Thus $\mathsf S_R$ has deficiency indices $(1,1)$.  If $r_w\in W_R$ denotes the Riesz vector of $\ell_w$, then
\begin{equation}\label{eq:deficiency-riesz}
 r_w\in\ker\!\left(\mathsf S_R^*-(\overline w-\tfrac12)\right).
\end{equation}
The involution $G$ preserves $W_R$ and anti-commutes with the dilation generator.  On source moments it sends $\Psi(s)$ to $\Psi(1-s)$, hence
\begin{equation}\label{eq:G-deficiency}
 Gr_w=r_{1-w}.
\end{equation}
The deficiency vectors are total: if $f\in W_R$ is orthogonal to $r_w$ on any open set of parameters avoiding zeta zeros, then the entire completed Mellin transform of $f$ vanishes on an open set and therefore vanishes identically.  Hence $\mathsf S_R$ is simple in the standard symmetric-operator sense.

Standard Liv\v{s}ic/von Neumann theory therefore supplies a one-parameter family of maximal skew-adjoint extensions of $\mathsf S_R$.  If $e_\pm$ are normalized generators of $\ker(\mathsf S_R^*\mp a)$ and
\[
 f=f_0+c_+e_++c_-e_-,\qquad f_0\in\mathcal D(\mathsf S_R),
\]
then the boundary form is
\begin{equation}\label{eq:boundary-form}
 \ip{\mathsf S_R^*f}{g}+\ip{f}{\mathsf S_R^*g}
 =2a\left(c_+\overline{d_+}-c_-\overline{d_-}\right).
\end{equation}
A maximal skew-adjoint extension is the restriction of $-\mathsf S_R^*$ (not $\mathsf S_R^*$) to the domain specified by one fixed unitary boundary relation $c_-=e^{i\vartheta}c_+$.  This statement is ordinary extension theory; it does not identify any such extension with the zeta zeros.

The zeta condition enters through an additional ambient incidence statement.  Fix an off-critical parameter $w$ and choose an admissible source $\eta_w\in\mathcal A_R$ with Mellin transform $H_w$ satisfying $H_w(w)\ne0$; the bump construction from the source-moment repair gives such a source.  Put
\[
 f_w=E(\eta_w)\in W_R,
 \qquad
 B_w(s)=\frac{s-(1-\overline w)}{s-w}.
\]
Then
\begin{theorem}[Exact zero-incidence criterion]\label{thm:zero-incidence}
For $0<\Re w<1$, $\Re w\ne1/2$,
\begin{equation}\label{eq:zero-incidence}
 \boxed{\zeta(w)=0\quad\Longleftrightarrow\quad B_wf_w\in H_R.}
\end{equation}
Equivalently, if $S_R^{(H)}$ denotes orthogonal projection onto $H_R$,
\begin{equation}\label{eq:zero-incidence-distance}
 \boxed{\zeta(w)=0\quad\Longleftrightarrow\quad
 \norm{(I-S_R^{(H)})B_wf_w}=0.}
\end{equation}
At a zero the non-arithmetic component therefore lies entirely in $Q_R$; away from a zero it has a nonzero $H_R^\perp$ component.
\end{theorem}
\begin{proof}
The completed Mellin transform of $f_w$ is
\[
 \mathcal M f_w(s)=c_0(s)\zeta(s)H_w(s).
\]
If $\zeta(w)=0$, then $\mathcal M f_w(w)=0$, and Burnol's ambient division theorem places $B_wf_w$ in $H_R$.  Conversely, if $B_wf_w\in H_R$, its completed Mellin transform is entire.  On the critical line it agrees with
\[
 \frac{s-(1-\overline w)}{s-w}\,c_0(s)\zeta(s)H_w(s),
\]
and hence with its meromorphic continuation.  Since $w\ne1-\overline w$ and $c_0(w)H_w(w)\ne0$, removability of the apparent pole at $s=w$ forces $\zeta(w)=0$.
\end{proof}

The theorem clarifies both the promise and the limitation of the boundary-triplet route.  The source generator $\mathsf S_R$ itself has the familiar deficiency pattern $(1,1)$ and therefore admits fixed maximal skew-adjoint boundary conditions.  But the zeta zeros are not presently characterized as eigenvalues of one of those extensions.  They are characterized by the external incidence condition \eqref{eq:zero-incidence}: the Cayley--Blaschke defect generated from $W_R$ must land in the larger Sonine space $H_R$, and at a zero its residual part is precisely the $Q_R$ leakage channel studied above.  The zero-evaluation vectors themselves live in $Q_R$, not in $W_R$.

Consequently a genuine Hilbert--P\'olya conclusion from this picture still requires one new theorem: identify the incidence condition \eqref{eq:zero-incidence}, or an equivalent condition on the $Q_R$ leakage, with a \emph{single zero-independent} maximal skew-adjoint boundary condition or with the characteristic function of one fixed conservative coupling.  No such identification is proved here.  In particular, choosing a boundary condition separately as a function of $w$ would simply repackage the zero problem and would be circular.

The correction above also explains the role of multiplicity.  The operator $\mathsf S_R$ has one-dimensional basic deficiency spaces.  A zero of multiplicity $m$ produces the codimension-$m$ closure of the algebraic source-core range of $(\mathsf A_{W,R}-(w-1/2))^m$ and the length-$m$ Jordan/source-jet chain of Theorem~\ref{thm:mleak}; it does not change the basic deficiency indices from $(1,1)$.


\section{A nested rank-one defect ladder and defect transfer}
The zero-incidence theorem admits a sharper quantitative form.  There are two nested compression defects associated with the same ambient Blaschke multiplier $B_w$: one for the Sonine space $H_R$, and one for its arithmetic subspace $W_R$.

Let $S_R^{(H)}$ and $R_R$ denote the orthogonal projections onto $H_R$ and $W_R$, respectively, and define
\[
 \Delta^{H}_{w,R}=(I-S_R^{(H)})B_w\big|_{H_R},
 \qquad
 \Delta^{W}_{w,R}=(I-R_R)B_w\big|_{W_R}.
\]
Burnol's division property implies that $\Delta^H_{w,R}$ vanishes on the codimension-one evaluation kernel
\[
 \{f\in H_R:\mathcal M f(w)=0\}.
\]
It is not the zero operator: if $k_w^R$ is the Riesz vector for completed Mellin evaluation, then $B_wk_w^R$ has an uncancelled pole at $w$ and therefore cannot belong to $H_R$.  Hence $\Delta^H_{w,R}$ has rank exactly one.  Put
\[
 d^H_{w,R}=\frac{\Delta^H_{w,R}k_w^R}{\norm{k_w^R}^2}.
\]
Then
\begin{equation}\label{eq:H-defect-factor}
 \boxed{\Delta^H_{w,R}f=\mathcal M f(w)\,d^H_{w,R}\qquad(f\in H_R).}
\end{equation}

Likewise the exact source-division identity shows that $\Delta^W_{w,R}$ vanishes on $\ker\ell_w$.  Thus it also has rank one.  If $r_w^R$ is the Riesz vector of the source-evaluation functional $\ell_w$ on $W_R$, set
\[
 d^W_{w,R}=\frac{\Delta^W_{w,R}r_w^R}{\norm{r_w^R}^2}.
\]
Then
\begin{equation}\label{eq:W-defect-factor}
 \boxed{\Delta^W_{w,R}f=\ell_w(f)\,d^W_{w,R}\qquad(f\in W_R).}
\end{equation}
For $f\in W_R$ the arithmetic Mellin factorization gives
\[
 \mathcal M f(w)=c_0(w)\zeta(w)\ell_w(f).
\]
Projecting \eqref{eq:W-defect-factor} further onto $H_R^\perp$ and comparing with \eqref{eq:H-defect-factor} therefore gives the exact vector identity
\begin{equation}\label{eq:defect-transfer-vector}
 \boxed{
 d^W_{w,R}=q_{w,R}+c_0(w)\zeta(w)d^H_{w,R},
 \qquad q_{w,R}=P_{Q_R}d^W_{w,R}\in Q_R.
 }
\end{equation}
The two summands are orthogonal because
\[
 W_R^\perp=Q_R\oplus H_R^\perp.
\]
Consequently
\begin{equation}\label{eq:defect-transfer-pythagoras}
 \boxed{
 \norm{d^W_{w,R}}^2
 =\norm{q_{w,R}}^2
 +|c_0(w)\zeta(w)|^2\norm{d^H_{w,R}}^2.
 }
\end{equation}

Equation \eqref{eq:defect-transfer-pythagoras} is a defect-conservation law.  The source generator always has one nontrivial Blaschke/Cayley defect direction off the critical line.  At a zeta zero its $H_R^\perp$ component vanishes and the entire defect lies internally in the arithmetic quotient $Q_R$; away from a zero a nonzero portion necessarily escapes the Sonine space.  Thus a zero does not annihilate the defect.  It transfers the defect from the outer Sonine channel to the inner arithmetic channel.

Theorem~\ref{thm:zero-incidence} can therefore be restated as the angle condition
\begin{equation}\label{eq:defect-angle-zero}
 \zeta(w)=0
 \quad\Longleftrightarrow\quad
 d^W_{w,R}\in Q_R
 \quad\Longleftrightarrow\quad
 \frac{|c_0(w)\zeta(w)|^2\norm{d^H_{w,R}}^2}{\norm{d^W_{w,R}}^2}=0.
\end{equation}
This resembles a scalar characteristic or transmission coefficient for the nested pair $W_R\subset H_R$, but no identification with the Liv\v{s}ic characteristic function of $\mathsf S_R$ is asserted here.  Establishing such an identification, with a zero-independent conservative colligation, would be a genuinely new theorem rather than a change of notation.


\section{Exact same-label returns are overstrong; the asymptotic regime is calibrated}
The complete/minimal zero systems suggest a formally natural same-label return between two scale fibers.  A necessary caution is that \emph{exact contractivity} of such a return is stronger than RH, whereas the asymptotic compression used earlier is correctly calibrated.

\begin{proposition}[Exact diagonal contractions force orthogonality]\label{prop:exact-return-overstrong}
Let $\mathcal H$ be a Hilbert space and let $(x_j)_{j\in J}$ have dense linear span.  Suppose a contraction $A$ satisfies
\[
 Ax_j=\lambda_jx_j,\qquad |\lambda_j|=1
\]
for every $j$.  Then $A$ is unitary.  In particular,
\[
 \lambda_i\ne\lambda_j\quad\Longrightarrow\quad \ip{x_i}{x_j}=0.
\]
\end{proposition}
\begin{proof}
For every $j$, equality of norms gives
\[
 0=\norm{x_j}^2-\norm{Ax_j}^2
   =\ip{(I-A^*A)x_j}{x_j}.
\]
Since $I-A^*A\succeq0$, this implies $(I-A^*A)x_j=0$.  Density and continuity give $A^*A=I$, so $A$ is an isometry.  Every $x_j$ belongs to $\operatorname{Ran}A$ because $x_j=\lambda_j^{-1}Ax_j$; hence the range is dense.  The range of an isometry is closed, so $A$ is onto and therefore unitary.  Distinct eigenspaces of a unitary operator are orthogonal.
\end{proof}

Under RH the desired prime multipliers $\lambda_\rho=p^{1/2-\rho}$ all have modulus one.  Thus an \emph{exact} contractive fixed-space operator diagonal on a complete zero system would force exact orthogonality between all zero vectors carrying distinct prime phases.  That orthogonality is an additional requirement not contained in RH.  Consequently exact same-label contractivity is not the correct target for the Burnol zero systems.

This is consistent with Burnol's later study of the biorthogonal family $\zeta(s)/(s-\rho)^k$: the system and its dual are complete and minimal, but bounded unconditional reconstruction is not supplied by completeness/minimality alone.  The hereditary-completeness problem is genuinely subtler; in the admissible mixed systems studied there the possible defect is controlled only up to codimension one.  Thus a formal biorthogonal return must not be treated as a bounded, let alone contractive, operator without a separate theorem.

Under Assumption~\ref{ass:evaluators}, the moving-cutoff construction avoids the preceding overconstraint because fixed critical-line zero kernels become asymptotically orthogonal. The following calculation records the consequence of those quantitative inputs, not a new proof of them.  For two fixed distinct critical-line parameters
\[
 w=\tfrac12+i\tau,\qquad z=\tfrac12+i\sigma,\qquad \tau\ne\sigma,
\]
let $Z_w^R,Z_z^R$ be Burnol's completed-Mellin evaluation vectors.  The raw Sonine representatives satisfy
\[
 C_1(v,w)=\gamma_+(w)v^{-w}+O_w(1),
 \qquad
 C_1(v,z)=\gamma_+(z)v^{-z}+O_z(1)
\]
near zero.  Hence their cross term has leading integral
\[
 \gamma_+(w)\overline{\gamma_+(z)}
 \int_{R^{-2}}^1 v^{-1-i(\tau-\sigma)}\,dv=O_{w,z}(1),
\]
while all remaining local and tail terms are bounded.  The Sonine projection changes each raw evaluator by $O_w(R^{-1})$, so
\begin{equation}\label{eq:critical-cross-bounded}
 \ip{Z_w^R}{Z_z^R}=O_{w,z}(1).
\end{equation}
On the other hand
\[
 \norm{Z_w^R}^2=C_w\log R+O_w(1),
 \qquad
 \norm{Z_z^R}^2=C_z\log R+O_z(1),
 \qquad C_w,C_z>0.
\]
Therefore
\begin{equation}\label{eq:critical-asymptotic-orthogonality}
 \boxed{
 \frac{\ip{Z_w^R}{Z_z^R}}
 {\norm{Z_w^R}\norm{Z_z^R}}
 =O_{w,z}\!\left(\frac1{\log R}\right).}
\end{equation}
For any fixed finite set of distinct critical-line parameters, the normalized Gram matrix therefore tends to the identity.

\begin{corollary}[Finite-mode asymptotic calibration]\label{cor:finite-mode-calibration}
Assume Assumption~\ref{ass:evaluators}. Fix a finite set $F$ of distinct critical-line parameters and unit-modulus numbers $\lambda_w$.  On
\[
 \mathcal V_R(F)=\operatorname{span}\{Z_w^R:w\in F\}
\]
define the diagonal operator $A_R(F)Z_w^R=\lambda_w Z_w^R$.  Then
\[
 \norm{A_R(F)}=1+O_F\!\left(\frac1{\log R}\right).
\]
In contrast, if a finite set contains a parameter with $|\lambda_w|>1$, then every operator having that vector as an exact eigenvector satisfies $\norm A\ge |\lambda_w|>1$.
\end{corollary}
\begin{proof}
After normalizing the basis vectors, equation \eqref{eq:critical-asymptotic-orthogonality} gives a Gram matrix $G_R=I+O_F(1/\log R)$.  The squared norm of the diagonal phase operator is the largest generalized eigenvalue of the pencil $(D^*G_RD,G_R)$.  Here $D=\operatorname{diag}(\lambda_w)$, so this eigenvalue equals $1+O_F(1/\log R)$.  The second assertion is immediate from the norm of an eigenvector.
\end{proof}

The correct fixed-space target is therefore \emph{asymptotic} rather than exact contractivity.  A useful non-circular formulation would construct, from the arithmetic source geometry and not from the zero list itself, operators $R_{R,p,T}$ on a fixed finite spectral window such that
\[
 \limsup_{R\to\infty}\norm{R_{R,p,T}}\le1
\]
for every fixed $T$.  Any hypothetical off-line zero inside the window would force a fixed eigenvalue modulus $p^{1/2-\Re\rho}>1$ and contradict this bound, while under RH the finite-mode Gram calibration above is compatible with norm $1+o(1)$.  The remaining difficulty is to define such a return by a zero-independent co-Poisson/prime formula.  Defining it spectrally from the zero vectors would merely encode the desired conclusion.


\section{Off-line resolvent localization and why the determinant shortcut is tautological}
The fixed Hilbert-space resolvent construction admits a useful extension away from the critical line.  This gives a zero-independent spectral window, but it also shows that determinant or spectral-radius inequalities on that window are endpoints equivalent to zero exclusion rather than independent mechanisms.

\begin{assumption}[Fixed closed arithmetic realization]\label{ass:closed-operator}
A fixed Hilbert space $\mathscr H_{\mathrm B}$ of functions holomorphic
in the nontrivial strip is specified, with continuous off-line evaluations
and an isometric critical-line realization. A closed densely defined
operator $A$ has graph formula
\[
 AF(s)=sF(s)-\eta_{\mathrm B}(F)\zeta(s)
\]
on its stated domain. The finite span of
$F_{\rho,j}=\zeta(s)/(s-\rho)^j$ is a graph core, with
$AF_{\rho,j}=\rho F_{\rho,j}+F_{\rho,j-1}$ for $j>1$ and
$AF_{\rho,1}=\rho F_{\rho,1}$.
The bounded divided-difference resolvents below preserve the realization,
and the resulting off-line meromorphic resolvent has no additional
spectral subspaces beyond the indicated finite root blocks.
\end{assumption}
This is an explicit realization hypothesis, not a completed operator
construction in this manuscript. The inherited critical-line norm, the
domain, closure, and graph-core theorem must be supplied together before
using the spectral claims that follow. Completeness and minimality of the
root system do not imply the graph-core assertion. Later finite root-space
sampling identities can be defined algebraically without this assumption;
calling their source space a Riesz window requires it.


Let $\mathscr H_{\mathrm B}$ denote the fixed right-Mellin Hilbert realization in Assumption~\ref{ass:closed-operator}, let $A$ be the closed operator defined there, and let
\[
 F_{\rho,j}(s)=\frac{\zeta(s)}{(s-\rho)^j},\qquad 1\le j\le m_\rho,
\]
be its complete minimal zero/jet system.  For a point $z$ satisfying
\[
 \Re z\ne\tfrac12,\qquad \zeta(z)\ne0,
\]
define
\begin{equation}\label{eq:complex-resolvent}
 (R_zF)(s)=\frac{F(s)-F(z)\zeta(s)/\zeta(z)}{s-z}.
\end{equation}
Point evaluation at $z$ is continuous on $\mathscr H_{\mathrm B}$.  Moreover, if
$\delta_z=|\Re z-1/2|$, then multiplication by $(s-z)^{-1}$ on the critical line has norm at most $\delta_z^{-1}$.  Since $\zeta(s)/s$ belongs to the ambient Mellin $L^2$ space and $s/(s-z)$ is bounded on the critical line, one also has
\[
 q_z(s):=\frac{\zeta(s)}{s-z}\in L^2.
\]
Consequently $R_z$ is bounded in the ambient norm, with the explicit coarse estimate
\begin{equation}\label{eq:complex-resolvent-bound}
 \norm{R_zF}
 \le\left(\delta_z^{-1}+\frac{\norm{k_z}\norm{q_z}}{|\zeta(z)|}\right)\norm F,
\end{equation}
where $k_z$ is the Riesz vector for evaluation at $z$.

On the finite zero/jet core the numerator in \eqref{eq:complex-resolvent} factors algebraically, so $R_z$ maps the core to itself and is the triangular Jordan resolvent.  Boundedness therefore extends $R_z$ to $\mathscr H_{\mathrm B}$.  The separately assumed graph-core property of the zero/jet system then upgrades the algebraic identities to
\[
 (A-zI)R_z=I,\qquad R_z(A-zI)=I
\]
on their natural domains.  Thus:

\begin{conditionalresult}[Off-line spectrum under a closed-operator realization]\label{thm:off-line-spectrum}
Assume Assumption~\ref{ass:closed-operator}. In the nontrivial strip away from the critical line,
\[
 \sigma(A)\cap\{0<\Re z<1:\Re z\ne\tfrac12\}
 =\{\rho:\zeta(\rho)=0,\ \Re\rho\ne\tfrac12\},
\]
with algebraic multiplicities matching the zeta multiplicities.
\end{conditionalresult}
\begin{proof}
If $\zeta(z)\ne0$, equations \eqref{eq:complex-resolvent}--\eqref{eq:complex-resolvent-bound} and the graph-core argument above show that $z\in\rho(A)$.  If $z=\rho$ is a nontrivial zeta zero of order $m_\rho$, $F_{\rho,1}$ is an eigenvector of $A$ and $F_{\rho,1},\ldots,F_{\rho,m_\rho}$ form the displayed Jordan chain.  The chain cannot be extended: if $G\in D(A)$ satisfied $(A-\rho I)G=F_{\rho,m_\rho}$, then the graph formula for $A$ would force
\[
 G(s)=\frac{\zeta(s)}{(s-\rho)^{m_\rho+1}}
      +c\,\frac{\zeta(s)}{s-\rho}
\]
for some scalar $c$.  The first term has a nonremovable pole at the nontrivial zero $\rho$, whereas elements of $\mathscr H_{\mathrm B}$ have no such pole.  Thus the root space has exactly the zeta multiplicity.  No assertion about spectrum exactly on the critical line is needed for this theorem.
\end{proof}

Under the same realization assumption, a compact contour $\Gamma=\partial\Omega$ lying strictly to one side of the critical line and avoiding zeta zeros defines a genuinely zero-independent Riesz projection
\begin{equation}\label{eq:riesz-window}
 \Pi_\Omega=\frac{1}{2\pi i}\oint_\Gamma (zI-A)^{-1}\,dz.
\end{equation}
Its range is precisely the finite-dimensional sum of the zero/jet blocks inside $\Omega$, without listing those zeros in the definition.

For a prime $p$ define on this window
\begin{equation}\label{eq:window-prime-calculus}
 U_{p,\Omega}
 =\frac{1}{2\pi i}\oint_\Gamma
 p^{1/2-z}(zI-A)^{-1}\,dz.
\end{equation}
On a zero block at $\rho$ this has the exact action
\[
 p^{1/2-\rho}\exp[-(\log p)N].
\]
Thus \eqref{eq:window-prime-calculus} is an exact fixed-space, zero-independent finite-window prime action.

It is tempting to seek $r(U_{p,\Omega})\le1$ or $|\det U_{p,\Omega}|\le1$ for a window strictly left of the critical line.  The determinant calculation shows why this is not by itself a new mechanism.  Holomorphic functional calculus gives
\begin{equation}\label{eq:window-det}
 \det\!\left(U_{p,\Omega}|_{\operatorname{Ran}\Pi_\Omega}\right)
 =\prod_{\rho\in\Omega}p^{m_\rho(1/2-\rho)},
\end{equation}
so
\begin{equation}\label{eq:window-det-modulus}
 \log\left|\det U_{p,\Omega}\right|
 =(\log p)\sum_{\rho\in\Omega}m_\rho(1/2-\Re\rho).
\end{equation}
Equivalently, when the contour avoids the pole and trivial zeros,
\begin{equation}\label{eq:window-det-contour}
 \log\det U_{p,\Omega}
 =\frac{\log p}{2\pi i}
 \oint_\Gamma\left(\frac12-z\right)\frac{\zeta'(z)}{\zeta(z)}\,dz,
\end{equation}
with a consistent branch on the finite-dimensional determinant.  For a left window every summand in \eqref{eq:window-det-modulus} is strictly positive.  Therefore the inequality $|\det U_{p,\Omega}|\le1$ is equivalent to the emptiness of that window.  The spectral-radius inequality has the same issue.  These quantities are useful diagnostics and exact spectral selectors, but a proof still requires an independent arithmetic estimate before taking the determinant.

\subsection*{A global hyponormality shortcut is impossible}
Let $T_{R,p}=P_RD_p|_{Q_R}$ and use Burnol's involution $G$, which preserves $Q_R$ and satisfies $GT_{R,p}G=T_{R,p}^*$.  Its self-commutator
\[
 C_{R,p}=T_{R,p}^*T_{R,p}-T_{R,p}T_{R,p}^*
\]
therefore obeys the exact anti-symmetry
\begin{equation}\label{eq:commutator-antisym}
 \boxed{GC_{R,p}G=-C_{R,p}.}
\end{equation}
Hence $C_{R,p}\succeq0$ or $C_{R,p}\preceq0$ on all of $Q_R$ would force $C_{R,p}=0$, i.e. normality of the compression.  Thus global hyponormality/cohyponormality cannot supply the missing one-sided zero exclusion unless one first proves the much stronger normality statement.  We do not claim a side-dependent asymptotic value for the diagonal of $C_{R,p}$: the existing inward-mode audit determines one of the two norm limits in each off-line half-plane, but not the complementary outward norm required for such a formula.

\section{Comparison with the finite-prime Weil-form program}
A recent and independently developed spectral program of Connes--Consani--Moscovici constructs, from the Weil quadratic form on $[\lambda^{-1},\lambda]$, self-adjoint finite-prime operators whose regularized determinants have only critical-line zeros.  Their finite form on this interval involves only prime-power terms $n\le\lambda^2$, the same multiplicative scale ratio that appears naturally in the Burnol source window $[1/R,R]$.  Their theorem produces a self-adjoint rank-one perturbation of the logarithmic scaling operator in a positive inner product obtained from the truncated Weil form, and hence an entire determinant with real zeros.  The stated missing step is convergence of those finite-prime determinants/spectra toward the Riemann $\Xi$ function; a rigorous proof of that convergence would establish RH \cite{ConnesConsaniMoscovici2025}.

\subsection*{A conditional obstruction to norm-resolvent convergence}
A unitary norm-resolvent comparison has the following necessary condition in the inherited Burnol norm.  Set
\[
 B=A-\frac12 I.
\]
On a critical-line zero block its eigenvalues are $i\gamma$.  If $D_n$ are self-adjoint finite-prime operators and $U_n$ are unitary identifications into one fixed Hilbert space, then $iU_nD_nU_n^*$ are skew-adjoint.  A norm-resolvent limit of skew-adjoint operators is again skew-adjoint.  But the exact boundary form from the resolvent audit gives
\[
 \ip{BF_\rho}{F_\sigma}+\ip{F_\rho}{BF_\sigma}
 =i(\gamma_\rho-\gamma_\sigma)\ip{F_\rho}{F_\sigma}
\]
for critical-line zero vectors.  If a pair of distinct critical ordinates satisfies
$\ip{F_\rho}{F_\sigma}\ne0$, this identity rules out skew-symmetry of $B$
and hence such a norm-resolvent limit to $B$ (equivalently a shifted limit
to $A$). The earlier draft referred to a certified two-zero calculation
but supplied no reproducible certificate. This revision does not assert
that numerical premise. The obstruction is conditional on verifying it
in the specified Hilbert realization. Finite-dimensional approximations
also require a precise common-space realization before norm-resolvent
convergence is a well-typed statement.

This does not rule out a comparison through a changing positive metric.  A nonnormal operator may be similar to a skew-adjoint one in an equivalent Hilbert norm, and the finite-prime Weil construction naturally provides such positive forms at each cutoff.

\subsection*{A finite-window Lyapunov metric criterion}
The preceding observation suggests a weaker and correctly typed bridge.  Let $\Omega$ be a compact Riesz window contained strictly in one half of the critical strip, let $\Pi_\Omega$ be \eqref{eq:riesz-window}, and keep $B=A-\frac12I$.

\begin{theorem}[Uniform metric defect excludes an off-line window]\label{thm:metric-window}
Assume the Riesz realization in Assumption~\ref{ass:closed-operator}. Suppose there are bounded positive operators $G_n$ on $\mathscr H_{\mathrm B}$ and numbers $c_\Omega>0$, $\varepsilon_n\to0$ such that for all sufficiently large $n$,
\begin{equation}\label{eq:metric-coercive}
 \ip{G_nf}{f}\ge c_\Omega\norm{f}^2,
 \qquad f\in\operatorname{Ran}\Pi_\Omega,
\end{equation}
and, on this finite-dimensional spectral subspace,
\begin{equation}\label{eq:metric-lyap}
 \left|\ip{G_nBf}{f}+\ip{G_nf}{Bf}\right|
 \le \varepsilon_n\norm{f}^2.
\end{equation}
Then $\operatorname{Ran}\Pi_\Omega=\{0\}$.
\end{theorem}
\begin{proof}
If $\rho\in\Omega$ is a zero and $F_\rho$ is a value eigenvector, then $BF_\rho=(\rho-\frac12)F_\rho$.  Substituting $f=F_\rho$ in \eqref{eq:metric-lyap} gives
\[
 2\left|\Re\rho-\frac12\right|\ip{G_nF_\rho}{F_\rho}
 \le\varepsilon_n\norm{F_\rho}^2.
\]
By \eqref{eq:metric-coercive},
\[
 2c_\Omega\left|\Re\rho-\frac12\right|\le\varepsilon_n,
\]
which is impossible for large $n$ because $\Omega$ stays a positive distance from the critical line.  Thus no zero block occurs in $\Omega$, and Theorem~\ref{thm:off-line-spectrum} gives $\Pi_\Omega=0$.
\end{proof}

The theorem is deliberately finite-window.  A single globally bounded positive metric making $B$ skew-adjoint would be substantially stronger and would impose global basis geometry on Burnol's complete/minimal zero system.  For RH it is enough to obtain \eqref{eq:metric-coercive}--\eqref{eq:metric-lyap} separately on each compact off-line window.

The finite-prime Weil forms make this target concrete.  On the Connes--Consani--Moscovici finite-dimensional space, under their stated hypotheses, including the required ground-state and boundary properties, subtracting the lowest form eigenvalue and quotienting its null direction makes the corrected logarithmic operator self-adjoint in the resulting positive form \cite{ConnesConsaniMoscovici2025}.  To turn that finite-prime positivity into Theorem~\ref{thm:metric-window}, one would need a zero-independent transport of those forms into the Burnol fixed space for which two properties survive simultaneously: a cutoff-uniform coercive lower bound on $\operatorname{Ran}\Pi_\Omega$, and a Lyapunov defect tending to zero.  This is weaker than direct norm-resolvent convergence and avoids the exact-orthogonality obstruction, but no such transport estimate is proved here.

This formulation also identifies what would count as genuine progress from the finite-prime side.  Mere convergence of eigenvalue lists, or convergence of the determinant after zeros are already matched, is not enough.  The required new input is an operator/form comparison that remains positive before any zero locations are inserted.

\subsection*{A finite-core criterion without a closed-operator assumption}
For zero exclusion the Riesz realization can be bypassed. Let $\mathcal V$
be the algebraic span of finitely many genuine zero chains
$F_{\rho,j}=\zeta(s)/(s-\rho)^j$, $1\le j\le m_\rho$. Distinct partial
fractions are linearly independent, and the formula
\[
 B F_{\rho,1}=(\rho-\tfrac12)F_{\rho,1},\qquad
 B F_{\rho,j}=(\rho-\tfrac12)F_{\rho,j}+F_{\rho,j-1}\quad(j>1)
\]
defines an operator on this finite-dimensional space without a global
domain or graph-core theorem. Equip $\mathcal V$ with any fixed positive
Hilbert norm. The source functional $\eta_{\mathrm B}$ is one on each
$F_{\rho,1}$ and zero on the higher members of each canonical chain.

\begin{theorem}[Finite-core metric criterion and exact transport budget]
\label{thm:v15-finite-metric}
Let $J_n:\mathcal V\to E_n$, let $D_n=D_n^*$, and let $S_n=S_n^*$
be finite matrices. Define
\[
 Y_n=iD_nJ_n-J_nB,\qquad G_n=J_n^*S_nJ_n,
 \qquad \mathcal L_n=G_nB+B^*G_n.
\]
If $G_n\succeq c_n I$ for some $c_n>0$ and
\begin{equation}\label{eq:v15-relative-criterion}
 \frac{\|\mathcal L_n\|}{c_n}\longrightarrow0,
\end{equation}
then every eigenvalue of $B|_{\mathcal V}$ has real part zero.
In particular a space consisting of off-line chains must be zero.
No global closed-operator or quantitative evaluator assumption is used.
The exact finite identity is
\begin{equation}\label{eq:v15-budget}
 \boxed{\mathcal L_n
 =-iJ_n^*[D_n,S_n]J_n-J_n^*S_nY_n-Y_n^*S_nJ_n.}
\end{equation}
Thus a sufficient condition for \eqref{eq:v15-relative-criterion} is
\[
 \|J_n^*[D_n,S_n]J_n\|+2\|J_n^*S_nY_n\|=o(c_n).
\]
Estimating the combined signed expression in \eqref{eq:v15-budget}
can be weaker than estimating its terms separately.
\end{theorem}
\begin{proof}
For a genuine eigenvector $Bf=\mu f$, first-slot linearity gives
\[
 \ip{\mathcal L_n f}{f}=2\Re\mu\,\ip{G_nf}{f}.
\]
Consequently $2|\Re\mu|c_n\le\|\mathcal L_n\|$. Every zero chain
contains its first, genuine eigenvector, so multiplicity creates no gap
in this argument. To obtain the budget, substitute
$J_nB=iD_nJ_n-Y_n$ into both terms of $G_nB+B^*G_n$ and use
$D_n=D_n^*$ and $S_n=S_n^*$. This proves the displayed identity with
the indicated sign and adjoints. The norm bound is the triangle inequality.
\end{proof}

The criterion is invariant under multiplication of $S_n$ by a positive
scalar. A vanishing unnormalized defect obtained by collapsing the metric
does not suffice. On the corrected sampler of
\eqref{v1:eq:graphdefect}, its total source is exactly
\begin{equation}\label{eq:v15-total-source}
 Y_n=z^{\mathrm c}_{L,2K}\eta_{\mathrm B}+\mathcal R_{L,K}.
\end{equation}
Both terms must enter \eqref{eq:v15-budget}. Weak G1 concerns particular
Weil residual pairings and does not bound this budget for an arbitrary
positive $S_n$. Applying the criterion to every hypothetical off-line
zero would prove RH, but the required small signed budget remains an
independent arithmetic obligation. The old Riesz theorem is still useful
for spectral projections; its realization assumption has not been proved.

\subsection*{A positive metric that annihilates the source and shell}
There is a concrete zero-list-independent benchmark for this criterion.
Use the physical scaling $t_n=2\pi n/L$, $D=D_L$, the centered output
$H_F=\kappa F$ of \eqref{v1:eq:hardy}, and the corrected sampler
$J=\Jr_{L,K}$ on $E_{2K}$. Fix $T>0$, write $P=P_{T,L}$ for the
orthogonal projection to $|t_n|\le T$, and choose
$K\ge\lfloor TL/(2\pi)\rfloor$ so the endpoint shell lies outside
this band. Put
\[
 z=z^{\mathrm c}_{L,2K},\qquad z_T=Pz,\qquad
 w=\frac{z_T}{\|z_T\|},\qquad
 S_{T,L}=P-|w\rangle\langle w|.
\]
For all sufficiently large $L$, $z_T\ne0$. Here the rank-one map is
$x\mapsto w\ip{x}{w}$, in the fixed first-slot-linear convention.
The construction uses values of $\kappa\zeta$ on a physical band, not a
list of nontrivial zeros. It is an analytic benchmark, not a newly proved
finite-prime Weil metric or positivity transfer.

\begin{proposition}[Source-annihilating metric with fixed-core coercivity]
\label{prop:v15-projector}
The matrix $S_{T,L}$ is an orthogonal projection and satisfies exactly
\begin{equation}\label{eq:v15-annihilation}
 S_{T,L}z=0,\qquad S_{T,L}\mathcal R_{L,K}=0,
 \qquad S_{T,L}J=S_{T,L}\Jc_{L,2K}.
\end{equation}
For every fixed finite root space $\mathcal V$ as above, there is
$c_{T,\mathcal V}>0$ such that, for sufficiently large $L$,
\begin{equation}\label{eq:v15-projector-coercivity}
 \ip{S_{T,L}JF}{JF}\ge c_{T,\mathcal V}\|F\|^2
 \qquad(F\in\mathcal V).
\end{equation}
These assertions are uniform in the ambient cutoff once its shell lies
outside the fixed band. The exact remaining Lyapunov defect is
\begin{equation}\label{eq:v15-commutator-budget}
 G_LB+B^*G_L=-iJ^*[D,S_{T,L}]J,
 \qquad G_L=J^*S_{T,L}J.
\end{equation}
\end{proposition}
\begin{proof}
Since $w\in\operatorname{Ran}P$ is a unit vector, $S_{T,L}$ is the
orthogonal projection onto $\operatorname{Ran}P\cap w^\perp$.
It annihilates $z$. Both $c_{L,K}$ and $D_Lc_{L,K}$ are supported on
$K<|n|\le2K$, so it also annihilates the two summands of
$\mathcal R_{L,K}F$ and the endpoint correction itself. This proves
\eqref{eq:v15-annihilation}, without any estimate on their sizes.

For $U,V\in\mathcal V+\mathbb C\zeta$, set
\[
 [U,V]_T=\frac1{2\pi}\int_{-T}^T
 H_U(\tfrac12+it)\overline{H_V(\tfrac12+it)}\,dt.
\]
The centered factors $(-1)^n$ cancel in ordinary Gram pairings.
Riemann sums with coefficient $L^{-1/2}$ therefore give, uniformly on
the fixed finite-dimensional unit sphere,
\begin{equation}\label{eq:v15-schur-gram}
 \ip{S_{T,L}JF}{JG}\longrightarrow
 [F,G]_T-\frac{[F,\zeta]_T[\zeta,G]_T}{[\zeta,\zeta]_T}.
\end{equation}
The denominator is positive because $H_\zeta$ is not identically zero
on an interval. Equality in the Cauchy--Schwarz inequality on the right
would give $H_F=aH_\zeta$ there. Since $F=\zeta R_F$ for a proper
rational function $R_F$, analytic continuation away from isolated zeros
forces $R_F\equiv a$. Its decay at infinity gives $a=0$ and $F=0$.
Thus the limiting Gram form is positive definite. Finite-dimensional
compactness proves \eqref{eq:v15-projector-coercivity}. Finally use
\eqref{eq:v15-budget} and \eqref{eq:v15-annihilation}.
\end{proof}

This disposes of the original source and shell terms for this particular
metric, not of the whole compatibility problem. The remaining
commutator can be computed exactly.

\begin{proposition}[Variance of the remaining commutator]
\label{prop:v15-variance}
Put $m_{T,L}=\ip{Dw}{w}\in\mathbb R$, $h=(D-m_{T,L})w$, and
$\varsigma_{T,L}=\|h\|$. Then
\begin{equation}\label{eq:v15-variance}
 [D,S_{T,L}]=-|h\rangle\langle w|+|w\rangle\langle h|,
 \qquad \|[D,S_{T,L}]\|=\varsigma_{T,L}.
\end{equation}
At each fixed $T>0$,
\begin{equation}\label{eq:v15-variance-limit}
 \varsigma_{T,L}^2\longrightarrow
 \frac{\int_{-T}^T(t-m_T)^2|H_\zeta(\tfrac12+it)|^2\,dt}
      {\int_{-T}^T|H_\zeta(\tfrac12+it)|^2\,dt}>0,
 \quad
 m_T=\frac{\int_{-T}^T t|H_\zeta(\tfrac12+it)|^2\,dt}
              {\int_{-T}^T |H_\zeta(\tfrac12+it)|^2\,dt}.
\end{equation}
Moreover the self-adjoint compression $\widehat D=S_{T,L}DS_{T,L}$
on $\operatorname{Ran}S_{T,L}$, with $\widehat J=S_{T,L}J$, has
the exact new graph source
\begin{equation}\label{eq:v15-projected-graph}
 i\widehat D\widehat JF-\widehat JBF
       =-i h\,\ip{JF}{w}.
\end{equation}
\end{proposition}
\begin{proof}
$D$ commutes with $P$ and $h\perp w$. Expanding the commutator gives
\eqref{eq:v15-variance}; on the span of $w$ and $h$ its two singular
values are $\|h\|$, and it vanishes on the orthogonal complement.
The Riemann-sum limits for the first two moments give
\eqref{eq:v15-variance-limit}. The variance is strictly positive since
the nonzero continuous density is not supported at a single frequency.
For \eqref{eq:v15-projected-graph}, first use $S_{T,L}Y=0$, then
\[
 i\widehat D\widehat J-\widehat JB
 =-i S_{T,L}D(I-S_{T,L})J.
\]
The part outside $P$ is killed by $S_{T,L}D$, and the part along $w$
gives $S_{T,L}Dw=h$, proving the formula.
\end{proof}

Thus the elementary operator-norm bound for this commutator does not
tend to zero. This does not by itself decide a restricted signed pairing.
For a hypothetical eigenvector $BF=\mu F$, however, the exact restriction is
\[
 -i\ip{[D,S_{T,L}]JF}{JF}
       =2\Re\mu\,\ip{S_{T,L}JF}{JF}.
\]
It has the required nonzero size under an off-line hypothesis; it is not
an independent upper estimate contradicting that hypothesis. Positivity
and source annihilation alone have simply moved the defect into the
commutator, or equivalently into \eqref{eq:v15-projected-graph}.

\subsection*{Why removing successive source derivatives loses coercivity}
Projecting away $Dz,D^2z,\ldots$ is a natural proposed repair of the last
commutator. On a fixed physical band it encounters an explicit obstruction.

\begin{proposition}[Polynomial source annihilation suppresses zero vectors]
\label{prop:v15-krylov}
Fix $T>0$ and a hypothetical off-line zero
$\rho=1/2+\mu$, where $\mu=d+i\gamma$ and $d\ne0$.
Let $Q_{m,L}$ be contractions supported on $\operatorname{Ran}P_{T,L}$
such that
\[
 Q_{m,L}D^jz_T=0\qquad(0\le j\le2m+1).
\]
Set $M^2=d^2+(T+|\gamma|)^2$ and
$r=1-d^2/M^2\in(0,1)$. For the genuine eigenvector
$F_{\rho,1}=\zeta/(s-\rho)$ one has
\begin{equation}\label{eq:v15-krylov-bound}
 \|Q_{m,L}JF_{\rho,1}\|
 \le \frac{r^{m+1}}{|d|}\,\|z_T\|.
\end{equation}
Consequently bounded projection metrics obtained by removing an increasing
number of these source derivatives cannot retain a uniform positive lower
bound on this vector when $m\to\infty$ at fixed $T$.
\end{proposition}
\begin{proof}
On $[-T,T]$ define $\xi(t)=1-(d^2+(t-\gamma)^2)/M^2\in[0,r]$.
The polynomial
\[
 p_m(t)=\frac{-d-i(t-\gamma)}{M^2}
                   \sum_{j=0}^m\xi(t)^j
\]
has degree $2m+1$. A finite geometric sum gives the exact remainder
\[
 \frac1{it-\mu}-p_m(t)
       =\frac{\xi(t)^{m+1}}{it-\mu},\qquad
 \sup_{|t|\le T}\left|\frac1{it-\mu}-p_m(t)\right|
       \le |d|^{-1}r^{m+1}.
\]
The centered source and root samples have the same half-period phase, so
$P_{T,L}JF_{\rho,1}=(iD-\mu)^{-1}z_T$ exactly; the high shell
does not meet the band. Applying $Q_{m,L}$ removes $p_m(D)z_T$.
The diagonal remainder estimate proves \eqref{eq:v15-krylov-bound}.
Finally $\|z_T\|$ stays bounded by its Riemann-sum limit.
\end{proof}

This is a limited obstruction to bounded polynomial-source projection
metrics, not to all positive metrics, changing bands, or signed arithmetic
estimates. Rescaling a collapsing metric must be assessed by the ratio in
\eqref{eq:v15-relative-criterion}; it cannot be judged by an unnormalized
commutator alone. Section 19 therefore now has a domain-free finite-core
criterion and a coercive positive benchmark with exact source removal,
but no independent small Lyapunov defect. G2 and RH remain open.


\subsection*{Why a sharp band projection does not inherit strip transfer}
The ordinary Gram convergence in Proposition~\ref{prop:v15-projector}
does not authorize applying the holomorphic Weil-transfer theorem after
the physical band projection. The following calculation identifies the
new cut terms on exactly that finite object.

\begin{lemma}[Transform of a sharp physical band]\label{lem:v16-band-transform}
Fix $T>0$, set $\Delta=2\pi/L$, $N=\lfloor T/\Delta\rfloor$ and
$T_L=N\Delta$. Let $a_L$ be functions with uniformly bounded $C^2$
norm on $[-T,T]$. In centered physical coordinates define
\[
 f_L(y)=\frac1L\sum_{n=-N}^N a_L(t_n)e^{it_ny},\quad |y|\le L/2,
 \qquad t_n=n\Delta,
\]
and extend by zero. For fixed $v\in\mathbb C$ with $\Re v\ne0$, put
$\mathcal H_L(v)=\int f_L(y)e^{-vy}\,dy$. Then exactly
\begin{equation}\label{eq:v16-exact-cut}
 \mathcal H_L(v)=-\frac{2\sinh(vL/2)}L
       \sum_{n=-N}^N\frac{(-1)^n a_L(t_n)}{it_n-v}.
\end{equation}
Moreover, uniformly when $v$ ranges over a compact subset of
$\{\Re v\ne0\}$,
\begin{equation}\label{eq:v16-cut-leading}
 \mathcal H_L(v)=-\frac{(-1)^N\sinh(vL/2)}L
 \left\{\frac{a_L(T_L)}{iT_L-v}
       +\frac{a_L(-T_L)}{-iT_L-v}+O(L^{-1})\right\}.
\end{equation}
If $a_L=a+O(L^{-1})$ in $C^2$, define
\[
 C_a(v)=\frac{a(T)}{iT-v}+\frac{a(-T)}{-iT-v}.
\]
Whenever $C_a(v)\ne0$,
\begin{equation}\label{eq:v16-cut-growth}
 |\mathcal H_L(v)|\sim\frac{|C_a(v)|}{2L}
                  e^{|\Re v|L/2}.
\end{equation}
This is compatible with a bounded ordinary norm, since
$\|f_L\|_2^2=L^{-1}\sum_{|n|\le N}|a_L(t_n)|^2=O(1)$.
\end{lemma}
\begin{proof}
Integrate each exponential over $[-L/2,L/2]$; both endpoint phases
are $(-1)^n$, which proves \eqref{eq:v16-exact-cut} with its sign
and factor two. Set $q(t)=a_L(t)/(it-v)$ and
$q_j=q(-T_L+j\Delta)$. An exact telescoping identity is
\[
 2\sum_{j=0}^{2N}(-1)^j q_j-q_0-q_{2N}
 =\sum_{k=0}^{N-1}(q_{2k}-2q_{2k+1}+q_{2k+2}).
\]
The right side has absolute value at most
$N\Delta^2\|q''\|_\infty=T_L\Delta\|q''\|_\infty$.
The denominators are uniformly separated from zero on the stated compact
sets of $v$, giving \eqref{eq:v16-cut-leading}. Now
$T_L=T+O(L^{-1})$ and the assumed $C^2$ convergence replace the
braces by $C_a(v)+O(L^{-1})$. Finally
$|\sinh(vL/2)|\sim\tfrac12 e^{|\Re v|L/2}$. Parseval on the
finite interval proves the ordinary norm assertion.
\end{proof}

\begin{proposition}[Source removal does not remove the sharp-cut growth]
\label{prop:v16-projected-cut}
Use the corrected sampler and the projection $S_{T,L}$ of
Proposition~\ref{prop:v15-projector}, with its endpoint shell outside
the fixed band. Suppose $\rho=\tfrac12+\mu$ is a hypothetical off-line
zero, take its genuine first eigenvector $F=\zeta/(s-\rho)$, and choose
$T>0$ such that $H_\zeta(\tfrac12\pm iT)\ne0$. Such bands exist
because the excluded endpoint zeros are isolated. Put
\[
 Z(t)=H_\zeta(\tfrac12+it),\qquad R(t)=\frac1{it-\mu},\qquad
 \alpha_T=\frac{\int_{-T}^T |Z(t)|^2R(t)\,dt}
                 {\int_{-T}^T |Z(t)|^2\,dt}.
\]
The centered physical representative of $S_{T,L}JF$ has the coefficients
of Lemma~\ref{lem:v16-band-transform} with
\[
 a_L(t)=Z(t)\{R(t)-\alpha_L\},\qquad
 \alpha_L=\frac{\sum_{|n|\le N}|Z(t_n)|^2R(t_n)}
                   {\sum_{|n|\le N}|Z(t_n)|^2}
 =\alpha_T+O(L^{-1}).
\]
Let $a(t)=Z(t)\{R(t)-\alpha_T\}$. For every fixed real $b>0$,
at least one of $C_a(b)$ and $C_a(-b)$ is nonzero. Thus at least one
of the two transforms at $s=\tfrac12\pm b$ grows as a nonzero
constant times $e^{bL/2}/L$. In particular these source-projected
vectors have no bounded full-strip transform limit, even on a compact
substrip containing those two points when $0<b<1/2$.
\end{proposition}
\begin{proof}
The projection coefficient is the first-slot-linear ratio
$\ip{PJF}{z_T}/\|z_T\|^2$, giving the displayed $\alpha_L$; the
half-period factors cancel in this ratio. The shell is absent from the
band. Since $\Re\mu\ne0$, $R$ is smooth there. Ordinary Riemann-sum
estimates give the $O(L^{-1})$ rate, including the positive limiting
denominator. Consequently $a_L=a+O(L^{-1})$ in $C^2$.

Write $A=a(T)$ and $B=a(-T)$. The two equations
$C_a(b)=C_a(-b)=0$ form a nonsingular two-by-two system in $A,B$:
its determinant is
\[
 -\frac{4ibT}{(T^2+b^2)^2}\ne0.
\]
They would force $A=B=0$. Since $Z(\pm T)\ne0$, this would imply
$R(T)=\alpha_T=R(-T)$, which is impossible for $T>0$.
Apply \eqref{eq:v16-cut-growth} to a nonzero coefficient.
\end{proof}

The functional-equation partner does not, by itself, cancel this edge
mechanism. The next statement concerns one reflected pair of summands in
the zero-side formula, not the complete Weil form.

\begin{proposition}[Reflected-pair edge contribution]
\label{prop:v17-pair-edge}
In Lemma~\ref{lem:v16-band-transform}, assume
$a_L=a+O(L^{-1})$ in $C^2[-T,T]$. Fix
$v=d+i\gamma$ with $d>0$, put $v^\sharp=-\overline v$, and suppose
the zeros $\tfrac12+v$ and $\tfrac12+v^\sharp$ have common
multiplicity $m$. The contribution of this reflected pair to the
value-level zero formula is
\[
 \mathcal P_{v,L}=m\left\{
 \mathcal H_L(v)\overline{\mathcal H_L(v^\sharp)}
 +\mathcal H_L(v^\sharp)\overline{\mathcal H_L(v)}\right\}.
\]
Writing
\[
 A_{a,T}(v)=C_a(v)\overline{C_a(v^\sharp)},
\]
one has
\begin{equation}\label{eq:v17-pair-edge}
 \boxed{
 \mathcal P_{v,L}
 =-\frac{m e^{dL}}{2L^2}
   \Re\!\left(e^{i\gamma L}A_{a,T}(v)\right)
   +O\!\left(\frac{e^{dL}}{L^3}\right).}
\end{equation}
Thus reflection converts the two growing transforms into a real
oscillatory term; it does not cancel them identically. If
$A_{a,T}(v)\ne0$ and $\gamma\ne0$, this pair contribution has
subsequences of both signs and magnitude comparable to $e^{dL}/L^2$.
For the source-projected root vector of
Proposition~\ref{prop:v16-projected-cut}, the coefficient is the explicit
product formed from $a(t)=Z(t)\{(it-\mu)^{-1}-\alpha_T\}$.
No assertion is made when one of these two edge coefficients vanishes.
\end{proposition}
\begin{proof}
Lemma~\ref{lem:v16-band-transform} gives, with
$\epsilon_L=(-1)^N$,
\[
 \mathcal H_L(v)=-\frac{\epsilon_L\sinh(vL/2)}L
                   \{C_a(v)+O(L^{-1})\}.
\]
Since $v^\sharp=-\overline v$,
\[
 \overline{\mathcal H_L(v^\sharp)}
 =\frac{\epsilon_L\sinh(vL/2)}L
       \{\overline{C_a(v^\sharp)}+O(L^{-1})\}.
\]
Their product therefore equals
\[
 -\frac{\sinh^2(vL/2)}{L^2}
  \{A_{a,T}(v)+O(L^{-1})\}.
\]
The other reflected summand is its complex conjugate. Finally
$\sinh^2(vL/2)=\tfrac14e^{vL}\{1+O(e^{-dL})\}$, which proves
\eqref{eq:v17-pair-edge}. Choosing the phase of $\gamma L$ near
$-\arg A_{a,T}(v)$ and near $\pi-\arg A_{a,T}(v)$ gives the two
subsequences.
\end{proof}

This does not prove divergence or negativity of the complete Weil form:
other zero summands may cancel the displayed contribution. The combined
prime, pole and archimedean expression is an equivalent representation
of that same form, not an additional contribution. Nor does this rule out a new
signed product estimate just on the zero set. It proves that the existing
bounded holomorphic-strip transfer cannot simply be inherited by these
sharply projected vectors. There is no conflict with weak G1, which
estimates a residual against band tests without asserting that those tests
retain the original Hardy output's strip bounds.

For a proposed correction $M=S_{T,L}(W+a_0I)S_{T,L}$, source and shell
annihilation still hold, but all three terms in
\[
 [D,M]=[D,S_{T,L}](W+a_0I)S_{T,L}
       +S_{T,L}[D,W]S_{T,L}
       +S_{T,L}(W+a_0I)[D,S_{T,L}]
\]
remain in the metric budget. Neither positivity of this compression nor
a small signed budget is proved. A valid next step must control its
actual finite prime, pole and archimedean terms together, or introduce
a different filter with independently justified product estimates.

\subsection*{Full-space projection with a proved strip-product bound}
The preceding obstruction concerns a fixed physical band. It can be
avoided without discarding the graph defect. Keep the original ambient
space $E_{2K}$, the physical $D$, the source $z=\Jc_{L,2K}\zeta$,
the shell $c=c_{L,K}$ and the corrected $J=\Jr_{L,K}$. Define
\[
 e_0=c/\|c\|,\qquad e_1=Dc/\|Dc\|,\qquad
 P_{\rm sh}=|e_0\rangle\langle e_0|+|e_1\rangle\langle e_1|,
 \qquad Q_0=I-P_{\rm sh}.
\]
Parity makes $e_0,e_1$ orthonormal. Put
\begin{equation}\label{eq:v18-full-projection}
 u=Q_0z,\qquad
 S_{L,K}=Q_0-\frac{|u\rangle\langle u|}{\|u\|^2},\qquad
 \alpha_{L,K}(F)=\frac{\ip{Q_0\Jc_{L,2K}F}{u}}{\|u\|^2}.
\end{equation}
The nonzero denominator for large $L$ is proved below. This construction
uses analytic samples of $\kappa\zeta$ and the two known shell vectors;
it does not use a zero list. It changes the metric, not $J$, its physical
endpoint, or the finite Weil matrix.

\begin{proposition}[Full-space source removal and Weil transfer]
\label{prop:v18-full-projection}
Let $\mathcal V$ be a fixed finite span of genuine root chains,
on either side of or on the critical line. Choose the fixed Hardy
smoothing so that, on $\mathcal U=\mathcal V+\mathbb C\zeta$,
\[
 |H_U(\tfrac12+it)|\le C\|U\|(1+|t|)^{-R},\quad R\ge3,
 \qquad |g_U^{(j)}(y)|\le C\|U\|e^{-p|y|},\quad j\le2,
\]
with $p>1/2$, as in \eqref{v1:eq:tails}. Let
$K\ge\lceil e^{2L}L^4\rceil$. Then $S=S_{L,K}$ is an orthogonal
projection and, exactly,
\begin{equation}\label{eq:v18-full-annihilation}
 Sz=Sc=SDc=0,\qquad S\mathcal R=0,\qquad
 SJF=Q_0\Jc_{L,2K}(F-\alpha_{L,K}(F)\zeta).
\end{equation}
If $\mathcal H_{L,K,F}$ is the transform of $SJF$ translated back
to $[-L/2,L/2]$ and zero extended, then throughout $0\le\Re s\le1$,
\begin{equation}\label{eq:v18-full-strip}
 \left|\mathcal H_{L,K,F}(s)-H_F(s)
             +\alpha_{L,K}(F)H_\zeta(s)\right|
 \le \frac{C\varepsilon_{L,K}\|F\|}{1+|\Im s|},
\end{equation}
where, uniformly over these cutoffs,
\[
 \varepsilon_{L,K}
 =e^{-(p-1/2)L/2}
       +Le^{L/4}(1+K/L)^{2-R}\longrightarrow0.
\]
Write $[U,V]_\infty=(2\pi)^{-1}\int_{\mathbb R}
H_U(1/2+it)\overline{H_V(1/2+it)}\,dt$. Uniformly on bounded
subsets of $\mathcal V$,
\begin{align}
 \ip{SJF}{SJG}&\longrightarrow
 [F,G]_\infty-
 \frac{[F,\zeta]_\infty[\zeta,G]_\infty}{[\zeta,\zeta]_\infty},
       \label{eq:v18-full-gram}\\
 \ip{WSJF}{SJG}&\longrightarrow\Q(g_F,g_G).
       \label{eq:v18-full-weil}
\end{align}
The first limiting form is positive definite on $\mathcal V$.
The second limit has error $O(\varepsilon_{L,K}+\varepsilon_{L,K}^2)$.
\end{proposition}
\begin{proof}
Let $X=1+K/L$. A centered sample vector $a=\Jc_{L,2K}U$ obeys
\[
 \|1_{K<|n|\le2K}a\|_2\le C\|U\|X^{1/2-R},\qquad
 \sum_n(1+|t_n|)|(P_{\rm sh}a)_n|
       \le C\sqrt L\|U\|X^{2-R}.
\]
For the second inequality use Cauchy--Schwarz on the $2K$ shell
coordinates, $\|P_{\rm sh}a\|\le\|1_{\rm sh}a\|$ and
$1+|t_n|\le CX$ there. The omitted coefficients of the infinite
centered Fourier series satisfy the same weighted sum bound, since
$R>2$. Riemann sums on the line and these tail bounds give
\[
 \|u\|^2\longrightarrow[\zeta,\zeta]_\infty>0,\qquad
 \alpha_{L,K}(F)\longrightarrow
       [F,\zeta]_\infty/[\zeta,\zeta]_\infty.
\]
In particular $\alpha_{L,K}$ is uniformly bounded on the fixed space.
The projection and annihilation identities now follow directly from
\eqref{eq:v18-full-projection}; $\mathcal R F$ lies in
$\operatorname{span}\{c,Dc\}$ exactly. Since $Sc=0$, the endpoint
correction in $J$ is annihilated, which proves the last identity in
\eqref{eq:v18-full-annihilation}.

Here is the strip estimate, including the physical endpoint jumps.
Set $U=F-\alpha_{L,K}(F)\zeta$. Compare the centered representative
of $Q_0\Jc_{L,2K}U$ with $1_{[-L/2,L/2]}P_Lg_U$.
Their difference $b$ consists of the omitted Fourier tail and the
negative shell projection. The weighted coefficient bounds above,
with the physical basis factor $L^{-1/2}$, give on this interval
\[
 \|b\|_\infty+\|b'\|_\infty
       \le C\|F\|X^{2-R}.
\]
For $|d|\le1/2$, zero extension therefore has
\[
 \|e^{-d\,\cdot}b\|_1+
 \operatorname{Var}(e^{-d\,\cdot}b)
       \le CL e^{L/4}\|F\|X^{2-R}.
\]
The variation includes both cut atoms. The $L^1$ estimate at bounded
ordinates and one integration by parts against this finite derivative
measure at large ordinates give the second term of
\eqref{eq:v18-full-strip}. Lemma~\ref{v1:lem:strip} supplies the
first term. At $R=3$ the second term is at most
$C e^{-7L/4}L^{-2}$ for the stated cutoffs. No fixed-band edge
asymptotic is used.

Ordinary Gram convergence and the projection formula prove
\eqref{eq:v18-full-gram}. Equality in its Cauchy--Schwarz bound
would give $H_F=aH_\zeta$ almost everywhere on the line.
Analyticity then forces the proper rational function $F/\zeta$
to be constant, hence zero. This proves positive definiteness and
thus a uniform lower bound on the fixed finite-dimensional unit sphere.

For the Weil limit, $H_\zeta(\rho)=0$ at every nontrivial zero;
the pole at $1$ has already been canceled by $\kappa$.
Thus subtracting $\alpha_{L,K}(F)H_\zeta$ changes none of the
global zero-form values. Equation~\eqref{eq:v18-full-strip} and
the bounded holomorphic transforms bound each product error by
\[
 C\|F\|\|G\|
 (\varepsilon_{L,K}+\varepsilon_{L,K}^2)
                       (1+|\Im\rho|)^{-2}.
\]
This is summable over all zeros, without RH. Apply
\eqref{v1:eq:zeros} and translate back to $[0,L]$; reflected
translation factors cancel in each product. This proves
\eqref{eq:v18-full-weil} on the actual matrix $W_{L,2K}$.
\end{proof}

The projection above resolves a transfer issue for a different metric;
it does not repair the fixed-band vectors of
Propositions~\ref{prop:v16-projected-cut}--\ref{prop:v17-pair-edge}.
Nor does it make the original graph defect small: it annihilates that
defect inside a particular metric. Its remaining signed budget is exactly
\begin{equation}\label{eq:v18-full-budget}
 G_LB+B^*G_L=-iJ^*[D,S_{L,K}]J,\qquad G_L=J^*S_{L,K}J.
\end{equation}
In particular $Q_0$ must not be commuted through $D$; removing the
two shell directions does not make their orthogonal complement invariant.

\begin{corollary}[A fixed scalar shift does not close the repaired metric]
\label{cor:v18-shift}
Let $\mathcal V$ be a nonempty hypothetical one-sided off-line root
space and fix $a_0>0$. With $S=S_{L,K}$ as above, put
\[
 M_L=S(W+a_0I)S,\quad A_L=J^*SWSJ,\quad
 G_L=J^*SJ,\quad H_L=J^*M_LJ.
\]
Then $A_L\to0$, $G_L\to G_\infty\succ0$ and
$H_L\to a_0G_\infty\succ0$, but
\begin{equation}\label{eq:v18-shift-budget}
 H_LB+B^*H_L\longrightarrow
                  a_0(G_\infty B+B^*G_\infty)\ne0.
\end{equation}
Consequently this candidate does not satisfy the relative defect
criterion of Theorem~\ref{thm:v15-finite-metric}.
\end{corollary}
\begin{proof}
The global Weil form vanishes on a one-sided root space, so
\eqref{eq:v18-full-weil} gives $A_L\to0$ in its finite matrix norm.
Since $B$ is fixed, $A_LB+B^*A_L\to0$ as well. The identity
$H_L=A_L+a_0G_L$ proves the limits. On a first root vector $f$,
$Bf=\mu f$ with $\Re\mu\ne0$, the limiting signed pairing is
$2a_0\Re\mu\,\ip{G_\infty f}{f}\ne0$. Its ratio to the
limiting metric pairing is $2\Re\mu$. The pullback metrics are
bounded above as well as below, so their positive lower bounds cannot
turn this nonzero limit into a vanishing relative defect.
\end{proof}

This is restricted coercivity, not positivity of $M_L$ on the entire
ambient space. The exact budget for $M_L$, whose source terms also
vanish, still contains all three terms
\[
 [D,M_L]=[D,S](W+a_0I)S+S[D,W]S+S(W+a_0I)[D,S].
\]
The new transfer theorem controls their combined pulled-back value via
\eqref{eq:v18-shift-budget}; it supplies no independent small signed
estimate. A successful G2 metric must do more than preserve the
one-sided zero Weil limit and add a fixed ordinary positive shift.

\subsection*{Norm-level Weil action and the squared-metric candidate}
Squaring the finite Weil matrix supplies positivity without a scalar
shift. Its behavior cannot be inferred from form convergence alone.
We first prove the needed norm estimate. The classical unconditional
zero count gives, with multiplicities,
\begin{equation}\label{eq:v19-local-count}
 \#\{\rho:j\le\Im\rho<j+1\}\le C\log(2+|j|),\qquad j\in\mathbb Z.
\end{equation}
For example, this follows by subtracting the estimates for $N(T)$ at
the two endpoints in Trudgian~\cite[Corollary 1]{Trudgian2012}; finitely
many small ordinates are absorbed into $C$. Neither simplicity nor RH
is used.

\begin{lemma}[Local synthesis from decaying zero coefficients]
\label{lem:v19-synthesis}
Suppose $|a_\rho|\le\epsilon/(1+|\Im\rho|)$ on the multiset of
nontrivial zeros. Then, for $A\ge1$, the series
\[
 r(y)=\sum_{\rho\in\Z}a_\rho e^{(\rho-1/2)y}
\]
converges in $L^2(-A,A)$, and
\begin{equation}\label{eq:v19-synthesis}
 \|r\|_{L^2(-A,A)}\le C\epsilon\sqrt{A+1}\,e^{A/2}.
\end{equation}
The constant is independent of $A$, $\epsilon$ and the coefficients.
\end{lemma}
\begin{proof}
Choose a nonnegative smooth function $\chi_A$ equal to one on
$[-A,A]$, zero outside $[-A-1,A+1]$, with its first two derivatives
bounded independently of $A$. Write $\rho=1/2+d_\rho+i\gamma_\rho$.
Since $|d_\rho|<1/2$, two integrations by parts, together with the
direct integral bound, give
\[
 \left|\int\chi_A(y)e^{(d_\rho+d_\tau)y}
                  e^{i(\gamma_\rho-\gamma_\tau)y}\,dy\right|
 \le\frac{C(A+1)e^A}{1+|\gamma_\rho-\gamma_\tau|^2}.
\]
For a finite partial sum put
$b_j=\sum_{j\le\gamma_\rho<j+1}|a_\rho|$. Grouping both
indices by these unit intervals bounds its squared norm by
\[
 C(A+1)e^A\sum_{j,k}\frac{b_jb_k}{1+|j-k|^2}
 \le C(A+1)e^A\sum_j b_j^2.
\]
The last inequality is the $\ell^1*\ell^2$ convolution bound.
Equation~\eqref{eq:v19-local-count} gives
$b_j\le C\epsilon\log(2+|j|)/(1+|j|)$, whose square is summable.
The same bound on omitted indices makes the partial sums Cauchy in
$L^2(-A,A)$ and proves \eqref{eq:v19-synthesis}. It also permits
arbitrary finite exhaustions of the zero multiset.
\end{proof}

Let $\mathcal T_L$ denote the unitary translation from the physical
interval $[0,L]$ to $[-L/2,L/2]$; on coefficient vectors it represents
$x$ as $L^{-1/2}\sum x_n(-1)^ne^{it_ny}$. Let $\Pi_{L,2K}$
be the ordinary orthogonal projection onto the corresponding centered
Fourier space. These operations do not change the finite matrix $W$.

\begin{proposition}[Weil-action transfer on a fixed root family]
\label{prop:v19-action}
Use the objects and cutoffs of Proposition~\ref{prop:v18-full-projection},
but choose the fixed smoothing with tail rate $p>1$. Define the finite sum
\[
 q_F(y)=\sum_{\rho\in\Z}H_F(\rho)e^{(\rho-1/2)y},
\]
where zeros are counted with multiplicity; only the finitely many
selected root locations can contribute. Then uniformly on the fixed
root family,
\begin{equation}\label{eq:v19-action}
 \|\mathcal T_LWSJF-\Pi_{L,2K}q_F\|_2
 \le C\Delta_{L,K}\|F\|,\qquad
 \Delta_{L,K}=\sqrt{L+1}\,e^{L/4}\varepsilon_{L,K}\longrightarrow0.
\end{equation}
In particular the first part of $\Delta_{L,K}$ is
$\sqrt{L+1}\,e^{-(p-1)L/2}$. The original condition $p>1/2$
alone does not make this bound tend to zero.
\end{proposition}
\begin{proof}
Let $x=SJF$ and let $\mathcal H_x$ be the transform of
$\mathcal T_Lx$. At a zero, Proposition~\ref{prop:v18-full-projection}
and $H_\zeta(\rho)=0$ give
\[
 |\mathcal H_x(\rho)-H_F(\rho)|
       \le C\varepsilon_{L,K}\|F\|/(1+|\Im\rho|).
\]
Lemma~\ref{lem:v19-synthesis} synthesizes these errors into a function
$r_F$ of norm at most the right side of \eqref{eq:v19-action}.
For each centered finite Fourier test $v$, the first-slot-linear
zero formula gives exactly
\[
 \Q(\mathcal T_Lx,v)=\sum_\rho\mathcal H_x(\rho)
          \overline{H_v(\rho^\sharp)}
       =\ip{q_F+r_F}{v}.
\]
Indeed $\overline{\rho^\sharp}-1/2=-(\rho-1/2)$, so the
exponential in the first slot is $e^{(\rho-1/2)y}$ with a positive
sign. The product series is absolutely convergent for these tests;
the synthesized series also converges in $L^2$. Translation invariance
of the paired form identifies the left side with
$\ip{\mathcal T_LWx}{v}$. Thus
$\mathcal T_LWx=\Pi_{L,2K}(q_F+r_F)$, and projection contractivity
proves the bound. At $R=3$ the cutoff contribution to $\Delta_{L,K}$
is $O(e^{-3L/2}L^{-3/2})$. The other contribution tends to zero
because $p>1$.
\end{proof}

The stronger tail margin is available, for example with fixed
$\sigma=2$, $a=4$ and sufficient $M$, taking $1<p\le3/2$ in
\eqref{v1:eq:tails}. Critical-line root chains are allowed here and
in Proposition~\ref{prop:v18-full-projection}: the zeros remove their
apparent transform poles, and the proper-rational independence proof
works away from the isolated zero locations. This extension assumes
genuine root functions, not inserted meromorphic poles.

\begin{corollary}[Mass of the squared-Weil metric]
\label{cor:v19-square}
On the same finite space set $M_L^{\square}=S W^2 S\succeq0$.
For a zero $\rho=1/2+\mu$ of multiplicity $m$, write
$F_{\rho,j}=\zeta/(s-\rho)^j$, $1\le j\le m$, and
\[
 C_\rho=m\kappa(\rho)\frac{\zeta^{(m)}(\rho)}{m!}\ne0.
\]
If $d=\Re\mu\ne0$, then
\begin{equation}\label{eq:v19-square-top}
 \ip{M_L^{\square}JF_{\rho,m}}{JF_{\rho,m}}
 =\|WSJF_{\rho,m}\|^2
 \sim |C_\rho|^2\frac{\sinh(|d|L)}{|d|}.
\end{equation}
For a critical-line zero the corresponding asymptotic is
$|C_\rho|^2L$. For every lower jet $j<m$,
\begin{equation}\label{eq:v19-square-lower}
 \ip{M_L^{\square}JF_{\rho,j}}{JF_{\rho,j}}
       =O(\Delta_{L,K}^2)\longrightarrow0.
\end{equation}
These assertions concern the displayed normalization; a rescaled metric
must still satisfy the relative criterion, not merely an absolute bound.
\end{corollary}
\begin{proof}
Self-adjointness of $W$ gives positivity and the squared norm identity.
For the top root, $q_F=C_\rho e^{\mu y}$; for lower jets $q_F=0$,
because all their zero values vanish. Moreover
\[
 \|e^{\mu y}\|_{L^2(-L/2,L/2)}^2
 =\begin{cases}\sinh(|d|L)/|d|,&d\ne0,\\L,&d=0.\end{cases}
\]
Its centered Fourier coefficient at index $n$ is
$2(-1)^n\sinh(\mu L/2)/(\sqrt L(\mu-it_n))$, with the removable
value used when $\mu=it_n$. For $|n|>2K$ this gives
\[
 \|(I-\Pi_{L,2K})e^{\mu\,\cdot}\|_2
       \le C_\mu e^{|d|L/2}\sqrt{L/K}\longrightarrow0.
\]
The cutoff dominates this exponential since $|d|<1/2$.
Apply Proposition~\ref{prop:v19-action} to obtain both asymptotics
and the lower-jet bound.
\end{proof}

This candidate is globally positive and annihilates the exact total
source, since $S(z\eta_{\mathrm B}+\mathcal R)=0$. Nevertheless
its full commutator is
\[
 [D,M_L^{\square}]=[D,S]W^2S
     +S\bigl([D,W]W+W[D,W]\bigr)S+SW^2[D,S].
\]
All terms belong to the signed G2 budget. At a hypothetical simple
off-line root $Bf=\mu f$, putting $G_L^{\square}=J^*M_L^{\square}J$
gives the exact quotient
\[
 \frac{\ip{(G_L^{\square}B+B^*G_L^{\square})f}{f}}
      {\ip{G_L^{\square}f}{f}}=2\Re\mu
\]
whenever its denominator is nonzero. At a multiple zero, the lower
jets additionally lose an unscaled coercive lower bound by
\eqref{eq:v19-square-lower}. Squaring has not supplied a small
relative defect. The norm-action theorem is a stronger transfer result,
not a proof of the independent signed arithmetic estimate or RH.

\section{Finite Weil residuals: compact-source radicality and G1 transfer}\label{sec:g1}
The finite-prime comparison can be narrowed further.  The following subsection isolates the part of the Connes--Consani--Moscovici comparison that is controlled by classical prolate asymptotics and the prime number theorem, separately from the still-open spectral-ordering problem.

\subsection*{A rank-two symmetry defect for an arbitrary even candidate}
Let $W_{\lambda,N}$ denote the finite Weil matrix, let $D_N$ be the Fourier-index matrix, and identify the equal-coefficient boundary covector with its Riesz vector $\eta_N$, writing $\eta_N(x)=\langle x,\eta_N\rangle$ in the first-slot-linear convention.  The exact CCM commutator has the form
\begin{equation}\label{eq:ccm-commutator}
 D_NW_{\lambda,N}-W_{\lambda,N}D_N
 =|\beta_N\rangle\langle\eta_N|-|\eta_N\rangle\langle\beta_N|,
\end{equation}
where $\beta_N$ is odd and $\eta_N$ is even.  Let $k_N$ be any real even vector satisfying $\eta_N(k_N)=1$, and define
\[
 \mathcal D_{k,N}=D_N-|D_Nk_N\rangle\langle\eta_N|.
\]
Then parity gives $\langle k_N,\beta_N\rangle=0$, and applying \eqref{eq:ccm-commutator} to $k_N$ gives
\[
 W_{\lambda,N}D_Nk_N=D_NW_{\lambda,N}k_N-\beta_N.
\]
Substitution yields the exact identity
\begin{equation}\label{eq:ranktwo-weil-defect}
 \boxed{
 W_{\lambda,N}\mathcal D_{k,N}
 -\mathcal D_{k,N}^{*}W_{\lambda,N}
 =-|D_NW_{\lambda,N}k_N\rangle\langle\eta_N|
  +|\eta_N\rangle\langle D_NW_{\lambda,N}k_N|.}
\end{equation}
Thus no ground-state vector is needed to identify the metric-symmetry defect: for an explicit even candidate it is rank at most two and is controlled by one linear residual.

It is useful to state the scale invariant version with the physical logarithmic generator $D_{\log}=(2\pi/L)D_N$, $L=2\log\lambda$, and the actual boundary-evaluation covector $\delta_N=L^{-1/2}\eta_N$.  On the fixed physical test band $|D_{\log}|\le T$, the restriction of the boundary-evaluation covector $\delta_N$ has norm bounded by a constant depending only on $T$, independently of how large the ambient section $E_N$ is.  Thus one may later choose a much larger diagonal section $N=N(\lambda)$ to resolve the prolate boundary layer without changing the fixed-window estimate.  Hence the relevant projective quantity is
\begin{equation}\label{eq:physical-projective-residual}
 \frac{\|P_TD_{\log}W_{\lambda,N}k_N\|}
 {|\delta_N(k_N)|},
\end{equation}
where $P_T$ is projection to the physical-frequency window $[-T,T]$.  This formulation avoids the factor-of-$\sqrt L$ ambiguity that results from simultaneously rescaling the covector and the candidate.

\subsection*{The compact-source radical passage}
The classical radical theorem applies to the even Schwartz class with
$h(0)=\widehat h(0)=0$~\cite[Section 3]{ConnesConsani2023}.
The following extension supplies the domain passage needed for the compact
prolate source. It is a statement about explicit-formula pairings and the
closed \emph{semilocal} form, not a claim of a positive global closed form.

\begin{theorem}[Radicality with a nonzero compact-source endpoint]
\label{thm:v14-radical}
Fix $b>0$. Let $h$ be even and smooth on $[-b,b]$, extended by zero
outside that interval, and suppose
\[
 h(0)=0,\qquad \int_0^b h(u)\,du=0.
\]
No vanishing of the interior trace $h(b^-)$ is required. In logarithmic
coordinates put
\[
 k(x)=\mathcal E(h)(e^x)=e^{x/2}\sum_{n\ge1}h(ne^x).
\]
For every $v$ obtained by zero-extending an $H^1$ function on a compact
interval $J$, all terms of the polarized explicit formula $QW(k,v)$ are
finite and $QW(k,v)=0$. In particular, for a compact interval $I$
containing the support of $v$, set $k_{\rm in}=1_I k$ and
$r=k-k_{\rm in}$. Then $k_{\rm in}$ and $v$ belong to the closed
semilocal form domain and
\[
 QW_I(k_{\rm in},v)=-QW(r,v).
\]
The assertion holds at each fixed $b$; its approximation constants need
not be uniform as $b$ varies.
\end{theorem}
\begin{proof}
\emph{A primitive gains one power.}
Write $\beta(t)=\tfrac12-\{t\}$ away from integers. Euler summation,
with $N=\lfloor b/u\rfloor$, gives for almost every sufficiently small $u$
\[
 \sum_{n=1}^N h(nu)
 =u^{-1}\int_0^{Nu}h(t)\,dt+
       \frac{h(Nu)-h(0)}2+O_h(u)
 =h(b^-)\beta(b/u)+O_h(u).
\]
For example, the composite trapezoid remainder is bounded by
$Cbu\|h''\|_\infty$ in the first equality; Taylor expansion of the
missing interval $[Nu,b]$ gives the second. Values at the arithmetic
cuts affect neither an integral nor the correlations below. Consequently
\begin{equation}\label{eq:v14-euler-tail}
 \mathcal E(h)(u)
 =h(b^-)u^{1/2}\beta(b/u)+O_h(u^{3/2}).
\end{equation}
In particular $k\in L^2(\mathbb R)$, it vanishes for $x>\log b$, and
it is of bounded variation on every compact logarithmic interval.
Let $K(x)=\int_{-\infty}^x k(t)\,dt$. Since $\beta$ has mean zero
and a bounded primitive,
\[
 \int_0^u v^{-1/2}\beta(b/v)\,dv
 =\sqrt b\int_{b/u}^{\infty}\beta(w)w^{-3/2}\,dw
 =O_b(u^{3/2}).
\]
The last bound follows by one integration by parts against that bounded
primitive. Integrating the remainder in \eqref{eq:v14-euler-tail} yields
\begin{equation}\label{eq:v14-primitive}
 K(x)=O_h(e^{3x/2})\qquad(x\longrightarrow-\infty).
\end{equation}
Both moments matter: without $h(0)=0$ an additional term
$-\tfrac12h(0)u^{1/2}$ generally prevents this gain.

\emph{Absolute convergence of the geometric side.}
The zero extension $v_0$ of $v$ has bounded variation, including its two
endpoint jumps. In the first-slot-linear convention its correlation is
\[
 C_{k,v}(y)=\int k(x)\overline{v_0(x-y)}\,dx
          =-\int K(x)\,d\overline{v_0(x-y)}.
\]
Thus \eqref{eq:v14-primitive} gives
$C_{k,v}(y)=O_{h,v}(e^{3y/2})$ as $y\to-\infty$; it is zero for
all sufficiently large positive $y$. Its reflected correlation has the
opposite support orientation. The even combinations therefore decay as
$O_{h,v}(e^{-3y/2})$ for $y\to+\infty$. Each growing pole integral
is absolutely convergent, and so is the prime series, since its tail is
bounded by a constant times $\sum_{n\ge2}(\log n)n^{-2}$.

At the archimedean singularity use, for $0<t\le1$,
\begin{equation}\label{eq:v14-translation}
 \|v_0(\cdot-t)-v_0\|_2
 \le C\left[t\|v'\|_2+
       \sqrt t\{|v(\inf J)|+|v(\sup J)|\}\right].
\end{equation}
This follows from the fundamental theorem of calculus on the overlap
and the two endpoint strips. Cauchy--Schwarz makes the subtracted
archimedean numerator $O(t^{1/2})$, which is integrable against its
$O(t^{-1})$ kernel. Its large-$t$ part is integrable by $L^2$ bounds
and exponential kernel decay.

\emph{Moment-preserving smoothing.}
Choose $\phi_\epsilon\ge0$ in $C_c^\infty((\epsilon,2\epsilon))$
with integral one, and define
\begin{equation}\label{eq:v14-smoothing}
 h_\epsilon(u)=\int\phi_\epsilon(t)e^{t/2}h(e^t u)\,dt,
 \qquad
 k_\epsilon(x)=\int\phi_\epsilon(t)k(x+t)\,dt.
\end{equation}
Multiplicative convolution smooths the endpoint; near zero it preserves
the smooth even germ. Hence $h_\epsilon$ is even and compactly supported
smooth, with support in $[-be^{-\epsilon},be^{-\epsilon}]$.
Both moments remain exactly zero: dilation multiplies the integral by
$e^{-t/2}$ and the central value is identically zero. Direct substitution
in the finite source sum proves
$\mathcal E(h_\epsilon)(e^x)=k_\epsilon(x)$, with all scale factors
as in \eqref{eq:v14-smoothing}. This is the scaling covariance used in
\cite[Lemma 6.2]{ConnesConsani2023}. Translation continuity gives
$k_\epsilon\to k$ in $L^2(\mathbb R)$, while
\[
 K_\epsilon(x)=\int\phi_\epsilon(t)K(x+t)\,dt
              =O_h(e^{3x/2})
\]
uniformly for sufficiently small $\epsilon$. In particular none of the
estimates needed for the passage uses derivatives of $h_\epsilon$ that
might blow up near its moving endpoint.

For a smooth compact logarithmic test $w$, the classical Schwartz
radical theorem gives $QW(k_\epsilon,w)=0$. To allow the stated $v$,
convolve its zero extension with a nonnegative smooth additive
approximate identity. These tests converge in $L^2$, have uniformly
bounded variation and support, and satisfy the same translation bound
\eqref{eq:v14-translation} by convolution contraction. The preceding
correlation bounds therefore dominate separately the prime series, both
pole integrals and the subtracted archimedean integral. Dominated
convergence gives $QW(k_\epsilon,v)=0$. The same argument, now using
the uniform primitive bound for $K_\epsilon$ and $L^2$ convergence,
allows $\epsilon\downarrow0$ and proves $QW(k,v)=0$.

\emph{The semilocal domain and truncation.}
On any fixed compact interval $k$ is a finite sum of smooth pieces with
finitely many jumps, so $k_{\rm in}$ is compactly supported BV. Such a
function has Fourier transform $O((1+|\xi|)^{-1})$, and hence
\[
 \int_{\mathbb R}(1+\log(2+|\xi|))
               |\widehat{k_{\rm in}}(\xi)|^2\,d\xi<\infty.
\]
The same holds for $v_0$. On a fixed compact support the archimedean
multiplier has logarithmic growth; the finitely many prime translations
and pole terms are bounded forms. This is the closed semilocal form
domain of the explicit formula~\cite{ConnesConsani2023}.
The smooth approximations converge there as well: on a fixed enlarged
interval they have uniform BV bounds and $L^2$ convergence, so a Fourier
tail split with $\log(2+|\xi|)/(1+|\xi|)^2$ proves convergence in the
logarithmic form norm. Compact truncation leaves the far lower primitive
unchanged, so $QW(r,v)$ is also finite. Linearity now gives the asserted
truncation identity, with the semilocal term equal to its closed-form
realization.
\end{proof}

The endpoint of the approximating sources need not equal that of $h$.
The theorem is proved at fixed $b$ \emph{before} normalizing by the
original source's separately evaluated one-sided endpoint. No endpoint
trace is inferred from form-norm convergence. This corrects the stronger,
unnecessary endpoint-preserving approximation requirement in version 1.3.

\subsection*{The focused PSWF cross-tail estimates}
The fixed-order radial estimates, the two-sided exterior-energy argument,
and Fuchs's concentration asymptotic below establish the source scale.
Together with Theorem~\ref{thm:v14-radical}, they close the fixed-band
and Sobolev-pairing G1 statements of this section for the explicit repaired
source. They do not give a uniform operator-norm bound on a growing deep
prolate block, spectral ordering, or the G2 graph/form comparison.


CCM's proposed source is the linear combination of the two fixed-order prolate modes $h_{0,\lambda},h_{4,\lambda}$ with vanishing integral.  The order-zero radial normalization can be matched exactly to the endpoint amplitude.  In Dunster's notation, for $m=0$ his matching constant satisfies
\[
 c_n^0(\gamma)=K_n^0(\gamma),
\]
where $K_n^0(\gamma)=Ps_n^0(1,\gamma^2)$.  His Bessel approximation is uniform for the full radial interval $1<x<\infty$.

\begin{proposition}[Endpoint-normalized radial bound]\label{prop:endpoint-radial}
Fix $n\in\{0,4\}$.  For all sufficiently large $\gamma$ there is a constant $C_n$, independent of $\gamma$, such that
\[
 |Ps_n^0(x,\gamma^2)|\le C_n\frac{|Ps_n^0(1,\gamma^2)|}{x},
 \qquad x\ge1.
\]
\end{proposition}
\begin{proof}
Dunster's formula (3.5) gives, uniformly for $1<x<\infty$,
\[
 Ps_n^0(x,\gamma^2)
 =c_n^0(\gamma)
 \left\{\frac{\eta}{(x^2-1)(x^2-\sigma^2)}\right\}^{1/4}
 \left[J_0(\gamma\eta^{1/2})
 +O(\gamma^{-1})\operatorname{env}J_0(\gamma\eta^{1/2})\right],
\]
with $\eta=\xi^2$.  His local relation
\[
 \eta=2(1-\sigma^2)(x-1)+O((x-1)^2)
\]
shows that the geometric prefactor tends to one as $x\downarrow1$; since $J_0(0)=1$ and $c_n^0(\gamma)=Ps_n^0(1,\gamma^2)$, the normalization is exactly the endpoint scale up to the uniform $O(\gamma^{-1})$ error.  For fixed $n$, Dunster also gives $\sigma^2=O_n(\gamma^{-1})$.

Use the standard envelope bound
\[
 \operatorname{env}J_0(y)\ll \min(1,y^{-1/2}).
\]
If $\gamma\xi\le1$, then $x-1=O_n(\gamma^{-2})$, so the geometric prefactor is uniformly bounded and $x^{-1}\asymp1$.  If $\gamma\xi\ge1$, the Bessel-envelope bound gives
\[
 \left\{\frac{\xi^2}{(x^2-1)(x^2-\sigma^2)}\right\}^{1/4}
 \operatorname{env}J_0(\gamma\xi)
 \ll_n
 \frac{\gamma^{-1/2}}{\{(x^2-1)(x^2-\sigma^2)\}^{1/4}}.
\]
On the boundary layer where $x-1=O(1)$ this is $O_n(1)$, while for $x$ bounded away from $1$ it is $O_n(x^{-1})$.  Enlarging the constant across the compact transition interval proves the claim.
\end{proof}

The preceding endpoint estimate can be sharpened away from the radial turning point, and the resulting gain is the key to summing the co-Poisson lattice.

\begin{lemma}[Away-boundary and exterior-energy bounds]\label{lem:pswf-away-energy}
Fix $n\in\{0,4\}$ and write
\[
 \Psi_{n,c}(x)=Ps_n^0(x,c^2),\qquad B_{n,c}=|\Psi_{n,c}(1)|.
\]
For all sufficiently large $c$,
\begin{align}
 |\Psi_{n,c}(x)|&\ll_n \frac{B_{n,c}}{\sqrt c\,x},\qquad x\ge2,\label{eq:pswf-away}\\
 \int_1^\infty |\Psi_{n,c}(x)|^2\,dx&\ll_n \frac{B_{n,c}^2}{c}.\label{eq:pswf-exterior-energy}
\end{align}
\end{lemma}
\begin{proof}
In Dunster's uniform formula (3.5), for $x\ge2$ one has $\xi(x)\asymp_n x$ and
$\{(x^2-1)(x^2-\sigma^2)\}^{1/4}\asymp_n x$.  The Bessel-envelope estimate therefore gives
\[
 \left\{\frac{\xi^2}{(x^2-1)(x^2-\sigma^2)}\right\}^{1/4}
 \operatorname{env}J_0(c\xi)
 \ll_n \frac1{\sqrt c\,x},
\]
which proves \eqref{eq:pswf-away}.

For $1<x<2$, use $d\xi/dx=((x^2-\sigma^2)/(x^2-1))^{1/2}$.  Squaring Dunster's formula and changing variables from $x$ to $\xi$ reduces the main term, up to bounded fixed-$n$ factors, to
\[
 B_{n,c}^2\int_0^{\xi(2)}
 \xi\,\operatorname{env}J_0(c\xi)^2\,d\xi.
\]
The bound $\operatorname{env}J_0(t)^2\ll\min(1,t^{-1})$ makes this $O_n(B_{n,c}^2/c)$; the uniform $O(c^{-1})$ remainder in Dunster's formula is smaller.  On $[2,\infty)$, \eqref{eq:pswf-away} is square integrable and contributes another $O_n(B_{n,c}^2/c)$.  This proves \eqref{eq:pswf-exterior-energy}.
\end{proof}

\subsection*{Normalized source coefficients and endpoint noncancellation}
We now settle the coefficient and repair-size parts of the former source
assumption. These estimates concern the actual two-mode source, not a
choice of independently rescaled endpoint values. The form-domain issue
is treated separately from the elementary correlation consequences below.

Let $\psi_{n,c}$ be the real angular PSWF with
$\|\psi_{n,c}\|_{L^2(-1,1)}=1$, extended to the real line by its finite
Fourier eigenrelation. For $n=0,4$ choose its sign so that
$\psi_{n,c}(0)>0$. Write $\vartheta_n(c)=\chi_n(c)^2$ for its
concentration eigenvalue, and put $b_n(c)=|\psi_{n,c}(1)|$.
These $\psi_{n,c}$ are scalar multiples of the preceding Dunster
normalization; its homogeneous radial estimates are unchanged.

\begin{lemma}[Endpoint size from exterior energy]
\label{lem:v12-endpoint-energy}
For each $n\in\{0,4\}$, as $c\to\infty$,
\begin{equation}\label{eq:v12-endpoint-energy}
 b_n(c)^2\asymp_n c\{1-\vartheta_n(c)\},\qquad
 b_n(c)\asymp_n c^{n/2+3/4}e^{-c}.
\end{equation}
In particular $b_4(c)/b_0(c)\asymp c^2$.
\end{lemma}
\begin{proof}
For clarity, the eigenrelation on $[-1,1]$ is
\[
 \int_{-1}^{1}e^{-icxy}\psi_{n,c}(y)\,dy
       =\mu_n(c)\psi_{n,c}(x),\qquad
 |\mu_n(c)|^2=\frac{2\pi}{c}\vartheta_n(c).
\]
Using the same formula to extend the function to all real $x$, Plancherel
gives the exact identity
\begin{equation}\label{eq:v12-exterior-identity}
 \int_1^\infty |\psi_{n,c}(x)|^2\,dx
     =\frac{1-\vartheta_n(c)}{2\vartheta_n(c)}.
\end{equation}
The upper bound for this integral by $C_nb_n(c)^2/c$ was proved in
Lemma~\ref{lem:pswf-away-energy}. We also need a lower bound; an upper
bound alone would not justify an endpoint comparison.

On the fixed interval $2\le x\le3$, Dunster's radial formula
\cite[eqs.~(3.5)--(3.7)]{Dunster2017} and the large-argument expansion
of $J_0$ give, with a real choice of the endpoint sign,
\[
 \psi_{n,c}(x)=\psi_{n,c}(1)\sqrt{\frac{2}{\pi c}}
 w_c(x)\left\{\cos(c\xi_c(x)-\pi/4)+O_n(c^{-1})\right\},
\]
where $w_c(x)=\{(x^2-1)(x^2-\sigma_n(c)^2)\}^{-1/4}$ and
$\xi_c'(x)=\{(x^2-\sigma_n(c)^2)/(x^2-1)\}^{1/2}$.
Here $\sigma_n(c)^2=O_n(c^{-1})$. On this interval $w_c$ and
$\xi_c'$ have uniform positive lower bounds and bounded derivatives.
One integration by parts therefore yields
\[
 \int_2^3 w_c(x)^2\cos^2(c\xi_c(x)-\pi/4)\,dx
      =\frac12\int_2^3w_c(x)^2\,dx+O_n(c^{-1})\gg_n1.
\]
The integrated cross and squared remainder terms are smaller by an
additional factor $c^{-1}$. Hence the exterior energy is also at least
$c_n b_n(c)^2/c$ for some $c_n>0$. Equation
\eqref{eq:v12-exterior-identity} and $\vartheta_n(c)\to1$ prove the
first comparison in \eqref{eq:v12-endpoint-energy}.

Fuchs's fixed-index asymptotic \cite[Theorem~1]{Fuchs1964}, in this
concentration normalization, is
\begin{equation}\label{eq:v12-fuchs}
 1-\vartheta_n(c)\sim
 \frac{4\sqrt\pi\,8^n}{n!}\,c^{n+1/2}e^{-2c}.
\end{equation}
Combining it with the two-sided energy comparison proves the second
comparison and the ratio assertion. No differentiation of an asymptotic
remainder is used.
\end{proof}

\begin{proposition}[An explicit normalized two-mode source]
\label{prop:v12-source}
Set $c=2\pi\lambda^2$ and
$h_{n,\lambda}(u)=\lambda^{-1/2}\psi_{n,c}(u/\lambda)$ on
$[-\lambda,\lambda]$, zero-extended for Fourier transformation.
Endpoint values denote interior traces. Define
\begin{align}
 a_{4,\lambda}&=\chi_0(c)h_{0,\lambda}(0),&
 a_{0,\lambda}&=-\chi_4(c)h_{4,\lambda}(0),\label{eq:v12-coefficients}\\
 h_\lambda^\circ&=a_{0,\lambda}h_{0,\lambda}
                 +a_{4,\lambda}h_{4,\lambda},&
 B_\lambda^+&=\sqrt\lambda\,h_\lambda^\circ(\lambda).
 \notag
\end{align}
The source has exactly zero integral. Both coefficients have finite
nonzero limits. For all sufficiently large $\lambda$,
\begin{align}
 |B_\lambda^+|&\asymp c^{11/4}e^{-c},\label{eq:v12-combined-endpoint}\\
 \frac{\sum_{n=0,4}|a_{n,\lambda}|\sqrt\lambda
                 |h_{n,\lambda}(\lambda)|}{|B_\lambda^+|}
     &=1+O(c^{-2}),\label{eq:v12-noncancellation}\\
 \frac{|h_\lambda^\circ(0)|}{|B_\lambda^+|}
     &=O(c^{7/4}e^{-c})=O_A(\lambda^{-A})
       \quad\text{for every fixed }A.\label{eq:v12-zero-value}
\end{align}
The last two conclusions are invariant under a common nonzero scalar
rescaling of the source.
\end{proposition}
\begin{proof}
The compressed physical Fourier eigenvalue is $\chi_n(c)>0$ for
$n=0,4$, so evaluation at zero gives exactly
\[
 \int_{-\lambda}^{\lambda}h_{n,\lambda}(u)\,du
       =\chi_n(c)h_{n,\lambda}(0).
\]
Thus \eqref{eq:v12-coefficients} cancels the integral. The fixed-order
Hermite limit, with the unit-norm scaling just specified, gives
\[
 h_{0,\lambda}(0)\longrightarrow2^{1/4},\qquad
 h_{4,\lambda}(0)\longrightarrow
       \frac{\sqrt3}{2\,2^{1/4}}.
\]
This central-value normalization follows, for example, from the
fixed-order limit used in \cite[Lemma~7.2]{ConnesConsaniMoscovici2025};
the uniform error there also makes the physical $L^2$ norm tend to one,
so unit-norm renormalization preserves these limits. Since $\chi_n\to1$,
both coefficients converge to the stated nonzero central constants,
with a minus sign for $a_{0,\lambda}$.

By Lemma~\ref{lem:v12-endpoint-energy}, the magnitude of the order-zero
endpoint contribution divided by that of order four is $O(c^{-2})$.
The triangle and reverse triangle inequalities prove
\eqref{eq:v12-combined-endpoint} and \eqref{eq:v12-noncancellation},
without needing a claim about their relative signs. Finally,
\[
 h_\lambda^\circ(0)
   =\{\chi_0(c)-\chi_4(c)\}
       h_{0,\lambda}(0)h_{4,\lambda}(0).
\]
Equation~\eqref{eq:v12-fuchs} and
$1-\chi_n=(1-\vartheta_n)/(1+\chi_n)$ give
$|h_\lambda^\circ(0)|=O(c^{9/2}e^{-2c})$.
Division by \eqref{eq:v12-combined-endpoint} gives
\eqref{eq:v12-zero-value}.
\end{proof}

\begin{corollary}[Size of the exact value-and-integral repair]
\label{cor:v12-repair}
Fix an even $\psi\in C_c^\infty((-1,1))$ with $\psi(0)=1$ and
$\int\psi=0$. For the source in Proposition~\ref{prop:v12-source},
\[
 \widetilde h_\lambda=h_\lambda^\circ-h_\lambda^\circ(0)\psi
\]
has zero value at zero, zero integral and the same endpoint as
$h_\lambda^\circ$. In lower-tail coordinates $u=\lambda^{-1}e^{-s}$,
for every fixed $A$ and every $S\ge0$ its co-Poisson correction obeys
\begin{equation}\label{eq:v12-repair-tail}
 \left\|1_{s\ge S}\,h_\lambda^\circ(0)
            \mathcal E(\psi)(\lambda^{-1}e^{-s})\right\|_2
     \le C_A |B_\lambda^+|\lambda^{-A}e^{-S/2}.
\end{equation}
This proves the repair's lower-tail size. Its repaired compact source
also meets Theorem~\ref{thm:v14-radical}, which supplies the radical and
semilocal form-domain passage used below.
\end{corollary}
\begin{proof}
The moment and endpoint statements follow directly from the fixed bump.
Poisson summation, retaining its zero-value term, gives
\[
 \mathcal E(\psi)(u)=-\frac12u^{1/2}
           +u^{-1/2}\sum_{m\ge1}\widehat\psi(m/u)=O(u^{1/2})
       \qquad(0<u\le1).
\]
The last bound follows from Schwartz decay, taking any fixed power
greater than one in the sum. The $L^2$ tail of $u^{1/2}$ in the
$s$ coordinate equals $\lambda^{-1/2}e^{-S/2}$. Apply
\eqref{eq:v12-zero-value}, increasing its arbitrary power if needed.
\end{proof}

The coefficient noncancellation and repair-size bounds are therefore no
longer assumptions for this explicitly normalized source. The tail-to-
correlation consequences and translated endpoint terms are proved below;
the global radical/form-domain passage is supplied by
Theorem~\ref{thm:v14-radical}.


The possible lattice resonance can be isolated entirely in a range where absolute summation is cheap.  Beyond that range the radial differential equation itself supplies a genuinely linear phase with a summable remainder.

\begin{lemma}[Uniform far field with controlled $c$-dependence]\label{lem:pswf-far-volterra}
For fixed $n\in\{0,4\}$ there is a scalar $A_{n,c}$ such that, uniformly for $x\ge c^2$,
\begin{equation}\label{eq:pswf-far-controlled}
 \Psi_{n,c}(x)
 =A_{n,c}\frac{\sin(cx-n\pi/2)}{x}
 +O_n\!\left(\frac{B_{n,c}\sqrt c}{x^2}\right),
 \qquad
 |A_{n,c}|\ll_n\frac{B_{n,c}}{\sqrt c}.
\end{equation}
\end{lemma}
\begin{proof}
Dunster's radial equation, with $m=0$, becomes after the standard substitution
$w(x)=(x^2-1)^{1/2}\Psi_{n,c}(x)$,
\[
 w''+c^2w=V_{n,c}(x)w,
 \qquad
 V_{n,c}(x)=\frac{\lambda_n^0(c^2)}{x^2-1}-\frac1{(x^2-1)^2}.
\]
His fixed-order relation $\sigma^2=2(n+\tfrac12)c^{-1}+O_n(c^{-2})$, together with
$\sigma^2=1+c^{-2}\lambda_n^0(c^2)$, gives
\[
 \lambda_n^0(c^2)=-c^2+(2n+1)c+O_n(1).
\]
Hence $|V_{n,c}(x)|\ll_n c^2x^{-2}$ for $x\ge2$ and
\[
 \int_{c^2}^\infty \frac{|V_{n,c}(t)|}{c}\,dt\ll_n c^{-1}.
\]
By Lemma~\ref{lem:pswf-away-energy}, $|w(x)|\ll_n B_{n,c}/\sqrt c$ for $x\ge2$.  On an interval of length $c^{-1}$ about $x_0=c^2$, the differential equation then gives
$|w''|\ll_n B_{n,c}c^{3/2}$; Taylor's formula yields
$|w'(x_0)|\ll_n B_{n,c}\sqrt c$.  Variation of constants for the perturbed oscillator and Gronwall therefore give
\[
 |w(x)|+c^{-1}|w'(x)|\ll_n B_{n,c}/\sqrt c,
 \qquad x\ge c^2.
\]
The usual incoming/outgoing amplitudes have derivatives bounded by
$|V_{n,c}(x)w(x)|/c$, which is integrable.  They therefore converge at infinity, and integration from $x$ to infinity gives an error
\[
 \ll_n \frac1c\int_x^\infty |V_{n,c}(t)w(t)|\,dt
 \ll_n \frac{B_{n,c}\sqrt c}{x}.
\]
After division by $(x^2-1)^{1/2}\asymp x$, this is the remainder in
\eqref{eq:pswf-far-controlled}.  Finally Dunster's exact radial phase at infinity, equation (1.24), identifies the limiting real combination with
$\sin(cx-n\pi/2)$, proving the stated form and amplitude bound.
\end{proof}

\begin{theorem}[Endpoint-normalized radial lattice bound]\label{thm:lattice-tail}
Put $c=2\pi\lambda^2$.  For fixed $n\in\{0,4\}$ and all $s\ge0$,
\begin{equation}\label{eq:lattice-log-target}
 \boxed{
 \left|\sum_{m\ge2}\Psi_{n,c}(me^s)\right|
 \ll_n
 \frac{B_{n,c}}{\sqrt c\,e^s}(1+\log c).}
\end{equation}
The estimate is uniform in $s$ and hence contains no unproved nonresonance hypothesis.
\end{theorem}
\begin{proof}
Set $x=e^s$ and
\[
 M=\max\!\left(2,\left\lceil\frac{c^2}{x}\right\rceil\right).
\]
For $2\le m<M$, Lemma~\ref{lem:pswf-away-energy} gives, by absolute summation,
\[
 \sum_{2\le m<M}|\Psi_{n,c}(mx)|
 \ll_n\frac{B_{n,c}}{\sqrt c\,x}\sum_{m<M}\frac1m
 \ll_n\frac{B_{n,c}(1+\log c)}{\sqrt c\,x}.
\]
For $m\ge M$, Lemma~\ref{lem:pswf-far-volterra} applies.  Since $n=0$ or $4$, the fixed phase $n\pi/2$ is an integral multiple of $2\pi$.  The elementary harmonic-sine bound
\[
 \sup_{\theta\in\mathbb R}\sup_{M\ge1}
 \left|\sum_{m=M}^{\infty}\frac{\sin(m\theta)}m\right|<\infty
\]
then gives
\[
 \left|\frac{A_{n,c}}x\sum_{m\ge M}\frac{\sin(cmx-n\pi/2)}m\right|
 \ll_n\frac{B_{n,c}}{\sqrt c\,x}.
\]
(The displayed uniform bound follows, for example, by splitting at $m\asymp |\theta|^{-1}$ and applying Dirichlet summation to the remaining tail.)
The remainder in \eqref{eq:pswf-far-controlled} contributes
\[
 \ll_n\frac{B_{n,c}\sqrt c}{x^2}\sum_{m\ge M}m^{-2}.
\]
If $M\asymp c^2/x$ this is $O_n(B_{n,c}/(c^{3/2}x))$; if $M=2$, then $x\gg c^2$ and it is even smaller than $B_{n,c}/(\sqrt c\,x)$.  Combining the two ranges proves \eqref{eq:lattice-log-target}.
\end{proof}

The $m=1$ Poisson sample is a boundary-layer term rather than a lattice-summation problem.  The finite-Fourier normalization introduces no additional power of $\lambda$: in the standard PSWF convention
\[
 \mathcal F_c\Psi_{n,c}=\mu_n(c)\Psi_{n,c},\qquad
 \frac{c}{2\pi}|\mu_n(c)|^2=\chi_n(c)^2,
\]
where $\chi_n^2$ is the concentration eigenvalue.  Since $c=2\pi\lambda^2$,
$|\mu_n(c)|=\chi_n(c)/\lambda$.  After scaling from $[-1,1]$ to
$[-\lambda,\lambda]$, if $u=\lambda^{-1}e^{-s}$ then the Fourier-lattice part of the lower tail is, up to the harmless fixed even-mode Fourier phase,
\begin{equation}\label{eq:scaled-poisson-radial}
 r_{n,\lambda}^{\rm lat}(s)
 =\chi_n(c)e^{s/2}\sum_{m\ge1}\Psi_{n,c}(me^s).
\end{equation}
For an even source with $\int h=0$ but $h(0)\ne0$, Poisson summation has in addition the elementary term
\begin{equation}\label{eq:poisson-zero-defect}
 \mathcal E(h)(u)
 =-\frac12u^{1/2}h(0)
   +u^{-1/2}\sum_{m\ge1}\widehat h(m/u).
\end{equation}
Thus the raw CCM two-mode source has a small zero-value defect in addition to the radial lattice term.  It will be estimated explicitly below rather than silently omitted.  In particular, at $s=0$ the $m=1$ lattice term is $\chi_n(c)$ times the upper-window boundary amplitude.

\begin{corollary}[Fixed-order co-Poisson tail and endpoint comparison]\label{cor:pswf-cross-tail}
Let $h_\lambda^{\circ}$ be the explicitly normalized two-mode source of
Proposition~\ref{prop:v12-source}, and let
$B_\lambda^{+}=\sqrt\lambda\,h_\lambda^{\circ}(\lambda)$ be its upper-window co-Poisson boundary amplitude.  Then the lower logarithmic tail $r_\lambda^{\circ}$ satisfies, for every $S\ge0$,
\begin{equation}\label{eq:lower-tail-l2}
 \left\|1_{\{s\ge S\}}r_\lambda^{\circ}\right\|_2
 \ll \frac{|B_\lambda^{+}|}{\lambda}(1+\log\lambda)e^{-S/2}.
\end{equation}
The same estimate holds for the exact value-and-integral repair. If
$B_\lambda$ is the common one-sided endpoint value of the inversion-even
finite-window proxy, then
\begin{equation}\label{eq:v13-endpoint-comparison}
 \frac{B_\lambda}{B_\lambda^+}
 =1+O\!\left(\frac{1+\log c}{\sqrt c}\right),
 \qquad c=2\pi\lambda^2.
\end{equation}
In particular $B_\lambda\asymp B_\lambda^+$ and it is nonzero for all
sufficiently large $\lambda$.
\end{corollary}
\begin{proof}
For the $m=1$ term in \eqref{eq:scaled-poisson-radial}, Lemma~\ref{lem:pswf-away-energy} gives the exact log-coordinate estimate
\[
 \int_0^\infty
 \left|\chi_n(c)e^{s/2}\Psi_{n,c}(e^s)\right|^2ds
 =\chi_n(c)^2\int_1^\infty|\Psi_{n,c}(x)|^2dx
 \ll_n\frac{B_{n,c}^2}{c}.
\]
For $s\ge\log2$, \eqref{eq:pswf-away} also gives the pointwise bound
$O_n(B_{n,c}c^{-1/2}e^{-s/2})$.  Thus its $L^2$ tail beyond depth $S$ is
$O_n(B_{n,c}c^{-1/2}e^{-S/2})$ for all $S\ge0$, after enlarging the constant on $0\le S\le\log2$.

For $m\ge2$, Theorem~\ref{thm:lattice-tail} and \eqref{eq:scaled-poisson-radial} give directly
\[
 |r_{n,\lambda}^{(\ge2)}(s)|
 \ll_n \frac{B_{n,c}(1+\log c)}{\sqrt c}e^{-s/2}.
\]
Since $\sqrt c\asymp\lambda$, this proves \eqref{eq:lower-tail-l2} mode by mode. The endpoint noncancellation bound in Proposition~\ref{prop:v12-source} controls the sum of the individual endpoint magnitudes by $|B_\lambda^+|$, so the modewise estimate passes to the actual normalized combination. This part no longer uses an unproved coefficient estimate.

For completeness, the exact radical source condition can be imposed without changing this scale.  Choose a fixed even $\psi\in C_c^\infty((-1,1))$ with $\psi(0)=1$ and $\int\psi=0$, and put
\[
 \widetilde h_\lambda=h_\lambda^{\circ}-h_\lambda^{\circ}(0)\psi.
\]
Then $\widetilde h_\lambda(0)=0$ and $\int\widetilde h_\lambda=0$, while its endpoint at $\pm\lambda$ is unchanged. Proposition~\ref{prop:v12-source} proves
$|h_\lambda^\circ(0)|/|B_\lambda^+|=O_A(\lambda^{-A})$ for every fixed $A$,
and Corollary~\ref{cor:v12-repair} supplies the lower-tail repair bound.
The repaired compact source satisfies Theorem~\ref{thm:v14-radical},
so its global radical and semilocal form-domain use is justified.
For the fixed bump one must retain its zero-value term: since $\widehat\psi(0)=\int\psi=0$ and $\psi(0)=1$, Poisson summation gives
\[
 \mathcal E(\psi)(u)
 =-\frac12u^{1/2}
   +u^{-1/2}\sum_{m\ge1}\widehat\psi(m/u).
\]
The second term is rapidly decreasing as $u\downarrow0$, while the first is $O(u^{1/2})$.  Because $|h_\lambda^{\circ}(0)|/|B_\lambda^{+}|=O_A(\lambda^{-A})$ for every fixed $A$, the whole correction is still $o(|B_\lambda^{+}|\lambda^{-1}e^{-s/2})$ at $u=\lambda^{-1}e^{-s}$ and is negligible in \eqref{eq:lower-tail-l2}.  The same estimate also absorbs the raw zero-defect term in \eqref{eq:poisson-zero-defect}.

It remains to identify the endpoint after inversion. Put
$f_\lambda=\mathcal E(\widetilde h_\lambda)$. At the upper endpoint only
the $m=1$ source summand remains, so
$f_\lambda(\lambda^-)=B_\lambda^+$. At the lower endpoint, Poisson
summation and \eqref{eq:scaled-poisson-radial} give
\[
 f_\lambda((\lambda^{-1})^+)
 =\sum_{n=0,4}a_{n,\lambda}\chi_n(c)
       \sqrt\lambda h_{n,\lambda}(\lambda)
 +O\!\left(\frac{1+\log c}{\sqrt c}
       \sum_{n=0,4}|a_{n,\lambda}|\sqrt\lambda
                    |h_{n,\lambda}(\lambda)|\right)
 +o(|B_\lambda^+|).
\]
The last term includes the fixed-bump repair. If $\lambda^2$ is an
integer, taking the indicated interior trace can also remove the single
source summand at the cut; its size is $O(|B_\lambda^+|/\lambda)$ and is
absorbed above. Since $\chi_n(c)=1+o(1)$ and
\eqref{eq:v12-noncancellation} controls the sum of the individual endpoint
magnitudes by $|B_\lambda^+|$, the displayed expression equals
$B_\lambda^+\{1+O((1+\log c)/\sqrt c)\}$. Finally
$p_\lambda=(f_\lambda+Jf_\lambda)/2$, so its two common one-sided endpoint
values are the average of the upper and lower traces. This proves
\eqref{eq:v13-endpoint-comparison}.
\end{proof}

\begin{lemma}[Uniform zero-extension correlations and the local term]
\label{lem:v13-correlations}
Let $L=2\log\lambda$, let $r_\lambda$ be either oriented logarithmic tail
of the repaired source, and zero-extend $v\in H^1([-L,0])$. For
\[
 C_{r,v}(y)=\int_{\mathbb R}r_\lambda(t)
                  \overline{v(t-y)}\,dt
\]
and for each reflected correlation occurring in the polarized even Weil
kernel, the interior-test derivative contribution satisfies
\begin{equation}\label{eq:pswf-cross-tail}
 \boxed{
 |C_{r,v,\mathrm{bulk}}'(y)|
 \le C\frac{|B_\lambda|}{\lambda}(1+\log\lambda)
       \|v'\|_2
 \begin{cases}
  1,&0\le y\le L,\\[2mm]
  e^{-(y-L)/2},&y>L.
 \end{cases}}
\end{equation}
Here $C_{r,v,\mathrm{bulk}}'(y)=-\int r_\lambda(t)
\overline{v'(t-y)}1_{(-L,0)}(t-y)\,dt$.
The full distributional derivative has, in addition, precisely the two terms
\begin{equation}\label{eq:v13-correlation-jumps}
 -\overline{v(-L)}\,r_\lambda(y-L)
 +\overline{v(0)}\,r_\lambda(y),
\end{equation}
up to the harmless sign and reflection dictated by the chosen orientation;
a term with a negative tail argument is zero. These translated-tail terms
are themselves locally integrable densities, not singular atoms in the
derivative of the correlation. They must not be omitted from its
absolutely continuous part. Moreover the local
archimedean functional of the even polarized sum $q_\lambda$ obeys
\begin{equation}\label{eq:v13-archimedean}
 |\mathcal A_\infty(q_\lambda)|
 \le C\frac{|B_\lambda|}{\lambda}(1+\log\lambda)
 \bigl(\|v\|_2+\|v'\|_2+|v(-L)|+|v(0)|\bigr).
\end{equation}
These estimates are uniform in $v$ and require no form-domain assumption.
\end{lemma}
\begin{proof}
Inside the support interval, differentiation in $y$ puts one derivative
on $v$ and never on the prolate tail. The overlap uses tail depths
$t\ge(y-L)_+$. Cauchy--Schwarz and \eqref{eq:lower-tail-l2}, followed by
\eqref{eq:v13-endpoint-comparison}, prove
\eqref{eq:pswf-cross-tail}. In distributions,
\[
 D(1_{[-L,0]}v)=1_{(-L,0)}v'
       +v(-L)\delta_{-L}-v(0)\delta_0.
\]
Substitution into the correlation, with the minus sign from differentiating
$v(t-y)$, gives \eqref{eq:v13-correlation-jumps}. Reflection gives the
other orientation and changes none of the bounds.

For the singularity at the origin use the elementary zero-extension
translation estimate, valid for $0<x\le1$,
\begin{equation}\label{eq:v13-translation}
 \|v(\cdot-x)-v\|_{L^2(\mathbb R)}
 \le C\left[x\|v'\|_2
   +\sqrt{x}\{ |v(-L)|+|v(0)|\}\right].
\end{equation}
It follows by applying the fundamental theorem of calculus on the overlap
and estimating the two boundary strips separately. Thus, with
$R_\lambda=|B_\lambda|(1+\log\lambda)/\lambda$,
\[
 |q_\lambda(x)+q_\lambda(-x)-2e^{-x/2}q_\lambda(0)|
 \le C R_\lambda\left[
  \sqrt{x}\{|v(-L)|+|v(0)|\}
  +x\{\|v'\|_2+\|v\|_2\}\right].
\]
The archimedean kernel is $O(x^{-1})$ on $(0,1]$, so this is integrable.
For $x\ge1$ the kernel is $O(e^{-x/2})$, while Cauchy--Schwarz and
\eqref{eq:lower-tail-l2} bound every correlation by
$CR_\lambda\|v\|_2$. Integration proves
\eqref{eq:v13-archimedean}.
\end{proof}

The truncated-Mellin identity remains a useful independent check on the powers and normalizations.  For an exact two-moment source,
\begin{equation}\label{eq:truncated-mellin-tail}
 \mathcal M^-_\lambda(\tau)
 =\int_0^{1/\lambda}\mathcal E(h_\lambda)(u)u^{-i\tau}\,\frac{du}{u},
\end{equation}
and Poisson summation gives
\begin{equation}\label{eq:truncated-mellin-poisson}
 \mathcal M^-_\lambda(\tau)
 =\sum_{m\ge1}m^{-1/2-i\tau}
   \int_{m\lambda}^{\infty}\widehat h_\lambda(x)
   x^{-1/2+i\tau}\,dx.
\end{equation}
The preceding argument unconditionally supplies the physical/logarithmic
tail and correlation estimates needed below, given the cited classical
fixed-order PSWF inputs. Theorem~\ref{thm:v14-radical} justifies passing
the global radical identity to the closed semilocal form at each fixed
$\lambda$; the quantitative estimates, not its approximation constants,
provide the required uniformity as $\lambda\to\infty$.


\begin{lemma}[Quantitative endpoint-jump bound]\label{lem:endpoint-jumps}
Let $E_{\lambda,T}$ be the periodic trigonometric space on $[-a,a]$, $a=\log\lambda$, with physical logarithmic frequencies in $[-T,T]$.  Extend $v\in E_{\lambda,T}$ by zero outside $[-a,a]$, and let $q_\lambda$ be its logarithmic correlation with the lower tail $r_\lambda$.  In the distributional derivative of $q_\lambda$, the two endpoint jumps contribute translates of $r_\lambda$ multiplied by the common periodic endpoint value $v(a)=v(-a)$.  There is a constant $c_1>0$, independent of $\lambda$, such that uniformly for $\|v\|_2=1$ their total contribution to the centered pole--prime Stieltjes pairing is
\begin{equation}\label{eq:endpoint-jump-rate}
 \ll_T |B_\lambda|\left((1+L)e^{-L/2}
 +(1+L)^{3/2}e^{-c_1\sqrt L}\right),
 \qquad L=2\log\lambda.
\end{equation}
In particular the endpoint-jump contribution is $o_T(|B_\lambda|)$.
\end{lemma}
\begin{proof}
The endpoint evaluation functional on $E_{\lambda,T}$ has norm $O_T(1)$ because the normalized Fourier basis has $O_T(L)$ modes and each basis vector has endpoint magnitude $L^{-1/2}$.  For the jump meeting the lower tail at depth $s=y$, split $0\le s\le\log2$ from $s\ge\log2$.  On the first interval, $V(s)=O(s)$ and Cauchy--Schwarz with \eqref{eq:lower-tail-l2} gives
\[
 O_T\!\left(|B_\lambda|\frac{1+\log\lambda}{\lambda}\right)
 =O_T\!\left(|B_\lambda|(1+L)e^{-L/2}\right).
\]
On the second, Theorem~\ref{thm:lattice-tail} together with the $m=1$ away-bound gives
\[
 |r_\lambda(s)|\ll |B_\lambda|\frac{1+\log\lambda}{\lambda}e^{-s/2},
 \qquad s\ge\log2,
\]
and \eqref{eq:centered-prime-cumulative} below supplies an integrable factor $e^{-c\sqrt s}$.  This contributes the same $O_T(|B_\lambda|(1+L)e^{-L/2})$ scale.

The second endpoint enters only after a shift $y=L+s$.  There
\[
 |V(L+s)|\ll \lambda e^{s/2}e^{-c\sqrt{L+s}}.
\]
For $0\le s\le\log2$, Cauchy--Schwarz and \eqref{eq:lower-tail-l2} cancel the factor $\lambda$ and give
$O_T(|B_\lambda|(1+L)e^{-c\sqrt L})$.  For $s\ge\log2$, the preceding pointwise tail estimate cancels $e^{s/2}$ and leaves
\[
 O_T\!\left(|B_\lambda|(1+L)
 \int_0^\infty e^{-c\sqrt{L+s}}\,ds\right).
\]
The integral is explicit:
\[
 \int_0^\infty e^{-c\sqrt{L+s}}\,ds
 =2e^{-c\sqrt L}\left(\frac{\sqrt L}{c}+\frac1{c^2}\right).
\]
After decreasing the exponential constant if necessary, these estimates give \eqref{eq:endpoint-jump-rate}.
\end{proof}

\subsection*{Pole--prime reduction and closure of the weak G1 estimate}
Let $k_\lambda$ be the exact-radical co-Poisson vector obtained from the corrected source, and write
\[
 k_\lambda=k_{\lambda,\mathrm{in}}+r_\lambda,
 \qquad k_{\lambda,\mathrm{in}}=1_{I_\lambda}k_\lambda.
\]
For every zero-extended $H^1$ test vector $v$ supported in $I_\lambda$,
Theorem~\ref{thm:v14-radical} gives
\begin{equation}\label{eq:radical-truncation}
 QW(k_{\lambda,\mathrm{in}},v)=-QW(r_\lambda,v).
\end{equation}

The dominant point-mass part and the prime part of the explicit formula pair the logarithmic correlation $q_\lambda$ with
\[
 e^{y/2}\,dy
 \quad\text{and}\quad
 \sum_{n\ge2}\frac{\Lambda(n)}{\sqrt n}\,\delta_{\log n},
\]
respectively; the additional $e^{-y/2}$ part is absolutely integrable.  Put
\[
 V(Y)=2(e^{Y/2}-1)-\sum_{2\le n\le e^Y}\frac{\Lambda(n)}{\sqrt n}.
\]
Partial summation together with the classical de la Vall\'ee-Poussin prime-number-theorem remainder gives constants $c,C>0$ such that
\begin{equation}\label{eq:centered-prime-cumulative}
 |V(Y)|\le C e^{Y/2}e^{-c\sqrt Y}\qquad(Y\ge1).
\end{equation}
Stieltjes integration by parts can be made quantitative.  Put
\[
 A_{\lambda,T}=C_T|B_\lambda|(1+L)e^{-L/2},
\]
so that \eqref{eq:pswf-cross-tail} bounds the interior-test derivative part by
$A_{\lambda,T}$ on $[0,L]$ and by
$A_{\lambda,T}e^{-(y-L)/2}$ on $[L,\infty)$.  On $1\le y\le L$ split at
$L/4$.  The first part is
$O_T(|B_\lambda|(1+L)L e^{-3L/8})$, while on $[L/4,L]$ the PNT factor gives
$e^{-c\sqrt y}\le e^{-c\sqrt L/2}$ and the $e^{y/2}$ integral cancels the
factor $e^{-L/2}$.  Beyond $L$ one obtains
\[
 A_{\lambda,T}e^{L/2}\int_L^\infty e^{-c\sqrt y}\,dy
 =2A_{\lambda,T}e^{L/2-c\sqrt L}
 \left(\frac{\sqrt L}{c}+\frac1{c^2}\right).
\]
Together with Lemma~\ref{lem:endpoint-jumps}, this yields constants
$c_1,C_T>0$ such that
\begin{equation}\label{eq:pole-prime-small}
 \left|\int_0^\infty q_\lambda(y)e^{y/2}\,dy
 -\sum_{n\ge2}\frac{\Lambda(n)}{\sqrt n}q_\lambda(\log n)\right|
 \le C_T|B_\lambda|\left((1+L)^{3/2}e^{-c_1\sqrt L}
 +(1+L)e^{-L/2}\right).
\end{equation}
The constants use only the classical unconditional zero-free-region/PNT remainder; no
RH assumption or information about individual nontrivial zeros is inserted.

For completeness, the archimedean local term may be estimated directly rather than hidden inside a multiplier statement.  In logarithmic coordinates the standard Weil formula contributes, after polarization,
\[
 \mathcal A_\infty(q)=-(\log 4\pi+\gamma_0)q(0)
 -\int_0^\infty
 \{q(x)+q(-x)-2e^{-x/2}q(0)\}
 \frac{e^{x/2}}{e^x-e^{-x}}\,dx .
\]
The kernel is $O(e^{-x/2})$ for $x\ge1$ and is $O(x^{-1})$ at the origin.
For zero-extended $H^1$ tests the numerator need only be $O(x^{1/2})$
because of the endpoint strips; this is still integrable. Lemma~\ref{lem:v13-correlations}
gives precisely the required uniform estimate:
\[
 |q(0)|\ll |B_\lambda|e^{-L/2}(1+\log\lambda),\qquad
 |\mathcal A_\infty(q_\lambda)|\ll_T |B_\lambda|e^{-L/2}(1+\log\lambda)
 =O_T\!\left(|B_\lambda|\frac{1+\log\lambda}{\lambda}\right).
\]
The remaining $e^{-y/2}$ pole contribution is absolutely integrable and obeys the same bound.

Define the explicit zero-independent envelope
\begin{equation}\label{eq:G1-envelope}
 \mathfrak r(\lambda)
 :=(1+L)^{3/2}e^{-c_1\sqrt L}+(1+L)e^{-L/2},
 \qquad L=2\log\lambda.
\end{equation}
Combining the pole--prime, archimedean, and decaying-pole pieces with global radicality gives the quantitative continuum G1 estimate
\begin{equation}\label{eq:G1-continuum}
 \boxed{
 \mathcal E_T(\lambda):=
 \sup_{\substack{\|v\|_2=1\\ v\in E_{\lambda,T}}}
 \frac{|QW(k_{\lambda,\mathrm{in}},v)|}{|B_\lambda|}
 \le C_T\mathfrak r(\lambda)\longrightarrow0.}
\end{equation}
The quantitative estimates use only fixed-order PSWF asymptotics, the
radial differential equation, Poisson summation, and the classical
unconditional prime-number-theorem remainder. The global radical and
semilocal domain passage is now Theorem~\ref{thm:v14-radical}; hence
\eqref{eq:G1-continuum} is proved for the explicitly repaired source,
with no remaining source assumption. The construction and all displayed
cutoff rules are independent of a nontrivial-zero list. This is the weak
G1 conclusion for a fixed physical band, not a uniform bound on growing
prolate blocks or an RH implication by itself.

\begin{lemma}[Inversion symmetrization preserves the G1 envelope]\label{lem:g1-inversion}
Let
\[
 p_\lambda=\tfrac12(I+J)k_{\lambda,\mathrm{in}},
 \qquad Jf(u)=f(u^{-1}).
\]
Then
\begin{equation}\label{eq:G1-sym-envelope}
 \sup_{\substack{\|v\|_2=1\\v\in E_{\lambda,T}}}
 \frac{|QW(p_\lambda,v)|}{|B_\lambda|}
 \le \mathcal E_T(\lambda)
 \le C_T\mathfrak r(\lambda).
\end{equation}
\end{lemma}
\begin{proof}
The Weil form is invariant under simultaneous inversion,
$QW(Jf,Jg)=QW(f,g)$, and $J$ is a unitary involution preserving the physical
frequency window $E_{\lambda,T}$; this is the semilocal form of the inversion
symmetry in Connes--Consani--Moscovici~\cite{ConnesConsaniMoscovici2025}.
Hence
\[
 QW(Jk_{\lambda,\mathrm{in}},v)
 =QW(k_{\lambda,\mathrm{in}},Jv).
\]
Taking the average with $QW(k_{\lambda,\mathrm{in}},v)$ and then the supremum
proves \eqref{eq:G1-sym-envelope}.
\end{proof}


\begin{proposition}[Sobolev form of the continuum G1 estimate]\label{prop:g1-sobolev}
Use the explicit repaired source above. Let $u$ be a periodic $H^1$ function on the logarithmic window $[-a,a]$,
$a=\log\lambda$, and interpret it by zero extension when it enters the
semilocal Weil form.  Put
\[
 \mathcal S_1(u):=\|u\|_2+\|D_{\log}u\|_2
 +|u(-a)|+|u(a)|.
\]
Then the same zero-independent envelope as in \eqref{eq:G1-continuum} obeys
\begin{equation}\label{eq:G1-sobolev}
 \boxed{
 \frac{|QW_\lambda(p_\lambda,u)|}{|B_\lambda|}
 \le C\,\mathfrak r(\lambda)\,\mathcal S_1(u).}
\end{equation}
The constant is independent of $\lambda$ and $u$.  In particular the
fixed-physical-band hypothesis in \eqref{eq:G1-sym-envelope} is a convenient
way to bound the Sobolev and trace factors, not an essential spectral
restriction on the continuum G1 pairing.
\end{proposition}
\begin{proof}
Repeat the proof of the cross-tail and pole--prime estimates without applying
a Bernstein inequality.  On the open logarithmic interval, differentiation
of the correlation is placed on $u$, so Cauchy--Schwarz gives the factor
$\|D_{\log}u\|_2$ in place of the earlier band constant $T$.  The two jumps
created by zero extension have coefficients exactly $u(-a)$ and $u(a)$, so
Lemma~\ref{lem:endpoint-jumps} carries over with the endpoint-evaluation norm
replaced by $|u(-a)|+|u(a)|$.  Finally
$|q(0)|\le\|r_\lambda\|_2\|u\|_2$, and the archimedean numerator is
controlled near the origin by the same $H^1$ correlation estimate.  The
Stieltjes integration-by-parts calculation leading to
\eqref{eq:pole-prime-small} is otherwise unchanged and gives
\[
 |QW_\lambda(k_{\lambda,\mathrm{in}},u)|
 \le C|B_\lambda|\mathfrak r(\lambda)\mathcal S_1(u).
\]
Simultaneous inversion preserves the $L^2$ norm, the logarithmic derivative
norm, and the two endpoint traces.  Averaging with the inverted vector as in
Lemma~\ref{lem:g1-inversion} proves \eqref{eq:G1-sobolev}.
\end{proof}

The ordinary Fourier finite-section issue requires more care.  The relevant regularity is supplied directly by the arithmetic finite-sum formula, rather than by a global continuity assumption.

\begin{lemma}[Arithmetic BV structure and absence of a periodic boundary jump]\label{lem:g1-bv-structure}
Let $\widetilde h_\lambda$ be the compactly supported exact-moment prolate source used above, and let
\[
 \mathcal E(\widetilde h_\lambda)(u)
 =u^{1/2}\sum_{n\ge1}\widetilde h_\lambda(nu)
\]
be its co-Poisson/Riemann-sum vector.  On the window
$I_\lambda=[\lambda^{-1},\lambda]$, the restriction of
$\mathcal E(\widetilde h_\lambda)$ is of bounded variation and has only finitely many possible interior jump points.  Consequently
\[
 p_\lambda=\tfrac12(I+J)k_{\lambda,\mathrm{in}},
 \qquad Jf(u)=f(u^{-1}),
\]
is also of bounded variation.  Moreover the inversion identity
$p_\lambda(u)=p_\lambda(u^{-1})$ implies
\begin{equation}\label{eq:g1-endpoint-limits}
 p_\lambda(\lambda^-)=p_\lambda((\lambda^{-1})^+)=:B_\lambda,
\end{equation}
so the periodic representative in logarithmic coordinates has no jump at the identified boundary.  No assertion of global continuity is needed.
\end{lemma}
\begin{proof}
The scale-invariant Riemann-sum formula
$\mathcal E(f)(u)=u^{1/2}\sum_{n>0}f(nu)$ is the standard co-Poisson map used in the prolate construction~\cite{ConnesConsani2023}.  Because
$\operatorname{supp}\widetilde h_\lambda\subset[-\lambda,\lambda]$, for
$u\in[\lambda^{-1},\lambda]$ only indices
$n\le\lambda/u\le\lambda^2$ can occur.  Thus the sum has only finitely many summands throughout the window.  Between the finitely many threshold points at which an argument $nu$ meets an endpoint or another piece boundary of the compact source, it is a finite sum of smooth functions; its one-sided limits are finite.  Hence its total variation on the compact window is finite.  The fixed smooth exact-moment repair used above does not alter this conclusion.

Inversion preserves bounded variation, so the symmetrized proxy $p_\lambda$ is BV.  Finally, by construction $p_\lambda=Jp_\lambda$.  Taking one-sided limits at the two ends of $I_\lambda$ gives \eqref{eq:g1-endpoint-limits}.  After the logarithmic endpoints are identified, this is exactly the absence of a periodic boundary jump.  Interior arithmetic jumps remain allowed.
\end{proof}

Thus for each \emph{fixed} cutoff ordinary Fourier projection converges both in the closed semilocal Weil form norm and at the endpoint.  This gives the following fixed-cutoff existence statement, but no quantitative diagonal compatible with the growing-band estimate is supplied.

\begin{proposition}[Fixed-cutoff boundary-preserving Galerkin approximation]\label{prop:boundary-galerkin}
Assume $p_\lambda$ is of bounded variation on $[-a,a]$ and that its one-sided endpoint limits agree, with common nonzero value
\[
 B_\lambda=p_\lambda(a^-)=p_\lambda((-a)^+).
\]
Let $P_N$ be ordinary Fourier projection onto the trigonometric space $E_N$.  For each fixed $\lambda$,
\[
 P_Np_\lambda\to p_\lambda
\]
in the form norm of the closed semilocal Weil form and
\[
 (P_Np_\lambda)(a)\to B_\lambda.
\]
Consequently, for every fixed $\lambda$ and every $\varepsilon>0$ there is an $N=N(\lambda,\varepsilon)$ such that
\[
 \|P_Np_\lambda-p_\lambda\|_{QW_\lambda}
 +|(P_Np_\lambda)(a)-B_\lambda|
 \le \varepsilon |B_\lambda|.
\]
No rate relating $N$ to $\lambda$ is asserted.
\end{proposition}
\begin{proof}
A bounded-variation function on the log interval has Fourier coefficients $\widehat p_\lambda(n)=O_\lambda(|n|^{-1})$.  The archimedean part of the semilocal Weil form has Fourier weight comparable to $1+\log(1+n^2)$, while at fixed $\lambda$ the pole term and the finite prime-power sum are bounded perturbations.  Hence
\[
 \sum_n(1+\log(1+n^2))|\widehat p_\lambda(n)|^2<\infty,
\]
and the Fourier tails tend to zero in the form norm.  Since the two one-sided endpoint limits agree, the periodic representative has no boundary jump; the Dirichlet--Jordan theorem therefore gives convergence at the identified endpoint.  Interior jumps do not affect this conclusion.
\end{proof}

The lack of a rate is not cosmetic.  The physical-band estimate proved above remains uniform only while the physical frequency grows sufficiently slowly compared with the quantitative PNT saving.  A generic Fourier approximation of a function with an exponentially small endpoint value may require a much larger cutoff.  Therefore Proposition~\ref{prop:boundary-galerkin} by itself does \emph{not} transfer \eqref{eq:G1-continuum} to a diagonal $N(\lambda)$.

There is nevertheless a zero-independent finite-dimensional transfer of the
\emph{fixed-band residual}, and the quantitative envelope above makes its exact
finite-section dependence transparent.  For $N$ large enough to contain the
physical band, set
\[
 k_{\lambda,N}=P_Np_\lambda,
\]
and define the relative boundary and band-form defects
\begin{align}
 \beta_{\lambda,N}
 &:=\frac{|\delta_N(k_{\lambda,N})-B_\lambda|}{|B_\lambda|},
 \label{eq:g1-beta}\\
 \gamma_{\lambda,N,T}
 &:=\frac1{|B_\lambda|}
 \sup_{\substack{v\in E_{\lambda,T}\\\|v\|_2=1}}
 |QW_\lambda(k_{\lambda,N}-p_\lambda,D_{\log}v)|.
 \label{eq:g1-gamma}
\end{align}
Both quantities are defined entirely from the prolate/co-Poisson proxy and the
finite Weil form; no zeta zero is used.

\begin{proposition}[Quantitative and cofinal finite-section transfer of G1]\label{prop:finite-g1-transfer}
Fix $T>0$.  Whenever $\beta_{\lambda,N}<1$,
\begin{equation}\label{eq:g1-master-bound}
 \boxed{
 \frac{\|P_TD_{\log}W_{\lambda,N}k_{\lambda,N}\|_2}
 {|\delta_N(k_{\lambda,N})|}
 \le
 \frac{T\mathcal E_T(\lambda)+\gamma_{\lambda,N,T}}
 {1-\beta_{\lambda,N}}.}
\end{equation}
For each fixed $\lambda$,
\begin{equation}\label{eq:g1-defects-fixed-lambda}
 \beta_{\lambda,N}\longrightarrow0,
 \qquad
 \gamma_{\lambda,N,T}\longrightarrow0
 \qquad(N\to\infty).
\end{equation}
Consequently, for all sufficiently large $\lambda$ there is a zero-independent
threshold $N_{G1}(\lambda,T)$ such that for \emph{every}
$N\ge N_{G1}(\lambda,T)$,
\begin{align}
 \beta_{\lambda,N}&\le\mathfrak r(\lambda),
 \label{eq:g1-threshold-beta}\\
 \gamma_{\lambda,N,T}&\le\mathfrak r(\lambda),
 \label{eq:g1-threshold-gamma}
\end{align}
and hence
\begin{equation}\label{eq:finite-g1-transfer}
 \boxed{
 \sup_{N\ge N_{G1}(\lambda,T)}
 \frac{\|P_TD_{\log}W_{\lambda,N}k_{\lambda,N}\|_2}
 {|\delta_N(k_{\lambda,N})|}
 \le C'_T\mathfrak r(\lambda)\longrightarrow0.}
\end{equation}
In particular one may choose any diagonal $N(\lambda)\ge N_{G1}(\lambda,T)$;
G1 does not require an economical or finely tuned Fourier cutoff.
\end{proposition}
\begin{proof}
Hermiticity of the finite Weil matrix and self-adjointness of $D_{\log}$ on the
periodic Fourier section give
\[
 \|P_TD_{\log}W_{\lambda,N}k_{\lambda,N}\|_2
 =\sup_{\substack{v\in E_{\lambda,T}\\\|v\|_2=1}}
 |QW_\lambda(k_{\lambda,N},D_{\log}v)|.
\]
Split $k_{\lambda,N}=p_\lambda+(k_{\lambda,N}-p_\lambda)$.  Since
$D_{\log}v\in E_{\lambda,T}$ and
$\|D_{\log}v\|_2\le T$, Lemma~\ref{lem:g1-inversion} bounds the first part by
$T|B_\lambda|\mathcal E_T(\lambda)$, while the second is exactly bounded by
$|B_\lambda|\gamma_{\lambda,N,T}$.  Also
\[
 |\delta_N(k_{\lambda,N})|
 \ge |B_\lambda|(1-\beta_{\lambda,N}).
\]
This proves \eqref{eq:g1-master-bound}.

For fixed $\lambda$, Proposition~\ref{prop:boundary-galerkin} gives the boundary
convergence and form-norm convergence of $P_Np_\lambda-p_\lambda$, proving the
first limit in \eqref{eq:g1-defects-fixed-lambda}.  To obtain the second without
an unspecified pointwise continuity constant, choose $m_\lambda$ so that
\[
 QW_\lambda(u,u)+m_\lambda\|u\|_2^2
\]
is a positive closed form.  Its Cauchy--Schwarz inequality, together with the
finite-dimensional boundedness of
$\{D_{\log}v:v\in E_{\lambda,T},\ \|v\|_2=1\}$ in the corresponding form norm,
shows that the supremum in \eqref{eq:g1-gamma} tends to zero with the form norm
of $P_Np_\lambda-p_\lambda$.

The convergence in \eqref{eq:g1-defects-fixed-lambda} is convergence of the
full sequence, so for each fixed $\lambda$ and the positive tolerance
$\mathfrak r(\lambda)$ there is a threshold after which
\eqref{eq:g1-threshold-beta}--\eqref{eq:g1-threshold-gamma} hold for all larger
$N$.  Since $\mathfrak r(\lambda)\to0$, it is $<1/2$ for large $\lambda$.
Insert these inequalities and
$\mathcal E_T(\lambda)\le C_T\mathfrak r(\lambda)$ into
\eqref{eq:g1-master-bound} to obtain \eqref{eq:finite-g1-transfer}.
\end{proof}


\subsection*{Quantitative high--low decay and an exact-boundary G1 repair}
Endpoint recovery of the ordinary Fourier partial sums in
Proposition~\ref{prop:finite-g1-transfer} was initially nonquantitative.
Proposition~\ref{prop:v114-polynomial-endpoint} below now supplies an
asymptotic polynomial cutoff for the fixed source. For the G1
residual itself one can also bypass endpoint selection.  The key observation is that the
explicit CCM matrix has an additional off-diagonal factor, so only bounded
variation of the co-Poisson proxy is needed.

Write $U_n=L^{-1/2}e^{2\pi inx/L}$ for the normalized logarithmic Fourier basis and $V_n=\kappa(U_n)$ for its multiplicative image, and set
\[
 \tau^{(\lambda)}_{n,m}:=QW_\lambda(V_n,V_m).
\]

\begin{lemma}[Zero-independent high--low decay of the Weil matrix]\label{lem:g1-high-low}
There is an absolute constant $C$ such that, for $\lambda\ge2$,
$L=2\log\lambda$, and $n\ne m$,
\begin{equation}\label{eq:g1-high-low}
 \boxed{
 |\tau^{(\lambda)}_{n,m}|
 \le C\frac{\lambda}{|n-m|}.}
\end{equation}
The constant uses no information about the nontrivial zeros of $\zeta$.
\end{lemma}
\begin{proof}
For $n\ne m$, the explicit convolution kernel of
Connes--Consani--Moscovici~\cite{ConnesConsaniMoscovici2025} is
\[
 q(U_n,U_m)(y)
 =\frac{\sin(2\pi my/L)-\sin(2\pi ny/L)}{\pi(n-m)},
 \qquad 0\le y\le L,
\]
so $q(U_n,U_m)(0)=0$ and
$\|q(U_n,U_m)\|_\infty\le 2/(\pi|n-m|)$.  The pole contribution is therefore
bounded by
\[
 \frac{C}{|n-m|}\int_0^L(e^{y/2}+e^{-y/2})\,dy
 \le C\frac{\lambda}{|n-m|}.
\]
For the nonarchimedean contribution, partial summation together with the
classical Chebyshev bound $\psi(x)\ll x$ gives
\[
 \sum_{1<k\le e^L}\frac{\Lambda(k)}{\sqrt{k}}
 \ll e^{L/2}=\lambda,
\]
and hence the same bound.

For the archimedean term set
\[
 \rho(y)=\frac{e^{y/2}}{e^y-e^{-y}},\qquad
 J_L(r)=\int_0^L\sin(2\pi r y/L)\rho(y)\,dy.
\]
Because the off-diagonal kernel vanishes at the origin, the CCM formula reduces
to
\[
 W_{\mathbb R}(V_n,V_m)
 =\frac{J_L(m)-J_L(n)}{\pi(n-m)}.
\]
Uniformly for $L\ge2\log2$ and integers $r$, $J_L(r)=O(1)$: on $0<y\le1$,
write $\rho(y)=(2y)^{-1}+O(1)$ and use the uniform boundedness of the sine
integral, while on $y\ge1$ use $\rho(y)\ll e^{-y/2}$.  Thus the archimedean
piece is $O(|n-m|^{-1})$.  Combining the three contributions proves
\eqref{eq:g1-high-low}.
\end{proof}

\begin{corollary}[Complete Fourier-matrix bound]\label{cor:g1-complete-matrix}
For all integers $n,m$ and $\lambda\ge2$,
\begin{equation}\label{eq:g1-complete-matrix}
 \boxed{
 |\tau^{(\lambda)}_{n,m}|
 \le
 \begin{cases}
 C\lambda/|n-m|,&n\ne m,\\[2mm]
 C\bigl(\lambda+1+\log(2+|n|/L)\bigr),&n=m.
 \end{cases}}
\end{equation}
The constant is absolute and uses no nontrivial-zero information.
\end{corollary}
\begin{proof}
The off-diagonal case is Lemma~\ref{lem:g1-high-low}.  The diagonal is a
separate CCM formula, not the limit of the off-diagonal divided difference:
for $0\le y\le L$,
\[
 q(U_n,U_n)(y)=2\left(1-\frac yL\right)
 \cos\frac{2\pi ny}{L},
 \qquad q(U_n,U_n)(0)=2.
\]
Thus the pole term is $O(\lambda)$ and, by the classical Chebyshev estimate,
the prime-power term is also $O(\lambda)$.

It remains to control the archimedean term.  Put $t_n=2\pi n/L$.  On
$0<x\le1$ the numerator in the logarithmic Weil kernel is, up to an absolute
constant,
\[
 \left(1-\frac xL\right)\cos(t_nx)-e^{-x/2}.
\]
Since $\rho(x)=e^{x/2}/(e^x-e^{-x})\ll x^{-1}$ there,
\[
 \left|\left(1-\frac xL\right)\cos(t_nx)-e^{-x/2}\right|
 \le C\{\min(1,t_n^2x^2)+x\}.
\]
Consequently
\[
 \int_0^1
 \left|\left(1-\frac xL\right)\cos(t_nx)-e^{-x/2}\right|
 \rho(x)\,dx
 \le C\bigl(1+\log(2+|t_n|)\bigr).
\]
For $x\ge1$ the kernel is $O(e^{-x/2})$ and the diagonal convolution is
bounded, so the remaining integral is $O(1)$.  Therefore the archimedean
contribution is
$O(1+\log(2+|n|/L))$, proving \eqref{eq:g1-complete-matrix}.
\end{proof}

\subsection*{A closed Weil operator and its omitted Fourier complement}
The entrywise bound above can be strengthened to an operator statement.
Here the closed operator is the canonical semilocal Weil operator; it is
not the arithmetic generator whose realization is conditional in Section~18.
Fix $\lambda>1$, put $L=2\log\lambda$, $t_n=2\pi n/L$, and write
\[
 P_\lambda=\sum_{1<m\le\lambda^2}\frac{\Lambda(m)}{\sqrt m},
 \qquad \rho(y)=\frac{e^{y/2}}{2\sinh y},
 \qquad h_L=\sinh^2(L/4).
\]
For later constants use the complete archimedean diagonal
\begin{align}
 A_n={}&\int_0^L\rho(y)
 \{2(1-y/L)\cos(t_ny)-2e^{-y/2}\}\,dy \notag\\
 &+\log(4\pi)+\gamma_{\rm E}+\log\tanh(L/2).
 \label{eq:v114-arch-diagonal}
\end{align}
The cancellation at zero makes this an ordinary convergent integral.

\begin{proposition}[Bounded commutator and the operator core]
\label{prop:v114-weil-core}
In the normalized logarithmic Fourier basis the canonical semilocal
Weil operator $\mathsf W_\lambda$ is
\begin{equation}\label{eq:v114-weil-decomposition}
 \mathsf W_\lambda=D_d+[M_b,H],\qquad
 H_{nm}=\begin{cases}(n-m)^{-1},&n\ne m,\\0,&n=m,\end{cases}
\end{equation}
where $D_d$ is real diagonal, $M_b$ is multiplication by the real odd
sequence
\begin{align}
 b_n={}&\frac{32Lh_Ln}{L^2+16\pi^2n^2}
 +\frac1\pi\left\{\int_0^L\rho(y)\sin(t_ny)\,dy
 +\sum_{1<m\le\lambda^2}\frac{\Lambda(m)}{\sqrt m}
                         \sin(t_n\log m)\right\},\notag\\
 d_n={}&\frac{32Lh_L(L^2-16\pi^2n^2)}{(L^2+16\pi^2n^2)^2}
 -2\sum_{1<m\le\lambda^2}\frac{\Lambda(m)}{\sqrt m}
        (1-\log m/L)\cos(t_n\log m)-A_n.
 \label{eq:v114-b-d}
\end{align}
In particular the diagonal is not filled in by a divided-difference limit.
One has
\begin{equation}\label{eq:v114-b-bound}
 \|b\|_\infty\le B_\lambda^{\rm mat}
 :=\frac{4h_L+P_\lambda+1+\pi/4}{\pi},\qquad
 \|[M_b,H]\|\le2\pi B_\lambda^{\rm mat}.
\end{equation}
At fixed $\lambda$, $d_n=\log|n|+O_\lambda(1)$ as $|n|\to\infty$.
Consequently
\begin{align}
 \mathcal D(\mathsf W_\lambda)
 &=\left\{x:\sum_n\log^2(2+|n|)|x_n|^2<\infty\right\},\notag\\
 \mathcal D[\mathsf W_\lambda]
 &=\left\{x:\sum_n\log(2+|n|)|x_n|^2<\infty\right\}.
 \label{eq:v114-weil-domains}
\end{align}
Finite Fourier sums form an operator core. The operator is lower bounded
and has compact resolvent.
\end{proposition}
\begin{proof}
The exact CCM entries give $\tau_{nm}=(b_n-b_m)/(n-m)$ for $n\ne m$
and $\tau_{nn}=d_n$; see also \eqref{eq:beta-bcoeff-direct}.
For $n>0$, expansion $\rho(y)=\sum_{k\ge0}e^{-(2k+1/2)y}$ and
$t_nL=2\pi n$ give
\[
 S_n:=\int_0^L\rho(y)\sin(t_ny)\,dy
 =\sum_{k\ge0}\frac{(1-e^{-(2k+1/2)L})t_n}
                         {(2k+1/2)^2+t_n^2}.
\]
Interchange is justified by integrability of
$\rho(y)|\sin(t_ny)|$. The summands are nonnegative. The $k=0$
term is at most $1$. For $k\ge1$ discard the factor
$1-e^{-(2k+1/2)L}\le1$ and compare the decreasing majorant
$t_n/\{(2k+1/2)^2+t_n^2\}$ with an integral; this gives at most
$\pi/4$. Thus
$0\le S_n\le1+\pi/4$. Oddness treats negative $n$, and $S_0=0$.
The pole part of $b_n$ has absolute value at most $4h_L/\pi$;
the prime sine sum is bounded by $P_\lambda$. This proves the first
bound in \eqref{eq:v114-b-bound}.

On $\ell^2(\mathbb Z)$, $H$ is the discrete Hilbert transform.
The Fourier series $\sum_{k\ne0}e^{ik\theta}/k=i(\pi-\theta)$
for $0<\theta<2\pi$ proves $H^*=-H$ and $\|H\|=\pi$.
Hence $[M_b,H]$ is bounded self-adjoint and has the asserted norm.
Moreover
\[
 A_n=A_0-\int_0^L w_L(y)(1-\cos(t_ny))\,dy,
 \qquad w_L(y)=2\rho(y)(1-y/L).
\]
Since $w_L(y)-1/y$ is integrable on $(0,L)$, the last integral is
$\log|t_n|+O_L(1)$. For example, split
$\int_0^{|t_n|L}(1-\cos v)\,dv/v$ at $1$ and integrate its
oscillatory tail by parts. The other terms in $d_n$ are bounded at
fixed $\lambda$, proving its logarithmic asymptotic.

The diagonal operator is self-adjoint on the first domain in
\eqref{eq:v114-weil-domains}, has finite sequences as an operator
core, and has compact resolvent. Bounded self-adjoint perturbation
preserves these statements and gives the displayed form domain.
This identifies the actual physical Weil realization: the finite sums
are a form core for its lower-bounded closed form by
\cite[Propositions~3.2--3.4]{ConnesConsaniMoscovici2025}, and both
forms have exactly the same entries there. Uniqueness of the form
closure proves equality, without imposing a new boundary condition.
\end{proof}

Every periodic BV proxy, including the repaired $p_\lambda$, therefore
belongs to $\mathcal D(\mathsf W_\lambda)$ at fixed $\lambda$:
its Fourier coefficients are $O_\lambda((1+|n|)^{-1})$.
Furthermore $P_Np_\lambda\to p_\lambda$ in this operator graph norm,
by the diagonal description and boundedness of the commutator. This
fixed-window statement supplies no endpoint-normalized uniform rate,
and does not remove the separate sampler graph defect.

\begin{proposition}[An explicit positive high-frequency tail]
\label{prop:v114-weil-tail}
Set $a=\min(1,L)$ and
\begin{align}
 C_\lambda^{\rm tail}
 :={}&|A_0|+4+\pi/2+a+2a/L+2P_\lambda
                   +2\sinh(L/2)-L,\notag\\
 \Gamma_{\lambda,N}
 :={}&\log\frac{2\pi(N+1)a}{L}-C_\lambda^{\rm tail}.
 \label{eq:v114-tail-constant}
\end{align}
For every vector $f$ in the form domain supported on Fourier indices
$|n|>N$,
\begin{equation}\label{eq:v114-tail-positive}
 QW_\lambda(f,f)\ge\Gamma_{\lambda,N}\|f\|_2^2.
\end{equation}
In particular this omitted complement is strictly positive whenever
$\Gamma_{\lambda,N}>0$. The constants use no nontrivial-zero data.
\end{proposition}
\begin{proof}
On $0<y<a$, $2\rho(y)\ge e^{-y/2}/y\ge1/y-1/2$, whence
\[
 w_L(y)\ge1/y-(1/2+1/L).
\]
Since $w_L\ge0$ on the full interval and
$F(T):=\int_0^T(1-\cos v)\,dv/v\ge\log T-2$ for $T>0$,
\[
 -A_n\ge\log(|t_n|a)-|A_0|-2-a-2a/L\qquad(n\ne0).
\]
The bound for $F$ is immediate for $T\le1$; for $T\ge1$ the
integral of $\cos v/v$ on $[1,T]$ has absolute value at most $2$
by integration by parts. The off-diagonal archimedean operator has
norm at most $2+\pi/2$, using $b_n^{\rm arch}=S_n/\pi$ in the
preceding commutator proof.

The complete prime operator is a sum of negative weighted truncated
translations and their adjoints, so its lower bound is $-2P_\lambda I$.
In centered logarithmic coordinates $x\in[-L/2,L/2]$, the pole form is
\[
 2|\ip f{\cosh(x/2)}|^2-2|\ip f{\sinh(x/2)}|^2
 \ge-\{2\sinh(L/2)-L\}\|f\|_2^2.
\]
Combining these three estimates gives
\eqref{eq:v114-tail-positive} first for finite Fourier sums and then
on the tail form domain by closure.
\end{proof}

\subsection*{A sharper positive Fourier complement at the certified window}
The previous bound loses constants by bounding each place on the full
space. On the omitted Fourier tail one can retain the limiting
archimedean commutator, the physical incidence of prime shifts, and the
decay of the negative pole vector. These are bounds for the same Weil
operator and its original Fourier projection.

For $L=2\log\lambda$ put
\begin{align}
 R_L&=\frac{e^{-L/2}}{1-e^{-2L}},&
 C_L&=\max\{1,1/4+R_L\},\notag\\
 E_L(t)&=\frac1{15t}+\frac7{2t^2}
       +\frac\pi{2Lt}+\frac{2+2R_L}{Lt^2},\qquad t>0.
 \label{eq:v115-arch-errors}
\end{align}
Define the weighted prime degree
\begin{equation}\label{eq:v115-prime-degree}
 M_\lambda=\mathop{\rm ess\,sup}_{0<x<L}
 \left\{\sum_{\substack{1<m<\lambda^2\\\log m\le x}}
                    \frac{\Lambda(m)}{\sqrt m}
       +\sum_{\substack{1<m<\lambda^2\\\log m\le L-x}}
                    \frac{\Lambda(m)}{\sqrt m}\right\}.
\end{equation}
The strict upper cut omits only a possible full-length translation,
which is zero on $L^2([0,L])$.

\begin{proposition}[A practical positive omitted complement]
\label{prop:v115-sharp-tail}
For $N\ge1$ let $t_*=2\pi(N+1)/L$, and set
\begin{align}
 \Gamma_{\lambda,N}^{\rm even}
   &=\log\frac{N+1}{L}-M_\lambda
                -\frac{2C_L}{t_*}-E_L(t_*),\notag\\
 \Gamma_{\lambda,N}^{\rm all}
   &=\Gamma_{\lambda,N}^{\rm even}-\frac\pi2
                   -\frac{4\sinh^2(L/4)L}{\pi^2N}.
 \label{eq:v115-sharp-tail}
\end{align}
For every form-domain vector $f$ whose Fourier coefficients vanish on
$|n|\le N$,
\[
 QW_\lambda(f,f)\ge\Gamma_{\lambda,N}^{\rm all}\|f\|_2^2.
\]
For an inversion-even $f$, the stronger lower bound with
$\Gamma_{\lambda,N}^{\rm even}$ holds. In particular, at the actual
window $\lambda=3$,
\begin{equation}\label{eq:v115-lambda3-tail}
 \Gamma_{3,256}^{\rm all}>0.0148,\qquad
 \Gamma_{3,256}^{\rm even}>1.5867,\qquad
 \Gamma_{3,512}^{\rm all}>0.7085.
\end{equation}
These inequalities control the entire omitted complement, not merely a
finite further block.
\end{proposition}
\begin{proof}
Write $a_k=2k+1/2$ and use the exact Fourier-grid series for $S_n$
in the proof of Proposition~\ref{prop:v114-weil-core}.
For $t>0$, the function $x\mapsto t/(x^2+t^2)$ decreases on $x>0$.
Comparing its values on $1/2+2\mathbb N$ with its integral gives
\[
 \frac\pi4-\frac1{4t}
 \le\sum_{k\ge0}\frac{t}{a_k^2+t^2}
 \le\frac\pi4+\frac1t.
\]
The exponentially weighted part of $S_n$ is nonnegative and at most
$R_L/|t_n|$. Oddness therefore proves
\begin{equation}\label{eq:v115-S-sharp}
 \left|S_n-\frac\pi4\operatorname{sgn}n\right|
       \le\frac{C_L}{|t_n|}\qquad(n\ne0).
\end{equation}
On the tail, the archimedean off-diagonal multiplier is consequently
$b_n^{\rm arch}=\operatorname{sgn}(n)/4+e_n$ with
$\|e\|_\infty\le C_L/(\pi t_*)$.
Since the compressed discrete Hilbert transform has norm at most $\pi$,
its commutator has norm at most $\pi/2+2C_L/t_*$.
For even coefficients the limiting commutator is instead positive:
under the normalized identification with positive indices its matrix is
\[
 K_{nm}=\frac1{2(n+m)},\qquad n,m>N.
\]
Positivity follows from
$1/(n+m)=\int_0^1x^{n+m-1}\,dx$. Thus only $2C_L/t_*$ is lost in
the even sector. This uses the unshifted Fourier basis of
\eqref{eq:v114-weil-decomposition}; translation to centered physical
coordinates conjugates by $(-1)^n$ and preserves both parity and the
positive quadratic form.

Next retain the complete archimedean diagonal. Its exact series, already
used in the independent matrix certificate, gives for
$z=1/4+it/2$ and $t=|t_n|>0$,
\begin{equation}\label{eq:v115-arch-exact}
 -A_n=\Re\psi_0(z)-\log\pi+\frac{\Re\psi_1(z)}{2L}
 -\frac2L\sum_{k\ge0}e^{-a_kL}
                   \frac{a_k^2-t^2}{(a_k^2+t^2)^2}.
\end{equation}
Here $\psi_0$ and $\psi_1$ are digamma and trigamma.
Binet's exact integral formula~\cite[eq.~5.9.15]{NISTDigamma} at
$w=z+1=5/4+it/2$ has remainder of modulus at most
\[
 2\int_0^\infty\frac{u\,du}
 {|u^2+w^2|(e^{2\pi u}-1)}
 \le\frac8{5t}\int_0^\infty\frac{u\,du}{e^{2\pi u}-1}
 =\frac1{15t},
\]
since $|u^2+w^2|\ge5t/4$. The recurrence
$\psi_0(z)=\psi_0(z+1)-1/z$, together with
$\Re\log w\ge\log(t/2)$,
$\Re(1/(2w))\le5/(2t^2)$, and $\Re(1/z)\le1/t^2$, yields
\[
 \Re\psi_0(z)\ge\log(t/2)-\frac1{15t}-\frac7{2t^2}.
\]
The convergent trigamma series and a decreasing-integrand comparison give
\[
 |\psi_1(z)|\le\sum_{k\ge0}\frac1{(k+1/4)^2+t^2/4}
                 \le\frac4{t^2}+\frac\pi t.
\]
The last term in \eqref{eq:v115-arch-exact} has modulus at most
$2R_L/(Lt^2)$. Hence
\begin{equation}\label{eq:v115-diagonal-sharp}
 -A_n\ge\log\frac{|n|}{L}-E_L(|t_n|).
\end{equation}
Its right side increases with $|n|$, so $|n|=N+1$ supplies the uniform
tail bound. No logarithmic constant or finite-$L$ diagonal term was
discarded from the exact identity.

For the prime part, set $w_m=\Lambda(m)/\sqrt m$ and $a_m=\log m$.
The elementary inequality $2\Re(z\overline w)\le|z|^2+|w|^2$ gives
\[
 -2\sum_m w_m\Re\int_0^{L-a_m}
      f(x+a_m)\overline{f(x)}\,dx
 \ge-\int_0^L d_\lambda(x)|f(x)|^2dx
 \ge-M_\lambda\|f\|_2^2,
\]
where $d_\lambda$ is the function inside braces in
\eqref{eq:v115-prime-degree}. This applies equally to complex vectors.

In centered physical coordinates the negative pole term is
$-2|\ip f{\sinh(x/2)}|^2$. Direct Fourier integration gives
\[
 |\widehat{\sinh(x/2)}(n)|^2
 =\frac{4\sinh^2(L/4)t_n^2}{L(t_n^2+1/4)^2}.
\]
Consequently
\[
 2\|Q_N\sinh(x/2)\|_2^2
 \le\frac{4\sinh^2(L/4)L}{\pi^2}
                         \sum_{n>N}n^{-2}
 \le\frac{4\sinh^2(L/4)L}{\pi^2N}.
\]
It vanishes identically on the even sector; the positive pole term can
be retained or omitted in a lower bound. Combining these estimates
proves \eqref{eq:v115-sharp-tail} first on finite Fourier sums and then
on the tail form domain by Proposition~\ref{prop:v114-weil-core}.

For the numerical constants at $\lambda=3$, symmetry reduces the prime
degree to $0<x<\log3$. Its breakpoints, expressed in $e^x$, are
$9/8,9/7,9/5,2,9/4$. With
$P=w_2+w_3+w_4+w_5+w_7+w_8$, the six interval values are
\[
 P,\quad P-w_8,\quad P-w_8-w_7,\quad
 w_2+w_3+w_4,\quad2w_2+w_3+w_4,\quad2w_2+w_3.
\]
All are at most $P$, using $w_5>w_2$ for the fifth value, and the first
occurs on an interval of positive length. Thus $M_3=P$ exactly.
Also $\sinh^2(L/4)=1/3$, $R_L=27/80$ and $C_L=1$.
Outward interval evaluation of the displayed elementary expressions
gives \eqref{eq:v115-lambda3-tail}; the reproducibility bundle contains
the 256-bit scalar certificate. These are not new finite-matrix
eigenvalue computations.
\end{proof}

\subsection*{A positive weight for the prime shifts and a diagonal inverse bound}
The flat degree bound can be sharpened without changing the Weil operator.
Let $\mathsf T_{\rm pr}$ denote the sum of the truncated translations and
their adjoints with weights $w_m=\Lambda(m)/\sqrt m$. Its coefficients
are nonnegative, and the prime contribution to the Weil form is
$-\mathsf T_{\rm pr}$. For every bounded function $\phi$ bounded away
from zero and every complex $f$,
\begin{equation}\label{eq:v116-weighted-prime-general}
 |\ip{\mathsf T_{\rm pr}f}{f}|
 \le\int_0^L\frac{(\mathsf T_{\rm pr}\phi)(x)}{\phi(x)}|f(x)|^2\,dx
 \le M_\phi\norm f_2^2,
 \qquad
 M_\phi:=\mathop{\rm ess\,sup}_{0<x<L}
       \frac{(\mathsf T_{\rm pr}\phi)(x)}{\phi(x)}.
\end{equation}
Here $\phi$ is required to be strictly positive. This follows edge by
edge from
\[
 2|f(x+s)f(x)|\le
 \frac{\phi(x)}{\phi(x+s)}|f(x+s)|^2
 +\frac{\phi(x+s)}{\phi(x)}|f(x)|^2.
\]
Thus the argument is a weighted Schur bound on the translation sum;
it assumes no sign for the complete Weil form.

\begin{proposition}[A certified weighted prime bound at $\lambda=3$]
\label{prop:v116-weighted-prime}
At $L=2\log3$, take $\phi(x)=\cosh(x-\log3)$. With
$w_m=\Lambda(m)/\sqrt m$, put
\begin{equation}\label{eq:v116-prime-constant}
 \mathfrak m_3=
 \frac{85}{116}w_2+\frac{65}{87}w_3+
 \frac{193}{232}w_4+\frac{137}{145}w_5+
 \frac{35}{29}w_7.
\end{equation}
Then $M_\phi=\mathfrak m_3$ and
\[
 2.6890557719319011796<\mathfrak m_3
       <2.6890557719319011797.
\]
In particular $\norm{\mathsf T_{\rm pr}}\le\mathfrak m_3$.
Replacing $M_3$ by $\mathfrak m_3$ in
Proposition~\ref{prop:v115-sharp-tail} gives
\begin{equation}\label{eq:v116-weighted-tail-values}
 \widetilde\Gamma_{3,256}^{\rm all}>0.4970,
 \qquad \widetilde\Gamma_{3,256}^{\rm even}>2.0690,
 \qquad \widetilde\Gamma_{3,512}^{\rm all}>1.1907.
\end{equation}
These are omitted-complement bounds, with no assertion yet about the
sign of the head/complement Schur matrix.
\end{proposition}
\begin{proof}
Write $u=e^x\in[1,9]$. The ratio for a forward shift by $\log m$ is
\[
 \frac{m u^2+9/m}{u^2+9},
\]
and for its backward shift it is $(u^2/m+9m)/(u^2+9)$.
On an interval with fixed active shifts, their weighted sum has the
form $A+B\tanh(x-\log3)$ and is monotone or constant. The endpoints
are $u=1,9,m,9/m$, for $m=2,3,4,5,7,8$.
By reflection it suffices to examine the six intervals in $1<u<3$,
whose successive upper endpoints are
\[
 9/8,\quad9/7,\quad9/5,\quad2,\quad9/4,\quad3.
\]
On each such interval the ratio is increasing. Its largest one-sided
value is approached from $u<9/7$ and is
\[
 \sum_{m=2,3,4,5,7}
       w_m\frac{9m^2+49}{58m}=\mathfrak m_3.
\]
The finite comparisons with the other endpoint values use rational
coefficients and are certified at 256-bit Arb precision by
\texttt{certify\_weighted\_prime\_tail.py}; its exact coefficient lists
and interval enclosures are included in the reproducibility bundle.
The $m=8$ term is retained on every interval where that shift is active.
The $m=9$ translation has full interval length and is zero in $L^2$.
This proves the asserted essential supremum and norm bound by
\eqref{eq:v116-weighted-prime-general}. Substitution in the explicit
formulas of Proposition~\ref{prop:v115-sharp-tail}, with
$h_L=1/3$, $R_L=27/80$ and $C_L=1$, gives
\eqref{eq:v116-weighted-tail-values}; the same scalar interval
certificate verifies the strict inequalities.
\end{proof}

\begin{proposition}[Diagonal domination of the tail inverse]
\label{prop:v116-diagonal-inverse}
Continue at $\lambda=3$ and retain $E_L$ and $C_L$ from
\eqref{eq:v115-arch-errors}. Fix $N\ge1$ and set
$t_*=2\pi(N+1)/L$. For the even sector define
\[
 \kappa_{\rm even}=\mathfrak m_3+2C_L/t_*;
\]
for the whole space, and hence also the odd sector, define
\[
 \kappa_{\rm all}=\kappa_{\rm even}+\pi/2
                         +4h_LL/(\pi^2N).
\]
In the selected sector put
\begin{equation}\label{eq:v116-diagonal-weight}
 g_n=\log\frac{|n|}{L}-E_L(2\pi|n|/L)-\kappa,
            \qquad |n|>N.
\end{equation}
If $g_{N+1}>0$, the canonical omitted-tail operator $T$ satisfies
the form inequality $T\succeq D_g$, where $D_g$ is this diagonal
operator, and therefore
\begin{equation}\label{eq:v116-inverse-order}
 T^{-1}\preceq D_g^{-1}.
\end{equation}
The same weights apply in the normalized positive-index parity basis.
\end{proposition}
\begin{proof}
The proof of Proposition~\ref{prop:v115-sharp-tail} bounds the
archimedean diagonal at each index by
$\log(|n|/L)-E_L(2\pi|n|/L)$ before any infimum is taken.
The remaining archimedean error has norm at most $2C_L/t_*$.
Its limiting commutator is nonnegative in the even sector and is
bounded below by $-\pi/2$ on the whole tail. Apply
Proposition~\ref{prop:v116-weighted-prime} to the prime term and the
existing negative-pole tail estimate to the whole-space case.
These bounds prove $T\succeq D_g$ on finite tail sums and hence on
the common logarithmic form domain. The weights $g_n$ increase with
$|n|$, since $E_L(t)$ is decreasing, so $D_g\succeq g_{N+1}I>0$.
The inverse inequality follows, for example, from the variational
identity
\[
 \ip{T^{-1}v}{v}
 =\sup_w\{2\Re\ip vw-\mathfrak t[w]\}
 \le\sup_w\left\{2\Re\ip vw-\sum_n g_n|w_n|^2\right\}
 =\sum_n\frac{|v_n|^2}{g_n}.
\]
The quadratic quantities are real in the fixed first-slot-linear
convention. Restriction to a parity sector is isometric and leaves
the diagonal weights unchanged.
\end{proof}


The preceding tail estimates alone do not establish full Weil positivity.
The complete signed certificate below also controls the effective head.
Its absolute arithmetic bound still gives a sufficient cutoff
exponential in $O(\lambda)$; it does not reach the natural prolate
scale. For the signed correction use $Q_N=I-P_N$ and let
$\Gamma_{\lambda,N}>0$ denote any of the proved tail bounds. Write
\[
 \mathsf W_\lambda=
 \begin{pmatrix}F&B^*\\B&T\end{pmatrix},\qquad
 F=P_N\mathsf W_\lambda P_N,\quad
 B=Q_N\mathsf W_\lambda P_N,\quad T\succeq\Gamma_{\lambda,N}I.
\]
The cross operator $B$ is bounded by \eqref{eq:v114-weil-decomposition}.
The remaining exact sign problem is the finite matrix
$K=F-B^*T^{-1}B$. There is a useful verified-solve alternative to
bounding $B$ alone. For any $Z:E_N\to\mathcal D(T)$ put
$R=B-TZ$ and $K_Z=F-B^*Z-Z^*B+Z^*TZ$. Then
\begin{equation}\label{eq:v114-schur-residual}
 K=K_Z-R^*T^{-1}R
 \succeq K_Z-\Gamma_{\lambda,N}^{-1}R^*R
 \succeq K_Z-\frac{\|R\|^2}{\Gamma_{\lambda,N}}I.
\end{equation}
Expansion proves the identity and $T^{-1}\preceq\Gamma_{\lambda,N}^{-1}I$
proves both bounds. Retaining the residual Gram matrix $R^*R$ preserves
directional cancellation lost by its scalar-norm majorant.
If the final finite lower bound is
$-\varepsilon I$, $\varepsilon\ge0$, completing the form square
gives $\mathsf W_\lambda\succeq-\varepsilon I$.
Proving such a signed bound along growing windows is still open;
positivity of $F$ alone is insufficient.

\subsection*{A directional certificate for the entire Fourier complement}

The residual Gram in \eqref{eq:v114-schur-residual} can be enclosed
without replacing its small directions by a scalar norm.  We retain
the actual matrix \eqref{eq:v114-b-d}, including its separate diagonal.
All adjoints below are the usual coordinate adjoints in the fixed
orthonormal Fourier basis; quadratic-form order is unchanged by our
first-slot-linear convention.

\begin{proposition}[Moment enclosure of the remote residual]
\label{prop:v116-remote-moment}
Fix $\lambda$, integers $1\le M<J$, and $r\ge1$.  Let
$G:\mathbb C^d\to E_M$ be any coefficient matrix, with row $G_m$ at
Fourier index $m$, and choose $B_*\ge\sup_n|b_n|$, for example
$B_*=B_\lambda^{\rm mat}$ from \eqref{eq:v114-b-bound}.
Define the $r$-by-$d$ moment matrices by
\begin{equation}\label{eq:v116-moments}
 S_j=\sum_{m=-M}^M(m/M)^jG_m,\qquad
 T_j=\sum_{m=-M}^M(m/M)^j b_mG_m,
 \qquad 0\le j<r,
\end{equation}
using $0^0=1$.  Set
\begin{align}
 H_{jk}&=M^{j+k}\{1+(-1)^{j+k}\}
                  \zeta(j+k+2,J+1),\notag\\
 c_{M,r}&=\sum_{m=-M}^M(|m|/M)^{2r},\notag\\
 \delta_{M,J,r}^{2}
 &=\frac{8B_*^2c_{M,r}M^{2r}\zeta(2r+2,J+1)}
                    {(1-M/(J+1))^2},
 \label{eq:v116-moment-constants}
\end{align}
where $\zeta(s,a)=\sum_{k\ge0}(k+a)^{-s}$ is the Hurwitz zeta
function.  For every $\tau>0$, the explicitly finite matrix
\begin{equation}\label{eq:v116-remote-bound}
 U_J=(1+\tau)\{2B_*^2S^*HS+2T^*HT\}
       +(1+\tau^{-1})\delta_{M,J,r}^{2}G^*G
\end{equation}
satisfies
\[
 (Q_J\mathsf W_\lambda G)^*(Q_J\mathsf W_\lambda G)\preceq U_J.
\]
The assertion also holds on either normalized parity sector, provided
the moments are those of the corresponding full Fourier coefficients.
\end{proposition}
\begin{proof}
Since $|n|>J>M\ge|m|$, no diagonal entry occurs.  The finite geometric
identity for $(n-m)^{-1}$ gives exactly
\begin{align*}
 (\mathsf W_\lambda G)_n&=A_n+E_n,\\
 A_n&=\sum_{j=0}^{r-1}\frac{M^j}{n^{j+1}}(b_nS_j-T_j),\\
 E_n&=\sum_{m=-M}^M\left(\frac mn\right)^r
                         \frac{b_n-b_m}{n-m}G_m.
\end{align*}
The matrix $H$ is the Gram matrix of the sequences
$M^j n^{-j-1}$ on $|n|>J$; pairing the two signs of $n$ proves the
formula in \eqref{eq:v116-moment-constants} and $H\succeq0$.
For a column $x$, apply
$|b_na_n-t_n|^2\le2B_*^2|a_n|^2+2|t_n|^2$ and sum over the remote
indices.  This bounds the Gram of $A$ by
$2B_*^2S^*HS+2T^*HT$.

For the remainder, Cauchy--Schwarz in $m$ yields
\[
 |E_nx|^2\le
 \frac{4B_*^2(M/|n|)^{2r}c_{M,r}}{(|n|-M)^2}\|Gx\|^2.
\]
Use $|n|-M\ge|n|(1-M/(J+1))$ and sum over both tails.
The resulting bound is
$\sum_{|n|>J}|E_nx|^2\le\delta_{M,J,r}^{2}\|Gx\|^2$;
this accounts for the factor $8$ in
\eqref{eq:v116-moment-constants}.  Finally apply
$|a+e|^2\le(1+\tau)|a|^2+(1+\tau^{-1})|e|^2$ before summing.
Every estimate is a quadratic-form inequality for arbitrary $x$, so
it proves the stated matrix order, not merely entrywise bounds.

For completeness, an even normalized parity vector has full
coordinates $G_0=v_0$, $G_m=G_{-m}=v_m/\sqrt2$ for $m>0$;
an odd vector has $G_0=0$, $G_m=-G_{-m}=v_m/\sqrt2$.
As $b$ is odd, $S_j=0$ for odd $j$ and $T_j=0$ for even $j$ in the
even sector, with the opposite zeros in the odd sector.
These zeros are exact.  The two tails are already counted in $H$;
no further factor of two is inserted in the parity calculation.
\end{proof}

\begin{proposition}[Diagonal-weighted Schur certificate]
\label{prop:v116-directional-schur}
Write $P=P_N$, $Q=I-P$ and
$F=P\mathsf W_\lambda P$, $B=Q\mathsf W_\lambda P$,
$T=Q\mathsf W_\lambda Q$.  Suppose an independent form estimate gives
\[
 T\succeq D_g:=\operatorname{diag}_{|n|>N}(g_{|n|})
       \succeq\gamma I>0,
\]
where $g_k$ is nondecreasing for $k>N$; the estimate of
Proposition~\ref{prop:v116-diagonal-inverse} is one source of such a
sequence.  Let $N<M<J$, let $Z:E_N\to E_M\cap Q\mathcal H$ be any
finite matrix, and put
\[
 G=I_{E_N\hookrightarrow E_M}-Z,\qquad
 K_Z=F-B^*Z-Z^*B+Z^*TZ.
\]
For $N<|n|\le J$ compute the rows
$R_n=(\mathsf W_\lambda G)_n$ using the complete actual matrix,
and form
\begin{equation}\label{eq:v116-directional-lower}
 \mathcal L_Z=K_Z-
     \sum_{N<|n|\le J}\frac{R_n^*R_n}{g_{|n|}}
                    -\frac{U_J}{g_{J+1}},
\end{equation}
with $U_J$ from Proposition~\ref{prop:v116-remote-moment}.
Then the exact Schur complement satisfies
$F-B^*T^{-1}B\succeq\mathcal L_Z$.
In particular, a certified inequality
$\mathcal L_Z\succeq-\varepsilon I$, $\varepsilon\ge0$, implies
$\mathsf W_\lambda\succeq-\varepsilon I$ on the entire form domain.
The same implication holds after restriction to either parity sector.
\end{proposition}
\begin{proof}
The cross operator $B$ is bounded by
Proposition~\ref{prop:v114-weil-core}; finite $Z$ has range in the
operator domain.  Thus $R=B-TZ=Q\mathsf W_\lambda G$ is well-defined.
Inverse order gives $T^{-1}\preceq D_g^{-1}$.  For unbounded
positive operators this follows, for example, by comparing the
variational formulas for their inverse quadratic forms; the positive
lower bound $\gamma$ makes both inverses bounded.  The exact
verified-solve identity therefore gives
\[
 F-B^*T^{-1}B=K_Z-R^*T^{-1}R
                  \succeq K_Z-R^*D_g^{-1}R.
\]
Split the last quadratic form into $N<|n|\le J$ and $|n|>J$.
Monotonicity gives $g_{|n|}^{-1}\le g_{J+1}^{-1}$ on the remote part,
where \eqref{eq:v116-remote-bound} applies.  This proves
\eqref{eq:v116-directional-lower}.  Writing $\mathfrak w_\lambda$ for the closed form of
$\mathsf W_\lambda$, completion of the square gives, for head $x$
and form-domain tail $y$,
\[
 \mathfrak w_\lambda[x+y]
 =\|T^{1/2}(y+T^{-1}Bx)\|^2
       +\langle(F-B^*T^{-1}B)x,x\rangle.
\]
Interpreted as the closed form on its form domain, this is at least
$-\varepsilon\|x\|^2\ge-\varepsilon\|x+y\|^2$.
For a calculation confined to one parity sector, all ingredients
are taken in that sector and the same proof applies.
\end{proof}

The use of $g_{|n|}$ on the intervening rows, and of the full moment
Grams on the remote rows, is substantive: replacing either by a
single operator norm can destroy the small head directions that the
certificate must preserve.  The finite rows in
\eqref{eq:v116-directional-lower} include the true diagonal $d_n$;
only the disjoint remote rows use the divided-difference expansion.
These propositions supply a finite sufficient test, not a
uniform-in-$\lambda$ positivity theorem.  A passed test at one window
does not close G2 or justify any transfer without its existing graph
defect.

\subsection*{A certificate for the complete fixed-window form}
The preceding estimates now permit a signed certificate that includes
the entire Fourier complement. This is a computer-assisted result for
one physical window; its finite arithmetic and its analytic tail bounds
are both part of the proof.

\begin{proposition}[Complete positivity at $\lambda=3$]
\label{prop:v116-window-positive}
For the canonical closed Weil form of
Proposition~\ref{prop:v114-weil-core}, at $\lambda=3$ there is a
constant $\epsilon_3>0$ such that
\[
 QW_3(f,f)\ge\epsilon_3\|f\|_2^2
\]
for every vector in its form domain. The result includes both parity
sectors and all complex vectors. No numerical value of $\epsilon_3$
is asserted here. In particular, the LDL pivots below are not lower
bounds for the smallest eigenvalue.
\end{proposition}
\begin{proof}[Analytic reduction and interval certificate]
Use the actual unshifted orthonormal Fourier basis on $[0,L]$,
$L=2\log3$, and its orthonormal parity combinations
$e_0$, $(e_m+e_{-m})/\sqrt2$ and
$(e_m-e_{-m})/\sqrt2$. The real matrix commutes with reflection, so
positivity in both sectors implies positivity for the whole complex
space. For each sector split the head at $N$, use an intermediate
inverse-action support through $M$, and compute residual rows through
$J$, with the following exact parameters:
\begin{center}
\begin{tabular}{lrrrrr}
\toprule
Sector & Head dimension & $N$ & $M$ & $J$ & Moment order $r$\\
\midrule
Even &257&256&512&4096&80\\
Odd  &512&512&1024&4096&100\\
\bottomrule
\end{tabular}
\end{center}
The tail lower functions $g_n$ are those of
Proposition~\ref{prop:v116-diagonal-inverse}, with the common base cut $N$
retained in every loss term. In particular their smallest values exceed
$2.0690$ in the even case and $1.1907$ in the odd case.

Let $T_M$ be the compression to $N<|n|\le M$, and let $B_M$ be the
corresponding head-to-tail block. A 768-bit Arb interval solve encloses
$T_M^{-1}B_M$; freeze its midpoint as the exact dyadic matrix $Z$.
The certificate then uses the complete expression
\[
 K_Z=F-B_M^*Z-Z^*B_M+Z^*T_MZ,
\]
and computes every residual row $N<|n|\le J$ from
$R_n=(\mathsf W_3G)_n$, where $G=I_{\rm head}-Z$.
This retains the nonzero intermediate residual from freezing $Z$.
The correct separate logarithmic diagonal is used whenever $n=m$.
For the remote rows, use
Proposition~\ref{prop:v116-remote-moment} with
\[
 B_* =\frac{4/3+P_3+1+\pi/4}{\pi},\qquad
 P_3=\sum_{1<m\le9}\frac{\Lambda(m)}{\sqrt m},\qquad
 \tau=10^{-6}.
\]
The optional full-length term in this conservative $B_*$ does not
enter the actual matrix. It only enlarges the absolute sequence bound.
The squared geometric remainder factors satisfy
\[
 \delta^2_{\rm even}<1.656\cdot10^{-148},\qquad
 \delta^2_{\rm odd}<3.288\cdot10^{-124}.
\]
The full matrix majorant $U_{\rm remote}$, including its leading
moment terms, is retained; these small remainder factors alone are not
estimates for the whole remote residual.

Proposition~\ref{prop:v116-directional-schur} gives the finite lower form
\begin{equation}\label{eq:v116-certified-lower}
 \mathcal L=K_Z-
 \sum_{N<|n|\le J}\frac{R_n^*R_n}{g_n}
 -\frac{U_{\rm remote}}{g_{J+1}}
 \preceq F-B^*T^{-1}B.
\end{equation}
In a parity basis the sum uses one normalized row per positive index;
the remote formula is evaluated in the isometric full-index embedding.
There is no additional factor two or omitted opposite-sign tail.
Interval LDL decomposition certifies every pivot of $\mathcal L$ to
be strictly positive. The smallest pivot enclosures are contained in
\[
 (9.17498,9.17499)\cdot10^{-9}\quad\hbox{(257 even pivots)},
 \qquad
 (9.68026,9.68027)\cdot10^{-8}\quad\hbox{(512 odd pivots)}.
\]
The reproducibility bundle records every pivot, the exact dyadic
witnesses, their hashes, the matrix assembly, and all analytic error
constants. Replaying both identical witnesses at 896-bit precision
again certifies all pivots positive.

For completeness, the coefficient evaluation does not truncate an
uncontrolled archimedean series. With $a_k=2k+1/2$, the exact
digamma and trigamma expressions from
\eqref{eq:v115-arch-exact} leave exponentially weighted sums.
After $k=0,\ldots,127$, put $a_{128}=256+1/2$ and
$e_{128}=3^{-513}$. The radii added to the sine sequence and
complete archimedean diagonal respectively are
\[
 \frac{e_{128}}{2a_{128}(1-1/81)},\qquad
 \frac{2e_{128}}{La_{128}^2(1-1/81)}.
\]
They follow from $t/(a^2+t^2)\le1/(2a)$ and
$|a^2-t^2|/(a^2+t^2)^2\le1/a^2$, uniformly in frequency.
The full-length $m=9$ prime translation is exactly zero; every
shorter prime-power shift is present. All remaining function evaluations
and matrix operations use outward ball enclosures.

Thus the exact Schur complement is positive definite in each finite
head and $T\succeq g_{N+1}I>0$ on its entire complement. Completing
the form square proves coercivity, since the triangular map
$(x,y)\mapsto(x,y+T^{-1}Bx)$ is boundedly invertible and the finite
positive Schur matrix has a positive smallest eigenvalue. The established
Fourier form core extends the result to the complete closed form domain.
No zero-location or Weil-positivity assumption is used.
\end{proof}

\subsection*{Fixed-window spectral ordering and ordinary ground overlap}

The fixed-window computer-assisted positivity certificate can be shifted to count its low
eigenvalues.  The thresholds in the two reflection sectors need not
coincide.  In particular, an even trial vector only sees the even spectral
gap when its overlap with the ground state is estimated.

\begin{proposition}[A fixed-window simple even ground state]
\label{prop:v117-ground-order}
Let $\mu_0\le\mu_1\le\cdots$ be the eigenvalues, counted with
multiplicity, of the complete canonical operator $\mathsf W_3$.
Then its ground state is simple and even, and
\begin{equation}\label{eq:v117-ground-order}
 0<\mu_0<3.644\cdot10^{-38},\qquad \mu_1>10^{-36}.
\end{equation}
Within the even sector alone, every eigenvalue other than $\mu_0$
exceeds $10^{-34}$.  These statements concern the complete operator,
including its infinite Fourier complement, at the single window
$\lambda=3$.
\end{proposition}
\begin{proof}[Shifted inertia and a matching trial direction]
For a threshold $a>0$ replace every diagonal $d_n$ by $d_n-a$ and
every tail lower weight $g_n$ by $g_n-a$.  Keep the same exact dyadic
witness $Z$ and the same head, intermediate and remote cutoffs as in
Proposition~\ref{prop:v116-window-positive}.  If $T_a=T-aI$, then
\[
 K_{Z,a}=K_{Z,0}-a(I+Z^*Z),\qquad
 R_a=B-T_aZ=R_0+aZ.
\]
Thus all intermediate residual rows change and are recomputed.  The
remote moment bound is unchanged, because those rows have no diagonal
entry and $G=\binom{I_{\rm head}}{-Z}$ is unchanged.  The diagonal inverse comparison gives
\[
 S_a:=F-aI-B^*T_a^{-1}B\succeq
 L_a:=K_{Z,a}-\sum_{N<n\le J}\frac{R_{a,n}^*R_{a,n}}{g_n-a}
                      -\frac{U_J}{g_{J+1}-a}.
\]
The sum uses normalized parity rows.  All denominators are strictly
positive by the earlier entire-tail bounds.

The signed interval LDL certificate for the even lower matrix at
$a_{\rm e}=10^{-34}$ has $257$ nonzero pivots, exactly one of which
is negative (zero-based index $24$).  For the odd lower matrix at
$a_{\rm o}=10^{-36}$ all $512$ pivots are strictly positive.
The complete signed elimination retains negative pivots with their
signs; it does not stop at the first negative pivot.  The exact dyadic
witnesses and every pivot enclosure are retained in the bundle;
\texttt{certify\_shifted\_inertia.py} recomputes all shifted residual rows.
The exact parameters are those of Proposition~\ref{prop:v116-window-positive}:
$(N,M,J,r)=(256,512,4096,80)$ for the even sector and
$(512,1024,4096,100)$ for the odd sector. Both certificates pass
at 768 bits and again at 896 bits using identical dyadic witnesses.
By finite-dimensional min--max (see also \cite[Section~4.3]{Teschl2009}), $S_{a_{\rm e}}$ has at most one
nonpositive eigenvalue, because the second eigenvalue of its lower
matrix is strictly positive.  The odd Schur matrix is positive definite.

A matching negative even direction is supplied by the existing
finite-supported candidate of
Proposition~\ref{prop:v112-finite-certificate}. More precisely, in that
$N=64$ certificate let $u$ be its normalized even vector, let $Q$ project
onto its even orthogonal complement. Set
$W_{64}=P_{64}\mathsf W_3P_{64}$, $\alpha=\ip{W_{64}u}{u}$,
$C=QW_{64}Q$ on that complement, and $r=QW_{64}u$. Write $A=C-\alpha I$, $z=A^{-1}r$, $q=\norm z$ and
$E=\ip zr$.  The vector
\[
 v=\frac{u-z}{\sqrt{1+q^2}}
\]
has norm one and belongs to the finite Fourier core.  Direct expansion
gives its exact complete-form Rayleigh value
\[
 \beta=QW_3(v,v)=\alpha-\frac{E}{1+q^2}.
\]
The previously certified enclosures imply
\[
 \alpha<5.312382\cdot10^{-38},\quad
 E>1.668982\cdot10^{-38},\quad q<0.000216,
 \qquad \beta<3.644\cdot10^{-38}.
\]
The last strict scalar comparison is separately checked by
\texttt{certify\_ground\_overlap\_constants.py}.  Because $v$ has
finite support, its finite-compression Rayleigh value equals its
complete-form value exactly.  No approximation of the omitted output
coefficients is used in this identity.  In particular $\beta<a_{\rm e}$.

On the closed form domain, square completion reads
\[
 q_a[x+y]=\ip{S_ax}{x}
       +\norm{T_a^{1/2}(y+T_a^{-1}Bx)}^2.
\]
The map $(x,y)\mapsto(x,y+T_a^{-1}Bx)$ is boundedly invertible and
preserves that domain.  It transfers the negative index to $S_a$;
also $\ker(\mathsf W_3-aI)$ is isomorphic to $\ker S_a$.
The negative even trial direction therefore makes the first even
Schur eigenvalue strictly negative, while its second is strictly
positive.  There is no even eigenvalue exactly at $a_{\rm e}$, and
exactly one even eigenvalue lies below it, counted with multiplicity.
Every odd eigenvalue exceeds $a_{\rm o}$.  The established compact
resolvent makes these full spectral statements.  Finally
Proposition~\ref{prop:v116-window-positive} supplies $\mu_0>0$,
and min--max applied to $v$ supplies $\mu_0\le\beta$.
Since $\beta<a_{\rm o}<a_{\rm e}$, the unique even low eigenvalue is
the simple ground state of the whole operator.  This proves
\eqref{eq:v117-ground-order} and the stronger even-sector threshold.
\end{proof}

\begin{corollary}[An ordinary ground-overlap bound at one window]
\label{cor:v117-ground-overlap}
For the normalized rational candidate $u$ of
Proposition~\ref{prop:v112-finite-certificate} and the complete unit
ground vector $\xi_0$ at $\lambda=3$,
\[
 \sin\angle(u,\xi_0)<0.01931.
\]
The same strict upper bound holds for the normalized exact $N=64$
source projection after the already certified normalization error is
included.  This is an ordinary Hilbert-space bound; it asserts no
endpoint comparison with $\xi_0$.
\end{corollary}
\begin{proof}
The improved trial $v$ in the proof above is even.  In its spectral
expansion, every non-ground component has energy greater than
$a_{\rm e}=10^{-34}$, and the ground energy is positive.  Hence
\[
 \sin\angle(v,\xi_0)
 <\sqrt{3.644\cdot10^{-38}/10^{-34}}.
\]
Also $z\perp u$, so the norm of the difference of the two rank-one
orthogonal projections onto $u$ and $v$ equals
$q/\sqrt{1+q^2}<0.000216$.  The triangle inequality for these
projectors proves
\[
 \sin\angle(u,\xi_0)
 <0.000216+\sqrt{3.644\cdot10^{-4}}
 <0.019305264<0.01931.
\]
For the unnormalized rational candidate $c$,
Proposition~\ref{prop:v113-source-certificate} gives
$\|P_{64}p_3-c\|<\varepsilon:=2.60\cdot10^{-36}$, while the
exact rational norm calculation gives $\|c\|>0.65172555$.
Writing $\rho=\varepsilon/\|c\|$, elementary ball geometry yields
\[
 \left\|\frac{P_{64}p_3}{\|P_{64}p_3\|}-\frac c{\|c\|}\right\|^2
 \le\frac{2\rho^2}{1+\sqrt{1-\rho^2}}<(4\cdot10^{-36})^2.
\]
Indeed the real inner product of these two unit vectors is positive
and at least $\sqrt{1-\rho^2}$, by minimizing the squared distance
from $c$ along the positive ray through $P_{64}p_3$.
This normalization error fits within the final margin.  The ordinary
norm controls projector distance but not the physical endpoint
functional; the earlier endpoint obstruction is therefore unchanged.
\end{proof}

This proves parity, simplicity and quantitative ordinary overlap at one
fixed window. It does not identify the unprojected source with the
ground vector. There is still no uniform family of spectral thresholds
or source overlap estimates along unbounded windows. The physical
endpoint is resolved below by a separate regularity argument; it does
not follow from this ordinary overlap bound. The sampler graph defect
and full-strip requirements remain separate; G2 and RH are not closed
by the fixed-window spectral result.

\subsection*{The physical endpoint of the complete Weil eigenfunctions}

The endpoint of an actual eigenfunction is not determined by ordinary
norm convergence of finite eigenvectors.  For the complete operator it
can instead be determined from the singular archimedean kernel.
Throughout this subsection $\Omega=(0,L)$, physical functions are
extended by zero outside $\Omega$, and Fourier frequency is normalized
as in $U_n(x)=L^{-1/2}e^{2\pi inx/L}$.

\begin{proposition}[Restricted logarithmic Laplacian decomposition]
\label{prop:v118-log-dirichlet}
Let $\mathcal A=\tfrac12L_\Delta^{\rm Dir}$ be the restricted,
exterior-Dirichlet logarithmic Laplacian on $\Omega$, whose full-line
symbol is $\log|t|$.  This is not the spectral logarithm of the
Dirichlet Laplacian.  The actual closed Weil operator satisfies
\begin{equation}\label{eq:v118-log-decomposition}
 \mathsf W_\lambda=\mathcal A+\mathcal K_\lambda,
 \qquad \mathcal D(\mathsf W_\lambda)=\mathcal D(\mathcal A),
\end{equation}
where, writing $k(y)=\rho(y)-(2y)^{-1}$, $a_m=\log m$, and
$w_m=\Lambda(m)/\sqrt m$,
\begin{align}
 (\mathcal K_\lambda f)(x)={}&-\log(2\pi)f(x)
       -\int_0^L k(|x-y|)f(y)\,dy\notag\\
 &-\sum_{1<m<\lambda^2}w_m
       \{\mathbf1_{x+a_m<L}f(x+a_m)
                   +\mathbf1_{x>a_m}f(x-a_m)\}\notag\\
 &+2\int_0^L\cosh((x-y)/2)f(y)\,dy.
 \label{eq:v118-bounded-physical-part}
\end{align}
The operator $\mathcal K_\lambda$ is real and self-adjoint on $L^2$.
For every $1\le p\le\infty$ it is bounded on $L^p(\Omega)$, with
\begin{equation}\label{eq:v118-K-bound}
 \|\mathcal K_\lambda\|_{p\to p}\le C_\lambda
 :=\log(2\pi)+2\int_0^L|k(y)|\,dy
                  +2P_\lambda+4\sinh(L/2).
\end{equation}
\end{proposition}
\begin{proof}
In dimension one the integral normalization of the logarithmic
Laplacian~\cite[Theorem~1.1]{ChenWethLog} gives
\[
 \mathcal A f(x)=\int_0^\infty
 \frac{2f(x)\mathbf1_{y<1}-f(x+y)-f(x-y)}{2y}\,dy
                   -\gamma_{\rm E}f(x)
\]
on compactly supported smooth tests.  The actual negative-archimedean
part of the Weil operator is
\[
 \int_0^\infty\rho(y)
 \{2e^{-y/2}f(x)-f(x+y)-f(x-y)\}\,dy
             -\{\log(4\pi)+\gamma_{\rm E}\}f(x).
\]
Its Fourier diagonal is precisely the negative of
\eqref{eq:v114-arch-diagonal}; the diagonal integral over $y>L$
equals $-\log\tanh(L/2)$.  Since
\[
 \int_0^\infty\left\{\frac1{\sinh y}
                -\frac{\mathbf1_{y<1}}y\right\}dy=\log2,
\]
subtracting $\mathcal A$ leaves $-\log(2\pi)I$ and the kernel $-k$.
The integral follows by evaluating $\log\tanh(y/2)$ with a positive
lower cutoff and then taking its limit.  In particular the constant
part of the logarithmic diagonal has been retained.
The prime shifts and the pole kernel $2\cosh((x-y)/2)$ give the
remaining terms of \eqref{eq:v118-bounded-physical-part}.
The latter kernel is the difference of the centered rank-one
$2\cosh\otimes\cosh$ and $2\sinh\otimes\sinh$ terms, so no pole
sign has been removed.

The function $k$ extends continuously to zero, with $k(0)=1/4$.
The integral-kernel row and column bounds, and the contraction of
each truncated translation, prove \eqref{eq:v118-K-bound}.
For the pole kernel the largest row integral is
$4\sinh(L/2)$.  Symmetry proves self-adjointness.

It remains to identify the closed forms, without imposing an
unproved endpoint condition.  Smooth compactly supported functions
are a form core for the restricted logarithmic
Laplacian~\cite[Theorem~3.1]{ChenWethLog}.  They are also a form core
for the operator in Proposition~\ref{prop:v114-weil-core}.
Indeed, approximate each finite periodic Fourier sum $f$ by smooth
cutoffs that remove endpoint strips of width $\epsilon$.
The discarded periodic function has $L^2$ norm $O(\epsilon^{1/2})$
and $H^1$ norm $O(\epsilon^{-1/2})$, hence $H^s$ norm
$O(\epsilon^{1/2-s})$ for any fixed $0<s<1/2$.
This tends to zero in the logarithmically weighted form norm.
The two closed lower-bounded forms therefore agree on a common form
core.  Bounded perturbation by $\mathcal K_\lambda$ proves
\eqref{eq:v118-log-decomposition}, including its operator-domain assertion.
\end{proof}

\begin{theorem}[Zero boundary trace of every complete Weil eigenfunction]
\label{thm:v118-eigen-endpoint-zero}
Fix $\lambda>1$.  Every $L^2$ eigenfunction
$\mathsf W_\lambda u=\mu u$ has a bounded representative whose zero
extension is continuous on $\mathbb R$.  With
\[
 \ell(r)=\frac1{|\log(\min\{r,0.1\})|},\qquad
 d(x)=\min\{x,L-x\},
\]
there is a finite constant $C$ such that
\begin{equation}\label{eq:v118-eigen-boundary-decay}
 |u(x)|\le C\sqrt{\ell(d(x))}\quad(0<x<L).
\end{equation}
In particular its actual physical endpoints satisfy $u(0)=u(L)=0$.
No uniform bound for $C$ as $\lambda$ grows is asserted.
\end{theorem}
\begin{proof}
First we prove boundedness, without assuming a positive Sobolev gain.
For $0<\epsilon<\min\{1,L\}$ let $\mathcal A_\epsilon$ be the
nonnegative killed small-jump generator
\[
 \mathcal A_\epsilon f(x)=\int_0^\epsilon
       \frac{2f(x)-f(x+y)-f(x-y)}{2y}\,dy
\]
in its closed Dirichlet-form realization, and put
\[
 J_\epsilon f(x)=\int_{\substack{y\in\Omega\\|x-y|\ge\epsilon}}
                    \frac{f(y)}{2|x-y|}\,dy.
\]
There is an exact decomposition
\[
 \mathcal A=\mathcal A_\epsilon+
       \{\log(1/\epsilon)-\gamma_{\rm E}\}I-J_\epsilon.
\]
The small-jump form is Markovian, including its nonnegative exterior
killing term.  Consequently its resolvent is consistent on $L^2$
and $L^\infty$, and
$\|(\mathcal A_\epsilon+aI)^{-1}\|_{p\to p}\le a^{-1}$ for
$a>0$, $p=2,\infty$.  One may obtain this directly by first cutting
off jumps smaller than a positive number, using the contraction
resolvent of that finite-activity killed jump generator, and then
passing to the increasing closed forms.  Also Cauchy--Schwarz gives
\[
 \|J_\epsilon\|_{2\to\infty}\le(2\epsilon)^{-1/2}.
\]
Choose $\epsilon$ so that
$a=\log(1/\epsilon)-\gamma_{\rm E}-\mu>C_\lambda$.
The eigen equation becomes
\[
 (\mathcal A_\epsilon+aI+\mathcal K_\lambda)u=J_\epsilon u.
\]
The inverse on the left exists on both $L^2$ and $L^\infty$ by the
Neumann series for
$[I+(\mathcal A_\epsilon+aI)^{-1}\mathcal K_\lambda]^{-1}
 (\mathcal A_\epsilon+aI)^{-1}$, with norm at most
$(a-C_\lambda)^{-1}$.  Since $J_\epsilon u\in L^2\cap L^\infty$,
the two series agree.  Uniqueness in $L^2$ identifies the bounded
solution with $u$, and proves
\[
 \|u\|_\infty\le
   \frac{\|u\|_2}{(a-C_\lambda)\sqrt{2\epsilon}}.
\]
Thus boundedness has not been inferred from the logarithmic operator
domain alone.

By Proposition~\ref{prop:v118-log-dirichlet}, the zero extension lies
in the logarithmic Dirichlet form space and satisfies
\[
 L_\Delta u=2(\mu u-\mathcal K_\lambda u)\in L^\infty(\Omega),
 \qquad u=0\quad\hbox{outside }\Omega.
\]
The interval satisfies an exterior uniform sphere condition.
The boundary regularity theorem of Hern\'andez-Santamar\'ia,
L\'opez R\'ios and Salda\~na~\cite[Theorem~1.1]{HLSLogBoundary}
therefore gives continuity of the zero extension and
\eqref{eq:v118-eigen-boundary-decay}.  For complex $u$, apply the
real theorem to its real and imaginary parts.  This proves the endpoint
assertion for every eigenfunction, independently of its spectral order.
\end{proof}

\begin{corollary}[Filtered endpoint decay and the nonzero-source obstruction]
\label{cor:v118-ground-fejer}
For an eigenfunction as above and $N\ge4$, define its Fej\'er endpoint
in the actual normalized Fourier coordinates by
\[
 \sigma_{N-1}u(0)=L^{-1/2}
       \sum_{|n|<N}(1-|n|/N)\widehat u(n).
\]
Then
\begin{equation}\label{eq:v118-fejer-endpoint}
 |\sigma_{N-1}u(0)|
 \le C\sqrt{\ell(L/\sqrt N)}+\frac{\|u\|_\infty}{2\sqrt N}
 =O_{\lambda,\mu,u}((\log N)^{-1/2}).
\end{equation}
For the exact repaired source with $B_\lambda=p_\lambda(0)\ne0$,
any complete-operator eigenfunction $u$ and scalar $c$ satisfy
\begin{equation}\label{eq:v118-ground-source-mismatch}
 \frac{|p_\lambda(0)-cu(0)|}{|B_\lambda|}=1.
\end{equation}
\end{corollary}
\begin{proof}
The positive Fej\'er kernel in physical coordinates is
$L^{-1}F_N(2\pi x/L)$, where
$F_N(\theta)=N^{-1}\{\sin(N\theta/2)/\sin(\theta/2)\}^2$;
it has integral one.  Put $h=L/\sqrt N\le L/2$.
On $d(x)<h$, \eqref{eq:v118-eigen-boundary-decay} bounds the
convolution by $C\sqrt{\ell(h)}$.
On $d(x)\ge h$, $\sin(\pi x/L)\ge2d(x)/L$ gives
\[
 L^{-1}F_N(2\pi x/L)\le\frac{L}{4N d(x)^2},\qquad
 \int_{d(x)\ge h}L^{-1}F_N(2\pi x/L)\,dx
       \le\frac{L}{2Nh}=\frac1{2\sqrt N}.
\]
This proves \eqref{eq:v118-fejer-endpoint}.
The mismatch identity follows from $u(0)=0$.
\end{proof}

Theorem~\ref{thm:v118-eigen-endpoint-zero} applies in particular to
the simple even ground at $\lambda=3$.  It disproves a nonzero physical
endpoint comparison between the repaired source and the complete
operator's ground, while leaving their ordinary-norm overlap intact.
It does not identify the endpoint of a finite-compression eigenvector
on a joint growing-window and growing-cutoff diagonal.
Nor does continuity alone prove absolute Fourier convergence or
convergence of sharp Fourier endpoint sums: the quantitative statement
\eqref{eq:v118-fejer-endpoint} concerns the stated positive filter.
The finite endpoint and its graph-defect terms therefore remain
separate from the complete eigenfunction's trace.

\begin{proposition}[The Weil graph norm does not control the physical trace]
\label{prop:v118-graph-trace}
At every fixed $\lambda>1$, physical endpoint evaluation on finite
Fourier polynomials has no continuous extension to
$\mathcal D(\mathsf W_\lambda)$ with its operator graph norm.
Indeed there are real even finite polynomials $h_N$ with
\[
 \delta(h_N)=1,\qquad
 \|h_N\|_2+\|\mathsf W_\lambda h_N\|_2
       =O_\lambda\!\left(\frac{\log(2+N)}{\sqrt N}\right)\longrightarrow0.
\]
The same failure holds with $\mathsf W_\lambda-\mu I$ for any fixed
$\mu\in\mathbb R$ and with the weaker closed-form norm.
\end{proposition}
\begin{proof}
In the physical logarithmic basis $U_n=L^{-1/2}e^{2\pi inx/L}$ set
\[
 h_N=\frac{\sqrt L}{2N}\sum_{N<|n|\le2N}U_n.
\]
There are exactly $2N$ terms. At either physical endpoint all their
phases are one, so $\delta(h_N)=1$ and
$\|h_N\|_2=\sqrt{L/(2N)}$. The exact decomposition in
Proposition~\ref{prop:v114-weil-core} gives
\[
 \|\mathsf W_\lambda h_N\|_2
 \le\left(\max_{N<|n|\le2N}|d_n|
              +2\pi B_\lambda^{\rm mat}\right)\sqrt{\frac L{2N}}.
\]
Since $d_n=\log|n|+O_\lambda(1)$, this proves the estimate.
A fixed scalar shift changes only the bounded constant. Graph convergence
implies form convergence for a lower-bounded self-adjoint operator.
In fact the trace is not closable on the polynomial core in the graph
norm: this sequence tends to zero while its trace tends to one.
\end{proof}

Let
\[
 \mathcal V_\lambda:=\operatorname{Var}(p_\lambda)
\]
be the total variation of the periodic logarithmic representative of the
inversion-even proxy.  Let
\[
 M_T(\lambda):=\max\{m\ge0:2\pi m/L\le T\}.
\]

\begin{proposition}[Explicit BV rate for the fixed-band Galerkin defect]\label{prop:g1-bv-rate}
Fix $T>0$.  If $N\ge2M_T(\lambda)+1$, then
\begin{equation}\label{eq:g1-gamma-rate}
 \boxed{
 \gamma_{\lambda,N,T}
 \le C_T\frac{\lambda L\,\mathcal V_\lambda}
 {|B_\lambda|\,N}.}
\end{equation}
In particular the form defect in G1 has a fully quantitative,
zero-independent rate; no $C^2$ or higher regularity of $p_\lambda$ is used.
\end{proposition}
\begin{proof}
Use the translated logarithmic coordinate $[0,L]$.  Since $p_\lambda$ has
bounded variation and matching periodic endpoint values, Stieltjes integration
by parts gives, for $n\ne0$,
\[
 |\widehat p_\lambda(n)|
 \le\frac{\sqrt L\,\mathcal V_\lambda}{2\pi|n|}.
\]
Let $g=D_{\log}v$ with $v\in E_{\lambda,T}$ and $\|v\|_2=1$.  Then
$\|g\|_2\le T$ and $g_m=0$ for $|m|>M_T(\lambda)$.  If $|n|>N\ge2M_T+1$,
Lemma~\ref{lem:g1-high-low} and Cauchy--Schwarz imply
\[
 \left|\sum_{|m|\le M_T}\tau^{(\lambda)}_{n,m}g_m\right|
 \le C_T\frac{\lambda\sqrt L}{|n|}.
\]
The resulting tail is absolutely summable, and hence
\[
 |QW_\lambda((P_N-I)p_\lambda,g)|
 \le C_T\lambda L\mathcal V_\lambda
       \sum_{|n|>N}\frac1{n^2}
 \le C_T\frac{\lambda L\mathcal V_\lambda}{N}.
\]
Divide by $|B_\lambda|$ and take the supremum over $v$.
\end{proof}

\subsection*{A compatibility-friendlier filtered finite section}
The exact-boundary repair below is useful for G1 alone, but it adds a new
high-frequency vector.  For the later same-object G1/G2 comparison it is
preferable to stay inside a canonical smoothing of the literal
prolate/co-Poisson proxy.  A de la Vall\'ee--Poussin filter does this while
preserving all sufficiently low Fourier modes exactly.

Work in the translated logarithmic coordinate $[0,L]$, with the endpoints
identified.  Let $\widehat p_\lambda(n)$ be the normalized Fourier
coefficients and define the multiplier
\begin{equation}\label{eq:g1-vp-multiplier}
 \nu_N(n):=
 \begin{cases}
  1,& |n|\le N,\\
  2-|n|/N,&N<|n|<2N,\\
  0,&|n|\ge2N.
 \end{cases}
\end{equation}
The de la Vall\'ee--Poussin mean is
\begin{equation}\label{eq:g1-vp-mean}
 \mathsf V_Np_\lambda
 :=\sum_{n\in\mathbb Z}\nu_N(n)\widehat p_\lambda(n)U_n\in E_{2N}.
\end{equation}
Equivalently its convolution kernel is the difference of two Fej\'er kernels,
$2F_{2N}-F_N$.  In particular its $L^1$ norm is bounded by an absolute
constant and, for $0<\delta\le L/2$, the kernel mass outside the boundary arc
$|x|<\delta$ is $O(L/(N\delta))$.

Define the endpoint modulus
\begin{equation}\label{eq:g1-edge-modulus}
 \omega_\lambda(\delta)
 :=\sup_{0\le t\le\delta}
 \max\{|p_\lambda(t)-B_\lambda|,
       |p_\lambda(L-t)-B_\lambda|\}.
\end{equation}
The no-boundary-jump property gives $\omega_\lambda(\delta)\to0$ as
$\delta\downarrow0$ for every fixed $\lambda$.

\begin{lemma}[Quantitative endpoint recovery by de la Vall\'ee--Poussin filtering]\label{lem:g1-vp-endpoint}
For every $N\ge1$ and $0<\delta\le L/2$,
\begin{equation}\label{eq:g1-vp-endpoint}
 \boxed{
 |\delta_{2N}(\mathsf V_Np_\lambda)-B_\lambda|
 \le C\left(
  \omega_\lambda(\delta)
  +\frac{\mathcal V_\lambda L}{N\delta}
 \right).}
\end{equation}
The constant is absolute.  Hence the endpoint recovery is quantitative in
terms of the explicit boundary modulus of the chosen prolate/co-Poisson
proxy and uses no zero information.
\end{lemma}
\begin{proof}
Let $K_N^{\mathrm{VP}}$ be the period-$L$ convolution kernel of
$\mathsf V_N$.  Since
$|K_N^{\mathrm{VP}}|\le2F_{2N}+F_N$, its $L^1$ norm is $O(1)$ and the
standard Fej\'er estimate gives
\[
 \int_{|x|\ge\delta}|K_N^{\mathrm{VP}}(x)|\,dx
 \le C\frac{L}{N\delta}.
\]
On the boundary arc use
$|p_\lambda(x)-B_\lambda|\le\omega_\lambda(\delta)$.  On its complement,
bounded variation and the absence of a periodic boundary jump give
$|p_\lambda(x)-B_\lambda|\le\mathcal V_\lambda$.  Splitting the convolution
integral proves \eqref{eq:g1-vp-endpoint}.
\end{proof}

For later bookkeeping it is useful to note that the endpoint modulus is itself
zero-independently computable from the source.  The compact support of the
prolate source makes the co-Poisson sum finite on the window; away from the
finitely many arithmetic thresholds it is a finite sum of smooth functions.
Thus there are explicit numbers $\Delta_\lambda>0$ and $H_\lambda<\infty$,
determined only by the chosen source and $\lambda$, for which
\begin{equation}\label{eq:g1-edge-lipschitz}
 \omega_\lambda(\delta)\le H_\lambda\delta,
 \qquad0<\delta\le\Delta_\lambda.
\end{equation}
For example one may take $\Delta_\lambda$ to be half the logarithmic distance
from the identified endpoint to the nearest interior arithmetic threshold;
$H_\lambda$ is then the supremum of the ordinary logarithmic derivative on
those two endpoint arcs.  No uniform-in-$\lambda$ estimate for $H_\lambda$ is
being asserted here.

\begin{proposition}[Filtered BV transfer with explicit boundary control]\label{prop:g1-vp-transfer}
Fix $T>0$ and suppose $N\ge2M_T(\lambda)+1$.  Set
\[
 k^{\mathrm{VP}}_{\lambda,N}:=\mathsf V_Np_\lambda\in E_{2N},
 \qquad
 \beta^{\mathrm{VP}}_{\lambda,N}
 :=\frac{|\delta_{2N}(k^{\mathrm{VP}}_{\lambda,N})-B_\lambda|}
 {|B_\lambda|}.
\]
Then
\begin{equation}\label{eq:g1-vp-form-defect}
 \boxed{
 \frac1{|B_\lambda|}
 \sup_{\substack{v\in E_{\lambda,T}\\\|v\|_2=1}}
 |QW_\lambda(k^{\mathrm{VP}}_{\lambda,N}-p_\lambda,
             D_{\log}v)|
 \le C_T\frac{\lambda L\mathcal V_\lambda}{|B_\lambda|N}.}
\end{equation}
Whenever $\beta^{\mathrm{VP}}_{\lambda,N}<1$,
\begin{equation}\label{eq:g1-vp-master}
 \boxed{
 \frac{\|P_TD_{\log}W_{\lambda,2N}
 k^{\mathrm{VP}}_{\lambda,N}\|_2}
 {|\delta_{2N}(k^{\mathrm{VP}}_{\lambda,N})|}
 \le
 \frac{T\mathcal E_T(\lambda)
 +C_T\lambda L\mathcal V_\lambda/(|B_\lambda|N)}
 {1-\beta^{\mathrm{VP}}_{\lambda,N}}.}
\end{equation}
Moreover, if $H_\lambda>0$ and
$N\ge C\mathcal V_\lambda L/(H_\lambda\Delta_\lambda^2)$, then
\begin{equation}\label{eq:g1-vp-beta-rate}
 \boxed{
 \beta^{\mathrm{VP}}_{\lambda,N}
 \le
 \frac{C}{|B_\lambda|}
 \sqrt{\frac{H_\lambda\mathcal V_\lambda L}{N}}.}
\end{equation}
Consequently one obtains a deterministic, zero-independent one-scale G1
diagonal by taking $N=N_\lambda$ large enough that
\begin{equation}\label{eq:g1-vp-cutoff}
 N_\lambda\ge C_T\max\left\{
 2M_T(\lambda)+1,
 \frac{\lambda L\mathcal V_\lambda}
      {|B_\lambda|\mathfrak r(\lambda)},
 \frac{H_\lambda\mathcal V_\lambda L}
      {|B_\lambda|^2\mathfrak r(\lambda)^2},
 \frac{\mathcal V_\lambda L}{H_\lambda\Delta_\lambda^2}
 \right\},
\end{equation}
with the last term omitted when it is vacuous.  Along such a diagonal,
\begin{equation}\label{eq:g1-vp-vanishing}
 \boxed{
 \frac{\|P_TD_{\log}W_{\lambda,2N_\lambda}
 k^{\mathrm{VP}}_{\lambda,N_\lambda}\|_2}
 {|\delta_{2N_\lambda}(k^{\mathrm{VP}}_{\lambda,N_\lambda})|}
 \le C'_T\mathfrak r(\lambda)\longrightarrow0.}
\end{equation}
\end{proposition}
\begin{proof}
The multiplier \eqref{eq:g1-vp-multiplier} is exactly one for $|n|\le N$.
Therefore the Fourier coefficients of
$k^{\mathrm{VP}}_{\lambda,N}-p_\lambda$ vanish in that range and, for
$|n|>N$, have modulus at most $|\widehat p_\lambda(n)|$.  The Stieltjes
integration-by-parts bound from Proposition~\ref{prop:g1-bv-rate} and the
high--low estimate of Lemma~\ref{lem:g1-high-low} therefore give
\eqref{eq:g1-vp-form-defect} with the same $O(N^{-1})$ tail as the ordinary
Galerkin projection.

Equation~\eqref{eq:g1-vp-master} is the same Hermitian fixed-band argument as
\eqref{eq:g1-master-bound}.  Lemma~\ref{lem:g1-vp-endpoint}, followed by
\eqref{eq:g1-edge-lipschitz} and the choice
$\delta\asymp(\mathcal V_\lambda L/(H_\lambda N))^{1/2}$, gives
\eqref{eq:g1-vp-beta-rate} once that $\delta$ lies inside the endpoint arc.
The cutoff \eqref{eq:g1-vp-cutoff} makes both the form defect and boundary
defect $O_T(\mathfrak r(\lambda))$; inserting the continuum envelope
\eqref{eq:G1-continuum} proves \eqref{eq:g1-vp-vanishing}.
\end{proof}

The filtered vector has another advantage over an external boundary repair:
it changes no Fourier coefficient below $N$.  Its derivative can still grow,
because the literal co-Poisson proxy is only BV across its arithmetic interior
thresholds, but the growth is intrinsic rather than injected by a separate
corrector.  The coefficient estimate used above gives the explicit bound
\begin{equation}\label{eq:g1-vp-derivative}
 \|D_{\log}k^{\mathrm{VP}}_{\lambda,N}\|_2
 \le C\mathcal V_\lambda\sqrt{\frac{N}{L}}.
\end{equation}

There is in fact a sharp obstruction to improving the derivative price by
pushing an exact boundary repair farther into the Fourier tail.

\begin{lemma}[High-frequency boundary repair has unavoidable square-root derivative cost]\label{lem:g1-repair-no-go}
Let
\[
 c(x)=\sum_{j=N+1}^{K}w_j\cos(2\pi jx/L),
 \qquad c(0)=\sum_{j=N+1}^{K}w_j=1.
\]
Then
\begin{equation}\label{eq:g1-repair-no-go}
 \boxed{
 \|D_{\log}c\|_2^2
 \ge
 \frac{2\pi^2}{L\displaystyle\sum_{j=N+1}^{K}j^{-2}}
 \ge\frac{2\pi^2N}{L}.}
\end{equation}
The first lower bound is sharp: equality is attained by weights
$w_j\propto j^{-2}$.
\end{lemma}
\begin{proof}
Orthogonality of the sine modes gives
\[
 \|D_{\log}c\|_2^2
 =\frac{2\pi^2}{L}\sum_{j=N+1}^{K}j^2|w_j|^2.
\]
Weighted Cauchy--Schwarz and $\sum w_j=1$ give
\[
 1\le
 \left(\sum j^2|w_j|^2\right)
 \left(\sum j^{-2}\right).
\]
This proves the first bound, with equality precisely for $w_j\propto j^{-2}$.
Finally
$\sum_{j>N}j^{-2}\le\int_N^\infty x^{-2}\,dx=N^{-1}$.
\end{proof}

Lemma~\ref{lem:g1-repair-no-go} shows that the square-root derivative growth
seen in the shell construction is not merely an artifact of equal weights: it
is unavoidable for every even exact-boundary corrector supported entirely
above frequency $N$.  For a norm-level same-object estimate, the filtered
literal proxy of Proposition~\ref{prop:g1-vp-transfer} is therefore still the
preferred G1 candidate.  This derivative lower bound is not, however, a
pairing-level no-go: the shell estimates in Section~\ref{sec:matching} below proves that the
shell's explicit G1 and scalar-shift pairings against the higher-order Hardy
sampler can be made arbitrarily small by moving the shell outward.

The endpoint defect can also be removed exactly without asking for a
quantitative Dirichlet--Jordan theorem.  Put
\[
 q_{\lambda,N}=P_Np_\lambda,
 \qquad
 d_{\lambda,N}=B_\lambda-q_{\lambda,N}(0),
\]
and, for an even integer $K>2N$, define the even high-frequency shell
corrector
\begin{equation}\label{eq:g1-shell-corrector}
 c^{\mathrm{sh}}_{\lambda,K}(x)
 :=\frac{2}{K}\sum_{j=K/2+1}^{K}\cos(2\pi jx/L)
 \quad\text{in logarithmic coordinates}.
\end{equation}
Every summand has boundary value one, so
$c^{\mathrm{sh}}_{\lambda,K}(0)=1$.  Moreover all of its Fourier modes lie
above $N$.  Then set
\begin{equation}\label{eq:g1-bc-candidate}
 k^{\mathrm{bc}}_{\lambda;N,K}
 :=q_{\lambda,N}+d_{\lambda,N}c^{\mathrm{sh}}_{\lambda,K}\in E_K.
\end{equation}
By construction it is real, inversion-even, and
\begin{equation}\label{eq:g1-bc-boundary}
 \delta_K(k^{\mathrm{bc}}_{\lambda;N,K})=B_\lambda
\end{equation}
exactly.

\begin{theorem}[Constructive finite fixed-band G1 closure]\label{thm:g1-boundary-corrected}
Use the explicit repaired source above and fix $T>0$. For $N\ge2M_T(\lambda)+1$ and an even integer
$K>2N$,
\begin{equation}\label{eq:g1-bc-master}
 \boxed{
 \begin{aligned}
 \frac{\|P_TD_{\log}W_{\lambda,K}
 k^{\mathrm{bc}}_{\lambda;N,K}\|_2}{|B_\lambda|}
 \le{}&T\mathcal E_T(\lambda)
 +C_T\frac{\lambda L}{N}
       \frac{\mathcal V_\lambda}{|B_\lambda|}\\
 &+C_T\frac{\lambda L}{K}
 \left(
 1+\sqrt{\frac{2N+1}{L}}
       \frac{\|p_\lambda\|_2}{|B_\lambda|}
 \right).
 \end{aligned}}
\end{equation}
Consequently there are completely zero-independent, deterministic cutoffs
$N_\lambda<K_\lambda$ for which
\begin{equation}\label{eq:g1-bc-vanishing}
 \boxed{
 \frac{\|P_TD_{\log}W_{\lambda,K_\lambda}
 k^{\mathrm{bc}}_{\lambda;N_\lambda,K_\lambda}\|_2}
 {|\delta_{K_\lambda}(k^{\mathrm{bc}}_{\lambda;N_\lambda,K_\lambda})|}
 \le C'_T\mathfrak r(\lambda)\longrightarrow0.}
\end{equation}
One admissible prescription is to choose, after enlarging $C_T$ if necessary,
\begin{align}
 N_\lambda&\ge
 \max\left\{2M_T(\lambda)+1,
 C_T\frac{\lambda L}{\mathfrak r(\lambda)}
       \frac{\mathcal V_\lambda}{|B_\lambda|}\right\},
 \label{eq:g1-N-explicit}\\
 K_\lambda&\text{ even},\qquad K_\lambda>2N_\lambda,\qquad
 K_\lambda\ge
 C_T\frac{\lambda L}{\mathfrak r(\lambda)}
 \left(1+\sqrt{\frac{2N_\lambda+1}{L}}
       \frac{\|p_\lambda\|_2}{|B_\lambda|}\right).
 \label{eq:g1-K-explicit}
\end{align}
All quantities on the right are explicit functions of the chosen
prolate/co-Poisson proxy and $\lambda$; no nontrivial zeta zero enters the
choice.
\end{theorem}
\begin{proof}
For $v\in E_{\lambda,T}$ with $\|v\|_2=1$, Hermiticity gives
\[
 |\langle P_TD_{\log}W_{\lambda,K}
 k^{\mathrm{bc}}_{\lambda;N,K},v\rangle|
 =|QW_\lambda(k^{\mathrm{bc}}_{\lambda;N,K},D_{\log}v)|.
\]
Split the candidate as
\[
 k^{\mathrm{bc}}_{\lambda;N,K}
 =p_\lambda+(P_N-I)p_\lambda+d_{\lambda,N}c^{\mathrm{sh}}_{\lambda,K}.
\]
The first term is at most $T|B_\lambda|\mathcal E_T(\lambda)$ by
Lemma~\ref{lem:g1-inversion}; Proposition~\ref{prop:g1-bv-rate} controls the
second.  For each mode in the shell, Lemma~\ref{lem:g1-high-low} and
Cauchy--Schwarz on the fixed physical band give a bound
$C_T\lambda L/j$; since $K/2<j\le K$ and the weights in
\eqref{eq:g1-shell-corrector} sum to one, averaging yields
\[
 |QW_\lambda(c^{\mathrm{sh}}_{\lambda,K},D_{\log}v)|
 \le C_T\frac{\lambda L}{K}.
\]
Finally the endpoint-evaluation norm on $E_N$ is
$\sqrt{(2N+1)/L}$, so
\[
 \frac{|d_{\lambda,N}|}{|B_\lambda|}
 \le1+\sqrt{\frac{2N+1}{L}}
       \frac{\|p_\lambda\|_2}{|B_\lambda|}.
\]
These estimates prove \eqref{eq:g1-bc-master}.  The choices
\eqref{eq:g1-N-explicit}--\eqref{eq:g1-K-explicit}, together with
$\mathcal E_T(\lambda)\le C_T\mathfrak r(\lambda)$, make each of the last two
terms $O_T(\mathfrak r(\lambda))$.  Equation~\eqref{eq:g1-bc-boundary} removes
the denominator defect exactly, proving \eqref{eq:g1-bc-vanishing}.
\end{proof}

Theorem~\ref{thm:g1-boundary-corrected} transfers the proved continuum G1
bound to a finite vector with exact endpoint normalization. It is a weak
fixed-band estimate. Increasing the cutoff does not make the full
derivative of the endpoint repair small: its squared norm is
\[
 |d_{\lambda,N}|^2\frac{8\pi^2}{K^2L}
 \sum_{j=K/2+1}^{K}j^2\le
 4\pi^2|d_{\lambda,N}|^2K/L.
\]
The growing derivative is relevant to any later graph argument.
The scalar-shift identity remains exact,
\[
 D_{\log}(W-\alpha I)k=D_{\log}Wk-\alpha D_{\log}k,
\]
but a weak bound on the first term supplies no norm bound on the second.

For filtered endpoint recovery, if the local Lipschitz constant is zero,
the near-endpoint error is identically zero; the cutoff is chosen from
the far-tail error alone. A formula that divides by that Lipschitz
constant is used only in its positive case. All other displayed cutoff
thresholds depend on the specified source, its variation, and the desired
tolerance, without using a list of zeta zeros.

The next section restores an explicit weighted-sampler estimate from the
earlier manuscript. The subsequent sampler construction then separates
ordinary norm, physical endpoint, Weil form, and graph behavior. These
four notions of compatibility cannot be interchanged.


\subsection*{Ordinary source mass and a full operator residual}
The Sobolev G1 estimate and the polynomial Fourier tail also give an
absolute ordinary-norm operator residual for this fixed source. This supplies the
Rayleigh and residual prerequisites in the spectral-overlap criterion;
it does not supply the overlap itself.

\begin{lemma}[A nonzero ordinary-norm limit of the repaired proxy]
\label{lem:v115-source-mass}
Put
\[
 h_\infty(u)=\frac{4\pi}{\sqrt3}u^2(2\pi u^2-3)e^{-\pi u^2},
 \qquad f_\infty=\mathcal E(h_\infty).
\]
Extend $p_\lambda$ by zero outside $I_\lambda$. Then
\begin{equation}\label{eq:v115-source-limit}
 \norm{p_\lambda-f_\infty}_{L^2(du/u)}=O(\lambda^{-1/2}),
 \qquad Jf_\infty=f_\infty.
\end{equation}
In particular $0<\norm{f_\infty}<\infty$. For
$N=N_\lambda=\lceil\lambda^8(1+L)\rceil$ and
$k_\lambda=P_Np_\lambda$, both $\norm{p_\lambda}$ and
$\norm{k_\lambda}$ tend to $\norm{f_\infty}$ and are bounded above
and away from zero. A valid eventual common lower bound is
\begin{equation}\label{eq:v115-mass-lower}
 m_*:=\frac14\left(\int_1^2h_\infty(u)^2\,du\right)^{1/2}>0.
\end{equation}
No numerical threshold in $\lambda$ is asserted.
\end{lemma}
\begin{proof}
The uniform fixed-mode Hermite approximation in
\cite[Lemma~7.2]{ConnesConsaniMoscovici2025} gives suitably normalized
$q_{j,\lambda}$ with
$\sup_{|u|\le\lambda}|q_{j,\lambda}(u)-H_j(u)|=O(\lambda^{-2})$,
$j=0,4$, where $H_j$ has unit $L^2(\mathbb R)$ norm. This rate survives
the manuscript's exact unit normalization: expanding the squared norm,
the cross term is $O(\lambda^{-2})\norm{H_j}_1$, the squared-error term
is $O(\lambda^{-3})$, and the omitted Hermite tail is Gaussian.
Thus $\norm{q_{j,\lambda}}=1+O(\lambda^{-2})$.

The limits of the coefficients in \eqref{eq:v12-coefficients} are
$-H_4(0)=-\sqrt3/(2\,2^{1/4})$ and $H_0(0)=2^{1/4}$, with error
$O(\lambda^{-2})$. Substitution of the normalized Hermite polynomials
gives exactly the displayed $h_\infty$, including its scalar factor.
The fixed repair has exponentially small amplitude by
\eqref{eq:v12-zero-value}. Consequently
\[
 \sup_{|u|\le\lambda}|\widetilde h_\lambda(u)-h_\infty(u)|
       \le C\lambda^{-2}.
\]
The limiting source is a Fourier-invariant combination of orders $0,4$,
with zero value at zero and zero integral. Poisson summation therefore
gives $Jf_\infty=f_\infty$. Its Gaussian decay at infinity, followed by
inversion at zero, gives square integrability.

On the physical window there are at most $\lambda/u$ compact-source
summands. Their accumulated approximation error is at most
$C\lambda^{-1}u^{-1/2}$. Its squared Haar norm is bounded by
\[
 C\lambda^{-2}\int_{1/\lambda}^\lambda u^{-2}\,du
       \le C\lambda^{-1}.
\]
The omitted Gaussian summands $nu>\lambda$ contribute an exponentially
smaller norm: bound their decreasing polynomial-Gaussian tail by its
first term plus $u^{-1}$ times its tail integral. The same estimate and
inversion control $f_\infty$ outside $I_\lambda$. Inversion averaging
is an $L^2$ contraction, proving \eqref{eq:v115-source-limit}.
For $u\ge1$, all summands $h_\infty(nu)$ are positive. Thus
$\norm{f_\infty}^2\ge\int_1^2h_\infty(u)^2\,du>0$.
Finally \eqref{eq:v114-high-tail} gives
\[
 \norm{(I-P_N)p_\lambda}
       \le C|B_\lambda|\lambda\sqrt{L/N}\longrightarrow0,
\]
so the projected vectors have the same nonzero norm limit. Choosing a
common sufficiently large threshold proves \eqref{eq:v115-mass-lower}.
\end{proof}

\begin{proposition}[A full residual after ordinary norm normalization]
\label{prop:v115-absolute-residual}
Use the same $N_\lambda,k_\lambda$ and the canonical closed operator
$\mathsf W_\lambda$ of Proposition~\ref{prop:v114-weil-core}.
For all sufficiently large $\lambda$, with $c=2\pi\lambda^2$,
\begin{equation}\label{eq:v115-absolute-residual}
 \norm{\mathsf W_\lambda p_\lambda}
 +\norm{\mathsf W_\lambda k_\lambda}
 +\norm{W_{\lambda,N_\lambda}k_\lambda}
       \le C\lambda^6 e^{-c/3}.
\end{equation}
For either source normalized to ordinary norm one, its Rayleigh value
and its full centered residual tend to zero at this rate. The finite
matrix assertion uses that same normalization and the original physical
Fourier cutoff.
\end{proposition}
\begin{proof}
All source vectors belong to $\mathcal D(\mathsf W_\lambda)$ by
Proposition~\ref{prop:v114-weil-core}. The weighted Chebyshev estimate
$P_\lambda=O(\lambda)$, together with $h_L=O(\lambda)$, gives
$\norm{[M_b,H]}=O(\lambda)$ in \eqref{eq:v114-b-bound}.
Corollary~\ref{cor:g1-complete-matrix}, including its separately proved
diagonal, gives uniformly
\[
 |d_n|\le C\{\lambda+1+\log(2+|n|/L)\}.
\]
Let $N=N_\lambda$ and $e=(I-P_N)p_\lambda$.
Using \eqref{eq:v114-high-tail} and comparing
$\sum_{n>N}\log^2(2+n/L)/n^2$ with its integral yields
\begin{align}
 \norm{\mathsf W_\lambda e}
 &\le C|B_\lambda|\lambda\sqrt{L/N}
        \{\lambda+1+\log(2+N/L)\}\notag\\
 &\le C|B_\lambda|\lambda^{-2}.
 \label{eq:v115-uniform-graph-tail}
\end{align}
Both the diagonal part and the bounded commutator are retained.

Next split the \emph{output} at an auxiliary integer $M>2N$.
This does not change the candidate or its finite matrix. For any
$v=P_Mv$ with $\norm v=1$,
\[
 \norm{D_{\log}v}\le2\pi M/L,\qquad
 |v(-a)|+|v(a)|\le2\sqrt{(2M+1)/L}.
\]
The full Sobolev estimate \eqref{eq:G1-sobolev}, whose constant is
independent of its test function, therefore implies
\begin{equation}\label{eq:v115-low-output}
 \norm{P_M\mathsf W_\lambda p_\lambda}
 \le C|B_\lambda|\mathfrak r(\lambda)
       \{1+M/L+\sqrt{M/L}\}.
\end{equation}
This includes the zero Fourier mode. It uses the already established
compact-source radicality and centered pole--prime cancellation, with
both source moments and inversion averaging intact. No fixed-band
constant is extrapolated to a growing band.

The diagonal has no cross block between $E_N$ and $E_M^\perp$.
For the complete off-diagonal matrix, Lemma~\ref{lem:g1-high-low}
gives the Hilbert--Schmidt bound
\[
 \norm{(I-P_M)\mathsf W_\lambda P_N}^2
 \le C\lambda^2\sum_{|m|\le N}\sum_{|n|>M}\frac1{(n-m)^2}
 \le C\lambda^2N/M.
\]
Lemma~\ref{lem:v115-source-mass} and
\eqref{eq:v115-uniform-graph-tail} now give
\begin{equation}\label{eq:v115-high-output}
 \norm{(I-P_M)\mathsf W_\lambda p_\lambda}
 \le C\lambda^5\sqrt{1+L}\,M^{-1/2}
         +C|B_\lambda|\lambda^{-2}.
\end{equation}
Thus all prime, pole and archimedean pieces remain in the actual matrix;
no divergent exterior pole or prime norm sum is estimated separately.

Take $M=\lceil e^{2c/3}\rceil$, which exceeds $2N$ eventually.
By \eqref{eq:v12-combined-endpoint} and
\eqref{eq:v13-endpoint-comparison},
$|B_\lambda|\le C\lambda^{11/2}e^{-c}$, while
$\mathfrak r(\lambda)$ is bounded. Equations
\eqref{eq:v115-low-output}--\eqref{eq:v115-high-output} are therefore
at most $C\lambda^6e^{-c/3}$ in total. This proves the first term of
\eqref{eq:v115-absolute-residual};
\eqref{eq:v115-uniform-graph-tail} proves the second, and compression
contractivity proves the third.
Finally ordinary normalization is legitimate by
\eqref{eq:v115-mass-lower}. For a unit vector $u$, letting
$\alpha=\ip{\mathsf W_\lambda u}u$, one has
$|\alpha|\le\norm{\mathsf W_\lambda u}$ and
$\norm{(\mathsf W_\lambda-\alpha)u}\le2\norm{\mathsf W_\lambda u}$.
The same argument applies to the compressed matrix.
\end{proof}

The preceding proposition supplies no endpoint-normalized operator residual:
dividing its bound by $|B_\lambda|$ leaves an exponentially growing
upper bound. Nor does it identify the spectral point near zero as the
lowest one. The elementary family $\operatorname{diag}(-1,\epsilon)$,
with candidate $(0,1)$, has Rayleigh value $\epsilon$ and centered
residual zero while its lowest level stays at $-1$. An independent
quantitative ground-space overlap or signed complementary lower bound
is still necessary. No estimate on the corrected sampler's graph defect,
or on a growing family of prolate modes, is inferred here.

\subsection*{G2 calibration from the screw-function realization}
Recent work of Suzuki~\cite{Suzuki2026} gives a useful, but more limited, structural calibration of G2.  Let $A_a$ denote the lower-bounded self-adjoint operator associated with the semilocal Weil form on $(-a,a)$, and write
\[
 \lambda_a=\inf\sigma(A_a).
\]
Suzuki proves that $A_a$ has discrete spectrum bounded from below and that $a\mapsto\lambda_a$ is continuous.  He also proves positivity, simplicity, and evenness of the ground state for \emph{sufficiently small} $a$ by a separate small-support asymptotic argument.  Those simplicity/parity conclusions must not be extrapolated to arbitrary growing support; the large-$a$ spectral ordering needed here remains open.

There is nevertheless an elementary global monotonicity which is important for claim discipline.  If $a_2>a_1$, extension by zero identifies the smaller compactly supported form domain with a subspace of the larger one, while the global Weil form is unchanged on such vectors.  The min--max principle therefore gives
\begin{equation*}
 \boxed{\lambda_{a_2}\le \lambda_{a_1}.}
 \tag{G2.1}\label{eq:g2-monotone}
\end{equation*}
In particular, once $\lambda_a$ becomes negative it remains bounded above by that same negative value at every larger support.  This is compatible with Suzuki's global continuity theorem and his observation that failure of RH is equivalent to the existence of a support with $\lambda_a<0$.

Weak fixed-band G1 alone does not test the form on the changing
proxy itself. The stronger operator estimate in
Proposition~\ref{prop:v115-absolute-residual} now supplies that missing
Rayleigh prerequisite for the actual normalized source and its literal
polynomial projection: $q_a(u_a,u_a)\to0$. Hence monotonicity gives
\[
 \lambda_a\to0\quad\Longleftrightarrow\quad
 \lambda_a\ge0\text{ for every finite }a.
\]
The forward implication uses only monotonicity; the reverse combines
positivity with the proved Rayleigh upper bound. The right side is the
compact-support Weil positivity criterion. The new source estimate
proves the upper bound, not positivity or the asserted limit of the
bottom level.

A quantitative spectral-overlap formulation isolates one possible extra ingredient without assuming simplicity.  Let $P_a^{(0)}$ be the orthogonal projection onto the (finite-dimensional) lowest eigenspace of $A_a$, let $p_a\in\mathcal D(A_a)$ be a normalized explicit proxy, and put
\[
 \alpha_a=\langle A_ap_a,p_a\rangle,
 \qquad
 r_a=(A_a-\alpha_a)p_a,
 \qquad
 \kappa_a=\|P_a^{(0)}p_a\|.
\]
Spectral expansion gives the exact bound
\[
 \boxed{
 |\lambda_a-\alpha_a|\,\kappa_a\le \|r_a\|.
 }
 \tag{G2.3}
\]
Proposition~\ref{prop:v115-absolute-residual} now proves both
$\alpha_a\to0$ and $\|r_a\|\to0$ for the normalized actual source,
with $\lambda=e^a$ and an upper envelope
\[
 \varepsilon_\lambda=C\lambda^6e^{-2\pi\lambda^2/3}.
\]
Therefore an independently proved overlap bound
$\kappa_a\ge\kappa_0>0$ would imply $\lambda_a\to0$ and close G2.
More generally it suffices that
$\varepsilon_\lambda/\kappa_a\to0$ along an unbounded sequence;
for example, $\kappa_a\ge e^{-\theta c}$ with fixed
$\theta<1/3$ and $c=2\pi\lambda^2$ would suffice.
No such lower bound is proved. If instead
$\lambda_a\le-\delta<0$, spectral projection of
$A_au_a$ gives $\delta\kappa_a\le\|A_au_a\|$.
Thus a negative ground state can coexist with the new residual estimate
by having exponentially small overlap with this source.

At small support Suzuki's positivity-improving analysis gives a simple
positive ground state, but it does not furnish a uniform large-support
overlap bound. Even positivity of both vectors would imply only a
nonzero overlap, without its required quantitative size. The model
$\operatorname{diag}(-1,\epsilon)$ with source $(0,1)$ shows directly
why a vanishing source residual does not determine the bottom sign.

This calibration is deliberately restrictive: continuity and monotonicity of the bottom level are structural, while large-support simplicity/parity and quantitative spectral overlap remain part of the open RH-strength content.  A successful continuation must therefore find new large-support structure, not extrapolate the small-support Dirichlet-form theorem.

\subsection*{A Schur-complement formulation of the large-support obstruction}
The preceding overlap criterion singles out one vector.  A more robust formulation does not require identification of the ground state and should be stated at the level of the closed semilocal Weil form $q_a$, rather than on the operator domain of $A_a$.  Let $P_a$ be any finite-rank orthogonal projection whose range is contained in the form domain, and decompose
\[
 L^2(-a,a)=\mathcal R_a\oplus\mathcal R_a^\perp,
 \qquad \mathcal R_a=\operatorname{Ran}P_a.
\]
The following elementary estimate isolates exactly the diagonal, cross, and complement quantities that a prolate near-radical construction would have to control.

\begin{proposition}[Form-Schur lower bound]\label{prop:g2-schur}
Suppose that for all $r\in\mathcal R_a$ and all $s\in\mathcal R_a^\perp$ in the form domain,
\[
 q_a(r,r)\ge-\alpha_a\norm r^2,
 \qquad
 |q_a(r,s)|\le\beta_a\norm r\norm s,
 \qquad
 q_a(s,s)\ge c_a\norm s^2
\]
with $c_a>0$, $\alpha_a\ge0$, and $\beta_a\ge0$. Then
\[
 \boxed{
 \lambda_a\ge -\alpha_a-\frac{\beta_a^2}{c_a}.
 }
 \tag{G2.4}
\]
In particular, $\alpha_a\to0$ and $\beta_a^2/c_a\to0$ imply $\liminf_{a\to\infty}\lambda_a\ge0$.
\end{proposition}
\begin{proof}
For $x=r+s$ the hypotheses give
\[
 q_a(x,x)
 \ge -\alpha_a\norm r^2-2\beta_a\norm r\norm s+c_a\norm s^2.
\]
Completing the square in $\norm s$ yields
\[
 q_a(x,x)
 \ge -\left(\alpha_a+\frac{\beta_a^2}{c_a}\right)\norm r^2
 \ge -\left(\alpha_a+\frac{\beta_a^2}{c_a}\right)\norm x^2.
\]
Taking the infimum of the Rayleigh quotient proves the claim.
\end{proof}

\subsection*{A complement bound without a positive spectral gap}
The ordinary operator residual permits a different reduction from the
coercive Schur estimate.  It does not require a positive gap on the
complement and therefore introduces no division by a collapsing level.

\begin{proposition}[A gap-free near-radical block bound]
\label{prop:v118-gap-free}
Let $A$ be a lower-bounded self-adjoint operator with closed form $q$,
and let $P$ be a finite-rank orthogonal projection with
$\operatorname{Ran}P\subset\mathcal D(A)$.  Put $Q=I-P$ and assume
$\norm{AP}\le\epsilon$.  Then $Q$ preserves the form domain and
\begin{equation}\label{eq:v118-gap-free-difference}
 \left|q(f,f)-q(Qf,Qf)\right|
 \le\frac{2}{\sqrt3}\epsilon\norm f^2
 \qquad(f\in\mathcal D(q)).
\end{equation}
Consequently, if
\[
 q(g,g)\ge-\eta\norm g^2\quad
 (g\in\operatorname{Ran}Q\cap\mathcal D(q)),\qquad \eta\ge0,
\]
then
\begin{equation}\label{eq:v118-gap-free-lower}
 \inf\sigma(A)\ge-\eta-\frac{2}{\sqrt3}\epsilon.
\end{equation}
The constant $2/\sqrt3$ in both statements is sharp.
\end{proposition}
\begin{proof}
Write $h=Pf$ and $g=Qf$.  The vector $h$ lies in the operator domain,
so $g$ belongs to the form domain and, with the fixed first-slot-linear
convention,
\[
 q(f,f)-q(g,g)
 =\ip{Ah}{h}+2\Re\ip{Ah}{g}
 =\Re\ip{Ah}{h+2g}.
\]
Since $h\perp g$, the absolute value is at most
\[
 \epsilon\norm h\sqrt{\norm h^2+4\norm g^2}
 \le\frac{2}{\sqrt3}\epsilon
           (\norm h^2+\norm g^2).
\]
For the last inequality, the maximum of
$t\sqrt{4-3t^2}$ on $0\le t\le1$ is $2/\sqrt3$.
This proves \eqref{eq:v118-gap-free-difference}; the complementary
lower bound and $\norm{Qf}\le\norm f$ prove
\eqref{eq:v118-gap-free-lower}.  Sharpness already holds on
$\mathbb C^2$: take $P$ onto the first coordinate and
\[
 A=-\frac{\epsilon}{\sqrt3}
   \begin{pmatrix}1&\sqrt2\\\sqrt2&0\end{pmatrix}.
\]
Then $\norm{AP}=\epsilon$, the complementary form is zero, and the
lowest eigenvalue is $-2\epsilon/\sqrt3$.
\end{proof}

For the rank-one projection onto either normalized source of
Proposition~\ref{prop:v115-absolute-residual},
$\norm{\mathsf W_\lambda P}=\norm{\mathsf W_\lambda u_\lambda}$
already tends to zero with the proved ordinary residual bound.
It therefore suffices to prove asymptotic nonnegativity
$q_\lambda|_{u_\lambda^\perp}\succeq-\eta_\lambda I$ with
$\eta_\lambda\to0$ along an unbounded window sequence.  Strict
complementary coercivity is unnecessary for this version of the
reduction.  For a growing source block the required uniform
$\norm{\mathsf W_\lambda P_\lambda}\to0$ remains a separate open
input.  Neither complementary sign estimate is proved here.

The missing sign cannot be concealed by removing the near-radical
block.  Indeed, suppose a fixed unit compact test $f$ has
$q(f,f)=-\delta<0$, and extend it by zero to every larger window.
Then \eqref{eq:v118-gap-free-difference} gives
\[
 q_\lambda((I-P_\lambda)f,(I-P_\lambda)f)
 \le-\delta+\frac{2}{\sqrt3}\epsilon_\lambda.
\]
Whenever $\epsilon_\lambda<\sqrt3\delta/2$, the projected vector
is nonzero and negative; its normalized Rayleigh value is bounded
above by the same negative quantity, since its norm is at most one.
Thus a negative direction persists in the complement of any uniformly
near-radical block, regardless of its rank.  By the existing cofinal
Weil-positivity criterion, the desired asymptotic complementary sign
for the proved rank-one source remains RH-strength.  The reduction
removes a gap denominator; it does not supply the arithmetic sign.

\subsection*{Strong cancellation persists in the exact source complement}

The proved ordinary source residual does not mean that significant
prime/non-prime cancellation is confined to one direction. The following
fixed-window certificate tests that possible shortcut on the literal
source, rather than on sampled basis or random vectors.

\begin{proposition}[A certified cancelling direction off the source]
\label{prop:v121-complement-cancellation}
At $\lambda=3$ there exists an inversion-even unit Fourier polynomial
$v\in E_{64}$ exactly orthogonal to the repaired source $p_3$ such that
\begin{equation}\label{eq:v121-complement-cancellation}
 \begin{gathered}
 \ip{v}{p_3}=0,\qquad
 0<QW_3(v,v)<3\cdot10^{-31},\\
 q_{\rm np}(v)>0.4,\qquad q_{\rm pr}(v)<-0.4.
 \end{gathered}
\end{equation}
Here $q_{\rm np}=q_{\rm ar}+q_{\rm po}$ includes both the pure
archimedean term and the pole term, and $q_{\rm pr}$ is the signed
prime contribution, so $QW_3=q_{\rm np}+q_{\rm pr}$.
\end{proposition}
\begin{proof}
Let $V$ have the two real dyadic columns supplied in
Appendix~\ref{app:v121-audit}, in the normalized even basis of $E_{64}$.
Let $c$ be the exact rational source candidate from
Proposition~\ref{prop:v113-source-certificate}, converted to the same
parity coordinates. Set $a=V^*c$ and $x=(-a_2,a_1)^t$. The exact
algebra gives $\ip{Vx}{c}=0$. The certified Gram matrix $V^*V$ is
positive definite and $\norm{Vx}>0$, so put $w=Vx/\norm{Vx}$.
All vectors are real here; the stated Hermitian convention is unchanged.

Outward interval evaluation of the separate finite matrices gives
\[
 \begin{aligned}
 2.5844134\cdot10^{-31}&<QW_3(w,w)<2.5844135\cdot10^{-31},\\
 q_{\rm np}(w)&>0.4173237590,\\
 q_{\rm pr}(w)&<-0.4173237590.
 \end{aligned}
\]
These are inequalities for exact dyadic/rational constructions;
approximate eigenvectors were used only to propose the dyadic columns.

Write $u=P_{64}p_3/\norm{P_{64}p_3}$ and $\widehat c=c/\norm c$.
The existing exact-source certificate, matched to $c$ by its file hash,
gives $\norm{P_u-P_{\widehat c}}<\delta=4\cdot10^{-36}$.
Since $w\perp c$, define
\[
 v=\frac{(I-P_u)w}{\norm{(I-P_u)w}}.
\]
Then $\norm{v-w}<2\delta$. Also $v\in E_{64}$, hence $v\perp u$
implies $v\perp p_3$, not just orthogonality to an approximate source.
For any finite Hermitian matrix $H$,
\[
 |\ip{Hv}{v}-\ip{Hw}{w}|
 \le2\norm H\norm{v-w}<4\delta\norm H.
\]
The certified all-row bounds for the compressed matrices are
\[
 \norm{W_{64}}<8.926,\qquad
 \norm{H_{\rm np,64}}<4.176,\qquad
 \norm{H_{\rm pr,64}}<5.490.
\]
Applying these bounds to the full interval enclosures, before rounding,
gives
\[
 2.5829853\cdot10^{-31}<QW_3(v,v)<2.5858415\cdot10^{-31}
\]
and preserves both strict $0.4$ bounds. A finite-support polynomial's
complete quadratic form equals its finite matrix quadratic form.
No invariant-subspace or full operator-residual assertion is required.
\end{proof}

Thus $QW_3(v,v)/q_{\rm np}(v)<7.5\cdot10^{-31}$ inside the exact
source complement. This rules out an explanation in which substantial
archimedean-plus-pole dominance eliminates the need for cancellation
everywhere off that source. It does not rule out positivity on the
complement, and exhibits neither a negative Weil direction nor an
additional exact null vector. Positive basis-vector or typical
random-vector energies cannot control these coherent linear combinations.
The pole term must be included explicitly: calling $q_{\rm np}$ the
pure archimedean contribution would change the asserted comparison.


\subsection*{Why a finite source block cannot reduce the unsigned prime norm}

The arithmetic bound in the preceding tail estimates is not merely a
poor choice of positive weight.  There is an exact recurrence obstruction
to replacing it by a small norm on the source complement.

\begin{proposition}[Prime norm survives every finite-dimensional removal]
\label{prop:v119-prime-essential}
Fix $\lambda>1$ and let $\mathsf T_{\rm pr}$ be the positive-coefficient
sum of truncated prime translations in
\eqref{eq:v116-weighted-prime-general}.  For every finite-rank
orthogonal projection $P$ on $L^2(0,L)$, $Q=I-P$, one has
\begin{equation}\label{eq:v119-prime-compression}
 \|Q\mathsf T_{\rm pr}Q\|=\|\mathsf T_{\rm pr}\|.
\end{equation}
For every compact operator $C$ on this space,
\begin{equation}\label{eq:v119-prime-compact}
 \|Q(\mathsf T_{\rm pr}-C)Q\|\ge\|\mathsf T_{\rm pr}\|.
\end{equation}
In particular the essential norm of $\mathsf T_{\rm pr}$ equals its
ordinary norm.  The projection need not be a Fourier projection and
may remove any prescribed finite source block.
\end{proposition}
\begin{proof}
For each positive integer $k$, let
$M_k f(x)=e^{2\pi ikx/L}f(x)$.  Simultaneous recurrence of the finitely
many phases gives integers $k_j\to\infty$ such that
\[
 e^{2\pi ik_j\log m/L}\longrightarrow1
 \quad\hbox{for every active }m.
\]
For completeness, partition the unit cube in the finitely many phase
coordinates into $H^d$ cubes and compare $H^d+1$ successive multiples.
Two lie in one cube, giving a positive return with coordinate errors
at most $1/H$.  If such returns have an unbounded subsequence, use it;
otherwise a fixed positive return has all coordinate errors zero,
and its increasing multiples give the required sequence.  No rational
independence of prime logarithms is assumed.

Conjugating a truncated translation by $M_k$ only multiplies it by
its corresponding phase.  Since each translation has norm at most one,
\begin{equation}\label{eq:v119-prime-recurrence}
 \|M_{k_j}^*\mathsf T_{\rm pr}M_{k_j}-\mathsf T_{\rm pr}\|
 \le2\sum_m w_m|e^{2\pi ik_j\log m/L}-1|\longrightarrow0.
\end{equation}
For every fixed $f\in L^2(0,L)$, $M_{k_j}f\rightharpoonup0$:
each scalar product is a Fourier coefficient of an $L^1$ function.
Thus $PM_{k_j}f\to0$, $PM_{k_j}\mathsf T_{\rm pr}f\to0$, and
$CM_{k_j}f\to0$ strongly.  Combining these statements with
\eqref{eq:v119-prime-recurrence} gives
\[
 Q(\mathsf T_{\rm pr}-C)Q M_{k_j}f
       -M_{k_j}\mathsf T_{\rm pr}f\longrightarrow0.
\]
For unit $f$, the input has norm at most one, so taking norms and then
the supremum over $f$ proves \eqref{eq:v119-prime-compact}.
Putting $C=0$ and using compression contractivity proves
\eqref{eq:v119-prime-compression}.  Taking the infimum over compact
$C$, with $P=0$, proves the essential-norm assertion.
\end{proof}

\begin{proposition}[The sharp growing-window scale of the unsigned prime norm]
\label{prop:v119-prime-scale}
With the actual weights $w_m=\Lambda(m)/\sqrt m$ and $L=2\log\lambda$,
\begin{equation}\label{eq:v119-prime-scale}
 \|\mathsf T_{\rm pr}\|=(1+o(1))\lambda
       \qquad(\lambda\longrightarrow\infty).
\end{equation}
Consequently \eqref{eq:v119-prime-scale} also holds for
$\|Q_\lambda\mathsf T_{\rm pr}Q_\lambda\|$ for any family of
finite-rank orthogonal projections $I-Q_\lambda$, with no restriction on their
finite ranks or Fourier cutoffs.
\end{proposition}
\begin{proof}
Write
\[
 a(x)=e^{-x/2},\quad b(x)=e^{-(L-x)/2},\quad\phi(x)=a(x)+b(x),
\]
and
\[
 \Psi(X)=\sum_{1<m<X}\Lambda(m),\quad
 S(X)=\sum_{1<m<X}\frac{\Lambda(m)}m.
\]
The strict cutoff agrees with the truncated translations almost
everywhere.  Direct evaluation gives
\begin{equation}\label{eq:v119-prime-weight-ratio}
 \frac{(\mathsf T_{\rm pr}\phi)(x)}{\phi(x)}
 =\frac{a(x)\{S(e^{L-x})+\Psi(e^x)\}
       +b(x)\{\Psi(e^{L-x})+S(e^x)\}}{a(x)+b(x)}.
\end{equation}
The exact identity
\[
 a(x)e^x+b(x)e^{L-x}=\lambda\{a(x)+b(x)\}
\]
identifies the leading constant.  The classical prime number
theorem~\cite{NISTPrimeAsymptotics}, equivalently $\Psi(X)\sim X$,
implies that for each $\delta>0$ there is $C_\delta<\infty$ with
\[
 |\Psi(X)-X|\le\delta X+C_\delta\qquad(X\ge1).
\]
The Chebyshev bound and partial summation give
$0\le S(X)\le C(1+\log X)$, with an absolute $C$.
Therefore, uniformly for $0<x<L$,
\[
 (1-\delta)\lambda-C_\delta
 \le\frac{\mathsf T_{\rm pr}\phi}{\phi}
 \le(1+\delta)\lambda+C_\delta+C(1+L).
\]
At each fixed window $\phi$ is bounded and bounded away from zero.
The upper estimate and the weighted Schur inequality
\eqref{eq:v116-weighted-prime-general} bound the operator norm from
above.  Integrating the lower estimate against $\phi^2$ gives the
same lower bound for its Rayleigh quotient and hence for the norm.
Divide by $\lambda$, let $\lambda\to\infty$, and then let
$\delta\downarrow0$. This proves \eqref{eq:v119-prime-scale}.
The projection assertion follows from
Proposition~\ref{prop:v119-prime-essential} at each window.
\end{proof}

These statements concern the unsigned prime operator, not the sign of
the full Weil form.  They make precise a limitation of the particular
place-by-place bound \eqref{eq:v115-sharp-tail}. Even replacing its
prime constant by the exact tail norm cannot certify positivity unless
\begin{equation}\label{eq:v119-unsigned-cutoff}
 \log\frac{N+1}{L}>\|\mathsf T_{\rm pr}\|
       =(1+o(1))\lambda.
\end{equation}
Here we refer to the whole-space bound, whose remaining subtracted
terms are nonnegative. Thus this method still requires an exponential
Fourier cutoff in $\lambda$, rather than the polynomial source cutoff.
This is a limitation of that sufficient scalar bound, not a necessary
cutoff for actual Weil positivity or for a frequency-weighted argument.

The physical pole operator has rank at most two.  Equation
\eqref{eq:v119-prime-compact} shows that combining it with the prime
operator does not make their ordinary operator norm smaller than
$\|\mathsf T_{\rm pr}\|$, even after a finite source removal.
The recurrence frequencies in the proof may be arbitrarily large;
there the archimedean logarithmic energy also grows.  No negative
direction for the complete Weil operator is inferred.  A scalable G2
estimate must control the arithmetic term together with that energy
or retain more precise signed correlations, rather than assert a small
unweighted norm on the complement.  The required signed estimate
remains open.

\subsection*{Retaining the archimedean energy in a one-sided estimate}

Choose a literal Fourier tail on which $a_n:=-A_n>0$ and write
$D=\operatorname{diag}(a_n)$ there. By
\eqref{eq:v114-weil-decomposition}, the compressed closed form has
$\mathsf W=D+R$, with $R$ bounded self-adjoint. It contains the
actual signed prime and pole terms and the archimedean off-diagonal;
no arithmetic diagonal is absorbed into $D$. Put
\begin{equation}\label{eq:v120-weighted-remainder}
 \mathcal B=D^{-1/2}R D^{-1/2}.
\end{equation}
The operator $\mathcal B$ is compact self-adjoint, since
$a_n\sim\log|n|$ at the fixed window. The form identity is
\[
 q_{\mathsf W}[f]
 =\ip{(I+\mathcal B)D^{1/2}f}{D^{1/2}f}.
\]
Thus the exact positivity target is $\mathcal B\succeq-I$;
the stronger two-sided condition $\|\mathcal B\|\le1$ can fail
because of a harmless positive direction.

\begin{proposition}[The scale of an approximate weighted sign]
\label{prop:v120-weighted-error}
Suppose $D\succeq d_0I>0$, $R\succeq-CI$, $C\ge0$, and
$D^{-1/2}RD^{-1/2}\succeq-(1+\eta)I$, $\eta\ge0$.
Then, on the entire form domain,
\begin{equation}\label{eq:v120-weighted-error}
 D+R\succeq-\frac{\eta C}{1+\eta}I.
\end{equation}
The constant is sharp. In particular, a dimensionless error
$\eta_\lambda\to0$ is not by itself an ordinary lower error
$o(1)$ when $C_\lambda$ grows.
\end{proposition}
\begin{proof}
The hypotheses give $D+R\succeq-\eta D$ and
$D+R\succeq D-CI$. Add these inequalities with weights
$1/(1+\eta)$ and $\eta/(1+\eta)$ respectively. For $C>0$,
the scalar example $D=C/(1+\eta)$, $R=-C$ attains equality.
Taking $C_\lambda=\lambda$ and $\eta_\lambda=1/\lambda$
shows why the missing scale factor cannot be omitted. A sufficient
conversion requires $\eta_\lambda C_\lambda\to0$, or a more
precise directional energy estimate.
\end{proof}

There is also a limitation on how the compact weighted remainder can
be estimated. A Hilbert--Schmidt budget for its whole infinite matrix
is unavailable, despite the valid finite-column residual Grams above.

\begin{proposition}[Weighted compactness without a square-summable matrix]
\label{prop:v120-not-schatten}
At any fixed $\lambda>\sqrt2$, let $D$ be a positive diagonal on a
fixed Fourier tail with $a_n\sim\log|n|$. The compact operator
$D^{-1/2}\mathsf T_{\rm pr}D^{-1/2}$ belongs to no Schatten class
$\mathcal S_p$, $0<p<\infty$. For the exact archimedean diagonal
$a_n=-A_n$, the same assertion holds for
$\mathcal B$ in \eqref{eq:v120-weighted-remainder}.
\end{proposition}
\begin{proof}
The prime diagonal is the nonzero finite trigonometric sequence
\[
 \tau_n=2\sum_{1<m<\lambda^2}
 \frac{\Lambda(m)}{\sqrt m}
 \left(1-\frac{\log m}{L}\right)
 \cos\frac{2\pi n\log m}{L}.
\]
Collect equal unit-circle frequencies to write
$\tau_n=\sum_{z\in F}c_z z^n$, with all $c_z>0$.
Finite geometric sums give
\[
 \frac1X\sum_{n=1}^X|\tau_n|^2\longrightarrow
 \sum_{z\in F}c_z^2>0.
\]
Boundedness then implies that $|\tau_n|$ exceeds a positive
constant on a set of positive lower density. Consequently, for each
$p\ge1$, $\sum_n|\tau_n/a_n|^p=\infty$: the partial sum on
that set is bounded below by a constant times $X/(\log X)^p$.
For self-adjoint $K\in\mathcal S_p$, Jensen's inequality in its
spectral measure gives
$\sum_n|\ip{KU_n}{U_n}|^p\le\operatorname{Tr}|K|^p$.
This is a contradiction. The inclusions
$\mathcal S_p\subset\mathcal S_1$ for $0<p<1$ settle the
remaining exponents. For the complete remainder its diagonal is
$(-\tau_n+O_\lambda(n^{-2}))/a_n$, since the exact pole diagonal
has this decay and the archimedean remainder has zero diagonal.
The same argument applies. None of this invalidates the separated
finite-column moment bounds, whose off-diagonal factor
$(n-m)^{-1}$ is square summable in the remote row index.
\end{proof}

\begin{proposition}[Sharp fixed-window decay of the weighted prime tail]
\label{prop:v120-weighted-tail-asymptotic}
Under the diagonal assumptions above, let $Q_J$ remove $|n|\le J$.
At each fixed window,
\begin{equation}\label{eq:v120-weighted-tail-asymptotic}
 \lim_{J\to\infty}(\log J)
 \|Q_JD^{-1/2}\mathsf T_{\rm pr}D^{-1/2}Q_J\|
       =\|\mathsf T_{\rm pr}\|.
\end{equation}
Subtracting any fixed compact operator from $\mathsf T_{\rm pr}$
before weighting leaves this limit unchanged. No uniformity in
$\lambda$ is asserted.
\end{proposition}
\begin{proof}
The upper bound is
$\|\mathsf T_{\rm pr}\|/\inf_{|n|>J}a_n$.
For the lower bound, fix $\epsilon>0$ and choose a unit Fourier
polynomial $f$, supported on $|m|\le M$, with
$|\ip{\mathsf T_{\rm pr}f}{f}|>\|\mathsf T_{\rm pr}\|-\epsilon$.
Returns of the finite vector of shift phases to any fixed neighborhood
of the identity are relatively dense. Indeed its cyclic closure is a
compact group; finitely many integer translates of that neighborhood
cover it. The largest difference between the corresponding translation
indices bounds the gap between return times. This also covers rational
phase relations.

Choose the neighborhood so that the conjugation error in
\eqref{eq:v119-prime-recurrence} is below $\epsilon$.
For every large $J$ there is therefore a return
$k\in[J+M+1,J+M+1+K]$, with $K$ independent of $J$.
Set $g=M_kf$ and $x=D^{1/2}g$. Both are supported outside $J$,
$\|x\|^2=(1+o(1))\log J$, and the weighted Rayleigh numerator
at $x$ equals $\ip{\mathsf T_{\rm pr}g}{g}$, whose absolute value
is greater than $\|\mathsf T_{\rm pr}\|-2\epsilon$.
Let $J\to\infty$ and then $\epsilon\downarrow0$.
For compact $C$ the change in the left side is at most
\[
 \frac{\log J}{\inf_{|n|>J}a_n}\|Q_JCQ_J\|\longrightarrow0.
\]
The return-length bound $K$ depends on the window and the tolerance.
Thus one may not combine \eqref{eq:v120-weighted-tail-asymptotic}
with $\|\mathsf T_{\rm pr}\|\sim\lambda$ to infer a bound on
an arbitrary joint scale $J=J(\lambda)$.
\end{proof}

\begin{proposition}[A signed energy bound at the larger window $\lambda=4$]
\label{prop:v120-lambda4-tail}
Let $L=2\log4$, retain the exact diagonal $a_n=-A_n$, and let
$Q_{16}$ be the literal projection onto $|n|>16$. For every complex
vector in its complete closed form domain,
\begin{equation}\label{eq:v120-lambda4-tail}
 QW_4(f,f)\ge10^{-8}\sum_{|n|>16}a_n|f_n|^2,
 \qquad f=Q_{16}f.
\end{equation}
Every weight in this sum is positive. This is an infinite-tail
computer-assisted bound, not a claim about the remaining low modes
or about an unbounded family of windows.
\end{proposition}
\begin{proof}
Apply the directional verified-solve proof to
$\mathsf W_4-cD$, $c=10^{-8}$, on $Q_{16}$.
Only the archimedean diagonal is changed, to $(1-c)a_n$;
the exact off-diagonal coefficients $b_n$, both pole signs, and every
prime power $2,3,4,5,7,8,9,11,13$ remain. The full-length $m=16$
translation is zero. The positive physical weight
$\phi(x)=e^{-x/2}+e^{-(L-x)/2}$ gives the certified Schur bound
\begin{equation}\label{eq:v120-lambda4-prime-bound}
 \|\mathsf T_{\rm pr}\|
 \le M_\phi<4.624042316766478.
\end{equation}
To verify it, set $u=e^x$. On each interval between
$1,16,m,16/m$ the ratio $(\mathsf T_{\rm pr}\phi)/\phi$
is a linear-fractional function of $u$, hence monotone or constant.
Every one-sided endpoint value is evaluated with outward rounding;
the chosen upper endpoint is checked against every candidate.

Use $E_L,C_L,h_L$ from \eqref{eq:v115-arch-errors}, with
$C_L=1$, $h_L=9/16$ and $R_L=64/255$. For $n>N$ the proved
remote inverse bound applies with
\[
 g_n=(1-c)\{\log(n/L)-E_L(2\pi n/L)\}-\kappa,
\]
where $\kappa=M_\phi+2C_L/t_*$ in the even sector and
$\kappa=M_\phi+2C_L/t_*+\pi/2+4h_LL/(\pi^2N)$ in the odd
sector, $t_*=2\pi(N+1)/L$. These weights are positive and
increasing for the choices below. Also the same diagonal estimate
at $n=17$ proves positivity of $a_n$ on the whole claimed tail.

The finite head now consists of $17\le n\le N$ in the normalized
parity basis; it does not include the omitted physical low modes.
The solve witness has support $N<n\le M$. With
$G=I-Z$, form the exact signed $K_Z$ and every residual row up to
$J$, and use the moment majorant $U_J$ beyond $J$. The lower matrix is
\begin{equation}\label{eq:v120-lambda4-lower}
 \mathcal L=K_Z-
  \sum_{N<n\le J}\frac{R_n^*R_n}{g_n}-\frac{U_J}{g_{J+1}}.
\end{equation}
The same full-index parity conversion and both infinite tails as in
Proposition~\ref{prop:v116-remote-moment} are used. Bounding its
geometric-remainder constant by including the absent indices
$1,\ldots,16$ is conservative. The certificate parameters are
\begin{center}
\begin{tabular}{lrrrrr}
\toprule
sector&head size&$N$&$M$&$J$&moment order\\
\midrule
even&496&512&1024&4096&48\\
odd&1520&1536&2048&4096&64\\
\bottomrule
\end{tabular}
\end{center}
Both $Z$ and an inverse-Cholesky congruence witness $C$ are frozen
to exact dyadic entries. Their numerical method of construction has
no proof role: outward interval evaluation verifies every row of
$C\mathcal L C^*$ has positive diagonal minus the sum of the
absolute off-diagonal entries. Thus $C\mathcal L C^*\succ0$
and $\mathcal L\succ0$. Square completion with the positive
remote operator proves \eqref{eq:v120-lambda4-tail} separately
in each parity sector and hence for arbitrary complex vectors.

The saved witnesses, analytic series radii, all finite residual rows
and the infinite moment bound are reproduced by
\texttt{assembly\_general.py} and
\texttt{certify\_weighted\_tail.py}. Replaying the same witnesses
at higher precision verifies the same strict inequalities.
These are internal computer-assisted certificates, not independent
external verification. No zero locations or RH assumptions enter them.
\end{proof}

\begin{proposition}[The two-sided weighted contraction fails at this cutoff]
\label{prop:v120-norm-counterexample}
For the same exact $D$ and $Q_{16}$ at $\lambda=4$,
\[
 \|D^{-1/2}Q_{16}(\mathsf W_4-D)Q_{16}D^{-1/2}\|>1.
\]
This does not contradict Proposition~\ref{prop:v120-lambda4-tail}.
\end{proposition}
\begin{proof}
The archived exact dyadic even vector $v$, supported on
$17\le|n|\le256$, has the rigorously enclosed quotient
\[
 2.05254409540345<\frac{QW_4(v,v)}{\ip{Dv}{v}}
                         <2.05254409540346.
\]
Its positive numerator excess over $2\ip{Dv}{v}$ exceeds
$0.0525$ in the saved normalization. The same analytic matrix
assembly, independent of the floating eigenvector that suggested
$v$, proves these interval inequalities. Testing the weighted
remainder on $D^{1/2}v$ gives a Rayleigh value greater than one.
This is a positive direction, not a negative Weil direction.
\end{proof}

The larger-window tail sign reduces full positivity at $\lambda=4$
to the signed effective matrix on
$E_{16}$:
\[
 F-B^*T^{-1}B,\qquad
 F=P_{16}\mathsf W_4P_{16},\quad B=Q_{16}\mathsf W_4P_{16},
 \quad T=Q_{16}\mathsf W_4Q_{16}.
\]
It has $17$ even and $16$ odd coordinates.
Proposition~\ref{prop:v125-even-complete} below settles the complete
even sector; the odd sector remains open. It is not the raw finite
Weil matrix: the inverse-tail correction must remain.
The next growing-window estimate must also satisfy the ordinary-error
scale in Proposition~\ref{prop:v120-weighted-error}.
A related recent preprint~\cite{Zhu2026Windows} studies certified
compact-window positivity and a pointwise prime-comb envelope barrier.
Our conclusion here concerns the actual Fourier-tail form and does
not use that preprint's computational certificates. Neither fixed-window
result supplies a uniform signed arithmetic estimate.

\subsection*{Reusing the signed tail certificate in the low-mode correction}
The scalar conclusion $T\succeq10^{-8}D$ discards most of the
information in Proposition~\ref{prop:v120-lambda4-tail}. The following
inverse bound retains its finite signed factorization. All adjoints
below use the fixed first-slot-linear convention.

\begin{proposition}[A structured inverse bound for a certified tail]
\label{prop:v120-structured-inverse}
Let a closed positive tail operator have a finite-head decomposition
\[
 T=\begin{pmatrix}F&B^*\\ B&U\end{pmatrix},
 \qquad U\succeq D_g\succeq\gamma I>0,
\]
where $B$ is bounded and $D_g$ is a positive diagonal operator.
Let $Z$ map the finite head into $\mathcal D(U)$, and set
\[
 R=B-UZ,\qquad K_Z=F-B^*Z-Z^*B+Z^*UZ.
\]
Suppose a finite Hermitian matrix $\underline L$ and invertible $C$
satisfy
\begin{equation}\label{eq:v120-inner-inverse-certificate}
 0\prec\underline L\preceq K_Z-R^*D_g^{-1}R,
 \qquad C\underline L C^*\succeq\mu I,\quad\mu>0.
\end{equation}
Then $T$ is bounded below by a strictly positive constant and, for
any tail right-hand side $r=(h,k)$,
\begin{equation}\label{eq:v120-structured-inverse}
 \ip{T^{-1}r}{r}
 \le \norm{D_g^{-1/2}k}^2+
 \mu^{-1}\norm{C(h-Z^*k-R^*D_g^{-1}k)}^2.
\end{equation}
The same statement holds as a Gram inequality for finitely many
right-hand sides simultaneously.
\end{proposition}
\begin{proof}
Write the original variable as $(x,y-Zx)$. Its exact form is
\[
 \ip{K_Zx}{x}+2\Re\ip{Rx}{y}+q_U[y],
\]
which is bounded below by
\[
 \ip{\underline Lx}{x}
       +\norm{D_g^{1/2}(y+D_g^{-1}Rx)}^2.
\]
The first change preserves the original form domain because
$Zx\in\mathcal D(U)$. The displayed lower form is defined on the
possibly larger domain consisting of the finite head and
$\mathcal D(D_g^{1/2})$. Its triangular change is bounded and
invertible there because $D_g^{-1}Rx\in\mathcal D(D_g)$ and the
head is finite dimensional. Thus that lower form is closed and
strictly positive; the original form dominates it on its own domain.
In its inverse variational formula the head part of the linear
functional is $h-Z^*k-R^*D_g^{-1}k$. Consequently its inverse
quadratic form equals the first term of
\eqref{eq:v120-structured-inverse} plus the quadratic form of
$\underline L^{-1}$ on that head vector. Inverse order and
\eqref{eq:v120-inner-inverse-certificate} give
$\underline L^{-1}\preceq\mu^{-1}C^*C$. This proves the bound.
Applying it to every linear combination of the right-hand sides
proves the Gram version. The argument concerns closed forms and
does not replace the logarithmic operator by a bounded matrix.
\end{proof}

If the saved factors instead certify $T-cD$, with $cD\succeq0$,
they still bound $T^{-1}$ because
$T^{-1}\preceq(T-cD)^{-1}$. The factors and residual in the
right side must then be those of $T-cD$ throughout; a shifted
residual is not an unshifted residual. In particular, the existing
$\lambda=4$ certificates are usable without losing their directional
information to the coarse upper bound $10^8D^{-1}$.

Here is a finite-row version that also retains every omitted mode.
When $Z$ has finite support, take $J$ past that support and split the inner far-tail rows
at $J$. For a matrix of right-hand sides define
\[
 H_J=h-Z^*k-\sum_{n\le J}R_n^*k_n/g_n,
 \qquad V_J=\sum_{n\le J}k_n^*k_n/g_n.
\]
The sums run over the inner far tail in normalized parity coordinates;
both Fourier signs are included by that decomposition. Suppose certified positive
matrices satisfy
\[
 \sum_{n>J}R_n^*R_n/g_n\preceq U_R,
 \qquad \sum_{n>J}k_n^*k_n/g_n\preceq U_k,
 \qquad \rho\ge\norm{CU_RC^*}.
\]
Then, for every $\tau>0$,
\begin{equation}\label{eq:v120-structured-finite-rows}
 r^*T^{-1}r\preceq V_J+U_k+
 \frac{(1+\tau)H_J^*C^*CH_J+
             (1+\tau^{-1})\rho U_k}{\mu}.
\end{equation}
Indeed, if $E=\sum_{n>J}R_n^*k_n/g_n$, then
$E^*C^*CE\preceq\rho U_k$ by the operator Cauchy--Schwarz
inequality. Apply the matrix Young inequality to $C(H_J-E)$.
The trace of $CU_RC^*$ is a permissible choice of $\rho$;
its possible growth with the inner head dimension must then be
retained. These are finite-column remote Grams, as in
Proposition~\ref{prop:v116-remote-moment}, not an invalid
Hilbert--Schmidt estimate for the whole weighted Weil operator.
A joint remote Gram can be sharper than separating this mixed term.

\begin{proposition}[A convergent arithmetic refinement of the far inverse]
\label{prop:v120-inverse-polynomial}
If $D_g\preceq U\preceq mD_g$, $D_g\succeq\gamma I>0$,
$m>1$, let $A=D_g^{-1/2}UD_g^{-1/2}$ denote the bounded
operator represented by the preconditioned form, and set
\[
 S=I-A/m,\qquad
 P_k(A)=\frac1m\sum_{j=0}^{k-1}S^j+S^k,\quad k\ge0.
\]
The empty sum gives $P_0=I$. Then
\begin{equation}\label{eq:v120-inverse-polynomial}
 A^{-1}\preceq P_{k+1}(A)\preceq P_k(A)\preceq I,
 \qquad
 0\preceq P_k(A)-A^{-1}\preceq(1-1/m)^k I.
\end{equation}
Thus $D_g^{-1/2}P_k(A)D_g^{-1/2}$ are decreasing upper
bounds for $U^{-1}$, with error at most
$(1-1/m)^kD_g^{-1}$ in form order. In particular,
\begin{equation}\label{eq:v120-first-inverse-refinement}
 U^{-1}\preceq D_g^{-1}
       -m^{-1}D_g^{-1}(U-D_g)D_g^{-1}.
\end{equation}
Products involving $U$ in this formula are understood through the
bounded preconditioned form.

For the actual far operators used in
Proposition~\ref{prop:v120-lambda4-tail}, including their shift
$\mathsf W_4-10^{-8}D$, one can take
\[
 m=67\quad\hbox{on the even tail }n>512,\qquad
 m=335\quad\hbox{on the odd tail }n>1536.
\]
These bounds cover all omitted Fourier modes at this fixed window.
\end{proposition}
\begin{proof}
The spectral theorem gives $I\preceq A\preceq mI$ and
$0\preceq S\preceq(1-1/m)I$. The geometric identity is
\[
 P_k(A)-A^{-1}=S^k(I-A^{-1}).
\]
Its factors commute and are positive. This proves the error bound;
$P_k-P_{k+1}=m^{-1}S^k(A-I)\succeq0$ proves monotonicity.
Conjugate by $D_g^{-1/2}$, using the inverse identity for the
positive closed form $U$. The case $k=1$ is
\eqref{eq:v120-first-inverse-refinement}.

To verify the two concrete constants, first strengthen
\eqref{eq:v115-diagonal-sharp} to
\begin{equation}\label{eq:v120-two-sided-arch}
 |a_n-\log(|n|/L)|\le E_L(|t_n|),\qquad n\ne0.
\end{equation}
For the upper bound use the same Binet formula at
$w=5/4+it/2$. It gives
\[
 \Re\log w-\log(t/2)
 =\tfrac12\log(1+25/(4t^2))\le25/(8t^2),
\]
while $-\Re(1/(2w))-\Re(1/z)\le0$, $z=1/4+it/2$.
The Binet remainder is at most $1/(15t)$ in absolute value.
Because $25/8<7/2$, the existing absolute trigamma and exponential
series bounds give the same $E_L$ for the upper error.

Write $\mathsf W_4=D+R_0$. From
\eqref{eq:v114-b-d}--\eqref{eq:v114-b-bound}, a valid bound is
\[
 \norm{R_0}\le C_R:=2\pi B_*+32h_L/L+2P_4.
\]
Here $P_4$ may omit the zero full-length translation $m=16$.
The pole diagonal bound is $32h_L/L$; its factor $L$ is essential.
The commutator bound is taken before the parity compression, so it
also bounds either normalized parity sector. With $c=10^{-8}$,
$g_n=(1-c)\{\log(n/L)-E_L(t_n)\}-\kappa$, and
$t_*=2\pi(N+1)/L$, the already proved lower bound and
\eqref{eq:v120-two-sided-arch} imply
\[
 D_g\preceq U\preceq
 D_g+\{\kappa+2(1-c)E_L(t_*)+C_R\}I.
\]
Since $g_n$ is positive and increasing, it suffices that
\begin{equation}\label{eq:v120-inverse-m-bound}
 m\ge1+\frac{\kappa+2(1-c)E_L(t_*)+C_R}{g_{N+1}}.
\end{equation}
Outward interval evaluation gives $C_R<34.450462$. The right side
of \eqref{eq:v120-inverse-m-bound} is below $66.76640$ for the
even $N=512$ choice, and below $334.71708$ for the odd $N=1536$
choice. The same prime bound, diagonal errors and inverse weights
as the tail certificate are used. Strict interval checks at two
precisions verify the stated integers.
\end{proof}

For a matrix of residuals $Y$, the error of the $k$th inverse bound
is at most
\[
 (1-1/m)^k Y^*D_g^{-1}Y.
\]
Along growing windows this part of the error is therefore controlled
if $\norm{D_{g,\lambda}^{-1/2}Y_\lambda}^2
 e^{-k_\lambda/m_\lambda}\to0$, with the actual values of
$m_\lambda$ retained. The explicit integers above apply only at
$\lambda=4$. Proposition~\ref{prop:v122-parity-inverse} below improves these
constants to $20$ and $90$ using the actual parity signs. Evaluating
the signed powers and their complete remote contributions, and
proving the remaining low-head sign, remain separate obligations. The inverse iteration proves neither of those signs.


\subsection*{Parity signs and certified evaluation of the inverse refinement}
The global remainder norm in the preceding inverse estimate loses useful
signs. Retaining them improves the complete far-operator sandwich;
it does not determine the sign of the low Schur matrix.

\begin{proposition}[Sharper inverse constants at the fixed window]
\label{prop:v122-parity-inverse}
For the actual shifted far operators and the unchanged weights $D_g$
in Proposition~\ref{prop:v120-lambda4-tail}, one has
\[
 D_g\preceq U\preceq20D_g\quad\text{in the even sector }n>512,
 \qquad
 D_g\preceq U\preceq90D_g\quad\text{in the odd sector }n>1536.
\]
Thus the constants in Proposition~\ref{prop:v120-inverse-polynomial}
can be replaced by $20$ and $90$, respectively, with error factors
$(19/20)^k$ and $(89/90)^k$. These are complete-operator bounds
at $\lambda=4$, with no assertion of growing-window uniformity.
\end{proposition}
\begin{proof}
Put $h=9/16$, $c=10^{-8}$, $L=2\log4$ and
$t_N=2\pi(N+1)/L$. Let $M_\phi$ be the certified prime bound
in \eqref{eq:v120-lambda4-prime-bound}. The limiting archimedean
off-diagonal in the normalized parity basis is the matrix
$1/[2(n+m)]$ in the even sector and its negative in the odd sector.
It is positive and has norm at most $\pi/2$; the remaining
archimedean commutator has norm at most $2/t_N$, by the proof of
Proposition~\ref{prop:v115-sharp-tail}.

In centered coordinates the pole form is
$2|\ip f{\cosh(x/2)}|^2-2|\ip f{\sinh(x/2)}|^2$.
Only the positive term acts on even vectors, and only the negative
term on odd vectors. Direct integration gives
\[
 |\widehat{\cosh(x/2)}(n)|^2
 =\frac{h}{L(t_n^2+1/4)^2},\qquad
 2\|Q_N\cosh(x/2)\|^2
 \le\frac{hL^3}{4\pi^4}\sum_{n>N}n^{-4}
 \le\frac{hL^3}{12\pi^4N^3}.
\]
The factor two for the form and both Fourier signs are included.
Consequently the remainder above the shifted archimedean diagonal
has the following upper bounds:
\begin{equation}\label{eq:v122-parity-upper}
 \begin{split}
 C_{\rm e}(N)&=M_\phi+2/t_N+\pi/2+hL^3/(12\pi^4N^3),\\
 C_{\rm o}(N)&=M_\phi+2/t_N.
 \end{split}
\end{equation}
The odd limiting Hilbert matrix and odd pole term are nonpositive
and may be omitted in an upper bound. The shift $cD$ changes only
the exact archimedean diagonal; none of these remainder terms is
multiplied by $1-c$.

For the same $\kappa$ and increasing $g_n$ as before, the two-sided
diagonal estimate \eqref{eq:v120-two-sided-arch} now yields
\[
 U\preceq D_g+
 \{\kappa+2(1-c)E_L(t_N)+C_{\rm e/o}(N)\}I.
\]
It suffices to bound one plus the braced quantity divided by
$g_{N+1}$. Outward interval evaluation gives values below
$19.215607$ for even $N=512$, and below $89.840960$ for odd
$N=1536$. The integer gates $20$ and $90$ pass at both 320 and
384 bits. The lower bounds and closed form domains are unchanged.
\end{proof}

\begin{proposition}[Complete residual bounds from a finite prefix]
\label{prop:v122-prefix-refinement}
Assume $D_g\preceq U\preceq mD_g$, with
$D_g\succeq\gamma I>0$. Put
$A=D_g^{-1/2}UD_g^{-1/2}$, $H=A-I$, and let $P$ be an
orthogonal prefix projection, $Q=I-P$. Suppose
$QHQ\preceq\nu Q$ with $\nu\ge0$. For a residual $k$ write
\[
 y=D_g^{-1/2}k=p+q,\quad p=Py,\quad q=Qy,\qquad
 v=\norm p^2,\quad a=\ip{Hp}{p},\quad \norm q^2\le E.
\]
With $[s]_+=\max(s,0)$, the following scalar bound includes
every omitted residual row:
\begin{equation}\label{eq:v122-prefix-scalar}
 \ip{U^{-1}k}{k}
 \le v+E-\frac1m[\sqrt a-\sqrt{\nu E}]_+^2.
\end{equation}
For a matrix of residual columns $Y=D_g^{-1/2}\mathcal R$, put
$Y_P=PY$, $Y_Q=QY$, $V=Y_P^*Y_P$ and
$A_P=Y_P^*HY_P$. If $Y_Q^*Y_Q\preceq\mathcal E$, then, for
every fixed $0<\theta\le1$,
\begin{equation}\label{eq:v122-prefix-matrix}
 \mathcal R^*U^{-1}\mathcal R\preceq
 V+\mathcal E-\frac1m
 \{(1-\theta)A_P-(\theta^{-1}-1)\nu\mathcal E\}.
\end{equation}
If a certified $0\le\rho<1$ instead satisfies
$\nu\mathcal E\preceq\rho^2A_P$, the braced lower bound can
be replaced by $(1-\rho)^2A_P$.
\end{proposition}
\begin{proof}
The bounded operator $H$ is positive, so the reverse triangle
inequality for $H^{1/2}p+H^{1/2}q$ gives
\[
 \ip{Hy}{y}\ge
 [\sqrt a-\|H^{1/2}q\|]_+^2
 \ge[\sqrt a-\sqrt{\nu E}]_+^2.
\]
Combine this with the first inverse polynomial
$A^{-1}\preceq I-H/m$ and
$\norm y^2=\norm p^2+\norm q^2\le v+E$.
For the matrix version, apply the lower Young inequality to
$H^{1/2}Y_Px+H^{1/2}Y_Qx$ for every head vector $x$.
Since $1-\theta^{-1}\le0$, the remote upper bound supplies the
stated lower form. The relative variant follows from the scalar
reverse triangle inequality in every head direction. No factor
equal to the number of columns is inserted. The scalar positive
part in \eqref{eq:v122-prefix-scalar} is not a matrix positive-part
operation; the latter does not follow from these arguments.
\end{proof}

In the actual parity tail based at $N$, the prefix can end at any
$J>N$. The same proof as \eqref{eq:v122-parity-upper} gives the
complete remote compression bound
\begin{equation}\label{eq:v122-remote-H}
 Q_JHQ_J\preceq\nu_J Q_J,\qquad
 \nu_J=\frac{\kappa_N+2(1-c)E_L(t_J)+C_{\rm e/o}(J)}{g_{J+1}}.
\end{equation}
Here $t_J=2\pi(J+1)/L$, whereas the constant $\kappa_N$ in
the original $D_g$ is kept fixed. The upper numerator decreases
and the positive denominator increases with the index. The moment
bound of Proposition~\ref{prop:v116-remote-moment}, divided by
$g_{J+1}$, supplies $\mathcal E$. Thus these formulas evaluate a
signed inverse improvement without setting the infinite residual
to zero. All occurrences of $U$ are interpreted through its
bounded preconditioned form.

\begin{proposition}[One complete far-residual comparison]
\label{prop:v122-directional-far}
At $\lambda=4$, take the archived exact dyadic even solve $X$
on modes $17,\ldots,256$, and the archived exact dyadic head
vector $v$ of length $17$ specified in
Appendix~\ref{app:v122-inverse}. Put
\[
 f=\binom{v}{-Xv},\qquad K=QW_4(f,f),\qquad
 k=Q_{512}\mathsf W_4 f.
\]
For the complete shifted even far operator
$U=Q_{512}(\mathsf W_4-10^{-8}D)Q_{512}$ one has
\begin{equation}\label{eq:v122-directional-far}
 \ip{U^{-1}k}{k}<5.758\cdot10^{-20}
 <6.469\cdot10^{-20}<K.
\end{equation}
This is one far-energy comparison. It does not include the mixed
term of \eqref{eq:v122-refined-structured}, and it is not a lower
bound for the full low Schur matrix.
\end{proposition}
\begin{proof}[Outward interval evaluation with a complete remote bound]
Use $m=20$ and $J=65536$ in
Proposition~\ref{prop:v122-prefix-refinement}. The frozen $f$ has
support through $256$. Thus its retained Weil energy and residual
rows are evaluated from the actual finite matrix coefficients,
without truncating $f$. The outer residual is unshifted; the far
inverse and its $D_g$ are those of the shifted tail certificate.
The following outward enclosures suffice:
\begin{center}
\begin{tabular}{@{}ll@{}}
\toprule Quantity & Certified enclosure\\\midrule
$K$ & $(6.4692843941,6.4692843942)\,10^{-20}$\\
$v_J=\|P_JD_g^{-1/2}k\|^2$ & $(6.9332416283,6.9332416284)\,10^{-20}$\\
$a_J=\ip{HP_JD_g^{-1/2}k}{P_JD_g^{-1/2}k}$
 & $(4.0338113722,4.0338113723)\,10^{-19}$\\
$\|Q_JD_g^{-1/2}k\|^2$ & $<3.4797080873\,10^{-21}$\\
$\|Q_JHQ_J\|$ & $<1.987323730$\\\bottomrule
\end{tabular}
\end{center}
The remote norm uses the complete order-$64$ moment majorant,
divided by $g_{J+1}$; the last row uses
\eqref{eq:v122-remote-H}. Substitution into
\eqref{eq:v122-prefix-scalar} gives the stricter computed upper
bound $5.757887789\cdot10^{-20}$, proving the displayed claim.
All gates replay with identical exact input witnesses at $768$ and
$896$ bits. The prefix matrix action uses exact Arb polynomial
products for its divided-difference and Hankel parts, as detailed
in the reproduction appendix. No floating-point sign decision,
zero-location assumption or omitted-residual cutoff enters the gate.
\end{proof}

\begin{proposition}[Refining both terms of the structured inverse]
\label{prop:v122-refined-structured}
Keep the factors $Z,R,C,\underline L,\mu$ of
Proposition~\ref{prop:v120-structured-inverse}, and assume also
$U\preceq mD_g$. Let $Q_*$ be a bounded positive operator such
that
\[
 U^{-1}\preceq Q_*\preceq D_g^{-1}.
\]
Then, for every $r=(h,k)$,
\begin{equation}\label{eq:v122-refined-structured}
 \ip{T^{-1}r}{r}
 \le\ip{Q_*k}{k}
       +\mu^{-1}\norm{C(h-Z^*k-R^*Q_*k)}^2.
\end{equation}
The same inequality holds for the simultaneous residual Gram.
In particular all inverse polynomials of
Proposition~\ref{prop:v120-inverse-polynomial} can be used while
retaining the original certified $C$ and $\mu$.
\end{proposition}
\begin{proof}
The assumed sandwich implies
$m^{-1}D_g^{-1}\preceq Q_*\preceq D_g^{-1}$.
Consequently $Q_*$ is injective and the closed form of $Q_*^{-1}$
has the same domain as $D_g^{1/2}$, with equivalent form norms.
Inverse order gives $U\succeq Q_*^{-1}$ as forms. Moreover
\[
 K_Z-R^*Q_*R\succeq K_Z-R^*D_g^{-1}R\succeq\underline L.
\]
Repeat the lower-form square completion in
Proposition~\ref{prop:v120-structured-inverse}, replacing the far
form by $Q_*^{-1}$ and its inverse by $Q_*$. Its exact inverse
contains $Q_*$ in both displayed terms, proving the assertion.
The change preserves the lower form domain since $Q_*Rx$ belongs
to the operator domain of $Q_*^{-1}$ for every finite head vector.
Applying the result to all linear combinations proves the Gram
version. If the saved factors certify $T-cD$, use inverse order
as before and retain those shifted factors throughout.
\end{proof}

Improving only the first term while leaving the old mixed term
unchanged is not justified. Positivity of the complete low Schur
matrix requires the resulting full correction, on the entire head,
to fit below $K_X$. The new estimates do not assert that comparison
or convergence along growing windows.

\begin{remark}[The simultaneous head acceptance test]
\label{rem:v122-head-ldl}
For a positive definite trial matrix $K_X$, put
$C_X=r^*T^{-1}r$. Then
\[
 \lambda_{\max}(K_X^{-1/2}C_XK_X^{-1/2})\le1
 \quad\Longleftrightarrow\quad K_X-C_X\succeq0.
\]
Thus an eigenvalue or inverse-square-root estimate is unnecessary.
It suffices to construct a Hermitian \emph{upper enclosure in form
order} $\overline C_X\succeq C_X$ for the complete residual Gram
and certify $K_X-\overline C_X\succeq0$, for example by interval
$LDL^*$ with strictly positive pivots. A fixed invertible dyadic
congruence may improve conditioning without changing the sign.
Strictly positive pivots prove the stronger strict comparison;
an exact zero pivot requires the corresponding semidefinite
factorization conditions, not an interval containing zero.
An entrywise upper bound is not by itself an upper bound in form
order. For $T_M=P_MT P_M$ and $r_M=P_Mr$, the variational identity
gives $r_M^*T_M^{-1}r_M\preceq C_X$, the opposite enclosure
direction. All mixed and remote contributions must be enclosed
before the $LDL^*$ test proves the complete Schur sign.
\end{remark}

\subsection*{Complete mixed correlations and a monotone head refinement}

The successful far-energy comparison in
Proposition~\ref{prop:v122-directional-far} leaves a mixed term.
The next estimate includes both the omitted input and the output
outside the retained Fourier rows. These are different contributions.

\begin{proposition}[A directional enclosure of the complete mixed term]
\label{prop:v123-mixed-tail}
Use the factors of Proposition~\ref{prop:v122-refined-structured},
and put
\[
 A=D_g^{-1/2}UD_g^{-1/2},\quad H=A-I,\quad
 Y=D_g^{-1/2}RC^*,\quad y=D_g^{-1/2}k,\quad d=C(h-Z^*k).
\]
Here $A$ is represented by its bounded closed form and
$0\preceq H\preceq(m-1)I$. Let $P$ be an orthogonal Fourier
prefix projection on the inner far tail, $Q=I-P$, $p=Py$ and
$q=Qy$. Suppose
\[
 \norm q^2\le E,\qquad 0\preceq QHQ\preceq\nu Q,
 \qquad 0\le\nu<m.
\]
Set $v=\norm p^2$, $a=\ip{Hp}{p}$ and
\[
 b=d-Y^*(I-H/m)y,\qquad
 b_P=d-(PY)^*(I-PHP/m)p.
\]
For any finite head vector $z$, let
$\eta_z=\ip{H PYz}{PYz}$ and suppose
$\norm{QYz}^2\le\rho_z$. Then
\begin{equation}\label{eq:v123-mixed-tail}
 |z^*(b-b_P)|\le
 \frac{\sqrt{\nu\eta_z E}+\sqrt{\nu a\rho_z}}{m}
       +\sqrt{\rho_z E}=:\epsilon_z.
\end{equation}
In addition the complete first-polynomial far energy satisfies
\begin{equation}\label{eq:v123-far-lower}
 \ip{(I-H/m)y}{y}\ge v-\frac{a}{m-\nu}.
\end{equation}
In particular, for $z\ne0$ the complete scalar-$\mu$ majorant
in \eqref{eq:v122-refined-structured} is at least
\begin{equation}\label{eq:v123-majorant-lower}
 v-\frac{a}{m-\nu}
 +\frac{[|z^*b_P|-\epsilon_z]_+^2}{\mu\norm z^2}.
\end{equation}
\end{proposition}
\begin{proof}
Write $f=Yz=f_P+f_Q$. The difference of the mixed pairings is
the sum of $m^{-1}f_P^*Hq$, $m^{-1}f_Q^*Hp$, and
$-f_Q^*(I-QHQ/m)q$. Positive-form Cauchy--Schwarz bounds the
first two by the corresponding square roots in
\eqref{eq:v123-mixed-tail}. Since $0\preceq QHQ\preceq\nu Q$
and $\nu<m$, the last factor is a contraction on the remote
subspace. This proves the mixed estimate without setting either
remote contribution to zero.

For $x=\norm q$, positivity of $H$ and its remote compression give
\[
 \ip{(I-H/m)y}{y}
 \ge v-a/m-2\sqrt{a\nu}\,x/m+(1-\nu/m)x^2.
\]
The minimum of this scalar quadratic over $x\ge0$ is
$v-a/(m-\nu)$. The directional lower bound on $\norm b$
then proves \eqref{eq:v123-majorant-lower}. All pairings are
bounded-operator pairings after preconditioning. No application
of the unbounded $U$ to an arbitrary Hilbert vector is used.
\end{proof}

\begin{proposition}[Failure of the first scalar-head majorant]
\label{prop:v123-first-majorant-fails}
On the same exact even head direction and frozen support-$256$
trial as Proposition~\ref{prop:v122-directional-far}, let
$\mathcal B_1$ denote the complete right side of
\eqref{eq:v122-refined-structured} with
$Q_*=D_g^{-1/2}(I-H/20)D_g^{-1/2}$ and the saved
$\mu=0.9999999999$. Then
\begin{equation}\label{eq:v123-first-majorant-fails}
 \mathcal B_1>6.7609\cdot10^{-20}
 >6.4693\cdot10^{-20}>K_X(v,v).
\end{equation}
Thus this particular complete majorant cannot certify that trial
direction. This is not a negative Weil vector and does not exclude
a better inverse polynomial or a stronger initial head certificate.
\end{proposition}
\begin{proof}[Complete interval certificate]
Retain $J=65536$, $m=20$ and the complete remote bounds of
Proposition~\ref{prop:v122-directional-far}. Choose the exact
$160$-bit dyadic head vector $z$ specified in the reproduction
appendix, obtained by rounding the direction of $b_P$. Its
normalization is enclosed, not assumed exact. The following
outward enclosures suffice in addition to the previous values of
$v,a,E,\nu$:
\begin{center}
\begin{tabular}{@{}ll@{}}
\toprule Quantity & Certified enclosure\\\midrule
$\eta_z$ & $(0.4276838592,0.4276838593)$\\
$\rho_z$ & $<0.0672057785$\\
$|z^*b_P|$ & $(1.7339321076,1.7339321077)\,10^{-10}$\\
$\norm z$ & $(1-10^{-45},1+10^{-45})$\\
$\epsilon_z$ & $<2.961708367\,10^{-11}$\\
complete far energy & $>4.693812399\,10^{-20}$\\
$\mu^{-1}\norm b^2$ & $>2.067157472\,10^{-20}$\\\bottomrule
\end{tabular}
\end{center}
The bound on $\rho_z$ uses the complete order-$64$ moment
majorant on the finite vector $[I;-Z]C^*z$, whose support ends
at $1024$, divided by $g_{J+1}$. It includes every row beyond
$J$. The scalar $\eta_z$ and the retained pairings use the exact
Arb polynomial action from the preceding appendix, with the
unchanged archimedean-series remainder and parity normalization.
Equation~\eqref{eq:v123-majorant-lower} gives
$\mathcal B_1>6.760969872\cdot10^{-20}$, whereas the already
certified $K_X(v,v)<6.469284395\cdot10^{-20}$.
The same frozen direction and solve replay at $768$ and $896$
bits. The outer residual is that of the unshifted Weil form;
the inner inverse factors consistently belong to the certified
shifted tail. Their inverse upper bound remains valid by order.
\end{proof}

\begin{proposition}[Retaining the strengthened head metric]
\label{prop:v123-monotone-head}
In the same setup, let a positive definite finite matrix $M_0$
satisfy
\[
 C K_Z C^*-Y^*Y\succeq M_0\succ0.
\]
The existing certificate permits $M_0=\mu I$. Let $P_j(A)$
be the inverse upper polynomials of
Proposition~\ref{prop:v120-inverse-polynomial}, and set
\begin{align*}
 M_j&=M_0+Y^*(I-P_j(A))Y,\\
 b_j&=d-Y^*P_j(A)y,\\
 \mathcal E_j(r)&=\ip{P_j(A)y}{y}+b_j^*M_j^{-1}b_j.
\end{align*}
Then these are complete simultaneous upper bounds with
\begin{equation}\label{eq:v123-monotone-head}
 \ip{T^{-1}r}{r}\le\mathcal E_{j+1}(r)
          \le\mathcal E_j(r).
\end{equation}
They converge as $j\to\infty$. Their limit need not equal
$\ip{T^{-1}r}{r}$: any loss in replacing
$CK_ZC^*-Y^*Y$ by $M_0$ remains.
For the first polynomial, the head improvement is exactly
$Y^*HY/m$. It must be retained together with the changed
mixed vector if monotonicity of the complete bound is required.
\end{proposition}
\begin{proof}
After the finite triangular substitution and head congruence,
the original form has blocks $CK_ZC^*$, $RC^*$ and $U$,
and the dual vector is $(d,k)$. Put
$Q_j=D_g^{-1/2}P_j(A)D_g^{-1/2}$ and consider the lower form
\[
 \mathcal L_j=
 \begin{pmatrix}
 M_0+Y^*Y & CR^*\\
 RC^* & Q_j^{-1}
 \end{pmatrix}.
\]
The assumed head comparison and $Q_j^{-1}\preceq U$ imply
that the transformed original form dominates $\mathcal L_j$.
Its Schur complement is $M_j\succeq M_0\succ0$.
Because $m^{-1}D_g^{-1}\preceq Q_j\preceq D_g^{-1}$,
all far forms have the same domain $\mathcal D(D_g^{1/2})$.
The finite head and bounded coupling give closed, coercive
lower forms by the same square completion as before.

Their inverse quadratic forms at $(d,k)$ are precisely
$\mathcal E_j(r)$. The polynomials decrease in operator order,
so $Q_j^{-1}$ increases. The head of $\mathcal L_j$ is fixed;
therefore $\mathcal L_{j+1}\succeq\mathcal L_j$, proving the
claimed inverse order. Norm convergence of $P_j(A)$ to $A^{-1}$,
boundedness of $Y$, and the common positive floor of $M_j$
give convergence of the displayed scalar and Gram expressions.
The possible head loss explains the stated limitation of the
limit. A certificate for the shifted tail is again transferred
only by the already stated inverse order.
\end{proof}

\begin{corollary}[The first head update does not rescue this trial]
\label{cor:v123-updated-first-fails}
For the same saved data, choose $M_0=\mu I$ in
Proposition~\ref{prop:v123-monotone-head}. Even its exact first
head update satisfies
\[
 \mathcal E_1(r)>6.6588\cdot10^{-20}>K_X(v,v).
\]
This excludes the first inverse degree with that initial certificate,
including its full positive head update. It does not exclude a
stronger initial $M_0$ or a higher inverse degree.
\end{corollary}
\begin{proof}
Use the same frozen $z$ and complete tail bounds. Positive-form
Cauchy gives
\[
 z^*M_1z\le\mu\norm z^2+
 \frac{(\sqrt{\eta_z}+\sqrt{\nu\rho_z})^2}{20}
       <1.051962238.
\]
Metric Cauchy--Schwarz therefore gives
\[
 b_1^*M_1^{-1}b_1
 \ge\frac{[|z^*b_P|-\epsilon_z]_+^2}{z^*M_1z}
       >1.965049122\cdot10^{-20}.
\]
Adding the complete far lower bound gives
$\mathcal E_1(r)>6.658861521\cdot10^{-20}$.
Both precision replays include this gate with the same exact
dyadic witnesses and all remote correlations.
\end{proof}

\begin{proposition}[Finite probes for the positive head improvement]
\label{prop:v123-head-probes}
Keep $H,Y,P,Q,\nu$ as above. Let $V$ have finitely many
finite-prefix columns, suppose $G=V^*HV\succ0$, and set
$B_P=V^*HPY$. If $Y^*QY\preceq E_Y$, then for every
$0<\theta<1$,
\begin{equation}\label{eq:v123-head-probes}
 Y^*HY\succeq
 (1-\theta)B_P^*G^{-1}B_P
       -(\theta^{-1}-1)\nu E_Y.
\end{equation}
Consequently the right side, divided by $m$ and added to $M_0$,
is a valid finite lower enclosure of the first improved head.
It may be inverted only after its positivity is certified.
\end{proposition}
\begin{proof}
The Gram matrix of $H^{1/2}V$ and $H^{1/2}Y$ gives
$Y^*HY\succeq B^*G^{-1}B$, where $B=V^*HY$.
Write $B=B_P+E_B$ with $E_B=V^*HQY$.
The projection onto the range of $H^{1/2}V$ shows
\[
 E_B^*G^{-1}E_B\preceq Y^*QHQY\preceq\nu E_Y.
\]
Young's inequality in the $G^{-1}$ metric proves
\eqref{eq:v123-head-probes}. This keeps the omitted probe
correlation explicitly. No unjustified matrix positive-part
operation is applied.
\end{proof}

These first-degree bounds do not settle the complete low-mode Schur
sign. The later simultaneous certificate proves the even-sector sign
at $\lambda=4$; the odd-sector sign remains open.
The first-degree obstruction now has a proof even with the
induced head update and the saved initial certificate. Increasing
the inverse degree, or strengthening the initial certificate,
requires a different complete enclosure.
Neither finite-prefix improvements nor convergence at this fixed
window establishes uniformity as $\lambda\to\infty$.

\subsection*{A complete Schur-direction certificate from a finite trial}

The first-degree obstruction above concerns a particular inverse
majorant and saved trial. A finite iterative calculation can instead
propose a new trial, after which the complete residual is verified
afresh. This uses the same verified-solve identity; it does not
identify a power of a finite compression with the full power.

\begin{proposition}[A posteriori verification without a full-power claim]
\label{prop:v124-posterior-identity}
Let the complete outer operator have finite head and coercive
tail $T$, with Schur form $S_\infty=F-B^*T^{-1}B$.
For a finite trial $g=(v,w)$ in its operator domain, set
$e=Bv+Tw$. Then
\begin{equation}\label{eq:v124-posterior-identity}
 S_\infty(v,v)=QW(g,g)-\ip{T^{-1}e}{e}.
\end{equation}
Thus any certified complete upper bound $\mathcal U(e)$ for the
last term gives
\[
 QW(g,g)-\mathcal U(e)\le S_\infty(v,v)\le QW(g,g).
\]
The algorithm that proposed the frozen finite coefficients of $w$
need not be an approximation to any particular infinite polynomial.
\end{proposition}
\begin{proof}
Expand $\ip{T^{-1}(Bv+Tw)}{Bv+Tw}$ and cancel the tail and
cross terms in $QW(g,g)$. Coercivity justifies the inverse.
Equivalently, for any residual $r$ and finite correction $z$,
\[
 \ip{T^{-1}r}{r}
 =2\Re\ip rz-\ip{Tz}{z}
       +\ip{T^{-1}(r-Tz)}{r-Tz}.
\]
If the inner certificate belongs to $T-cD$, use
$T^{-1}\preceq(T-cD)^{-1}$ only in bounding the new residual
term. The energy and $e=P_{\rm tail}Wg$ remain unshifted.
This avoids an erroneous omission of a $cDz$ residual or a shift
energy when the two operators are mixed.
\end{proof}

For completeness, let the new residual split as $e=(h,k)$ in
the already certified inner decomposition. Past the support of
the saved $Z$, retain rows through $J$ and put
\[
 V_J=\sum_{n\le J}|k_n|^2/g_n,\qquad
 b_J=C\left(h-Z^*k-\sum_{n\le J}R_n^*k_n/g_n\right).
\]
The sums run over the inner far rows. If complete remote bounds
give $\sum_{n>J}|k_n|^2/g_n\le E$ and
$\|D_g^{-1/2}Q_JRC^*\|^2\le\rho$, then the old structured
inverse certificate supplies
\begin{equation}\label{eq:v124-posterior-upper}
 \mathcal U(e)=V_J+E+
       \mu^{-1}\bigl(\norm{b_J}+\sqrt{\rho E}\bigr)^2.
\end{equation}
Every row beyond $J$ is included. For a matrix of trial columns,
the simultaneous version retains their Gram matrices and uses
\eqref{eq:v120-structured-finite-rows}; separate positive scalar
tests are insufficient for positivity on their span.

\begin{proposition}[One complete positive Schur direction at $\lambda=4$]
\label{prop:v124-positive-direction}
Let $v$ be the exact archived even head vector used in
Propositions~\ref{prop:v122-directional-far} and
\ref{prop:v123-first-majorant-fails}. For the \emph{complete}
even Schur form with head $0\le n\le16$,
\begin{equation}\label{eq:v124-positive-direction}
 3.0930\cdot10^{-20}<S_\infty(v,v)<3.2356\cdot10^{-20}.
\end{equation}
This certifies one fixed direction. It does not assert that the
whole even head, the odd head, or the full $\lambda=4$ form is
positive, and it provides no growing-window uniformity.
\end{proposition}
\begin{proof}[Frozen finite trial and complete residual certificate]
Start with the previous support-$256$ vector $f=[v;-Xv]$.
Subtract the exact dyadic tail correction $z$, supported on
$17\le n\le4096$, specified in the reproduction appendix.
The new vector $g=f-(0,z)$ has exactly the same head $v$.
A finite preconditioned iteration proposed $z$; neither its
convergence nor its recursive residual is used in this proof.

Evaluate $QW_4(g,g)$ and the unshifted residual
$e=Q_{16}W_4g$ through $J=65536$. Retain the original shifted
inner factors, $\mu=0.9999999999$, and the previously certified
$\rho<0.165102333$. Direct interval evaluation gives
\begin{center}
\begin{tabular}{@{}ll@{}}
\toprule Quantity & Certified enclosure\\\midrule
$QW_4(g,g)$ & $(3.2355843481,3.2355843482)\,10^{-20}$\\
$V_J$ & $<3.742924324\,10^{-22}$\\
$\norm{b_J}^2$ & $<1.617782100\,10^{-23}$\\
$E$ & $<8.080848756\,10^{-22}$\\
$\mathcal U(e)$ & $<1.424889\,10^{-21}$\\\bottomrule
\end{tabular}
\end{center}
The bound $E$ is the complete order-$64$ moment majorant applied
to the finite vector $g$, with its actual support through $4096$,
divided by $g_{J+1}$. The saved remote $\rho$ belongs to the same
inner witness and cutoff. Substitution in
\eqref{eq:v124-posterior-upper} and
\eqref{eq:v124-posterior-identity} proves the lower bound.
The upper bound follows from $T^{-1}\succeq0$.

The frozen correction uses exact dyadic coefficients produced at
$256$ bits; the proof evaluates them anew at $768$ and $896$
bits. Both checks retain the zero-mode $\sqrt2$ factor, normalized
parity, logarithmic diagonal and archimedean-series remainder.
Selected rows, including zero and the support endpoint, are also
checked by dense coefficient assembly. The independently expanded
energy identity agrees with the direct energy. No finite power is
used as a substitute for an infinite operator power.
\end{proof}

For the \emph{old} trial in this same direction the complete inverse
correction is consequently bounded by
\[
 3.2337\cdot10^{-20}<\ip{T^{-1}r}{r}
       <3.3762\cdot10^{-20}.
\]
In particular $S_\infty(v,v)>0.4781K_X(v,v)$ on this one
direction. This explains why the earlier majorant could fail
without identifying an intrinsically difficult Schur direction.
The next target is the simultaneous matrix comparison, including
the remaining directions of small energy. This directional
certificate cannot be extrapolated to that matrix or to RH.

\subsection*{A simultaneous complete Schur certificate}

A finite trial may be chosen in coordinates adapted to its head energy;
the final certificate must still use its actual head and every residual
cross term.  This distinction is essential when some ordinary head
energies are extremely small.

\begin{proposition}[Simultaneous trial and ordinary-error certificate]
\label{prop:v125-simultaneous}
Let the complete closed Weil form have finite head and coercive tail
$T$, and let $G$ be a matrix of finite Fourier trial columns.  Write
$V=P_{16}G$, $E=Q_{16}\mathsf W_4G$ and $K=G^*\mathsf W_4G$.
Then
\begin{equation}\label{eq:v125-simultaneous}
 V^*S_\infty V=K-E^*T^{-1}E.
\end{equation}
Suppose $V$ is square and invertible, and a complete matrix upper bound
$\mathcal U\succeq E^*T^{-1}E$ is certified.  If
$K-\mathcal U\succ0$, the entire parity form is strictly positive.
More generally, if $\varepsilon\ge0$ and
\[
 K-\mathcal U+\varepsilon V^*V\succeq0,
\]
then $S_\infty\succeq-\varepsilon I$ and the entire parity form is
bounded below by $-\varepsilon I$ in ordinary norm.
\end{proposition}
\begin{proof}
Apply Proposition~\ref{prop:v124-posterior-identity} to every linear
combination of the columns.  Invertibility of the actual head $V$
transfers the resulting congruent inequality to $S_\infty$.
For a head vector $x$ and tail vector $y$, the closed form identity is
\[
 q_{\mathsf W_4}[x,y]
 =\ip{S_\infty x}{x}
   +\norm{T^{1/2}(y+T^{-1}Bx)}^2.
\]
The head-to-tail map $B$ is bounded because the head is finite
and its Fourier vectors belong to the operator domain.  Thus the
completion is valid on the form domain.  The tail square is
nonnegative, and $\norm{x}\le\norm{(x,y)}$.  Strict head and tail
positivity also imply fixed-window coercivity by the bounded
invertible triangular change in this identity.  No quantitative
ordinary spectral gap is identified with an $LDL^*$ pivot.
\end{proof}

\begin{proposition}[Retaining a joint remote Gram]
\label{prop:v125-joint-remote}
In the finite-row structured inverse bound, put
$Y=D_g^{-1/2}Q_JRC^*$ and $y=D_g^{-1/2}Q_Jk$, and let
$H=C(h-Z^*k-R_J^*D_g^{-1}k_J)$ be the retained mixed matrix.
Suppose one joint enclosure and a lower bound satisfy
\[
 \Gamma=\begin{pmatrix}\Gamma_{YY}&\Gamma_{Yk}\\
 \Gamma_{Yk}^*&\Gamma_{kk}\end{pmatrix}
 \succeq [Y,y]^*[Y,y],\qquad 0\preceq L_Y\preceq Y^*Y,
 \qquad A_0=\mu I+L_Y-\Gamma_{YY}\succ0.
\]
Then the complete inverse correction is bounded by
\begin{equation}\label{eq:v125-joint-remote}
 E^*T^{-1}E\preceq V_J+\Gamma_{kk}
 +(H-\Gamma_{Yk})^*A_0^{-1}(H-\Gamma_{Yk}).
\end{equation}
Here $V_J=k_J^*D_g^{-1}k_J$.  One may take $L_Y=0$.
If $\Gamma_{YY}\preceq\rho I$ with $\rho<\mu$, replacing
$A_0^{-1}$ by $(\mu-\rho)^{-1}I$ gives a simpler valid bound.
\end{proposition}
\begin{proof}
For an outer coefficient vector $x$, the remote portion of the old
structured majorant is
\[
 \norm{yx}^2+\mu^{-1}\norm{Hx-Y^*yx}^2
 =\sup_z\{\norm{yx-Yz}^2-\norm{Yz}^2
       +2\Re\ip{Hx}{z}-\mu\norm z^2\}.
\]
Apply the joint upper Gram to $(-z,x)$ and the lower bound $L_Y$
to $\norm{Yz}^2$, then complete the finite $A_0$ square.
This proves the claim for every $x$.  The bounded finite-column maps
justify all remote pairings.  The cross block must come from the
\emph{same joint} Gram; unrelated diagonal upper bounds do not
justify an invented mixed term.
\end{proof}

\begin{proposition}[The complete even form at $\lambda=4$]
\label{prop:v125-even-complete}
The canonical complete Weil form $\mathsf W_4$ restricted to its
even parity space is strictly positive on every nonzero vector
in its form domain, and is bounded below by a positive constant
at this fixed window.
\end{proposition}
\begin{proof}[Frozen matrix certificate]
Keep the outer head $0\le n\le16$ and the complete tail certified
in Proposition~\ref{prop:v120-lambda4-tail}.  The archived matrix
trial has $17$ columns and support through $M=4096$.  Its exact
dyadic base and correction coefficients define $G$; the actual
head $V=P_{16}G$ is certified invertible.  The proposal first scales
the old finite trial by its head-energy $LDL^*$ factors, rounds
that base to exact dyadics, and uses $16$ preconditioned finite CG
steps at $512$ bits for each column.  Neither CG convergence nor
its recursive residual enters the proof.

Recompute the unshifted $K=G^*\mathsf W_4G$ and
$E=Q_{16}\mathsf W_4G$ at $768$ and $896$ bits from the same
frozen inputs.  Retain all cross-Grams through $J=65536$ and apply
the order-$64$ moment bound to the entire trial for every omitted
row.  The saved shifted inner factors have $N=512$, support $1024$,
$\mu=0.9999999999$ and complete remote bound
$\rho<0.165102333$.  Inverse order is used only for the residual
inverse, so the trial energy and residual stay unshifted.
For $\tau=1/10$, equation~\eqref{eq:v120-structured-finite-rows}
gives the complete matrix upper bound $\mathcal U_\tau$.
All $17$ interval $LDL^*$ pivots of $K-\mathcal U_\tau$ are positive;
the smallest exceeds $0.7645944101$ in these saved head coordinates.
Both precision replays give this strict gate.  This number certifies
sign and is not an ordinary-norm eigenvalue bound.

The verifier checks the actual head determinant, exact dyadic
reconstruction, inner-certificate provenance, positivity of the
far diagonal, and nine independently assembled physical rows for
every column.  The complete remote moment includes the zero-mode
normalization and both Fourier signs through the parity basis.
Proposition~\ref{prop:v125-simultaneous} therefore proves positivity
of the entire even form, not only a Fourier compression.
\end{proof}

A common remote Gram for the padded inner source $[I;-Z]C^*$ and
this trial also supplies the cross centre in
Proposition~\ref{prop:v125-joint-remote}.  Its scalar-denominator
certificate has all $17$ pivots positive, the smallest greater than
$0.7892519172$; using the full finite denominator improves that
smallest pivot to more than $0.7892519657$.  This is an additional
fixed-witness check, not needed for the preceding even sign proof.
The whole-operator conclusion comes from a simultaneous matrix
inequality and complete tail control; separately positive directions
would not suffice.

The comparison also has a certified dimensionless margin:
\begin{equation}\label{eq:v125-relative-margin}
 \mathcal U_{1/10}\preceq
 \left(1-\frac{62629}{100000}\right)K.
\end{equation}
Interval $LDL^*$ verifies the equivalent positive matrix at both
precisions.  Thus the generalized relative margin exceeds $0.62629$;
unlike a smallest ordinary eigenvalue or a coordinate pivot, this
ratio is invariant under an invertible head congruence.  On the
original physical head direction from
Proposition~\ref{prop:v124-positive-direction}, this new matrix trial
has certified headroom greater than $12.1284$.  Its trial is different
from the earlier single-direction trial, whose headroom was about
$22.71$; changing the trial changes this diagnostic while leaving
$S_\infty(v,v)$ unchanged.

The analogous odd-head bounds tested here are inconclusive.  The
exact statement established so far is positivity of the complete
even sector at $\lambda=4$; it is not yet
$\mathsf W_4\succeq0$ on both sectors.  No finite list of such
fixed-window statements implies the cofinal estimate needed for RH.

\subsection*{What the present diagonal majorant costs as the window grows}

\begin{proposition}[Exponential cutoff cost of the scalar far bound]
\label{prop:v125-cutoff-cost}
Fix $0\le c<1$, put $L=2\log\lambda$, and use the present far
weights
\[
 g_{N+1}=(1-c)\{\log((N+1)/L)-E_L(2\pi(N+1)/L)\}
               -\kappa_{\lambda,N},
 \qquad \kappa_{\lambda,N}\ge M_{\phi,\lambda}
                    \ge\|\mathsf T_{{\rm pr},\lambda}\|.
\]
If $g_{N+1}>0$, then
\begin{equation}\label{eq:v125-cutoff-cost}
 N+1>L\exp\!\left(\frac{M_{\phi,\lambda}}{1-c}\right).
\end{equation}
Consequently, along any unbounded family using these positive weights,
\[
 \liminf_{\lambda\to\infty}
 \frac{\log((N(\lambda)+1)/L)}{\lambda}\ge\frac1{1-c}.
\]
This is a cost of this particular scalar diagonal majorant, not a
lower bound on the cutoff required by every possible argument.
\end{proposition}
\begin{proof}
The error $E_L$ is positive, so positivity of $g_{N+1}$ gives
$(1-c)\log((N+1)/L)>\kappa_{\lambda,N}\ge M_{\phi,\lambda}$.
Use Proposition~\ref{prop:v119-prime-scale}, which proves
$\|\mathsf T_{{\rm pr},\lambda}\|/\lambda\to1$, for the limit.
\end{proof}

For $c=10^{-8}$, freshly evaluating the physical weighted prime bound
and the parity constants at each window gives the following smallest
integer $N$ for which the stated $g_{N+1}$ is positive.  These are
certificates for the remote block only; they do not certify the
intervening inner head or its inverse factors.
\begin{center}
\begin{tabular}{rrrr}
\toprule
$\lambda$ & $M_{\phi,\lambda}$ (rounded) & even $N$ & odd $N$\\
\midrule
3&2.980079&44&209\\
4&4.624042&283&1360\\
5&6.145139&1502&7224\\
6&7.644648&7488&36021\\
8&10.306338&124442&598623\\
\bottomrule
\end{tabular}
\end{center}
Each threshold has positive interval margin, while its preceding
integer has negative margin; monotonicity in $N$ proves minimality
for this formula.  The displayed prime values are rounded summaries,
not the outward endpoints used by the verifier.

The appropriate whole-head diagnostic is
$1-\lambda_{\max}(\mathcal U,K)$, with $K\succ0$ and a
\emph{complete} inverse upper bound $\mathcal U$.  The separately
useful ratios $\lambda_{\min}(K-\mathcal U)/\lambda_{\min}(K)$ and
$\|\mathcal U\|/\lambda_{\min}(K)$ depend on the chosen coordinates;
they should be recorded with the head dimension and condition of
its congruence.  A strong single direction cannot distinguish all
possible causes of a weak whole-head margin.  For example,
\[
 K=\operatorname{diag}(1,\epsilon),\qquad
 U=\operatorname{diag}(\delta,\epsilon(1-\eta)),
 \quad 0<\epsilon,\delta,\eta<1,
\]
has generalized margin $\min(1-\delta,\eta)$, while the headroom
in its first direction is $1/\delta$ and the condition of a
normalizing congruence is $\epsilon^{-1/2}$.  These parameters
may vary independently.  The actual error terms
$\alpha_\lambda$, $e_\lambda^2$ and $t_\lambda^2/\mu_\lambda$
still require complete window-specific witnesses.  The same
numerical cutoff or moment order may be used at different windows
only after its inequalities are verified afresh.


\begin{proposition}[An ordinary-error target for growing windows]
\label{prop:v120-structured-uniform}
For each window, let the complete form have an outer finite head
and positive complement $T_\lambda$, with bounded coupling
$B_\lambda$. For a finite-support solve $X_\lambda$ put
\[
 \begin{split}
 K_{X,\lambda}&=F_\lambda-B_\lambda^*X_\lambda
       -X_\lambda^*B_\lambda+X_\lambda^*T_\lambda X_\lambda,\\
 r_\lambda&=B_\lambda-T_\lambda X_\lambda.
 \end{split}
\]
Assume the inner tail admits the factors of
Proposition~\ref{prop:v120-structured-inverse}, and write
$r_\lambda=(h_\lambda,k_\lambda)$ in that inner decomposition.
If, in operator norm on the entire outer head,
\[
 \begin{gathered}
 K_{X,\lambda}\succeq-\alpha_\lambda I,\qquad
 \norm{D_{g,\lambda}^{-1/2}k_\lambda}\le e_\lambda,\\
 \norm{C_\lambda(h_\lambda-Z_\lambda^*k_\lambda
       -R_\lambda^*D_{g,\lambda}^{-1}k_\lambda)}\le t_\lambda,
 \qquad \alpha_\lambda\ge0,
 \end{gathered}
\]
then the whole closed form satisfies
\begin{equation}\label{eq:v120-structured-uniform}
 \mathsf W_\lambda\succeq
 -\left(\alpha_\lambda+e_\lambda^2+
               t_\lambda^2/\mu_\lambda\right)I.
\end{equation}
Thus convergence of the displayed error to zero along a cofinal
window family is sufficient for the large-support sign criterion.
No such convergence is asserted here.
\end{proposition}
\begin{proof}
The exact outer Schur matrix is
\[
 S_\lambda=K_{X,\lambda}-r_\lambda^*T_\lambda^{-1}r_\lambda.
\]
Apply \eqref{eq:v120-structured-inverse} to each outer head vector.
Completing the outer form square leaves $S_\lambda$ on the head
and a nonnegative complementary square. Since the head projection
is orthogonal, its norm is at most the full ordinary norm, proving
\eqref{eq:v120-structured-uniform}.
\end{proof}

Two scales cannot be suppressed in applying this criterion.
If only individual bounds
$\norm{D^{-1/2}r e_j}\le\delta$ are known for a head of
dimension $d$, the valid conclusion is
$\norm{D^{-1/2}r}^2\le d\delta^2$. Equal residual columns
attain this factor. In particular, a single source estimate or
columnwise decay does not prove the required growing-head bound.
Likewise, $C\mathcal L C^*\succeq-\eta I$ implies only
\[
 \mathcal L\succeq-\eta C^{-1}C^{-*}
             \succeq-\eta\norm{C^{-1}}^2I.
\]
Congruence preserves positivity, but it does not preserve the size
of an ordinary negative error. Using merely $T\succeq cD$
yields the coarser correction
$c^{-1}\norm{D^{-1/2}r}^2$, and its division by $c$ is sharp
already for a scalar tail $T=cD$.

For clarity, the remote truncation itself has an explicit uniform
budget. If $G$ has Fourier support $|m|\le M$, $J>M$, and
$|b_n|\le B_*$, the actual off-diagonal divided difference gives
\begin{equation}\label{eq:v120-separated-uniform-budget}
 (Q_J\mathsf W_\lambda G)^*(Q_J\mathsf W_\lambda G)
 \preceq\frac{8B_*^2(2M+1)}{J-M}\,G^*G.
\end{equation}
Its proof bounds the squared Hilbert--Schmidt norm of the separated
row/column block by
$4B_*^2(2M+1)\sum_{|n|>J}(|n|-M)^{-2}$ and uses the
integral bound $\sum_{n>J}(n-M)^{-2}\le1/(J-M)$.
The diagonal is absent. If the diagonal lower weights $g_n$ are increasing,
dividing by $g_{J+1}$ gives the weighted Gram bound. For
$G=[I;-X]$ its exact final factor is $I+X^*X$, with no extra
number-of-columns factor. The moment bound is usually much sharper,
but a growing moment order controls only its geometric remainder;
it does not force the leading moment Gram to vanish.

These estimates control certification of the low-mode correction.
The simultaneous certificate above settles the complete even form at
$\lambda=4$. The odd sign and the corresponding growing-window
head/inverse estimates remain unproved; positivity of a remote tail
alone does not supply them.


The time--frequency geometry suggests a canonical family of candidate projections $P_a$.  Put $\lambda=e^a$ and consider the Slepian concentration operator at time--bandwidth parameter $c\asymp\lambda^2$.  Connes--Consani construct corrected co-Poisson ``prolate vectors'' from the highly concentrated PSWFs and observe that the number of extremely small semilocal Weil eigenvalues grows on the same $\asymp\lambda^2$ scale \cite{ConnesConsani2023}.  Quantitative Slepian eigenvalue bounds sharpen the geometric picture.  If $\theta_j(c)$ are the concentration eigenvalues and one chooses
\[
 \delta_\lambda=\lambda^{-K},
\]
then the number of indices in the plunge region
\[
 \delta_\lambda<\theta_j(c)<1-\delta_\lambda
\]
is at most
\[
 O\!\left(\log c\,\log\frac1{\delta_\lambda}\right)=O_K(a^2),
 \tag{G2.5}
\]
using the standard quantitative transition estimates for one-dimensional time--frequency limiting operators \cite{SlepianTransition}.  Thus the concentration geometry has three distinct pieces:
\[
 \boxed{
 \mathcal H_a
 =\mathcal R_a^{\rm deep}\oplus
   \mathcal T_a^{\rm plunge}\oplus
   \mathcal C_a^{\rm low},
 }
 \tag{G2.6}
\]
where $\dim\mathcal R_a^{\rm deep}\asymp\lambda^2$, the plunge block has dimension $O_K(a^2)$, and $\mathcal C_a^{\rm low}$ is an infinite-dimensional low-concentration complement.

Equation \eqref{eq:G1-continuum}, now proved using Theorem~\ref{thm:v14-radical}, controls fixed-band pairings for the fixed-order proxy treated above. It does \emph{not} prove the operator-norm estimate $\norm{A_aP_a}\to0$ uniformly on the full deep block. A possible route would combine the exact global-radical identity with the concentration defect, but it needs a uniform bound on the truncation cost in the relevant form or graph norm, as well as normalized endpoint control for the growing family. Neither follows from the fixed-source approximation theorem, whose constants may depend on the source. Thus the growing-block estimate remains an open subgate, not a consequence of the weak G1 closure.

More importantly, the small plunge dimension does \emph{not} dispose of the complement.  Time--frequency concentration alone supplies no sign for the Weil form on $\mathcal C_a^{\rm low}$.  One sufficient G2 route seeks a coercivity estimate
\[
 \boxed{
 P_{\mathcal C_a^{\rm low}}A_aP_{\mathcal C_a^{\rm low}}\succeq c_a I
 }
 \tag{G2.7}
\]
together with control of the couplings to the deep and plunge blocks.
Proposition~\ref{prop:v118-gap-free} gives an alternative: a uniform
operator residual on the removed block and asymptotic nonnegativity of
its whole complement suffice, without a strictly positive gap.  Propositions~\ref{prop:v119-prime-essential}--\ref{prop:v119-prime-scale}
make one limitation exact: finite source removal leaves the unsigned
prime norm unchanged, and that norm is $(1+o(1))e^a$.
Subtracting this scalar from the minimum archimedean tail diagonal
forces the expensive cutoff \eqref{eq:v119-unsigned-cutoff}.
This does not exclude a frequency-weighted or directional estimate;
the recurrence vectors restoring the prime norm may have much larger
archimedean energy than that minimum.

This three-block formulation is useful mainly because it localizes the unresolved content.  Classical prolate theory controls the geometry and shows that the genuine transition region is only polynomial in $a$; the RH-strength issue is the semilocal Weil coercivity of the low-concentration complement and the uniform residual control of the entire deep near-radical block.  Neither statement follows from Slepian concentration by itself.

\subsection*{Finite trial positivity and the exact omitted-tail correction}

The trial-head term in Proposition~\ref{prop:v120-structured-uniform}
must be distinguished from the exact Schur complement. A finite-support
certificate can settle the former without settling the latter.

\begin{proposition}[Finite-support trial forms and Galerkin Schur bounds]
\label{prop:v121-finite-trial}
Let the outer head be finite dimensional, let $T\succeq cI$ with
$c>0$ be a positive self-adjoint tail operator, and let $B$ be the
bounded head-to-tail coupling. For $X$ with range in $\mathcal D(T)$,
define $K_X,r$ as in Proposition~\ref{prop:v120-structured-uniform}
and put $S_\infty=F-B^*T^{-1}B$. Then
\begin{align}
 K_X-S_\infty
  &=(X-T^{-1}B)^*T(X-T^{-1}B)\succeq0,
       \label{eq:v121-square}\\
 S_\infty&=K_X-r^*T^{-1}r.\label{eq:v121-residual}
\end{align}
Let $P_M$ be the orthogonal projection onto a finite tail space in $\mathcal D(T)$, and
write $T_M=P_MT|_{\operatorname{Ran}P_M}$, $B_M=P_MB$ and
$S_M=F-B_M^*T_M^{-1}B_M$. If $X=P_MX$, then
\begin{equation}\label{eq:v121-finite-trial}
 K_X=\begin{bmatrix}I\\-X\end{bmatrix}^{*}
       \begin{bmatrix}F&B_M^*\\B_M&T_M\end{bmatrix}
       \begin{bmatrix}I\\-X\end{bmatrix}.
\end{equation}
In particular, $S_M\succeq0$ implies $K_X\succeq0$ for every such
$X$, and strict positivity transfers as well. This is the complete
operator's $K_X$, not a truncation of that trial form.

For the exact finite solve $X_M=T_M^{-1}B_M$, embedded in the complete
tail, one has $K_{X_M}=S_M$ and
\begin{equation}\label{eq:v121-galerkin-defect}
 \boxed{S_\infty=S_M-r_M^*T^{-1}r_M,
 \qquad r_M=B-TX_M.}
\end{equation}
If the finite tail spaces are nested and their union is a form core
for $T$, then $S_M\downarrow S_\infty$ in matrix norm.
\end{proposition}
\begin{proof}
Expanding the first square gives $K_X-F+B^*T^{-1}B$.
Since $r=-T(X-T^{-1}B)$, its value is $r^*T^{-1}r$,
which proves both identities. In \eqref{eq:v121-finite-trial}
every trial lies in the head plus the retained finite tail space,
so compression changes none of the quadratic terms. Finite square completion gives
\[
 K_X=S_M+(X-T_M^{-1}B_M)^*T_M(X-T_M^{-1}B_M).
\]
Thus $K_X\succeq S_M$, and the graph map $h\mapsto(h,-Xh)$ is
injective. Substitution of the exact finite solve proves
\eqref{eq:v121-galerkin-defect}; its retained residual rows vanish,
but its omitted rows need not vanish.

For each head vector $h$, the variational identity is
\[
 \ip{B^*T^{-1}Bh}{h}
 =\sup_{y\in\mathcal D(T^{1/2})}
 \{2\Re\ip{Bh}{y}-\norm{T^{1/2}y}^2\}.
\]
Restricting the supremum to $\operatorname{Ran}P_M$ gives
$\ip{B_M^*T_M^{-1}B_Mh}{h}$. Nestedness and form-core density make
these quantities increase to the full supremum. Polarization and
finite dimensionality of the head give matrix-norm convergence of
the Schur complements in the stated decreasing order.
\end{proof}

A finite list of positive $S_M$ therefore does not prove positivity
of $S_\infty$. Positivity for every member of a cofinal sequence at
a fixed window would imply $S_\infty\succeq0$, but strict positivity
would still require a positive limiting margin. The structured inverse
bounds concern this same $T$; a comparison metric or a shifted lower
operator does not redefine the target Schur complement.

\begin{remark}[The semidefinite range condition]\label{rem:v121-range}
For $T\succeq0$, the notation $T^{-1}B$ requires an additional
condition. A form version is available if
$\operatorname{Ran}B\subset\operatorname{Ran}T^{1/2}$. Define the
reduced inverse $C=T^{-1/2}B$, with values orthogonal to $\ker T$.
For $X$ with range in the form domain, use
$(T^{1/2}X)^*T^{1/2}X$ in the definition of $K_X$. Then
\[
 K_X-(F-C^*C)=(T^{1/2}X-C)^*(T^{1/2}X-C)\succeq0.
\]
Orthogonality to $\ker T$, or membership in the closure of its range,
does not suffice when the range is not closed. At the certified
$\lambda=4$ split there is no such issue:
Proposition~\ref{prop:v120-lambda4-tail} gives
$T\succeq10^{-8}D\succeq1.7940\cdot10^{-8}I$ on $Q_{16}$.
\end{remark}

\begin{proposition}[A certified local zero head-error term]
\label{prop:v121-alpha-zero}
For the actual $\lambda=4$ Weil form with outer head $E_{16}$,
$K_X\succ0$ in both parity sectors for $X=0$ and for the frozen
dyadic solves supported on modes $17,\ldots,256$ specified in
Appendix~\ref{app:v121-audit}. Consequently $\alpha_4=0$ is
admissible in Proposition~\ref{prop:v120-structured-uniform} for
each of these choices, on the complete operator.
\end{proposition}
\begin{proof}
Assemble the exact finite entries with the separate logarithmic
diagonal, normalized parity basis and $L=2\log4$. Outward interval
$LDL^*$ factorization of $K_X$ verifies all $17$ even pivots and
all $16$ odd pivots strictly positive, both for $X=0$ and for the
specified dyadics. The identical witnesses pass at $768$ and $896$
bits; Appendix~\ref{app:v121-audit} gives the files and gates.
Real symmetric positivity also proves positivity on the complex
parity spaces. Equation~\eqref{eq:v121-finite-trial} identifies these
matrices with the complete trial forms. No infinite Schur sign or
small residual is used in this argument.
\end{proof}

This result does not certify $\alpha_\lambda=0$ at untested windows
or for larger unsupported solves. Positivity of an exact Schur
complement would guarantee $K_X\succeq0$ for every trial; without
that premise the least admissible $\alpha$ can depend on $X$.
For example, $F=0$, $B=T=1$ gives $S=-1$, $K_0=0$ and $K_1=-1$.
Even $K_0\succ0$ need not help the signed target: its residual is
the entire coupling $B$.

With a proved $\alpha_\lambda=0$ along a growing family, the precise
remaining sufficient estimates would be
\begin{equation}\label{eq:v121-scaled-target}
 e_\lambda\longrightarrow0,
 \qquad t_\lambda/\sqrt{\mu_\lambda}\longrightarrow0.
\end{equation}
The second condition is stronger than $t_\lambda\to0$ when
$\mu_\lambda\to0$. Geometric convergence of an inverse approximation
improves an enclosure of $r_\lambda^*T_\lambda^{-1}r_\lambda$;
it does not force that residual energy to vanish for a fixed trial.
These growing-window estimates remain unproved. Failure of this
sufficient bound would not itself refute Weil positivity.


\subsection*{An inverse-weighted finite spectral certificate}
The complement estimate can also identify a single candidate without
dividing its entire residual by the smallest spectral gap.
The following finite statement is independent of any PSWF approximation.

\begin{proposition}[Inverse-weighted finite certificate]
\label{prop:v112-schur-angle}
Let $W$ be a finite Hermitian matrix, $u$ a unit vector, and $P$ the
orthogonal projection onto $u^\perp$. Put
\[
 \alpha=\ip{Wu}{u},\quad r=PWu,\quad
 C=PWP|_{u^\perp},\quad A=C-\alpha I.
\]
If $A\succeq\gamma I$ for some $\gamma>0$, then $W$ has a simple
lowest eigenvalue $\lambda_0\le\alpha$. Set
$q=\norm{A^{-1}r}$ and $E=\ip{A^{-1}r}{r}$. For the angle $\theta$
between $u$ and the lowest eigenspace,
\begin{equation}\label{eq:v112-angle}
 \tan\theta\le q,\qquad
 \sin\theta\le\frac{q}{\sqrt{1+q^2}},\qquad
 \alpha-E\le\lambda_0\le\alpha.
\end{equation}
If $W$ commutes with reflection and $u$ is even, its ground state is even.
For any approximate solve $z\in u^\perp$ one also has
\begin{equation}\label{eq:v112-solve}
 q\le\norm z+\frac{\norm{r-Az}}{\gamma}.
\end{equation}
\end{proposition}
\begin{proof}
Min--max gives $\lambda_0\le\alpha$ and
$\lambda_1\ge\inf\sigma(C)>\alpha$. Write a ground vector as $bu+h$,
where $h\perp u$ and $b\ne0$. The complementary equation gives
$h=-b(C-\lambda_0 I)^{-1}r$. With $t=\alpha-\lambda_0\ge0$, the
scalar equation is $t=\ip{(A+tI)^{-1}r}{r}$. Spectral calculus for
the same positive operator $A$ gives $t\le E$ and
$\norm{(A+tI)^{-1}r}\le\norm{A^{-1}r}$. These prove
\eqref{eq:v112-angle}. An odd ground state would be orthogonal to the
even $u$, contradicting the complement bound. Finally
$A^{-1}r=z+A^{-1}(r-Az)$ proves \eqref{eq:v112-solve}.
\end{proof}

For certified approximations $\widetilde A,\widetilde r$, suppose
$\norm{A-\widetilde A}\le\delta_A$,
$\norm{r-\widetilde r}\le\delta_r$ and
$\widetilde A\succeq\widetilde\gamma I$ with
$\widetilde\gamma>\delta_A$. Then the usable error budget is
\[
 q\le\norm z+
 \frac{\delta_r+\delta_A\norm z+
       \norm{\widetilde r-\widetilde A z}}
      {\widetilde\gamma-\delta_A}.
\]
For a fixed exact candidate and matrix error $\norm{W-\widetilde W}\le\eta$,
one may take $\delta_r=\eta$ and $\delta_A=2\eta$.
True-source approximation error is a separate input. The condition
$\gamma>0$ is sufficient, not necessary, for source-to-ground
convergence: in $W=\operatorname{diag}(0,d,1)$, the candidate
$u=\sqrt{1-\tau^2}e_0+\tau e_2$ has angle tending to zero but
$\gamma\le d-\tau^2<0$ when $0<d<\tau^2$.

Appendix~\ref{app:v112-certificate} certifies the hypotheses for one
explicit rational candidate at $\lambda=3$, $N=64$, using the actual
arithmetic Weil matrix. That certificate alone proves no uniform complement theorem or
exact-source error bound. The following argument supplies the latter
at its fixed parameters; endpoint and graph transfer remain separate.


\subsection*{From a finite prolate approximation to the exact source}
The preceding certificate treats its input vector as exact. To identify
that vector with a projection of the manuscript's source, the omitted
Legendre tail must also be controlled. Here this can be done at fixed
parameters, including pointwise values rather than only an $L^2$ error.

\begin{lemma}[Certified angular eigenfunction error]
\label{lem:v113-angular}
On $L^2(-1,1)$ let
$e_\ell(x)=\sqrt{(2\ell+1)/2}\,P_\ell(x)$,
$De_\ell=\ell(\ell+1)e_\ell$, and
$\mathsf J_c=D+c^2X^2$, where $Xf(x)=xf(x)$ and
$\mathcal D(\mathsf J_c)=\mathcal D(D)$.
Restrict to the even subspace. Let $J_M$ be its compression to
$e_0,e_2,\ldots,e_{2M-2}$, let $b_M$ be the link from its last
coordinate to $e_{2M}$, and put $d_M=2M(2M+1)$.
For $x<d_M$, the number of eigenvalues of $\mathsf J_c$ below $x$
is between the negative inertia counts of
\begin{equation}\label{eq:v113-inertia}
 J_M-xI,\qquad
 J_M-xI-\frac{b_M^2}{d_M-x}E_{M-1,M-1}.
\end{equation}
If both matrices are nonsingular with the same count, that count is exact.

Suppose these counts certify exactly one eigenvalue $\nu$ in
$(\xi-g,\xi+g)$ and no other spectrum there, with $g>0$.
Let $\xi\in\mathbb R$, let $v$ be an exact normalized real vector in
this even finite span, and suppose
$\norm{(\mathsf J_c-\xi)v}\le\epsilon<g$.
For the normalized eigenfunction $\psi$ at $\nu$, with sign chosen
so that $\ip v\psi>0$, put $\delta=\sqrt2\epsilon/g$.
Then
\begin{align}
 |\nu-\xi|&\le\epsilon,& \norm{v-\psi}_2&\le\delta,\notag\\
 \norm{v-\psi}_\infty
 &\le 2\epsilon+(1+|\xi|+c^2)\delta.
 \label{eq:v113-pointwise}
\end{align}
The residual includes the first omitted coordinate $b_Mv_{M-1}$.
The gap $g$ is measured from $\xi$ to the other exact eigenvalues.
\end{lemma}
\begin{proof}
The diagonal operator $D$ is self-adjoint with compact resolvent;
$c^2X^2$ is bounded, positive and self-adjoint. Its angular realization
is the regular PSWF operator (the differential spectral parameter is
shifted by $c^2$ relative to the convention in \cite{DLMFPSWF}).
The omitted compression $T$ satisfies $T\succeq d_MI$. For $x<d_M$
its positive inverse gives the finite Schur complement
\[
 J_M-xI-b_M^2\ip{(T-xI)^{-1}e_{2M}}{e_{2M}}E_{M-1,M-1}.
\]
The scalar lies between $0$ and $(d_M-x)^{-1}$. Inertia and min--max
therefore prove \eqref{eq:v113-inertia}. Completing the square in the
quadratic form justifies the Schur argument for the unbounded $T$.

The spectral theorem places an eigenvalue within $\epsilon$ of $\xi$;
the certified interval identifies it as $\nu$. All other coefficients
of $v$ have squared mass at most $(\epsilon/g)^2$. Choosing positive
overlap gives $\norm{v-\psi}_2\le\sqrt2\epsilon/g$.
For $e=v-\psi$,
\[
 \mathsf J_ce=\xi e+(\xi-\nu)\psi+(\mathsf J_c-\xi)v,
\]
so $\norm{(I+D)e}_2\le2\epsilon+(1+|\xi|+c^2)\delta$.
Finally $|P_\ell(x)|\le1$ on $[-1,1]$ and
\[
 \sum_{\ell\ge0}\frac{(2\ell+1)/2}{[1+\ell(\ell+1)]^2}
 \le\frac12+\frac12\sum_{\ell\ge1}
             \left(\frac1{\ell^2}-\frac1{(\ell+1)^2}\right)=1.
\]
Cauchy--Schwarz gives uniform convergence of the Legendre expansion and
$\norm e_\infty\le\norm{(I+D)e}_2$, including both endpoint traces.
\end{proof}

Here is the source error budget used below. Set $\lambda=3$ and let
$\widehat h_j=\lambda^{-1/2}v_j(\cdot/\lambda)$ approximate
$h_{j,\lambda}$ for $j=0,4$. Let $d_j$ and $u_j$ bound their physical
$L^2$ and uniform errors. Define
\[
 \widehat a_0=-\int\widehat h_4,\quad
 \widehat a_4=\int\widehat h_0,\quad
 \widehat h^\circ=\widehat a_0\widehat h_0+
                  \widehat a_4\widehat h_4.
\]
Both this polynomial and the exact $h^\circ$ have zero integral.
With $d_{a_0}=\sqrt{2\lambda}\,d_4$ and
$d_{a_4}=\sqrt{2\lambda}\,d_0$, a valid uniform error is
\begin{equation}\label{eq:v113-source-budget}
 \eta=\sum_{j=0,4}
 \{(|\widehat a_j|+d_{a_j})u_j+
                d_{a_j}\norm{\widehat h_j}_\infty\}.
\end{equation}
Use the fixed repair with $b=3/4$,
$\phi(t)=\exp(1-1/(1-(t/b)^2))$ for $|t|<b$, zero outside, and
$\psi(t)=\phi(t)(1-Ct^2)$ where
$C=\int_0^b\phi/\int_0^b t^2\phi$.
Since $\phi\le1$ everywhere and $\phi\ge1/2$ on $[b/4,b/2]$,
$Cb^2\le128$ and $\norm\psi_\infty\le129$.
Consequently, with $Z=|\widehat h^\circ(0)|$,
\begin{equation}\label{eq:v113-repair-budget}
 \norm{\widetilde h-\widehat h^\circ}_\infty
 \le U:=\eta+129(Z+\eta).
\end{equation}
This bounds the entire exact repair without assuming a numerical
quadrature value for $C$.

Let $\widehat p$ be the inversion-even co-Poisson proxy of
$\widehat h^\circ$ on the same window. At almost every
$u\in[\lambda^{-1},\lambda]$, at most $\lambda/u$ summands occur,
so the physical factor $\sqrt u$ gives
$|\mathcal E(\widetilde h-\widehat h^\circ)(u)|
 \le\lambda^{3/2}U$.
Inversion averaging is an $L^2(du/u)$ contraction. Bessel's inequality
and the strict interior endpoint formula therefore give
\begin{align}
 \norm{P_Np_3-P_N\widehat p}_2&\le\sqrt{\lambda^3L}\,U,
       &L&=2\log3,\label{eq:v113-projection-budget}\\
 |B_3-\widehat B|&\le\tfrac12\left(\sqrt\lambda\,\eta+
                                     8\lambda^{-1/2}U\right),
 \label{eq:v113-endpoint-budget}
\end{align}
where
$\widehat B=\tfrac12\{\sqrt\lambda\,\widehat h^\circ(\lambda)
+\lambda^{-1/2}\sum_{m=1}^{8}\widehat h^\circ(m/\lambda)\}$.
The $m=9$ term is excluded from this lower interior trace; the repair
vanishes at the upper endpoint. These are fixed-parameter estimates,
not a growing-window spectral theorem.

\begin{proposition}[Exact-source angle and endpoint at one window]
\label{prop:v113-source-certificate}
Use the exact repaired source of Proposition~\ref{prop:v12-source}
with the bump just specified. At $\lambda=3,N=64$, let $c^*$ be the
unnormalized rational coefficient vector in
Proposition~\ref{prop:v112-finite-certificate}, and let $v^*$ be the
coefficient vector of the exact $P_{64}p_3$ in the same translated
Fourier basis. The interval certificate in
Appendix~\ref{app:v113-source} proves
\[
 \norm{v^*-c^*}_2<2.60\,10^{-36},\qquad
 \sin\angle(v^*,\text{ground state of }W_{3,64})<0.000216.
\]
It also proves
\begin{equation}\label{eq:v113-endpoint-failure}
 5.58669\,10^{-19}<B_3<5.58670\,10^{-19},\qquad
 1.374<\frac{|\delta_{64}(P_{64}p_3)-B_3|}{B_3}<1.375.
\end{equation}
Thus the small ordinary ground-state angle coexists with a relative
physical-endpoint error greater than one for this exact projected source.
\end{proposition}
\begin{proof}
The appendix certifies all error budgets above, including the full
infinite angular residuals, and the exact polynomial Fourier moments.
It gives the first bound and the displayed endpoint intervals.
For nonzero vectors $v^*,c^*$ with error at most $e<\norm{c^*}$,
\[
 \norm{P_{\mathbb Cv^*}-P_{\mathbb Cc^*}}
 =\sin\angle(v^*,c^*)\le\frac{e}{\norm{c^*}-e}.
\]
This is below $4\,10^{-36}$ here. The triangle inequality for rank-one
orthogonal projections and the sharper enclosed sine bound from
Proposition~\ref{prop:v112-finite-certificate} prove the angle claim.
No new positive-complement estimate for the perturbed candidate is
assumed. Equation~\eqref{eq:v113-endpoint-failure} concerns physical
endpoint evaluation, not the equal-coefficient functional without its
$L^{-1/2}$ factor.
\end{proof}

\section{Weighted G1 transfer and the arithmetic jump train}\label{sec:weighted}
\begin{proposition}[Cofinal finite G1 bound on a weighted sampler class]\label{prop:g1-weighted-sampler}
Use the explicit repaired source of Section~\ref{sec:g1}.
Let $R>3$, let $K\ge2N$, and write
\[
 x=\sum_{|m|\le K}x_mU_m\in E_K,
 \qquad
 \|x\|_{\mathcal H_{R,L}}
 :=\sup_{|m|\le K}
 \sqrt L\left(1+\frac{|m|}{L}\right)^R|x_m|.
\]
If $N\ge2L$, then the filtered proxy of
Proposition~\ref{prop:g1-vp-transfer} satisfies
\begin{equation}\label{eq:g1-weighted-sampler}
 \boxed{
 \frac{|\langle D_{\log}W_{\lambda,K}
 k^{\mathrm{VP}}_{\lambda,N},x\rangle|}{|B_\lambda|}
 \le C_R\|x\|_{\mathcal H_{R,L}}
 \left[
 \mathfrak r(\lambda)
 +\frac{\lambda\mathcal V_\lambda}{|B_\lambda|}
 \left\{
 \frac{L}{N}
 +\left(\frac{L}{N}\right)^{R-1}(1+\log(2N))
 \right\}
 \right].}
\end{equation}
The bound is uniform in the ambient cutoff $K$.  Thus, once the arithmetic
core cutoff $N$ has been fixed, enlarging the Hardy-sampler cutoff cannot
spoil this G1 estimate.
\end{proposition}
\begin{proof}
Put $e_N=k^{\mathrm{VP}}_{\lambda,N}-p_\lambda$.  Because the
Vallee--Poussin multiplier is one for $|n|\le N$, one has
$\widehat e_N(n)=0$ there, while for $|n|>N$ the BV estimate gives
\[
 |\widehat e_N(n)|
 \le |\widehat p_\lambda(n)|
 \le \frac{\sqrt L\,\mathcal V_\lambda}{2\pi|n|}.
\]
The coefficient hypothesis and $R>3$ imply, by comparison with Riemann sums,
\[
 \|D_{\log}x\|_2+\|D_{\log}^2x\|_2
 +|(D_{\log}x)(-a)|+|(D_{\log}x)(a)|
 \le C_R\|x\|_{\mathcal H_{R,L}}.
\]
Indeed the two $L^2$ sums reduce to
$L^{-1}\sum_m (|m|/L)^{2j}(1+|m|/L)^{-2R}$, $j=1,2$, and each endpoint
trace to $L^{-1}\sum_m (|m|/L)(1+|m|/L)^{-R}$.
Proposition~\ref{prop:g1-sobolev}, applied to $D_{\log}x$, therefore gives
\[
 |QW_\lambda(p_\lambda,D_{\log}x)|
 \le C_R|B_\lambda|\mathfrak r(\lambda)
       \|x\|_{\mathcal H_{R,L}}.
\]

For the projection error separate the off-diagonal and diagonal matrix
entries.  Lemma~\ref{lem:g1-high-low} gives, uniformly in $K$,
\[
 \begin{split}
 |QW_\lambda(e_N,D_{\log}x)|_{\mathrm{off}}
 \le{}& C_R\lambda\mathcal V_\lambda
 \|x\|_{\mathcal H_{R,L}}
 \sum_{|n|>N}\frac1{|n|}
 \sum_{m\ne n}\frac{a_m}{|n-m|},\\
 a_m:={}&\frac{|m|}{L}
       \left(1+\frac{|m|}{L}\right)^{-R}.
 \end{split}
\]
For nonzero integer $n-m$ one may replace $|n-m|^{-1}$ by
$2(1+|n-m|)^{-1}$.  For $|n|>N\ge2L$, split the inner sum into
$|m|\le|n|/2$, $|n|/2<|m|\le2|n|$, and $|m|>2|n|$.
Since $\sum_m a_m\le C_RL$, the first range is $O_R(L/|n|)$.
On the middle range,
$a_m\le C_R(|n|/L)^{1-R}$ and the harmonic sum contributes
$O(1+\log(2|n|))$.  On the last range the denominator is $\gg|m|$, and
summation of $m^{-R}$ gives $O_R((|n|/L)^{1-R})$.  Hence
\[
 \sum_{m\ne n}\frac{a_m}{|n-m|}
 \le C_R\left[
 \frac{L}{|n|}
 +\left(\frac{|n|}{L}\right)^{1-R}
  (1+\log(2|n|))\right],
\]
and therefore
\begin{equation}\label{eq:g1-weighted-offdiag}
 |QW_\lambda(e_N,D_{\log}x)|_{\mathrm{off}}
 \le C_R\lambda\mathcal V_\lambda\|x\|_{\mathcal H_{R,L}}
 \left[
 \frac{L}{N}
 +\left(\frac{L}{N}\right)^{R-1}(1+\log(2N))\right].
\end{equation}

The diagonal must be handled separately.  Corollary~\ref{cor:g1-complete-matrix},
the BV coefficient estimate, and the definition of $\mathcal H_{R,L}$ give
\[
 \begin{split}
 |QW_\lambda(e_N,D_{\log}x)|_{\mathrm{diag}}
 \le{}& C\frac{\mathcal V_\lambda}{L}
 \|x\|_{\mathcal H_{R,L}}\\
 &\times\sum_{N<|n|\le K}
 \bigl(\lambda+1+\log(2+|n|/L)\bigr)
 \left(1+\frac{|n|}{L}\right)^{-R}.
 \end{split}
\]
Comparison with the corresponding decreasing integral, using $N\ge2L$ and
$R>3$, yields
\begin{equation}\label{eq:g1-weighted-diag}
 |QW_\lambda(e_N,D_{\log}x)|_{\mathrm{diag}}
 \le C_R\mathcal V_\lambda\|x\|_{\mathcal H_{R,L}}
 \bigl(\lambda+1+\log(2N)\bigr)
 \left(\frac{L}{N}\right)^{R-1}.
\end{equation}
Because $\lambda\ge2$, \eqref{eq:g1-weighted-diag} is absorbed by the
second term on the right of \eqref{eq:g1-weighted-offdiag}.  Thus the
projection-error contribution has exactly the rate stated in
\eqref{eq:g1-weighted-sampler}, despite the logarithmic diagonal growth.
Finally
$\langle D_{\log}W_{\lambda,K}k^{\mathrm{VP}},x\rangle
=QW_\lambda(k^{\mathrm{VP}},D_{\log}x)$ because both finite vectors lie in
$E_K$.  Combining the two pieces proves \eqref{eq:g1-weighted-sampler}.
\end{proof}

\subsection*{Exact extraction of the arithmetic jump train}
The BV estimate above treats all variation of the prolate/co-Poisson proxy in
one lump.  That is unnecessarily expensive after normalization by the tiny
endpoint amplitude.  The actual discontinuities have a much more rigid
arithmetic form and can be separated exactly.

Let $y=\log(\lambda u)\in[0,L]$ and let $B_\lambda^{+}
=\sqrt\lambda\,\widetilde h_\lambda(\lambda^-)$ denote the upper source
boundary amplitude.  Put
\[
 \mathcal M_\lambda:=\{m\in\mathbb N:1<m<\lambda^2\}.
\]
The fixed smooth exact-moment repair used above is supported strictly inside
the source interval, so it creates no endpoint jumps.

\begin{lemma}[Exact arithmetic jump train]\label{lem:g1-jump-train}
There is a zero-mean periodic BV function $\jmath_\lambda$ on the logarithmic
circle such that its distributional derivative is
\begin{equation}\label{eq:g1-jump-measure}
 d\jmath_\lambda
 =\frac{B_\lambda^{+}}2
 \sum_{m\in\mathcal M_\lambda}m^{-1/2}
 \bigl(\delta_{\log m}-\delta_{L-\log m}\bigr).
\end{equation}
Moreover
\[
 c_\lambda:=p_\lambda-\jmath_\lambda
\]
is continuous and periodic, is piecewise smooth, and its first logarithmic
derivative has bounded variation.  If
\[
 t_n:=\frac{2\pi n}{L},\qquad
 \mathcal D_\lambda(t):=
 \sum_{m\in\mathcal M_\lambda}m^{-1/2}\sin(t\log m),
\]
then for every $n\ne0$,
\begin{equation}\label{eq:g1-jump-fourier-exact}
 \boxed{
 \widehat{\jmath_\lambda}(n)
 =-\frac{B_\lambda^{+}}{t_n\sqrt L}\,
 \mathcal D_\lambda(t_n).}
\end{equation}
Consequently, for all sufficiently large $\lambda$,
\begin{equation}\label{eq:g1-jump-fourier-trivial}
 \boxed{
 |\widehat{\jmath_\lambda}(n)|
 \le C|B_\lambda|\frac{\lambda\sqrt L}{|n|}.}
\end{equation}
Finally, if
\[
 \mathcal W_\lambda
 :=\bigl\|D_y^2c_\lambda\bigr\|_{\mathcal M(\mathbb T_L)}
\]
denotes the total variation of the distributional second derivative of the
continuous remainder, then
\begin{equation}\label{eq:g1-continuous-fourier}
 \boxed{
 |\widehat{c_\lambda}(n)|
 \le \frac{L^{3/2}\mathcal W_\lambda}{4\pi^2n^2},
 \qquad n\ne0.}
\end{equation}
All quantities in the statement are determined by the explicit source and
$\lambda$; no zeta zero is used.
\end{lemma}
\begin{proof}
For
$f_\lambda(u)=\mathcal E(\widetilde h_\lambda)(u)$, the summand with index
$m$ disappears as $u$ increases through $u=\lambda/m$.  Its one-sided jump is
therefore
\[
 -\left(\frac\lambda m\right)^{1/2}
  \widetilde h_\lambda(\lambda^-)
 =-\frac{B_\lambda^{+}}{\sqrt m}.
\]
In the $y$ coordinate this occurs at $L-\log m$.  The inverted function
$f_\lambda(u^{-1})$ has the opposite jump at $\log m$.  Since
$p_\lambda=\frac12(f_\lambda+Jf_\lambda)$ on the window, its complete
interior jump measure is exactly \eqref{eq:g1-jump-measure}.  The total mass
of that measure is zero, so it is the derivative of a periodic BV function;
fixing zero mean makes $\jmath_\lambda$ unique.  Subtracting it removes every
function jump.  The arithmetic sum has only finitely many thresholds and is
smooth between them, so $c_\lambda$ is continuous, piecewise smooth, and
$D_yc_\lambda$ is BV.

For $n\ne0$, periodic Stieltjes integration by parts gives
\[
 \widehat{\jmath_\lambda}(n)
 =\frac1{it_n\sqrt L}
 \int_0^Le^{-it_ny}\,d\jmath_\lambda(y).
\]
Because $e^{-it_nL}=1$, \eqref{eq:g1-jump-measure} yields
\[
 \int_0^Le^{-it_ny}\,d\jmath_\lambda(y)
 =\frac{B_\lambda^{+}}2
 \sum_{m\in\mathcal M_\lambda}m^{-1/2}
 \{m^{-it_n}-m^{it_n}\}
 =-iB_\lambda^{+}\mathcal D_\lambda(t_n),
\]
which is \eqref{eq:g1-jump-fourier-exact}.  The elementary estimate
\[
 |\mathcal D_\lambda(t)|
 \le\sum_{m<\lambda^2}m^{-1/2}\le2\lambda
\]
and $B_\lambda^{+}\asymp B_\lambda$ from
Corollary~\ref{cor:pswf-cross-tail} prove
\eqref{eq:g1-jump-fourier-trivial}.

For the continuous remainder there is no first-order boundary term.  A second
periodic Stieltjes integration by parts, now against the finite measure
$D_y^2c_\lambda$, gives
\[
 |\widehat{c_\lambda}(n)|
 \le\frac{\mathcal W_\lambda}{\sqrt L\,t_n^2}
 =\frac{L^{3/2}\mathcal W_\lambda}{4\pi^2n^2},
\]
proving \eqref{eq:g1-continuous-fourier}.
\end{proof}

The exact sine sum in \eqref{eq:g1-jump-fourier-exact} is potentially useful
in its own right: any unconditional exponential-sum improvement immediately
sharpens the cutoff below.  Even the trivial bound already removes the
factor $\mathcal V_\lambda/|B_\lambda|$ from the discontinuous arithmetic
part.

\begin{proposition}[Jump/continuous split for the weighted G1 transfer]\label{prop:g1-jump-split}
Let $R>3$, $K\ge2N$, $N\ge2L$, and let $x\in E_K$ be measured by
$\|x\|_{\mathcal H_{R,L}}$ as in
Proposition~\ref{prop:g1-weighted-sampler}.  Then
\begin{align}
 &\frac{
 |QW_\lambda((\mathsf V_N-I)p_\lambda,D_{\log}x)|
 }{|B_\lambda|}
 \notag\\
 &\qquad\le C_R\|x\|_{\mathcal H_{R,L}}
 \Bigg[
 \lambda^2\left\{
 \frac LN+
 \left(\frac LN\right)^{R-1}(1+\log(2N))
 \right\}
 \label{eq:g1-jump-split}\\
 &\hspace{43mm}
 +\frac{\lambda\mathcal W_\lambda}{|B_\lambda|}
 \left\{
 \left(\frac LN\right)^2+
 \left(\frac LN\right)^R(1+\log(2N))
 \right\}
 \Bigg].\notag
\end{align}
The first bracket is the entire contribution of the arithmetic function jumps
and has no small-boundary denominator.  The only surviving endpoint-normalized
regularity cost belongs to the continuous remainder and gains one additional
power of $N^{-1}$.
\end{proposition}
\begin{proof}
Write
$p_\lambda=\jmath_\lambda+c_\lambda$ and use that the
Vallee--Poussin multiplier is exactly one on $|n|\le N$.  For the jump part,
\eqref{eq:g1-jump-fourier-trivial} replaces the BV coefficient estimate used
in Proposition~\ref{prop:g1-weighted-sampler} by
\[
 |\widehat{(\mathsf V_N-I)\jmath_\lambda}(n)|
 \le C|B_\lambda|\frac{\lambda\sqrt L}{|n|},
 \qquad |n|>N.
\]
Repeating the off-diagonal convolution estimate there therefore replaces
$\lambda\mathcal V_\lambda/|B_\lambda|$ by $C\lambda^2$.  The diagonal
bound in Corollary~\ref{cor:g1-complete-matrix} is absorbed by the same
second term because $\lambda\ge2$.

For the continuous remainder use \eqref{eq:g1-continuous-fourier}.  With
\[
 a_m=\frac{|m|}{L}\left(1+\frac{|m|}{L}\right)^{-R},
\]
the off-diagonal part is bounded by
\[
 C_R\lambda L\mathcal W_\lambda
 \|x\|_{\mathcal H_{R,L}}
 \sum_{|n|>N}\frac1{n^2}
 \left[
 \frac L{|n|}+
 \left(\frac{|n|}{L}\right)^{1-R}
 (1+\log(2|n|))
 \right].
\]
Comparison with decreasing integrals gives
\[
 C_R\lambda\mathcal W_\lambda
 \|x\|_{\mathcal H_{R,L}}
 \left[
 \left(\frac LN\right)^2+
 \left(\frac LN\right)^R(1+\log(2N))
 \right].
\]
Using the diagonal estimate from
Corollary~\ref{cor:g1-complete-matrix} with the $n^{-2}$ coefficient bound
produces no larger term.  Division by $|B_\lambda|$ proves
\eqref{eq:g1-jump-split}.
\end{proof}

\begin{corollary}[The arithmetic jumps are resolvable at near-Slepian scale]\label{cor:g1-jump-natural-scale}
Fix $R>3$ and put
\begin{equation}\label{eq:g1-jump-cutoff}
 N_\lambda^{\rm jump}
 :=\left\lceil
 C_R\frac{\lambda^2L}{\mathfrak r(\lambda)}
 \right\rceil.
\end{equation}
For $C_R$ sufficiently large, the jump contribution in
\eqref{eq:g1-jump-split} is $O_R(\mathfrak r(\lambda))$ uniformly for every
ambient cutoff $K\ge2N_\lambda^{\rm jump}$.  Moreover
\begin{equation}\label{eq:g1-jump-near-slepian}
 \boxed{
 N_\lambda^{\rm jump}=\lambda^{2+o(1)}.}
\end{equation}
At the same cutoff, the continuous-remainder contribution is
\begin{equation}\label{eq:g1-continuous-at-jump-scale}
 O_R\!\left(
 \frac{\mathcal W_\lambda}{|B_\lambda|}
 \frac{\mathfrak r(\lambda)^2}{\lambda^3}
 \right)
\end{equation}
(up to a smaller $R$-dependent logarithmic term).  Thus the concrete
zero-independent condition
\begin{equation}\label{eq:g1-continuous-subgate}
 \boxed{
 \frac{\mathcal W_\lambda}{|B_\lambda|}
 \ll_R\frac{\lambda^3}{\mathfrak r(\lambda)} }
\end{equation}
would upgrade the filtered G1 transfer to the natural
$\lambda^{2+o(1)}$ Fourier scale.  Condition
\eqref{eq:g1-continuous-subgate} is \emph{not} proved here; it is a newly
isolated regularity subgate rather than an RH assumption.
\end{corollary}
\begin{proof}
The first term in \eqref{eq:g1-jump-split} is made
$O_R(\mathfrak r(\lambda))$ by \eqref{eq:g1-jump-cutoff}; the second
high-order term is smaller because $R>3$ and $\mathfrak r(\lambda)<1$ for
large $\lambda$.  Since
\[
 \log\frac1{\mathfrak r(\lambda)}=O(\sqrt L+\log L)=o(L)
\]
while $\lambda^2=e^L$, one has
$\log N_\lambda^{\rm jump}=L+o(L)$, proving
\eqref{eq:g1-jump-near-slepian}.  Finally
$L/N_\lambda^{\rm jump}\asymp
\mathfrak r(\lambda)/\lambda^2$; substitution in the continuous term of
\eqref{eq:g1-jump-split} gives
\eqref{eq:g1-continuous-at-jump-scale}.  Requiring it to be
$O_R(\mathfrak r(\lambda))$ gives
\eqref{eq:g1-continuous-subgate}.
\end{proof}


\subsection*{A polynomial cutoff for the literal source}
The small physical endpoint need not force an exponentially large Fourier
cutoff for this fixed two-mode source.  The additional ingredient is an
exterior norm estimate for a finite entire extension of the arithmetic
sum.  It controls the endpoint derivative jets without bounding the whole
interior variation by the endpoint amplitude.

\begin{lemma}[Polynomial approximation with a growing derivative order]
\label{lem:v114-entire-polynomial}
Fix a compact real interval $I$ of positive length. Suppose
\[
 F(z)=\int_{-\Omega}^{\Omega}e^{izt}\,d\mu(t),\qquad
 \Omega\le A_0c,\quad \norm\mu\le c^{A_1}.
\]
For each fixed $H>0$, there is $A>0$ such that a polynomial $T$ of degree
$D=\lceil Ac\rceil$ satisfies, simultaneously for $0\le j\le\lceil c\rceil$,
\begin{equation}\label{eq:v114-polynomial-error}
 \norm{(F-T)^{(j)}}_{L^\infty(I)}
       \le e^{-Hc}(Cc)^j.
\end{equation}
For all sufficiently large $c$, the constants are independent of $c$ and $j$. Moreover,
\[
 \norm T_\infty\le C(D+1)\norm T_2,\qquad
 \norm{T^{(j)}}_\infty\le(CD^2)^j\norm T_\infty
\]
on $I$.
\end{lemma}
\begin{proof}
Take the degree-$D$ Taylor polynomial about the midpoint of $I$ under the
finite measure. If $d$ is the half-length of $I$, its differentiated
remainder is at most
\[
 \norm\mu\,\Omega^j e^{\Omega d}
       \frac{(\Omega d)^{D+1-j}}{(D+1-j)!}.
\]
The inequality $k!\ge(k/e)^k$, with $A$ sufficiently large, proves
\eqref{eq:v114-polynomial-error} for the entire stated range of $j$.
The first polynomial inequality follows from normalized Legendre expansion,
$|P_k|\le1$, and $\sum_{k=0}^D(2k+1)=(D+1)^2$. The second is Markov's
inequality iterated on the derivatives, with the fixed interval scaling.
Real and imaginary parts give the same bounds for complex polynomials.
\end{proof}

\begin{lemma}[Endpoint jets of the finite arithmetic sum]
\label{lem:v114-endpoint-jets}
Put $c=2\pi\lambda^2$, $r=\lceil c\rceil$, and let $p_\lambda^\circ$
be the inversion-even proxy before the fixed value-at-zero repair.
If $J_j$ is the sum of absolute jumps of its $j$th ordinary logarithmic
derivative, including the jump at the identified periodic endpoint, then
\begin{equation}\label{eq:v114-jet-bound}
 J_j\le C|B_\lambda^+|\lambda^3(C\lambda^4)^j,
          \qquad 1\le j<r.
\end{equation}
Its zeroth-order periodic boundary jump is zero. On its smooth pieces,
for fixed constants $A,C,\rho>0$,
\begin{equation}\label{eq:v114-smooth-remainder}
 \norm{(p_\lambda^\circ)^{(r)}_{\rm ac}}_1
       \le CL\lambda^3 r!\rho^{-r}e^{Ac}.
\end{equation}
\end{lemma}
\begin{proof}
Write
\[
 g(x)=\sqrt\lambda h_\lambda^\circ(\lambda x)
     =\sum_{j=0,4}a_{j,\lambda}\psi_{j,c}(x),\qquad g(1)=B_\lambda^+.
\]
The exact finite Fourier eigenrelation used in
Lemma~\ref{lem:v12-endpoint-energy} represents $g$ by a measure supported
on $[-c,c]$ with total variation $O(\lambda)$; this is also the
bandlimited extension in \cite[\S30.15]{DLMFPSWF}.
Proposition~\ref{prop:endpoint-radial} and endpoint noncancellation give
$|g(x)|\le C|B_\lambda^+|$ for $x\ge1$.

Set $M=\lceil2\lambda^2\rceil$ and
$G(v)=\sum_{m=1}^M g(mv/\lambda^2)$. Its measure has bandwidth
$2\pi M=O(c)$ and total variation $O(\lambda^3)$. For $1/2\le v\le1$,
\[
 G(v)=\lambda v^{-1/2}f_\lambda^\circ(v/\lambda)
       +\sum_{\substack{m\le M\\mv\ge\lambda^2}}
                       g(mv/\lambda^2)
\]
almost everywhere. The last sum is $O(|B_\lambda^+|\lambda^2)$.
Corollary~\ref{cor:pswf-cross-tail}, which includes the unrepaired source,
and $dv/v=du/u$ therefore give
\[
 \norm G_{L^2(1/2,1)}\le C|B_\lambda^+|\lambda^2.
\]
Since $|B_\lambda^+|\asymp c^{11/4}e^{-c}$, apply
Lemma~\ref{lem:v114-entire-polynomial} with fixed sufficiently large $H$.
Its polynomial inequalities imply
\begin{align*}
 |G^{(j)}(1)|&\le C|B_\lambda^+|\lambda^4(C\lambda^4)^j,\\
 \sup_{1\le x\le3}|g^{(j)}(x)|
       &\le C|B_\lambda^+|(C\lambda^4)^j,
          \qquad 0\le j\le r.
\end{align*}
For the second bound approximate $g$ on $[1,3]$ and use its already
bounded sup norm, so the extra Legendre degree factor is unnecessary.

At the strict interior lower trace $v=1+$, subtract from $G$ precisely the
terms with $m\ge\lambda^2$, including equality when $\lambda^2$ is an
integer, and multiply by $\sqrt v/\lambda$. Their arguments lie in
$[1,3]$, their number is $O(\lambda^2)$, and differentiation costs at most
$3^j$. Thus the logarithmic trace jets are bounded by
$C|B_\lambda^+|\lambda^3(C\lambda^4)^j$.
To justify the conversion uniformly in $j$, use
$(v\partial_v)^j=\sum_k S(j,k)v^k\partial_v^k$ and
$S(j,k)\le\binom jk j^{j-k}$; hence
$\sum_k S(j,k)A^k\le(A+j)^j$. Here $j=O(\lambda^2)$ and
$A=C\lambda^4$. For the product factor, $|(\sqrt v)^{(k)}|_{v=1}\le k!$ and
$\sum_{k=0}^j\binom jk k!A^{j-k}\le2A^j$ when $A\ge2j$;
thus its cost is uniform in $j$ as well.
At the upper trace only the first summand survives, giving the stronger
bound without $\lambda^3$.

Each interior derivative jump is a source endpoint jet times
$m^{-1/2}$, with the reflected sign. Summing $m<\lambda^2$ costs
$O(\lambda)$; the preceding trace estimates control the periodic jump.
This proves \eqref{eq:v114-jet-bound}. Inversion symmetry gives the zero
function jump at the periodic boundary.
Finally extend each smooth piece using its fixed finite collection of
entire summands. On a complex logarithmic disk of fixed small radius
$\rho$, every active source argument has imaginary part bounded by a
fixed constant. The Fourier representation bounds the extension by
$C\lambda^3e^{Ac}$. Cauchy's estimate, integrated over total length $L$,
proves \eqref{eq:v114-smooth-remainder}. No asymptotic remainder was
differentiated, and the smooth compact repair is not used in this step.
\end{proof}

\begin{proposition}[Polynomial physical-endpoint recovery]
\label{prop:v114-polynomial-endpoint}
For the source of Proposition~\ref{prop:v12-source}, repaired as in
Corollary~\ref{cor:v12-repair}, set
\[
 N_\lambda=\left\lceil\lambda^8(1+L)\right\rceil,
       \qquad k_\lambda=P_{N_\lambda}p_\lambda.
\]
For all sufficiently large $\lambda$,
\begin{align}
 \frac{|\delta_{N_\lambda}(k_\lambda)-B_\lambda|}{|B_\lambda|}
       &\le\frac C\lambda,\label{eq:v114-polynomial-endpoint}\\
 |\widehat p_\lambda(n)|
       &\le C|B_\lambda|\frac{\lambda\sqrt L}{|n|},
            \qquad |n|>N_\lambda.\label{eq:v114-high-tail}
\end{align}
This is an asymptotic estimate with no certified numerical threshold or
optimized exponent. The vector is the literal ordinary Fourier projection;
no endpoint shell is added.
\end{proposition}
\begin{proof}
First omit the compact repair. The periodic proxy is BV and continuous
at the identified endpoint, so the Dirichlet--Jordan theorem identifies
its symmetric Fourier sum there with its common one-sided trace. We may
therefore estimate the actual endpoint error by its symmetric Fourier tail.
In the jump train of
Lemma~\ref{lem:g1-jump-train}, the endpoint tail contributed by a threshold
$m$ is, up to sign,
\[
 \frac{B_\lambda^+}{\pi\sqrt m}
       \sum_{n>N}\frac{\sin(2\pi n\log m/L)}n.
\]
Dirichlet summation and the bounded sine-series partial sums give
\[
 \left|\sum_{n>N}\frac{\sin(n\theta)}n\right|
       \le C\min\left(1,\frac1{N|\sin(\theta/2)|}\right).
\]
Separate the largest integer $m_0<\lambda^2$: its contribution is
$O(|B_\lambda^+|/\lambda)$ even when $\lambda^2$ approaches an integer.
For every remaining $m$, $\lambda^2-m\ge1$. Splitting at $m=\lambda$
and $m=\lambda^2/2$ gives
\[
 \sum_{\substack{1<m<\lambda^2\\m\ne m_0}}
 \frac{m^{-1/2}}{|\sin(\pi\log m/L)|}
       \le C\lambda L(1+\log\lambda).
\]
Near the upper end use
$\log(\lambda^2/m)\ge(\lambda^2-m)/\lambda^2$ and sum the harmonic
denominator. Thus the entire function-jump endpoint error divided by
$|B_\lambda^+|$ is at most
\begin{equation}\label{eq:v114-resonance-bound}
 C\left\{\lambda^{-1}+\frac{\lambda L(1+\log\lambda)}N\right\}.
\end{equation}
This retains, rather than assumes away, near-threshold arithmetic resonance.

Integrate by parts $r=\lceil c\rceil$ times on the smooth pieces, retaining
the jumps of every derivative and the periodic derivative jumps. With
$t_n=2\pi n/L$, the $j$th jump term in a normalized Fourier coefficient
has magnitude at most $J_j/(\sqrt L\,|t_n|^{j+1})$. Its physical-endpoint
tail is at most
\[
 \frac{2J_j}{L}\sum_{n>N}|t_n|^{-j-1}
       \le \frac{C J_j}{j}\left(\frac L{2\pi N}\right)^j.
\]
By Lemma~\ref{lem:v114-endpoint-jets}, summing $1\le j<r$ gives
$C|B_\lambda^+|\lambda^7L/N$ whenever $C\lambda^4L/N<1/2$.
The final smooth-piece remainder, after endpoint summation and division by
$|B_\lambda^+|$, is bounded by
\begin{equation}\label{eq:v114-analytic-tail}
 C\lambda^A(1+L)^Ae^{Ac}r!
                 \left(\frac{CL}{N}\right)^{r-1}.
\end{equation}
Here $A,C$ are fixed and independent of $r,\lambda$.
At $N=N_\lambda$, $r!\le(C\lambda^2)^r$, so the logarithm of this bound
is at most $-6c\log\lambda+O(c+\log\lambda+\log L)$.
It is smaller than every fixed inverse power of $\lambda$.
Combining this with \eqref{eq:v114-resonance-bound} proves the first
claim for $p_\lambda^\circ$.

The same expansion gives, at every $|n|>N_\lambda$,
\[
 |\widehat p_\lambda^\circ(n)|
 \le \frac{C|B_\lambda^+|}{\sqrt L}
       \left(\frac\lambda{|t_n|}+\frac{\lambda^7}{|t_n|^2}\right)
       +|R_n|.
\]
The second term and remainder are absorbed by the first. For the remainder,
its ratio to $|B_\lambda^+|\lambda/(\sqrt L|t_n|)$ decreases with $|n|$
and is already smaller than every fixed inverse power at $N_\lambda$.
This proves \eqref{eq:v114-high-tail} before repair.

Finally let $e=p_\lambda-p_\lambda^\circ$ and
$a=h_\lambda^\circ(0)$. Direct counting gives
$\norm e_\infty\le C\lambda^{3/2}|a|$.
The Fourier Lebesgue constant is $O(\log(N+2))$, so
\[
 \frac{|(P_N-I)e(0)|}{|B_\lambda|}
 \le C\lambda^{3/2}\log(N+2)\frac{|a|}{|B_\lambda|}
       =O_A(\lambda^{-A})
\]
along this polynomial diagonal, by \eqref{eq:v12-zero-value}.
For the coefficient assertion a direct variation estimate is stronger:
if the fixed repair bump is supported in $[-b,b]$, then
\[
 \int_{1/\lambda}^{\lambda}|D_{\log}\mathcal E(\psi)(u)|\frac{du}{u}
 \le \sum_{m\le b\lambda}m^{-1/2}
   \int_0^b x^{-1/2}\bigl(\tfrac12|\psi(x)|+x|\psi'(x)|\bigr)\,dx
 \le C\sqrt\lambda.
\]
Inversion averaging preserves this bound and has no periodic function
jump. Hence $|\widehat e(n)|\le C|a|\sqrt{\lambda L}/|n|$, which is
negligible. Use $B_\lambda^+\asymp B_\lambda$ to finish.
\end{proof}

\begin{corollary}[Weak G1 on the same polynomial-cutoff vector]
\label{cor:v114-polynomial-g1}
Let $R>3$, $N=N_\lambda$, $K\ge N$, and $x\in E_K$, measured by
the weighted norm of Proposition~\ref{prop:g1-weighted-sampler}.
For all sufficiently large $\lambda$,
\begin{align}
 \frac{|\ip{D_{\log}W_{\lambda,K}k_\lambda}{x}|}
      {|\delta_K(k_\lambda)|}
 &\le C_R\norm x_{\mathcal H_{R,L}}
 \left[\mathfrak r(\lambda)+\lambda^2
 \left\{\frac LN+\left(\frac LN\right)^{R-1}
                   (1+\log(2N))\right\}\right].
 \label{eq:v114-polynomial-g1}
\end{align}
The bound is uniform in $K$ and tends to zero on bounded weighted test
families. This is a fixed-source weak pairing statement, not control
in operator norm on a growing source block.
\end{corollary}
\begin{proof}
The coefficients of $(P_N-I)p_\lambda$ vanish for $|n|\le N$.
Insert \eqref{eq:v114-high-tail} into the absolute convolution proof
of Proposition~\ref{prop:g1-weighted-sampler}: on the only frequencies
used there, $\mathcal V_\lambda$ is replaced by $C|B_\lambda|\lambda$.
Both its off-diagonal estimate and its separately treated logarithmic
diagonal therefore give the displayed $\lambda^2$ error. The sums
converge absolutely for $R>3$; the fixed-window domain identification
in Proposition~\ref{prop:v114-weil-core} also justifies the passage
from finite Fourier sums to $p_\lambda$.
The continuum term is Proposition~\ref{prop:g1-sobolev}. Finally
$|\delta_K(k_\lambda)|\ge(1-C/\lambda)|B_\lambda|$ and, with the
first-slot-linear convention,
$\ip{D_{\log}W_{\lambda,K}k_\lambda}{x}
=QW_\lambda(k_\lambda,D_{\log}x)$.
Unlike the filtered vector, $P_Np_\lambda$ lies in $E_N$, so $K\ge N$
suffices.
\end{proof}

These results assert no positive Weil complement along this diagonal,
a small inverse-weighted spectral residual, or control of the corrected
sampler's graph defect. Those obligations remain separate. The fixed
orders $0,4$ are essential; no uniform claim over a growing prolate block
is made.

\subsection*{Use on the centered and corrected sampler}
The weighted estimate is unchanged by the factors $(-1)^n$, because its
proof bounds coefficient magnitudes. For any fixed finite root family,
sufficient Hardy smoothing therefore controls the centered sampler and
any fixed finite number of its derivatives. The source input is now
proved in Section~\ref{sec:g1}; constants may depend on the test family.
It supplies no uniformity for a derivative order growing with $L$.

For the later even endpoint shell, the zeroth residual pairing is preserved
exactly by parity: an even candidate $k$ has odd $D_{\log}Wk$ and hence
zero pairing with that shell. Higher residual jets do not all have the
same parity, so no automatic preservation is claimed for every jet.
A fixed finite Gramian made from pairings that tend to zero also tends to
zero, by summing their squared magnitudes. Such a Gramian cannot by itself
supply a uniformly positive metric on that root space.

\begin{lemma}[A legitimate common-cutoff selection]\label{lem:diagonal-selection}
Suppose a normalized G1 error is at most $\epsilon$ for every
$N\ge N_0(L,\epsilon)$ and every ambient cutoff $K\ge2N$.
Suppose finitely many additional errors are at most $\epsilon$ for every
$K\ge K_0(L,N,\epsilon)$. Then one can choose a common cutoff satisfying
all these bounds. The same statement holds with nested eventual thresholds
$N\ge N_0$, $K\ge K_0(N)$ when that is what has actually been proved.
\end{lemma}
\begin{proof}
Choose an admissible core cutoff first, then choose the ambient cutoff above
all the finitely many thresholds. Apply this with a positive tolerance
sequence tending to zero. Each error must be normalized by the actual
endpoint or scale used in the target expression before taking thresholds.
Merely knowing that two sets of acceptable cutoffs are unbounded would not
suffice: two unbounded sets can be disjoint.
\end{proof}
The explicit corrected-sampler rule below can be enlarged to meet additional
proved eventual thresholds: its error bounds still vanish when $K$ grows
beyond that rule. A formula can be independent of the zero list while a
proof of its convergence on a fixed hypothetical window has window-dependent
constants. Neither observation proves a missing graph-pairing estimate.
\section{A consistent finite sampling problem}\label{sec:sampling}

The immediate question is whether a finite sampler can simultaneously retain
the Cauchy evaluation, use the actual CCM endpoint functional, and converge
in the same finite Weil form. The graph identity required by the wider
G1/G2 argument is a fourth property. It is kept separate here because it
does not follow from the other three.

Fix $L\ge2$, $\lambda=e^{L/2}$, $A=L/2$, and
\[
 t_n=2\pi n/L,\qquad s_n=\tfrac12+it_n.
\]
In the coordinate $x=\log(\lambda u)\in[0,L]$, use
\[
 V_n(x)=L^{-1/2}e^{it_nx},\qquad D_LV_n=t_nV_n.
\]
Interval vectors are extended by zero outside the interval when evaluating the
global Weil form. Write $W_{L,N}$ for its compression to $|n|\le N$.
This is the same compression previously denoted $W_{\lambda,N}$, with $\lambda=e^{L/2}$; $\Q=QW$ denotes the same global Weil form in the fixed first-slot-linear convention. The coordinate change is unitary. The Hermitian product is linear in its first slot. The physical endpoint is
\begin{equation}\label{v1:eq:delta}
 \delta_{L,N}\Big(\sum a_nV_n\Big)=L^{-1/2}\sum a_n,
 \qquad \eta_{L,N}=\sqrt L\,\delta_{L,N}.
\end{equation}
This is the equal-coefficient endpoint of the original matrix, not an
interior evaluation. The basis and matrix realization are those of
\cite{ConnesConsaniMoscovici2025}; its opposite sesquilinear convention is translated here.

To remove the Fourier-sign ambiguity, define
\begin{equation}\label{v1:eq:transform}
 H_F(s)=\int_{\mathbb R}g_F(y)e^{-(s-1/2)y}\,\dd y,
 \qquad
 g_F(y)=\frac1{2\pi}\int_{\mathbb R}H_F(\tfrac12+it)e^{ity}\,\dd t.
\end{equation}
Thus the critical-line transform has the usual negative Fourier sign.
Reversing this sign reverses physical tails and Fourier indices, but does
not remove the half-period character identified below.

For the manuscript's Hardy output put
\begin{equation}\label{v1:eq:hardy}
 \begin{split}
 H_F(s)&=\kappa(s)F(s),\\
 \kappa(s)&=\frac{\sigma(\sigma+a)^M}{\sigma-1}
       \frac{s-1}{(s+a)^M(\sigma-s)},
 \qquad \E(F)=\sigma F(\sigma).
 \end{split}
\end{equation}
Here $F$ ranges over a fixed finite-dimensional space of root functions
$\zeta(s)/(s-\rho)^j$, with the denominator canceled by a genuine zeta
zero of order at least $j$. For the one-sided application all these zeros
have real part greater than $1/2$. This is a conditional model: their
existence is not asserted. The additional output $H=\kappa\zeta$ will
provide an unconditional numerical test of the sampling mechanism.

Choosing $M$ sufficiently large gives uniform polynomial vertical decay
of any fixed order needed below on the relevant bounded strip. Contour
shifts, after cancellation at $s=1$ and at the selected zeros, imply the
following useful facts. They also follow directly by Fourier inversion:
\begin{equation}\label{v1:eq:tails}
 g_F(y)=\E(F)e^{qy}\quad(y<0),\qquad q=\sigma-\tfrac12,
 \qquad |g_F^{(j)}(y)|\le C e^{-ry}\quad(y>0),\quad j\le2,
\end{equation}
where any fixed $r<a+1/2$ can be used after sufficient smoothing.
Constants are uniform on a finite-dimensional unit sphere. In particular,
$g_F(0)=\E(F)$. To verify the negative tail, close the inverse contour to
the right for $y<0$. The only pole crossed is $\sigma$, whose residue is
$-\sigma F(\sigma)$; the contour orientation gives \eqref{v1:eq:tails}.
For the positive tail shift left to $\Re s=1/2-r>-a$.
Using a rate strictly below $a+1/2$ avoids hiding the polynomial associated
with the order-$M$ pole at $-a$.

Both transfer and the uncentered nonconvergence theorem use $q,r>1/2$.
These margins are available for every fixed $\sigma>1$ and $a>0$, after
choosing $r<a+1/2$ and taking $M$ sufficiently large. No conclusion below
requires moving these parameters with $L$.

In the convention \eqref{v1:eq:transform}, the first-slot-linear zero formula is
\begin{equation}\label{v1:eq:zeros}
 \Q(f,g)=\sum_{\rho\in\Z}
 H_f(\rho)\overline{H_g(\rho^\sharp)},
 \qquad \rho^\sharp=1-\overline\rho,
\end{equation}
with nontrivial zeros counted with multiplicity. The symmetries of the
zero set make this consistent with the Fourier form of the explicit
formula. All sums used below are absolutely convergent. The standard
bound $N(T)=O(T\log T)$ implies
$\sum_{\rho}(1+|\Im\rho|)^{-2}<\infty$.

\section{The two samplers are different finite vectors}

Let $P_Lg(y)=\sum_{k\in\mathbb Z}g(y+kL)$. Its Fourier series is
$L^{-1}\sum_n H(s_n)e^{it_ny}$. Define
\begin{align}
 \Ju_{L,N}F&=L^{-1/2}\sum_{|n|\le N}H_F(s_n)V_n,
       \label{v1:eq:raw}\\
 \Jc_{L,N}F&=L^{-1/2}\sum_{|n|\le N}(-1)^nH_F(s_n)V_n.
       \label{v1:eq:centered}
\end{align}
The superscripts mean uncentered and centered, respectively.
For fixed $L$, their full Fourier limits on $[0,L]$ are
\[
 u_{L,F}(x)=P_Lg_F(x),\qquad
 c_{L,F}(x)=P_Lg_F(x-A).
\]
Translation of $c_{L,F}$ back to $[-A,A]$ gives the centered fundamental
representative. Translation invariance of the global form therefore
applies to $c_{L,F}$; it does not identify $u_{L,F}$ with that representative.

\begin{proposition}[Endpoint and interior evaluations]\label{v1:prop:endpoints}
With Fourier truncations taken sufficiently far,
\[
 \delta(\Ju F)\longrightarrow \E(F),\qquad
 \delta(\Jc F)\longrightarrow0\quad(L\to\infty).
\]
On the centered sampler the evaluation at $x=A$, rather than at the cut,
converges to $\E(F)$.
\end{proposition}
\begin{proof}
For fixed $L$, Fourier convergence gives
$\delta(\Ju F)\to P_Lg_F(0)$ and
$\delta(\Jc F)\to P_Lg_F(A)$.
Equation \eqref{v1:eq:tails} gives
$P_Lg_F(0)=\E(F)+O(e^{-\min(q,r)L})$ and
$P_Lg_F(A)=O(e^{-\min(q,r)L/2})$.
Evaluation of \eqref{v1:eq:centered} at $A$ cancels $(-1)^n$.
\end{proof}

The diagonal unitary $P=\operatorname{diag}((-1)^n)$ commutes with $D_L$.
A pure change of basis must replace $W$ by $PWP$ and the boundary covector
by its transformed covector. Applying $P$ to the samples while retaining
the original $W$ and $\delta$ is instead a different physical vector.
That is precisely the change used in \eqref{v1:eq:centered}.

\section{Centered Weil transfer with the boundary terms retained}

\begin{lemma}[A sufficient strip estimate]\label{v1:lem:strip}
Suppose $|g^{(j)}(y)|\le C e^{-\mu|y|}$ for $j\le2$, with $\mu>1/2$.
Let $h_L=1_{[-A,A]}P_Lg$ and let $H_L$ be its transform. Then, uniformly
for $0\le\Re s\le1$,
\begin{equation}\label{v1:eq:stripbound}
 |H_L(s)-H(s)|\le
 \frac{C e^{-(\mu-1/2)L/2}}{1+|\Im s|}.
\end{equation}
The estimate includes the jumps of the zero-extended representative.
\end{lemma}
\begin{proof}
Set $e_L=h_L-g$ and $d=\Re s-1/2$.
Outside $[-A,A]$ use the tails of $g$; inside use the nonzero periodic
images. Integrating the resulting geometric series, including their first
derivatives, gives
\[
 \norm{e_L e^{-d\,\cdot}}_1+
 \operatorname{Var}(e_L e^{-d\,\cdot})
 \le C e^{-(\mu-1/2)A},\qquad |d|\le1/2.
\]
The variation includes both endpoint jumps, whose weighted sizes obey
the same bound. The $L^1$ bound handles bounded ordinates. One integration
by parts in the sense of a finite derivative measure handles large
ordinates, giving \eqref{v1:eq:stripbound}.
\end{proof}

Arbitrary inverse powers of the ordinate are not claimed. A nonzero
endpoint jump normally produces a $1/t$ term. That decay is sufficient:
each Weil summand is a product of two transforms.

\begin{theorem}[Centered form transfer]\label{v1:thm:center}
On any fixed finite-dimensional family satisfying the preceding decay
bounds and $|H_F(s)|\le C(1+|\Im s|)^{-1}$ on the strip,
\[
 \Q(c_{L,F},c_{L,G})\longrightarrow\Q(g_F,g_G)
\]
uniformly on bounded sets of $F,G$. The same conclusion holds for
$\ip{W_{L,N}\Jc_{L,N}F}{\Jc_{L,N}G}$ with a sufficiently large common
Fourier cutoff. On a one-sided root space the self-form limit is zero.
For reflected top-root outputs, the nonzero global cross-form also transfers.
\end{theorem}
\begin{proof}
Let $\epsilon_L=e^{-(\mu-1/2)L/2}$. In \eqref{v1:eq:zeros}, replace
each transform by $H_F+E_{L,F}$. Lemma~\ref{v1:lem:strip} bounds each difference
by $C\epsilon_L/(1+|\Im\rho|)$, and the unchanged transform has the same
ordinate decay without $\epsilon_L$. Hence the difference of the forms is
at most
$C(\epsilon_L+\epsilon_L^2)\sum_\rho(1+|\Im\rho|)^{-2}$.
Translation of $h_L$ to $[0,L]$ does not change the paired form.

For completeness, Fourier truncation can be justified directly, without
inferring convergence of these particular projections from abstract core
density. The correct CCM matrix bounds are
\begin{equation}\label{v1:eq:matrix}
 |W_{nm}|\le C\lambda/|n-m|\ (n\ne m),\qquad
 |W_{nn}|\le C\{\lambda+1+\log(2+|n|/L)\}.
\end{equation}
The diagonal uses $2(1-y/L)\cos(t_ny)$, not a limit of the off-diagonal
kernel. Near zero its archimedean numerator is bounded by
$C\{\min(1,t_n^2y^2)+y\}$; integration against $O(1/y)$ gives the
logarithm in \eqref{v1:eq:matrix}. The prime and pole pieces are $O(\lambda)$.

If $|H_F(s_n)|\le C(1+|n|/L)^{-R}$ with $R\ge2$, then the normalized
coefficients $a_n$ have
\[
 \sum_n|a_n|\le C\sqrt L,\qquad
 \sum_{|n|>N}|a_n|\le C\sqrt L(1+N/L)^{1-R}.
\]
The absolute double matrix series therefore converges. Bounding off-diagonal
entries by $C\lambda$ and treating the logarithmic diagonal separately gives,
for $N\ge L$,
\begin{equation}\label{v1:eq:galerkin}
 |\Q(c_{L,F},c_{L,G})-
 \ip{W_{L,N}\Jc_{L,N}F}{\Jc_{L,N}G}|
 \le C\lambda L(1+N/L)^{1-R}.
\end{equation}
The constants are uniform on the fixed family. This also identifies the
matrix series with the semilocal form, using its closed realization.

For a one-sided root space, $H_F$ is nonzero at at most finitely many zeros,
all on one side, so each paired summand in its self-form is zero. If $F$ and
$G$ are the top root functions at $\rho$ and $\rho^\sharp$, their cross-form is
\[
 m_\rho\,\kappa(\rho)\frac{\zeta^{(m_\rho)}(\rho)}{m_\rho!}
 \overline{\kappa(\rho^\sharp)
 \frac{\zeta^{(m_\rho)}(\rho^\sharp)}{m_\rho!}},
\]
which is nonzero. This is the first-slot-linear convention.
\end{proof}

\section{The uncentered sampler has an explicit cut defect}

The failure of the original proof is more than an omitted phase in its
notation. It is possible to compute the form created by the different cut.
In this section assume $q,r>1/2$, and abbreviate $E=\E(F)$,
$H=H_F$, $z_\rho=\rho-1/2$. Let $g_+=1_{(0,\infty)}g$ and
$g_-=1_{(-\infty,0)}g$. Set
\[
 f_L=g_++g_-(\,\cdot-L).
\]
The negative half has been moved by $L$. Since
$H_-(s)=E/(\sigma-s)$, its full transform is
\begin{equation}\label{v1:eq:cuttransform}
 \widehat f_L(s)=H(s)+
       E\frac{e^{-(s-1/2)L}-1}{\sigma-s}.
\end{equation}

\begin{lemma}[Comparison with the actual periodic cut]\label{v1:lem:uncut}
For $u_L=1_{[0,L]}P_Lg$ and some $\nu>0$,
\[
 \Q(u_L,u_L)=\Q(f_L,f_L)+O(e^{-\nu L}).
\]
One may take $\nu=\min(q,r)-1/2>0$.
\end{lemma}
\begin{proof}
Write $s=1/2+d+it$, $|d|\le1/2$, and $p_0=\min(q,r)$.
On $[0,L]$ the negative images sum exactly to
$E e^{q(x-L)}/(1-e^{-qL})$. Their transform differs from the transform
of $g_-(\,\cdot-L)$ by exactly
\[
 \frac{E e^{-qL}}{1-e^{-qL}}
       \frac{e^{-(s-1/2)L}-1}{\sigma-s}.
\]
The extra positive images are $O(e^{-rL}e^{-rx})$ on $[0,L]$; the removed
positive tail lies in $x>L$. Their weighted $L^1$ norms and variations,
including endpoint jumps, give the same type of bound. Consequently
\[
 |\widehat u_L(s)-\widehat f_L(s)|
 \le \frac{C e^{-p_0L}(1+e^{-dL})}{1+|t|},\qquad
 |\widehat f_L(s)|\le \frac{C(1+e^{-dL})}{1+|t|}.
\]
It matters that the partner has real displacement $-d$. Products in the
form error contain
$(1+e^{-dL})(1+e^{dL})\le4e^{L/2}$, rather than two independent
worst-case factors. Expanding the zero sum therefore bounds the error by
$C e^{-(p_0-1/2)L}\sum_\rho(1+|\Im\rho|)^{-2}$, as required.
\end{proof}

For the one-sided root family, let $S$ be the finite set of zeros at which
$H$ does not vanish. Sums over $S$ below include their multiplicities.
For $H=\kappa\zeta$, take $S$ empty.

\begin{proposition}[Cut-defect formula]\label{v1:prop:cut}
Under the stated hypotheses,
\begin{align}
 \Q(u_L,u_L)={}&\Q(g,g)
 +2\Re\sum_{\rho\in S}
 \frac{H(\rho)\overline E}{q+z_\rho}
          (e^{z_\rho L}-1)\notag\\
 &+2|E|^2\sum_{\rho\in\Z}
       \frac{1-e^{z_\rho L}}{q^2-z_\rho^2}
 +O(e^{-\nu L}).\label{v1:eq:cutdefect}
\end{align}
Both displayed sums are well defined without RH; the second is absolutely
convergent for each fixed $L$.
\end{proposition}
\begin{proof}
Insert \eqref{v1:eq:cuttransform} into \eqref{v1:eq:zeros}. The two mixed sums
are complex conjugates after replacing $\rho$ by $\rho^\sharp$.
The product of the two added terms has numerator
$2-e^{z_\rho L}-e^{-z_\rho L}$ and denominator $q^2-z_\rho^2$.
The symmetry $\rho\mapsto1-\rho$ makes the two exponential sums equal,
giving \eqref{v1:eq:cutdefect}. Lemma~\ref{v1:lem:uncut} handles the actual cut.
Absolute convergence follows from $|\Re z_\rho|<1/2$ and the quadratic
denominator at large ordinates.
\end{proof}

\begin{theorem}[Failure of the claimed full-scale raw transfer]\label{v1:thm:rawfailure}
Suppose $E\ne0$, and $H$ vanishes at every critical-line zero and at all
but finitely many other nontrivial zeros. Under the preceding decay
hypotheses, $\Q(u_L,u_L)$ does not converge to any finite limit as real
$L\to\infty$. In particular it cannot converge to $\Q(g,g)=0$ for a
one-sided root output with nonzero $E$.
\end{theorem}
\begin{proof}
Choose a critical-line zero $\rho_0=1/2+i\gamma_0$ of multiplicity $m_0$,
with $\gamma_0\ne0$. Such zeros exist unconditionally by Hardy's theorem
\cite{Hardy}. The coefficient of $e^{i\gamma_0L}$ in the infinite sum
in \eqref{v1:eq:cutdefect} is
\begin{equation}\label{v1:eq:residue}
 -\frac{2|E|^2m_0}{q^2+\gamma_0^2}\ne0.
\end{equation}
The finite mixed terms have exponents strictly off the imaginary axis,
so do not contribute at this exponent.

Here is an argument that does not assume away possible off-line zeros.
Take the Laplace transform in $L\ge L_0$ of \eqref{v1:eq:cutdefect}, initially
for $\Re p>1/2$. Its series of rational-exponential terms continues
meromorphically across the imaginary axis: the denominators at large
ordinates are $O(|\Im\rho|^3)$, giving locally normal convergence away
from the discrete poles. The exponentially decaying error is analytic
in $\Re p>-\nu$. At $p=i\gamma_0$ the continued transform has residue
\eqref{v1:eq:residue}; the lower integration limit does not change the residue.

If $\Q(u_L,u_L)$ had a finite limit, it would be bounded and its Laplace
transform would be holomorphic for $\Re p>0$. All apparent poles in that
half-plane would therefore have to cancel. Moreover bounded convergence,
after subtracting the finite limit, would give
\[
 \lim_{\epsilon\downarrow0}\epsilon
 \int_{L_0}^\infty e^{-(\epsilon+i\gamma_0)L}
           \Q(u_L,u_L)\,\dd L=0.
\]
For the constant part the same limit is zero since $\gamma_0\ne0$.
But the meromorphic expression gives the nonzero residue
\eqref{v1:eq:residue}. This is a contradiction.
\end{proof}

\begin{remark}[Exact scope]
The theorem concerns the full real-scale limit. It does not exclude
specially chosen subsequences $L_j$ on which the form happens to approach
zero. Nor does it exclude different sampling weights or a changed
matrix. It applies to Galerkin approximations whenever their error relative
to this continuum cut tends to zero; choosing a larger faithful Fourier
cutoff cannot remove the continuum defect.

For any nonempty hypothetical one-sided root window, a basis vector has
$E=\sigma\zeta(\sigma)/(\sigma-\rho)^j\ne0$. Thus the original full-scale
transfer and the original unphased sampling convention are incompatible
in that window. This is a conditional diagnosis, not an assertion that an
off-line zero exists or an unconditional disproof of an implication with
a potentially empty premise. The output $H=\kappa\zeta$ is an unconditional
example of the same failure in the larger analytic output class.

On $\ker\E$, the negative half in \eqref{v1:eq:tails} vanishes, and this cut
obstruction disappears. That codimension-one restriction does not recover
the whole nonempty root window.
\end{remark}

\section{A single finite object with form transfer and exact endpoint}

For an integer $K\ge1$ use the ambient space $|n|\le2K$, and define the
even endpoint shell
\begin{equation}\label{v1:eq:shell}
 c_{L,K}=\frac{\sqrt L}{2K}\sum_{K<|n|\le2K}V_n.
\end{equation}
There are exactly $2K$ terms. Consequently
\begin{equation}\label{v1:eq:shellnorms}
 \delta(c_{L,K})=1,\qquad
 \norm{c_{L,K}}^2=\frac L{2K},\qquad
 \norm{D_Lc_{L,K}}^2=\frac{\pi^2}{3L}(14K+9+K^{-1}).
\end{equation}
Define a linear functional and a corrected sampler by
\begin{equation}\label{v1:eq:repair}
 d_{L,K}(F)=\E(F)-\delta(\Jc_{L,2K}F),\qquad
 \Jr_{L,K}F=\Jc_{L,2K}F+c_{L,K}d_{L,K}(F).
\end{equation}

\begin{theorem}[Three-property compatibility]\label{v1:thm:repair}
On a fixed finite-dimensional Hardy-output family as above, choose
\[
 K(L)=\lceil\lambda^4L^4\rceil,\qquad \lambda=e^{L/2},
\]
and choose $M$ sufficiently large to give $R\ge2$ in
\eqref{v1:eq:galerkin}. The same finite matrix, sampler, and endpoint obey
\begin{align}
 \delta_{L,2K}(\Jr_{L,K}F)&=\sigma F(\sigma)
                       &&\text{exactly},\label{v1:eq:exactendpoint}\\
 \norm{\Jr_{L,K}F-\Jc_{L,2K}F}&\longrightarrow0,\label{v1:eq:normrepair}\\
 \ip{W_{L,2K}\Jr_{L,K}F}{\Jr_{L,K}G}
            &\longrightarrow\Q(g_F,g_G),\label{v1:eq:formrepair}\\
 \ip{\Jr_{L,K}F}{\Jr_{L,K}G}
            &\longrightarrow\ip{g_F}{g_G}.\label{v1:eq:hilbertrepair}
\end{align}
The limits are uniform on bounded subsets of the fixed family. The
construction requires no zero list in the matrix, shell, cutoff rule, or
formula for the endpoint correction. It is not an economical cutoff claim.
\end{theorem}
\begin{proof}
The endpoint identity follows immediately from \eqref{v1:eq:shellnorms}.
Uniform decay of the critical-line samples gives
$\norm{\Jc F}\le C\norm F$, and the corresponding Riemann sum converges
to $\norm{g_F}^2$. The same coefficient bound and the endpoint result give
$\norm{d_{L,K}}\le C$ for large $L$ and the stated cutoffs. Therefore
$\norm{\Jr F-\Jc F}\le C\sqrt{L/K}\norm F$.

The Schur row-sum estimate applied to the correct matrix bounds yields
\begin{equation}\label{v1:eq:opnorm}
 \norm{W_{L,2K}}\le C\lambda(1+\log(4K+1)).
\end{equation}
Expansion of the two corrected arguments bounds the change in their
form pairing by
\begin{equation}\label{v1:eq:formprice}
 C\lambda(1+\log(4K+1))
       \left(\sqrt{L/K}+L/K\right)\norm F\norm G.
\end{equation}
For the displayed cutoff this is $O(\lambda^{-1}L^{-1/2})$.
The centered Fourier-truncation error in \eqref{v1:eq:galerkin} is at most
$O(\lambda^{-3}L^{-2})$ when $R=2$, and smaller for larger $R$.
Theorem~\ref{v1:thm:center} proves \eqref{v1:eq:formrepair}.
The ordinary norm correction tends to zero, so the Riemann-sum limit also
proves \eqref{v1:eq:hilbertrepair}.
\end{proof}

If $F\mapsto g_F$ is injective on the fixed finite-dimensional space, the
last limit gives a uniform lower ordinary-norm bound there. In particular,
on a one-sided root space the corrected self-form tends to zero, and adding
any fixed positive scalar to $W$ gives restricted coercivity for sufficiently
large cutoffs. The nonzero reflected cross-form also transfers to this
corrected sampler. These are legitimate consequences of this particular
construction; they do not identify its graph transport with the old one.

\section{The graph identity is the remaining compatibility cost}

Suppose the root-space operator satisfies
\[
 BF=(s-\tfrac12)F-\eta_{\mathrm B}(F)\zeta.
\]
The centered samples satisfy the exact coefficientwise identity
\begin{equation}\label{v1:eq:centergraph}
 iD_L\Jc-\Jc B=z^{\mathrm c}_{L,2K}\eta_{\mathrm B},\qquad
 z^{\mathrm c}_{L,2K}=L^{-1/2}\sum_{|n|\le2K}
        (-1)^n\kappa(s_n)\zeta(s_n)V_n.
\end{equation}
For the corrected sampler, however,
\begin{equation}\label{v1:eq:graphdefect}
 iD_L\Jr-\Jr B=z^{\mathrm c}_{L,2K}\eta_{\mathrm B}+\mathcal R_{L,K},
 \qquad
 \mathcal R_{L,K}F=iD_Lc_{L,K}\,d_{L,K}(F)
                         -c_{L,K}\,d_{L,K}(BF).
\end{equation}
This is an exact rank-at-most-two defect. Parity gives
$\ip{D_Lc_{L,K}}{c_{L,K}}=0$, so
\begin{equation}\label{v1:eq:graphnorm}
 \norm{\mathcal R_{L,K}F}^2
 =|d_{L,K}(F)|^2\frac{\pi^2(14K+9+K^{-1})}{3L}
       +|d_{L,K}(BF)|^2\frac L{2K}.
\end{equation}
For $\E(F)\ne0$, this diverges along the displayed construction.
Smallness of a correction in ordinary norm or Weil form therefore does
not make it small in the graph identity.

\begin{proposition}[Rigidity and a general trace obstruction]\label{v1:prop:graph}
For a fixed one-sided off-line root space, the equation
$iD_LJ-JB=z\eta_{\mathrm B}$ has at most one solution $J$ for each prescribed $z$.
Moreover, any family of corrections $T_L$ with $\norm{T_L}\to0$ and
$\delta(T_LF)\to\E(F)\ne0$ for some fixed $F$ has
$\norm{iD_LT_LF-T_LBF}\to\infty$.
\end{proposition}
\begin{proof}
All eigenvalues of $B$ have positive real part, whereas those of $iD_L$
are purely imaginary. Row by row, the homogeneous equation is
$\ell_n(it_nI-B)=0$, and $it_nI-B$ is invertible. This proves uniqueness.

For a periodic $H^1$ function $b$ on $[0,L]$, choose a point where
$|b|^2$ does not exceed its average and integrate the derivative of
$|b|^2$. Cauchy--Schwarz gives
\begin{equation}\label{v1:eq:trace}
 |b(0)|^2\le L^{-1}\norm b^2+2\norm b\,\norm{D_Lb}.
\end{equation}
Apply this to $b=T_LF$. Its norm tends to zero while its endpoint tends
to a nonzero value, so $\norm{D_LT_LF}\to\infty$.
Since $B$ is bounded on the fixed finite-dimensional space,
$\norm{T_LBF}\to0$. The reverse triangle inequality proves the claim.
\end{proof}

There is also a useful exact description at the weaker pairing level.
Here $\eta$ denotes the boundary covector $\eta_{L,2K}$, distinct from
the root-space functional $\eta_{\mathrm B}$. In physical scaling the finite commutator is
\[
 [D_L,W]=|\beta\rangle\eta- |\eta\rangle\beta^*,
 \qquad \eta(c_{L,K})=\sqrt L,\qquad \beta^*(c_{L,K})=0,
\]
where $\beta=\beta^{\log}$ includes the physical factor $2\pi/L$ and is odd and $c_{L,K}$ is even. Here the rank-one notation
means $|a\rangle b^*:x\mapsto a\ip{x}{b}$, and the vector representing
$\eta$ is understood. Therefore for every finite test $v$,
\begin{equation}\label{v1:eq:paireddefect}
 \begin{split}
 \ip{W\mathcal R_{L,K}F}{v}
 ={}&i\,d_{L,K}(F)\ip{D_LWc_{L,K}}{v}
       -i\sqrt L\,d_{L,K}(F)\ip{\beta}{v}\\
    &-d_{L,K}(BF)\ip{Wc_{L,K}}{v}.
 \end{split}
\end{equation}
Thus even where the first and last terms can be made small, the boundary
commutator term has not disappeared. No unproved asymptotic value for this
term is used in this manuscript.

One part of G1 is preserved particularly simply. The original finite matrix
preserves parity; if the G1 candidate $k$ is even, $D_LWk$ is odd. Hence
\begin{equation}\label{v1:eq:g1parity}
 \ip{D_LWk}{\Jr F}=\ip{D_LWk}{\Jc F}.
\end{equation}
This is exact and requires no estimate on the shell derivative. Establishing
the centered right-hand G1 bound still requires the independently audited
continuum G1 inputs. Equation \eqref{v1:eq:g1parity} does not cancel
\eqref{v1:eq:paireddefect} in the later graph calculation.

\section{The rebuilt finite metric and reflected-partner argument}\label{sec:g2}
Throughout this section $J=\Jr_{L,K}$, $W=W_{L,2K}$, $D=D_L$,
$z=z^{\mathrm c}_{L,2K}$, and $\mathcal R=\mathcal R_{L,K}$ are the
objects of \eqref{v1:eq:repair} and \eqref{v1:eq:graphdefect} on the
\emph{same} finite space. Write $\eta_\partial=\sqrt L\delta$ and
$\beta=\beta^{\log}$. Then
\[
 [D,W]=|\beta\rangle\eta_\partial-|\eta_\partial\rangle\beta^*.
\]
The rank-one convention is $|u\rangle v^*:x\mapsto u\ip{x}{v}$.
In particular the root functional $\eta_{\mathrm B}$ is not identified
with either finite endpoint functional.
Theorem~\ref{thm:v15-finite-metric} gives the full finite-core metric
budget without a global closed-operator assumption. The source-annihilating
projection of Proposition~\ref{prop:v15-projector} is a separate positive
benchmark on this same finite object, not a replacement for $W$ in the
identities below or a claim of small metric defect. Proposition~\ref{prop:v16-projected-cut}
also shows why inserting that sharp band projection into a Weil pairing
does not automatically preserve the earlier holomorphic-strip transfer.
Proposition~\ref{prop:v17-pair-edge} further shows that reflection alone
does not cancel the leading edge term; control of the entire zero sum is
still required.
Proposition~\ref{prop:v18-full-projection} supplies a different,
full-space projection that kills the source and both shell directions
while preserving ordinary coercivity and Weil transfer through a proved
full-strip product bound. Its exact commutator remains in
\eqref{eq:v18-full-budget}. Corollary~\ref{cor:v18-shift} shows why
adding a fixed positive scalar to its projected Weil form still fails
the small-defect criterion. Neither result licenses meromorphic
resolvent transfer or changes the endpoint of $J$.
With the stronger fixed tail margin $p>1$,
Proposition~\ref{prop:v19-action} now controls $WSJF$ in ordinary norm
by an explicit finite exponential sum. Corollary~\ref{cor:v19-square}
audits the globally positive candidate $SW^2S$: top-root mass grows,
while lower jets at multiple zeros lose unscaled coercivity. The full
squared-metric commutator is still required.

\begin{proposition}[Restricted positive shifts after corrected transfer]\label{prop:v1-shift}
On any fixed one-sided root space, for every fixed $a_0>0$ there is
$c>0$ such that, for sufficiently large $L$,
\[
 \ip{(W+a_0I)JF}{JF}\ge c\|F\|^2.
\]
This is restricted coercivity, not positivity on the full finite space.
\end{proposition}
\begin{proof}
The output map $F\mapsto g_F$ is injective, so its Gram matrix has a
positive minimum eigenvalue on the fixed space. Theorem~\ref{v1:thm:repair}
transfers that ordinary Gram matrix, while the Weil Gram matrix tends to
zero by one-sided isotropy. The positive scalar term therefore dominates.
\end{proof}
If a negative compact Weil test is used to bound the lowest finite form
eigenvalue away from zero under failure of RH, its particular Fourier
approximants must also converge in the closed form. For a smooth compact
test this follows from the correct matrix bound and rapid coefficient
decay. This optional calibration is not needed for the proposition.
In either formulation coercivity alone does not imply a small Lyapunov defect.

\begin{proposition}[Lyapunov identity with the graph defect retained]\label{prop:v1-lyapunov}
Let $BF=\mu F$, $d=\Re\mu>0$, $\|F\|=1$, and put
\[
 x=JF,\quad e=\eta_{\mathrm B}(F),\quad
 y=e z+\mathcal RF,\quad b=\eta_\partial(x).
\]
For every self-adjoint finite matrix $S$,
\begin{equation}\label{eq:v1-general-metric}
 d\ip{Sx}{x}
 =-\Re\ip{Sy}{x}-\frac{i}{2}\ip{[D,S]x}{x}.
\end{equation}
For $T=W+a_0I$ this becomes
\begin{equation}\label{eq:v1-residual}
 2d\ip{Tx}{x}=-2\Re\ip{T(ez+\mathcal RF)+ib\beta}{x}.
\end{equation}
If $\ip{Tx}{x}\ge c$ and $\|x\|\le C$, then
\[
 \|T(ez+\mathcal RF)+ib\beta\|\ge dc/C.
\]
\end{proposition}
\begin{proof}
The exact graph equation is $iDx-\mu x=y$. Pair it in both slots
with $Sx$. Self-adjointness gives
$2d\ip{Sx}{x}=-i\ip{[D,S]x}{x}-2\Re\ip{Sy}{x}$.
For $S=T$ the scalar part commutes with $D$. The rank-two commutator
has expectation $b\ip{\beta}{x}-\overline{b\ip{\beta}{x}}$,
which gives \eqref{eq:v1-residual}. Apply Cauchy--Schwarz for the bound.
\end{proof}
For this sampler $b=\sqrt L\,\mathscr E(F)$ exactly. The vector
$\mathcal RF$ cannot be omitted on the ground that $J-\Jc$ tends to zero
in ordinary norm: \eqref{v1:eq:graphnorm} proves the opposite behavior
for its derivative. On a left root window the signs reverse, with the
same absolute-value conclusion.

\begin{proposition}[Candidate comparison and scalar cancellation]\label{prop:v1-candidate}
For an even candidate $k$ with $a=\eta_\partial(k)\ne0$, put
$r=DWk$, $q=ib/a$, and $m=y-qDk$. Then
\begin{align}
 T y+ib\beta&=T m+q(r+a_0Dk),\label{eq:v1-candidate}\\
 d\ip{Wx}{x}&=-\Re\ip{Wm}{x}-\Re\bigl(q\ip{r}{x}\bigr).
 \label{eq:v1-unshifted}
\end{align}
More generally, a self-adjoint correction commuting with $D$ satisfies
$d\ip{Sx}{x}=-\Re\ip{S y}{x}$ exactly.
\end{proposition}
\begin{proof}
Parity gives $[D,W]k=a\beta$, hence $WDk=r-a\beta$.
Since $qa=ib$, expansion proves the first identity. Use
Proposition~\ref{prop:v1-lyapunov} with $S=W$ for the second.
The final statement is \eqref{eq:v1-general-metric} with zero commutator.
\end{proof}
Thus scalar coercivity balances its own source and graph terms. It does
not independently exclude an off-line root. For a noncommuting positive
$S$, the exact combined quantity
\[
 \mathfrak X_S=\Re\ip{S(ez+\mathcal RF)}{x}
       +\frac i2\ip{[D,S]x}{x}=-d\ip{Sx}{x}
\]
still needs an independent upper estimate before its lower bound can be
used in a contradiction. With $S=W^2$, $[D,W^2]=[D,W]W+W[D,W]$ has
rank at most four. Form smallness alone does not give norm smallness.

\begin{proposition}[Reflected top-root lower bound for the corrected sampler]\label{prop:v1-reflected}
Suppose $\rho$ is an off-line zero of exact multiplicity $m$ and
$\Re\rho>1/2$. Set $\rho^\sharp=1-\overline\rho$,
$F_+=\zeta/(s-\rho)^m$, and $F_-=\zeta/(s-\rho^\sharp)^m$.
On the common finite construction, with $x=JF_+$ and $x_-=JF_-$,
\[
 \ip{Wx}{x}\to0,\qquad \ip{Wx}{x_-}\to C_{\rho,m}\ne0,
 \qquad \liminf\|Wx\|>0,
\]
where the first-slot-linear constant is
\begin{equation}\label{eq:v1-crossconstant}
 C_{\rho,m}=m\,\kappa(\rho)\frac{\zeta^{(m)}(\rho)}{m!}
 \overline{\kappa(\rho^\sharp)
                   \frac{\zeta^{(m)}(\rho^\sharp)}{m!}}.
\end{equation}
\end{proposition}
\begin{proof}
The functional equation and conjugation preserve zero multiplicity, so
$\rho^\sharp$ also has exact order $m$. The top root output is nonzero
only at its selected zero in the global value-level zero sum. Thus the
global cross form is exactly \eqref{eq:v1-crossconstant}; its factors
are nonzero. Both outputs meet the holomorphic strip hypotheses of
Theorem~\ref{v1:thm:repair}, which transfers the self and cross forms.
Ordinary norms of $x_-$ are bounded. The inequality
$\|Wx\|\|x_-\|\ge|\ip{Wx}{x_-}|$ proves the lower bound.
\end{proof}
For $m>1$ this is a bound on the top generalized vector. The genuine
eigenvector $\zeta/(s-\rho)$ vanishes at every zero when $m>1$, and the
value-level Weil formula supplies no analogous bound on it. No claim
of coercivity on an entire multiple root block follows from this proposition.

\subsection*{Boundary-annihilating tests: the exact surviving identity}
Write $\mu=\rho-1/2$, $A_\rho=iD-\mu I$, and $x_j=JF_{\rho,j}$.
With $x_0=0$, the corrected graph relation gives
\[
 A_\rho x_m=u_m,\qquad
 u_m=x_{m-1}+\eta_{\mathrm B}(F_{\rho,m})z+\mathcal R F_{\rho,m}.
\]
In the canonical chain $\eta_{\mathrm B}(F_{\rho,1})=1$ and the
functional vanishes on higher chain vectors. Thus the old radical
predecessor now has the additional term $\mathcal R F_{\rho,m}$.
Define $v^\pm=(A_\rho^*)^{-1}x_\pm$, with $x_+=x_m$.
The inverse has norm at most $1/d$, where $d=\Re\rho-1/2$.
If $\delta(v^+)\ne0$, set
\[
 \alpha=\delta(v^-)/\delta(v^+),\qquad
 \widetilde v=v^- -\alpha v^+,\qquad
 \widetilde y=x_- -\alpha x_m=A_\rho^*\widetilde v.
\]
These are finite definitions, not assertions of a bounded limiting ratio.
For $b_m=\eta_\partial(x_m)$, commutation gives exactly
\begin{equation}\label{eq:v1-boundarytest}
 \ip{Wu_m+ib_m\beta}{\widetilde v}
       =\ip{Wx_m}{\widetilde y}.
\end{equation}
Indeed $A_\rho Wx_m=Wu_m+i[D,W]x_m$, and
$\eta_\partial(\widetilde v)=0$ removes the other boundary channel.
If the ratio $\alpha$ is bounded along a chosen sequence, the right side
tends to $C_{\rho,m}$ by Proposition~\ref{prop:v1-reflected} and
conjugate-linearity in the second slot. Corollary~\ref{cor:v11-ratio}
below shows that this bounded-ratio premise actually fails for the
corrected sampler. Even a hypothetical bounded-ratio construction would
also require, before isolating the beta term,
\begin{equation}\label{eq:v1-needed-pole-transfer}
 \ip{W(x_{m-1}+\eta_{\mathrm B}(F_{\rho,m})z)}{\widetilde v}
       +\ip{W\mathcal R F_{\rho,m}}{\widetilde v}\longrightarrow0.
\end{equation}
Neither term is disposed of by the holomorphic form-transfer theorem.
The first involves a potentially meromorphic adjoint-resolvent test; the
second is the new graph correction. Equation~\eqref{v1:eq:paireddefect}
shows its explicit boundary contribution.

\paragraph{What is withdrawn.}
The previous pole-aware transfer proof controlled the singular zero but
did not bound the product error over all zeros. Its conclusion is not used
here. A valid replacement must give a summable majorant for the paired
transforms on the entire strip, including cut jumps and slow reflected
tails, on the actual finite test family. Consequently the old nonzero
limit for $\sqrt L\ip{\beta}{\widetilde v}$ is not an established result
of this version. The exact arithmetic formula for that quantity remains
valid for every finite test, as shown in Section~\ref{sec:arithmetic}.

\subsection*{The endpoint ratio and a limitation on fixed test spaces}
The bounded-ratio premise above can now be decided for the corrected
sampler. Fix a hypothetical zero $\rho=1/2+\mu$, where
$\mu=d+i\gamma$ and $0<d<1/2$, and set
$p=\rho^\sharp=1/2-\overline\mu$. Let $\mathcal F$ be a fixed
finite-dimensional space of Hardy inputs satisfying the holomorphy,
tail and vertical-decay hypotheses of Sections~\ref{sec:sampling}--\ref{sec:g2}.
One may include any fixed finite set of genuine root functions and
$\zeta$ itself, choosing the smoothing order sufficiently large.
Fix any norm on $\mathcal F$. Use the common cutoff
\[
 K\ge\left\lceil e^{2L}L^4\right\rceil,
 \qquad J=\Jr_{L,K},\qquad
 V_LG=(A_\rho^*)^{-1}JG.
\]
The letter $V_L$ here denotes a resolvent map; $V_n$ still denotes a
Fourier basis vector. Choose $r$ as in \eqref{v1:eq:tails} and put
$\nu=\min(q,r,1)>1/2$. All constants below may depend on the fixed
space and on $\rho$, but not on $L$, $K$, or a bounded input.

\begin{lemma}[Endpoint of the corrected adjoint resolvent]
\label{lem:v11-resolvent-endpoint}
Uniformly for $G\in\mathcal F$,
\begin{equation}\label{eq:v11-endpoint-asymptotic}
 \delta V_LG
 =-H_G(p)\frac{e^{-\overline\mu L/2}}{1-e^{-\overline\mu L}}
       +O(e^{-\nu L/2})\|G\|.
\end{equation}
Consequently the rescaled endpoint functional satisfies, in the dual
norm of this fixed space,
\begin{equation}\label{eq:v11-rescaled-endpoint}
 \ell_L(G):=e^{\overline\mu L/2}(1-e^{-\overline\mu L})\delta V_LG
 =-H_G(p)+O(e^{-(\nu-d)L/2})\|G\|.
\end{equation}
\end{lemma}
\begin{proof}
The multiplier of $(A_\rho^*)^{-1}$ is
$(-it-\overline\mu)^{-1}$. Thus the full centered Fourier series is
the periodization of the inverse transform of
\[
 H_{v,G}(s)=-\frac{H_G(s)}{s-p}.
\]
This transform is used only for Fourier inversion, not in a Weil
zero-sum formula. Its possible pole at $p$ has residue $-H_G(p)$.
Contour shifts with the same decay bounds as \eqref{v1:eq:tails} give
\[
 h_G(y)=
 \begin{cases}
 -\E(G)(\sigma-p)^{-1}e^{qy},&y<0,\\
 -H_G(p)e^{-\overline\mu y}+O(e^{-ry})\|G\|,&y>0.
 \end{cases}
\]
Summing these tails at the half-period points $(j+1/2)L$ gives
\begin{align*}
 \delta V_L^{\mathrm c,\infty}G
 ={}&-H_G(p)\frac{e^{-\overline\mu L/2}}{1-e^{-\overline\mu L}}\\
 &-\frac{\E(G)}{\sigma-p}
          \frac{e^{-qL/2}}{1-e^{-qL}}+O(e^{-rL/2})\|G\|.
\end{align*}
If $|H_G(1/2+it)|\le C(1+|t|)^{-R}\|G\|$, $R\ge2$,
the omitted Fourier endpoint tail is at most
$C(1+K/L)^{-R}\|G\|$, because the resolvent supplies one further
inverse power. The endpoint correction has coefficient
$d_{L,K}(G)=\E(G)-\delta\Jc_{L,2K}G=O(\|G\|)$.
For its shell $c=c_{L,K}$, pairing the positive and negative indices
gives the exact identity
\begin{equation}\label{eq:v11-shell-resolvent}
 \begin{split}
 \delta(A_\rho^*)^{-1}c
 &=\frac1{2K}\sum_{K<|n|\le2K}\frac1{-it_n-\overline\mu}\\
 &=-\frac{\overline\mu}{K}
       \sum_{n=K+1}^{2K}\frac1{t_n^2+\overline\mu^{\,2}}
   =O_\mu(L^2/K^2).
 \end{split}
\end{equation}
For large $L$, $K/L$ exceeds every fixed multiple of $|\mu|$,
which justifies the last bound despite the complex denominator.
Both cutoff errors are $O(e^{-4L}L^{-6})$ at the stated minimum
cutoff and decrease upon enlarging it. They are absorbed by
$O(e^{-\nu L/2})$. This proves \eqref{eq:v11-endpoint-asymptotic};
multiplication by the displayed scale proves
\eqref{eq:v11-rescaled-endpoint}.
\end{proof}

\begin{lemma}[Uniform lower norm bound for the resolvent map]
\label{lem:v11-resolvent-gram}
There are $c,C>0$ such that for all sufficiently large $L$,
\[
 c\|G\|\le\|V_LG\|\le C\|G\|\qquad(G\in\mathcal F).
\]
\end{lemma}
\begin{proof}
The Gram form of the centered resolvent samples is a Riemann sum
converging, uniformly on the fixed finite-dimensional space, to
\[
 \frac1{2\pi}\int_{\mathbb R}
       \frac{|H_G(1/2+it)|^2}{|-it-\overline\mu|^2}\,dt.
\]
This form is positive definite: a zero integral forces the analytic
output, and hence $G$, to vanish. Its minimum eigenvalue is positive.
The shell changes the resolvent vector by at most
$d^{-1}|d_{L,K}(G)|\sqrt{L/(2K)}=o(\|G\|)$.
Finite-dimensional compactness proves both bounds.
\end{proof}

\begin{corollary}[Divergence of the proposed endpoint ratio]
\label{cor:v11-ratio}
For the top-root pair in Proposition~\ref{prop:v1-reflected}, put
$v^\pm=V_LF_\pm$. Wherever $\delta v^+\ne0$ and $L$ is sufficiently
large, the proposed ratio obeys
\[
 \left|\frac{\delta v^-}{\delta v^+}\right|
       \ge c e^{(\nu-d)L/2}.
\]
Thus it cannot be bounded along any sequence tending to infinity on
which it is defined. At zeros of the denominator the proposed formula
for the test is undefined. The vectors
$v^- -(\delta v^-/\delta v^+)v^+$ have unbounded ordinary norm.
\end{corollary}
\begin{proof}
Genuine zero cancellation gives
\[
 H_{F_+}(p)=0,\qquad
 H_{F_-}(p)=\kappa(p)\frac{\zeta^{(m)}(p)}{m!}\ne0.
\]
The preceding endpoint lemma bounds the denominator by $Ce^{-\nu L/2}$
and the numerator below by $ce^{-dL/2}$. The inputs $F_+$ and $F_-$
are linearly independent, and Lemma~\ref{lem:v11-resolvent-gram}
then gives norm divergence of the displayed combination.
\end{proof}

\begin{theorem}[Loss of the reflected signal on a fixed test space]
\label{thm:v11-fixed-tests}
Include $F_+$ in $\mathcal F$ and let $G_L\in\mathcal F$.
Suppose $v_L=V_LG_L$ is bounded in ordinary norm and $\delta v_L=0$.
Then
\begin{equation}\label{eq:v11-signal-loss}
 \ip{WJF_+}{A_\rho^*v_L}
       =\ip{WJF_+}{JG_L}\longrightarrow0.
\end{equation}
This conclusion uses holomorphic form transfer only. No pole-aware
Weil transfer is assumed.
\end{theorem}
\begin{proof}
Lemma~\ref{lem:v11-resolvent-gram} makes $G_L$ bounded.
Equation~\eqref{eq:v11-rescaled-endpoint} and $\delta v_L=0$
imply $H_{G_L}(p)\to0$. On the other hand, the global value-level
zero formula and the top-root cancellation give
\[
 \Q(g_{F_+},g_G)=mH_{F_+}(\rho)\overline{H_G(p)}.
\]
Theorem~\ref{v1:thm:repair} transfers this identity uniformly for
bounded $G\in\mathcal F$. Substitution of $G_L$ proves the result.
Equivalently, for every bounded family of inputs,
\begin{equation}\label{eq:v11-endpoint-cross-relation}
 \ip{WJF_+}{JG_L}
   =-mH_{F_+}(\rho)\overline{\ell_L(G_L)}+o(1).
\end{equation}
\end{proof}

In particular, normalizing the divergent two-vector test does not retain
the desired nonzero cross limit. The endpoint constraint suppresses the
very reflected evaluation detected by the Weil form. This is a limitation
of a fixed finite-dimensional input space with the specified resolvent
and sampler. It does not exclude test spaces that grow or change with
$L$, nor other ways to handle the boundary channel.

\paragraph{A remaining alternative and its cost.}
A finite test can instead be corrected directly:
\[
 w_L=v^- -\delta(v^-)c_{L,K},\qquad \delta w_L=0.
\]
It is bounded and differs from $v^-$ by a norm-vanishing vector, but
$A_\rho^*w_L=x_- -\delta(v^-)A_\rho^*c_{L,K}$ is no longer the
image of a bounded input in the same fixed test space. The exact price is
\[
 \ip{WJF_+}{A_\rho^*w_L}
 =\ip{WJF_+}{x_-}
  -\overline{\delta(v^-)}\ip{WJF_+}{A_\rho^*c_{L,K}}.
\]
For this particular large shell,
$\|A_\rho^*c_{L,K}\|^2=\|Dc_{L,K}\|^2+
|\mu|^2\|c_{L,K}\|^2$ by symmetry. Consequently its graph
correction has size comparable to
$e^{-dL/2}\sqrt{K/L}$, which diverges at the specified cutoff.
A different endpoint correction or a smaller auxiliary shell inside the
same ambient space may change that cost. Either choice requires an
independent bound for the displayed Weil pairing and the complete graph
source. The present calculation establishes no such bound.


\subsection*{A parity benchmark that uses ordinary norm transfer only}
Let $\mathsf P V_n=V_{-n}$ and $P_+=(I+\mathsf P)/2$.
Then $P_+\succeq0$ and $[D,P_+]=D\mathsf P$.
For the fixed one-sided root family the even projection of the Hardy
output is injective. To see this, a vanishing even part implies
$\kappa(s)F(s)+\kappa(1-s)F(1-s)=0$ by analytic continuation.
Writing $F=\zeta R_F$ and using the functional equation gives a
meromorphic identity whose rational principal parts lie in disjoint
right and left windows. The coefficients at each pole on the right
must vanish, so $F=0$. Finite-dimensional compactness gives a positive
lower norm bound. Ordinary Riemann-sum convergence transfers it to
$\Jc$; the centered phase commutes with reflection, and the norm-small
endpoint correction preserves the bound for $J$.
Thus $P_+$ is a bounded noncommuting positive benchmark. Its exact
Lyapunov source is $P_+(ez+\mathcal RF)$, however, and the new graph
term is still present. Positivity by itself supplies no small upper bound.

\section{The arithmetic boundary observable}\label{sec:arithmetic}
\begin{proposition}[Exact arithmetic realization of the CCM boundary channel]\label{prop:beta-arithmetic-channel}
Let $L=2\log\lambda$, $t_n=2\pi n/L$, and let
$\mathfrak b_\lambda(n)$ for $n\in\mathbb Z$ denote the real odd coefficient in the CCM
rank--two commutator before physical rescaling, so that
\begin{equation}\label{eq:beta-bcoeff-expansion}
 \beta^{\log}_{\lambda,N}
 =\frac{2\pi}{L}\sum_{|n|\le N}\mathfrak b_\lambda(n)V_n.
\end{equation}
With
\[
 \rho_\infty(y)=\frac{e^{y/2}}{e^y-e^{-y}},
 \qquad
 \mathcal V(y)=2(e^{y/2}-1)
 -\sum_{2\le k\le e^y}\frac{\Lambda(k)}{\sqrt k},
 \qquad 0\le y\le L,
\]
where $\mathcal V$ is exactly the centered prime cumulative already used in
\eqref{eq:centered-prime-cumulative}, one has the finite, zero-independent
identity
\begin{align}
 \mathfrak b_\lambda(n)
 ={}&\frac{32L\sinh^2(L/4)\,n}
          {L^2+16\pi^2n^2}
 +\frac1\pi\int_0^L\sin(t_ny)\rho_\infty(y)\,dy \notag\\
 &+\frac1\pi\sum_{1<k\le e^L}
   \frac{\Lambda(k)}{\sqrt k}\sin(t_n\log k)                                      \label{eq:beta-bcoeff-direct}\\
 ={}&\boxed{\frac{t_n}{\pi}\int_0^L
       \mathcal V(y)\cos(t_ny)\,dy
 +\frac1\pi\int_0^L
   \bigl(\rho_\infty(y)-e^{-y/2}\bigr)\sin(t_ny)\,dy.}             \label{eq:beta-bcoeff-V}
\end{align}
Equivalently, if
\[
 \mathcal E_L(t)
 :=\sum_{1<k\le e^L}\frac{\Lambda(k)}{\sqrt k}\sin(t\log k)
   -\int_0^L e^{y/2}\sin(ty)\,dy,
\]
then at the lattice frequencies
\begin{equation}\label{eq:beta-bcoeff-E}
 \boxed{
 \mathfrak b_\lambda(n)
 =\frac{\mathcal E_L(t_n)}{\pi}
  +\frac1\pi\int_0^L\sin(t_ny)\rho_\infty(y)\,dy
  -\frac{4t_n(1-\lambda^{-1})}{\pi(1+4t_n^2)}.}
\end{equation}
Thus the same centered pole--prime discrepancy that drives G1 also generates
the boundary vector appearing in the graph calculation; no nontrivial-zero data occur in
\eqref{eq:beta-bcoeff-direct}--\eqref{eq:beta-bcoeff-E}.

More precisely, write an arbitrary finite vector as
\[
 v=L^{-1/2}\sum_{|n|\le N}h_nV_n,
 \qquad
 H_v(y)=\frac1L\sum_{|n|\le N}h_ne^{it_ny},
 \qquad
 A_v(y)=\overline{H_v(-y)}-\overline{H_v(y)}.
\]
Then, for the first-slot-linear convention used throughout this manuscript,
\begin{equation}\label{eq:beta-arithmetic-trace}
 \boxed{
 i\sqrt L\,\langle\beta^{\log}_{\lambda,N},v\rangle
 =\int_0^L\mathcal V(y)A_v'(y)\,dy
  +\int_0^L\bigl(\rho_\infty(y)-e^{-y/2}\bigr)A_v(y)\,dy.}
\end{equation}
Since $A_v$ is $L$-periodic and odd, one also has the exact folded form
\begin{align}
 i\sqrt L\,\langle\beta^{\log}_{\lambda,N},v\rangle
 ={}&\int_0^{L/2}
   \bigl(\mathcal V(y)+\mathcal V(L-y)\bigr)A_v'(y)\,dy \notag\\
 &+\int_0^{L/2}
 \Bigl[
  \rho_\infty(y)-e^{-y/2}
  -\rho_\infty(L-y)+e^{-(L-y)/2}
 \Bigr]A_v(y)\,dy.                                      \label{eq:beta-arithmetic-folded}
\end{align}
\end{proposition}
\begin{proof}
Connes--Consani--Moscovici show that the truncated Weil matrix has
$\tau_{ij}=(b_i-b_j)/(i-j)$ off the diagonal and that its commutator with the
integer Fourier generator is the corresponding rank--two matrix
\cite{ConnesConsaniMoscovici2025}.  Their point contribution
\[
 W_{0,2}(V_n,V_m)
 =\frac{32L\sinh^2(L/4)(L^2-16\pi^2mn)}
 {(L^2+16\pi^2m^2)(L^2+16\pi^2n^2)}
\]
therefore contributes
$32L\sinh^2(L/4)n/(L^2+16\pi^2n^2)$ to $b_n$.
For $n\ne m$, the kernel already used in Lemma~\ref{lem:g1-high-low} is
\[
 q(U_n,U_m)(y)
 =\frac{\sin(t_my)-\sin(t_ny)}{\pi(n-m)}.
\]
The signs in the semilocal Weil decomposition then give the positive prime
sine sum and the positive archimedean sine integral in
\eqref{eq:beta-bcoeff-direct}.  This proves the first identity directly from
the finite Weil matrix.

The point term also has the integral representation
\[
 \frac{32L\sinh^2(L/4)n}{L^2+16\pi^2n^2}
 =-\frac1\pi\int_0^L
   (e^{y/2}+e^{-y/2})\sin(t_ny)\,dy.
\]
In Stieltjes form,
\[
 d\mathcal V(y)=e^{y/2}\,dy
 -\sum_{1<k\le e^L}\frac{\Lambda(k)}{\sqrt k}\,\delta_{\log k}(dy).
\]
Substitution in \eqref{eq:beta-bcoeff-direct} gives
\[
 \mathfrak b_\lambda(n)
 =-\frac1\pi\int_0^L\sin(t_ny)\,d\mathcal V(y)
 +\frac1\pi\int_0^L
   (\rho_\infty(y)-e^{-y/2})\sin(t_ny)\,dy.
\]
The boundary term in Stieltjes integration by parts vanishes because
$t_nL=2\pi n$, giving \eqref{eq:beta-bcoeff-V}.  Formula
\eqref{eq:beta-bcoeff-E} follows from
\[
 \int_0^Le^{-y/2}\sin(t_ny)\,dy
 =\frac{4t_n(1-\lambda^{-1})}{1+4t_n^2}.
\]

Finally, \eqref{eq:beta-bcoeff-expansion} and the first-slot-linear inner
product give
\[
 \sqrt L\,\langle\beta^{\log}_{\lambda,N},v\rangle
 =\frac{2\pi}{L}\sum_{|n|\le N}
   \mathfrak b_\lambda(n)\overline{h_n}.
\]
But
\[
 A_v(y)=\frac{2i}{L}\sum_{|n|\le N}
 \overline{h_n}\sin(t_ny),
 \qquad
 A_v'(y)=\frac{2i}{L}\sum_{|n|\le N}
 t_n\overline{h_n}\cos(t_ny),
\]
which proves \eqref{eq:beta-arithmetic-trace}.  Periodicity gives
$A_v(L-y)=-A_v(y)$ and $A_v'(L-y)=A_v'(y)$; splitting at $L/2$ proves
\eqref{eq:beta-arithmetic-folded}.
\end{proof}
\begin{proposition}[The present absolute PNT remainder is quantitatively insufficient for the G2 boundary channel]\label{prop:pnt-beta-barrier}
Fix $T>0$.  The unconditional estimate
\eqref{eq:centered-prime-cumulative}, inserted absolutely into the exact
formula \eqref{eq:beta-bcoeff-V}, gives constants $c_T,C_T>0$ such that
\begin{equation}\label{eq:pnt-beta-barrier}
 \boxed{
 \sup_{|t_n|\le T}|\mathfrak b_\lambda(n)|
 \le C_T\bigl(1+\lambda e^{-c_T\sqrt L}\bigr).}
\end{equation}
In particular this already-audited PNT input, used only by absolute values,
does not imply that the boundary observable
$\sqrt L\langle\beta^{\log},v\rangle$ on a varying test family tends to zero.  The
extra factor $e^{-L/2}=\lambda^{-1}$ that closes the G1 pole--prime estimate
comes from the special co-Poisson cross-tail kernel and is absent from the
periodic Hardy boundary kernel in \eqref{eq:beta-arithmetic-trace}.
This records the limitation of this absolute-value estimate. It is not a lower bound on the actual error, nor a proof that signed cancellation is impossible.
\end{proposition}
\begin{proof}
The second integral in \eqref{eq:beta-bcoeff-V} is $O_T(1)$: near zero
$\rho_\infty(y)=(2y)^{-1}+O(1)$ while
$\sin(t_ny)=O_T(y)$, and for $y\ge1$ the kernel decays exponentially.
For the first integral, \eqref{eq:centered-prime-cumulative} gives
\[
 |t_n|\int_1^L|\mathcal V(y)|\,dy
 \ll_T\int_1^Le^{y/2-c\sqrt y}\,dy
 \ll_T\lambda e^{-c_T\sqrt L},
\]
after splitting at $L/4$ and decreasing $c_T$ if necessary.  The bounded
initial interval is absorbed into the constant.  This proves
\eqref{eq:pnt-beta-barrier}.
\end{proof}
\section{Analytic filters, moving phases, and the reflected pole}\label{sec:filters}
The following continuum identities remain useful diagnostics. They do
not identify the endpoint or graph behavior of the corrected finite
sampler without a further transfer argument.

\subsection*{The exact meromorphic normal form}
For an off-line zero $\rho$ of order $m$, the critical-line adjoint
resolvent multiplier is $-(s-\rho^\sharp)^{-1}$. Define
\[
 H_{v^+}(s)=-\frac{\kappa(s)\zeta(s)}{(s-\rho)^m(s-\rho^\sharp)},
 \qquad
 H_{v^-}(s)=-\frac{\kappa(s)\zeta(s)}{(s-\rho^\sharp)^{m+1}}.
\]
Their continuum values at the uncentered cut follow by the right-hand
Cauchy pole at $\sigma$. They give
\[
 \alpha_\infty=
 \left(\frac{\sigma-\rho}{\sigma-\rho^\sharp}\right)^m,
 \qquad
 H_{\widetilde v}(s)=\frac{\kappa(s)\zeta(s)}{s-\rho^\sharp}
 \left[\frac{\alpha_\infty}{(s-\rho)^m}
              -\frac1{(s-\rho^\sharp)^m}\right].
\]
The bracket vanishes at $\sigma$, canceling the Hardy Cauchy pole.
At $\rho$ the zeta zero cancels the apparent pole. At $\rho^\sharp$
there remains exactly one simple pole, with residue
\begin{equation}\label{eq:v1-pole-residue}
 -\kappa(\rho^\sharp)\frac{\zeta^{(m)}(\rho^\sharp)}{m!}\ne0.
\end{equation}
This is an algebraic statement about the uncentered continuum transform.
For the corrected finite resolvent tests, Corollary~\ref{cor:v11-ratio}
proves that the endpoint ratio instead diverges in magnitude; it cannot
tend to this finite $\alpha_\infty$.

Adding a transform holomorphic at $\rho^\sharp$ cannot change this
residue. If such a correction $w$ is cross-invisible to the top root,
then, wherever the paired form is well-defined,
\[
 QW(g_{\rho,m},A_\rho^*(c v^-+w))
       =\overline c\,C_{\rho,m}.
\]
The conjugate on $c$ is required by the first-slot-linear convention.
This corrects the opposite convention used in the previous proof.

\begin{proposition}[The bounded-filter tradeoff]\label{prop:v1-filter}
Let $\Phi_L$ be bounded and holomorphic on $\Re s>1/2$ and admit
continuation near $\rho^\sharp$. Put
$M_L=\|\Phi_L\|_{H^\infty(\Re s>1/2)}$.
If $|\Phi_L(\rho)\overline{\Phi_L(\rho^\sharp)}|\ge c_0>0$, then
\[
 |\Phi_L(\rho^\sharp)|\ge c_0/M_L.
\]
In particular residue damping by $e^{-\epsilon L}$ while retaining
that product forces $M_L\ge c_0e^{\epsilon L}$.
\end{proposition}
\begin{proof}
Since $\rho$ is in the right half-plane, $|\Phi_L(\rho)|\le M_L$.
The conclusion follows by dividing the product lower bound by $M_L$.
\end{proof}
To interpret the product as a global Weil cross form, impose in addition
global transform admissibility and the strip/tail estimates required for
that form and its transfer. Local continuation at two points is not
sufficient. A boundary-annihilation ratio also requires
$\Phi_L(\sigma)\ne0$. Under these extra conditions the cross constant
is multiplied by $\Phi_L(\rho)\overline{\Phi_L(\rho^\sharp)}$ and
the residue by $\Phi_L(\rho^\sharp)$.

The absolute PNT envelope applied to a reflected image of size
$|\Phi_L(\rho^\sharp)|e^{-dL}$ yields an upper estimate with scale
$|\Phi_L(\rho^\sharp)|e^{(1/2-d)L-c\sqrt L}$.
Making this particular estimate tend to zero by amplitude damping
alone calls for exponentially small residue and hence an exponentially
large available Hardy norm. The existing G1 estimate pays uniform
Sobolev constants and does not cover that regime. This is a limitation
of this proof strategy. A large upper bound on an error does not prove
that the actual error is large, and the calculation does not exclude
signed cancellation or a new estimate that avoids that norm cost.

\subsection*{Moving phases and the omitted periodic image}
For real $\tau$ the multiplier
$H_F^{(\tau)}(s)=e^{-\tau(s-1/2)}H_F(s)$ translates the inverse
transform to $g_F(y-\tau)$. Global cross forms are invariant, because
the two reflected factors cancel. The uncentered Cauchy value at zero
is $e^{-q\tau}\mathscr E(F)$ for $\tau>0$.
These are exact continuum statements, and the coefficientwise graph
identity changes its source by the same phase.

The earlier claimed periodic kernel bound proportional to $e^{-\tau/2}$
omitted the nearest reflected-pole image. A pole at real displacement
$-d$, $0<d<1/2$, contributes a term on the scale
$e^{-d(L-\tau)}$ at the opposite end of the period. For example,
$d=0.1$ and $\tau=0.26L$ give $e^{-0.074L}$, larger than
$e^{-0.13L}$. Restricting $\tau<L/2$ does not repair that comparison.

The historical threshold comparison
\[
 \frac{d}{d+q}<\frac1{2q+1}\quad\Longleftrightarrow\quad d<1/2
 \qquad(q>0)
\]
is algebraically correct. It proves a conflict only if both proposed
asymptotic thresholds have first been justified, with remainders small
relative to the relevant periodic image. An absolute $O(e^{-cL})$
remainder is not enough if the image decays faster. The previous
unconditional exclusion of the whole moving-phase route is therefore
withdrawn. Moving $\sigma$ likewise requires uniform smoothing, cutoff,
and transfer estimates; the scalar relation
$\sigma F(\sigma)\to\eta_{\mathrm B}(F)$ on a fixed root space is
not such a uniform theorem.

\section{Candidate changes, exact matching, and their scope}\label{sec:matching}
\subsection*{Projective candidate invariance is an exact finite identity}
For any candidate $k$ with $a=\eta_\partial(k)\ne0$, parity is not
needed for the full commutator identity
\[
 W\frac{Dk}{a}=\frac{DWk}{a}-\beta
             +\frac{\beta^*(k)}{a}\eta_\partial.
\]
When paired with a fixed test $v$ having $\eta_\partial(v)=0$, the
last term vanishes. For two such candidates,
\[
 \ip{W(Dk_1/a_1-Dk_2/a_2)}{v}
 =\ip{DWk_1/a_1-DWk_2/a_2}{v}.
\]
Consequently a candidate change cannot alter that fixed paired mismatch
when both normalized G1 pairings tend to zero. This does not assert a
nonzero limit, and it does not cover changes to the tests, matrix, or
sampler. Replacing the prolate candidate by an endpoint shell preserves
this exact identity but does not import a disputed G2 lower limit.

\subsection*{An exactly matching sampler}
On an algebraic one-sided root space put
\[
 R_F(s)=F(s)/\zeta(s),\qquad
 R_{BF}(s)=(s-1/2)R_F(s)-\eta_{\mathrm B}(F).
\]
For an even finite candidate $k$ with physical endpoint $B_k\ne0$,
let $z^{(k)}=iDk/B_k$ and define
\begin{equation}\label{eq:v1-exactmatcher}
 [J^{(k)}F]_n=[z^{(k)}]_nR_F(s_n).
\end{equation}
Multiplication of the rational graph equation by $[z^{(k)}]_n$ proves
\[
 iDJ^{(k)}-J^{(k)}B=z^{(k)}\eta_{\mathrm B}.
\]
With $x=J^{(k)}F$ and $q=i\eta_\partial(x)/\eta_\partial(k)$,
\[
 \eta_{\mathrm B}(F)z^{(k)}-qDk
 =\frac{iDk}{B_k}\bigl(\eta_{\mathrm B}(F)-\delta(x)\bigr).
\]
This removes a vector mismatch algebraically. Endpoint recovery is only
one remaining requirement: injectivity, uniform ordinary norm bounds,
Weil transfer, and weak G1 on these tests also need proof. No coercivity
follows from the displayed coefficient identity alone.

\subsection*{Endpoint recovery is a periodic Laplace problem}
Let $k$ be a trigonometric polynomial on $[0,L]$, with
$k(0)=k(L)=B_k$, and $\Re\alpha>0$. Put
\[
 E_{L,k}(\alpha)=\frac{\alpha}{B_k(1-e^{-\alpha L})}
                   \int_0^Le^{-\alpha y}k(y)\,dy.
\]
Solving $(\partial_y-\alpha)x=k'/B_k$ with periodic boundary condition
gives $x(0)=1-E_{L,k}(\alpha)$: multiply the equation by
$e^{-\alpha y}$, integrate, and integrate its right side by parts.
Since $\partial_\alpha(\partial_y-\alpha)^{-1}
=(\partial_y-\alpha)^{-2}$, higher Jordan coordinates satisfy
\[
 \eta_{\mathrm B}(F_{\rho,j})-\delta(J^{(k)}F_{\rho,j})
 =\frac{\partial_\alpha^{j-1}E_{L,k}(\alpha)}{(j-1)!},
 \qquad \rho=1/2+\alpha.
\]
This formula is exact for each finite candidate.

For the endpoint shell $k=c_{L,K}$ of \eqref{v1:eq:shell}, the
coefficient of the first adapted root is
$\sqrt L(2K)^{-1}it_n/(it_n-\alpha)$ on its $2K$ modes.
If $K/L\to\infty$, its physical endpoint tends to one while its squared
ordinary norm is $O(L/K)\to0$. Higher root endpoints have the corresponding
vanishing limits and their ordinary masses also vanish. Thus near-perfect
endpoint matching can coexist with collapse of the whole sampled norm.
The shell exact matcher is not a coercive transport.

\subsection*{Global divisibility and endpoint-jet expansions}
For a smooth exact-moment source $h$ supported up to $\lambda$, the
Mellin identity is $M_0(Eh)=\zeta M_0h$. Set
$a=\log\lambda$, $\varphi(y)=e^{y/2}h(e^y)$ for $y\le a$, and
$B^+=\varphi(a)\ne0$. With the positive Mellin convention used in
the Sonine sections,
\[
 it\,M_0h(1/2+it)/B^+
 =e^{ita}-\frac1{B^+}\int_{-\infty}^a\varphi'(y)e^{ity}\,dy.
\]
This is integration by parts. The negative Fourier convention of
Section~\ref{sec:sampling} represents the same formula after $y\mapsto-y$;
the signs must be changed together. If $\varphi'\in L^1$, the bulk
integral tends to zero for fixed $\lambda$ and $|t|\to\infty$.
Thus the global zeta quotient cancels but a nondecaying endpoint character
remains. If $\varphi'\in L^2$, Plancherel gives bulk squared norm
$\|\varphi'\|_2^2/|B^+|^2$.

For $r\ge1$, assuming one-sided derivatives, vanishing lower traces,
and $\varphi^{(r)}\in L^1$, repeated integration gives
\[
 it\,M_0h(1/2+it)/B^+
 =e^{ita}\sum_{j=0}^{r-1}
   \frac{(-1)^j\varphi^{(j)}(a)}{B^+(it)^j}
   +O\!\left(\frac{\|\varphi^{(r)}\|_1}
                       {|B^+||t|^{r-1}}\right).
\]
This is a fixed-cutoff high-frequency expansion. The terms $(it)^{-j}$
cannot be subtracted through $t=0$ without a low-frequency prescription.
On the lattice $t_n=2\pi n/L$, $L=2a$, the character is $(-1)^n$.
Its appearance is consistent with the centered/unphased distinction and
does not license treating the two samplers as one vector.

\subsection*{Shell estimates and the limited scalar-shift obstruction}
The corrected logarithmic diagonal is sufficient for shell weak-pairing
estimates. For the explicit shell and a bounded ordinary-norm test,
\[
 |\ip{Wc_{L,K}}{v}|
 \le C\lambda(1+\log(4K+1))\sqrt{L/K}\,\|v\|.
\]
For a Hardy test with adequate weighted coefficient decay, the weighted
convolution estimate controls $\ip{DWc}{v}=\ip{Wc}{Dv}$ similarly.
One may enlarge a common ambient cutoff to make these shell pairings
small, but \eqref{v1:eq:paireddefect} still contains its separate beta
term. The old shell calculation is not a proof of graph compatibility.

Finally a scalar-shift shortcut has a limited, elementary obstruction.
If on the same normalized candidate a full residual is $o(1)$, while a
proposed shift has $|\alpha|\ge\alpha_0>0$ and normalized derivative
norm bounded below, the reverse triangle inequality prevents the shifted
full residual from also tending to zero. Weak fixed-band G1 supplies
neither of those full-norm premises. The exact cancellation in
Proposition~\ref{prop:v1-candidate} is the appropriate statement for the
weaker paired argument.

\section{The remaining fixed-space criterion and research goals}\label{sec:goals}
The original relative boundary criterion is unchanged. If, for every
fixed nontrivial zero and prime $p$,
\[
 \frac{\|P_RD_p(v_\rho^{pR}-v_\rho^R)\|}{\|v_\rho^R\|}
 \longrightarrow0,
\]
then contractivity of $T_{R,p}$ and \eqref{eq:boundary} give
$|p^{1/2-\rho}|\le1$. Reflection gives RH. The critical-line rate
derived from Assumption~\ref{ass:evaluators} is consistent with this
criterion; it does not establish the off-line estimate. Equivalently,
the exact diagonal defect asks for an independent arithmetic sign on
each zero evaluator. Full positivity on their span is an unnecessarily
strong sufficient condition.
Proposition~\ref{prop:v110-normalized-shell} now removes a technical
prerequisite from this alternative: the Fredholm remainder and its signed
moment correction vanish after evaluator normalization without invoking
Assumption~\ref{ass:evaluators}. The precise outstanding target is
\eqref{eq:v110-sign-target} at the single prime $2$. The elementary
off-line upper bound suffices for this reduction; a new arithmetic
argument must still orient the leading shell imbalance. The reduction
does not infer that sign from the already known prime eigenvalue.
Proposition~\ref{prop:v111-adjoint} now writes the arithmetic constraint
as the local distributional equation $\mathscr Ay=A_y+B_y/t$ on the
expanded source interval. The integral-corrected cutoff converges with
the explicit bound in \eqref{eq:v111-adjoint-error}; the raw cutoff can
instead diverge in norm. Corollary~\ref{cor:v111-gram} justifies fixed
finite source-Gram calculations, not inversion of a limiting source
Gram operator. The next step is a signed quadratic estimate exploiting
this equation together with both Sonine support conditions. Neither
linear source orthogonality alone nor the power response
\eqref{eq:v111-power-response} supplies that estimate.

The weak G1 source chain is now closed: Theorem~\ref{thm:v14-radical}
supplies compact-source radicality and the semilocal domain passage;
Proposition~\ref{prop:v12-source} and Corollary~\ref{cor:v12-repair}
give the normalized source and exact repair; the subsequent tail,
endpoint and correlation estimates give \eqref{eq:G1-continuum} and
\eqref{eq:G1-sobolev}. The deterministic finite transfer in
Theorem~\ref{thm:g1-boundary-corrected} is consequently unconditional
for this source and a fixed physical test band.

Three distinct research tasks remain for the finite-Weil route. First,
seek genuinely uniform estimates for a growing deep prolate block, or
the separately stated continuous-remainder improvement needed for an
economical near-Slepian cutoff. Neither is required for the proved weak
G1 limit, and neither follows from its fixed-source domain passage.
Second, estimate the
corrected graph pairings on a new, fully defined reflected test family.
The finite object is fixed by \eqref{v1:eq:repair}, but the former
bounded endpoint ratio is ruled out by Corollary~\ref{cor:v11-ratio}.
Theorem~\ref{thm:v11-fixed-tests} further shows that bounded resolvent
tests from any fixed finite-dimensional input space lose the reflected
cross signal under exact endpoint annihilation. A viable replacement must
address that obstruction, for example through a changing test space or
a direct endpoint correction with its additional Weil pairing controlled.
For such a replacement, the full graph source, the relevant meromorphic
transform errors, and any entire zero-sum estimate remain obligations.
Third, prove a positive
metric comparison whose Lyapunov defect tends to zero independently of
the assumed off-line roots. Restricted scalar coercivity is already a
consequence of form transfer and cannot serve as that independent input.
Section 19 now sharpens this task: Theorem~\ref{thm:v15-finite-metric}
supplies a finite-core criterion, and Proposition~\ref{prop:v15-projector}
constructs a positive, uniformly coercive metric that kills the zeta source
and endpoint-shell graph terms exactly. Its remaining commutator has the
strictly positive variance limit in Proposition~\ref{prop:v15-variance};
no independent vanishing signed pairing follows. Removing more source
derivatives by bounded projections loses coercivity by
Proposition~\ref{prop:v15-krylov}. A sharp fixed-band Weil correction
faces a separate transfer cost: Proposition~\ref{prop:v16-projected-cut}
proves exponential off-line transform growth after source projection.
Proposition~\ref{prop:v17-pair-edge} computes the reflected-pair term:
unless its explicit endpoint product vanishes, it remains exponentially
large and oscillatory. Any proposed cancellation must therefore be checked
over the entire zero sum or on the actual finite prime/pole/archimedean
form; bounded ordinary norms and pairwise reflection do not supply it.
There is now a transfer-preserving alternative:
Proposition~\ref{prop:v18-full-projection} removes the source and the
two shell directions on the full expanding Fourier space. It proves a
uniform summable strip-product estimate, fixed-core coercivity and the
unchanged global Weil pullback. The next metric target is its remaining
signed commutator, or an independently justified arithmetic correction
to it. Reusing the same projected Weil form plus a fixed scalar is
excluded by Corollary~\ref{cor:v18-shift}: the Weil pullback tends to
zero on the one-sided root space and cannot cancel the nonzero scalar
metric budget. A proposed correction must be tested against the full
relative criterion before any positivity claim is made.
Proposition~\ref{prop:v19-action} now supplies a stronger diagnostic:
at tail rate $p>1$, the actual finite Weil action converges in norm to
the projected finite exponential sum of zero values. Consequently
the positive square $SW^2S$ has the top-root growth and lower-jet
collapse in Corollary~\ref{cor:v19-square}; it does not furnish the
missing signed estimate. Further metrics should be checked for both
the relative defect and multiplicity-sensitive coercivity, with any
scale change accounted for explicitly.

The fixed-space route has separate foundational obligations: prove the
quantitative evaluator estimates and give a complete closed-operator
realization if the Riesz-window formulation is retained. These are not
prerequisites for the finite-core exclusion criterion in
Theorem~\ref{thm:v15-finite-metric}, nor for the direct complete-Weil
route below. Bypassing them does not establish them for the Riesz
formulation. The shell-sign, uniform metric, large-support
spectral-overlap, and Schur-complement formulations are alternative
targets; none has been shown to follow from weak G1.

\begin{proposition}[The direct cofinal Weil route]\label{prop:v121-cofinal-rh}
Suppose $\lambda_j\to\infty$ and the complete, unshifted semilocal
Weil forms satisfy
\[
 QW_{\lambda_j}(f,f)\ge-\epsilon_j\norm f^2
 \quad(f\in\mathcal D(QW_{\lambda_j})),
 \qquad \epsilon_j\ge0,\quad\epsilon_j\to0.
\]
Then the Weil form is nonnegative on every smooth compactly supported
test function, and RH follows by Weil's positivity criterion.
\end{proposition}
\begin{proof}
Fix such a test $f$. For all sufficiently large $j$, its multiplicative
support lies in $[\lambda_j^{-1},\lambda_j]$. Extension by zero
preserves both its ordinary norm and the global Weil value, so
\[
 QW(f,f)=QW_{\lambda_j}(f,f)\ge-\epsilon_j\norm f^2.
\]
Sending $j$ to infinity proves nonnegativity. Apply the compact-support
Weil criterion recalled in \cite{ConnesConsaniMoscovici2025}.
\end{proof}

Consequently the full cofinal conclusion of
Proposition~\ref{prop:v120-structured-uniform}, if its hypotheses and
vanishing error are proved, would already yield RH by this direct route.
The conditional Riesz realization and evaluator assumptions belong to
other routes and are not additional premises of
Proposition~\ref{prop:v121-cofinal-rh}. A fixed-window certificate,
a bound on a finite compression alone, or an incomplete transport
estimate does not meet this proposition. The implication isolates the
strength of the missing estimate; it does not establish that estimate.



Proposition~\ref{prop:v112-schur-angle} isolates an inverse-weighted
residual that can be far smaller than the full residual divided by the
smallest gap. Proposition~\ref{prop:v112-finite-certificate} certifies
one finite case. Proposition~\ref{prop:v113-source-certificate} now
identifies the exact projected source there, including its angular tail,
and proves a relative physical-endpoint mismatch greater than one.
Proposition~\ref{prop:v114-polynomial-endpoint} now recovers the exact
source endpoint, with relative error $O(\lambda^{-1})$ at the unchanged
polynomial Fourier cutoff. Corollary~\ref{cor:v114-polynomial-g1}
puts weak G1 on this same finite vector, uniformly over the ambient
Fourier cutoff in the weighted test norm. These results supply no
practical certified threshold
at $\lambda=3$, and no endpoint comparison to an actual ground vector.
Proposition~\ref{prop:v114-weil-core} closes the semilocal Weil
operator-core issue and establishes fixed-window graph convergence for
BV sources. This is distinct from the open arithmetic realization in
Section~18. Proposition~\ref{prop:v115-sharp-tail} sharpens the omitted-tail
estimate to $N=256$ at $\lambda=3$. Proposition~\ref{prop:v115-absolute-residual}
additionally proves a small full ordinary residual and Rayleigh value
for the actual source and its polynomial projection. It gives no
endpoint-normalized operator estimate and no lowest-state overlap.
Proposition~\ref{prop:v116-window-positive} now closes the signed
head/complement test at $\lambda=3$ for both parity sectors and the
entire closed form. Positive weighted prime bounds and the increasing
tail diagonal are essential to its verified inverse-action estimate;
the residual's directional moments control every omitted row.
Proposition~\ref{prop:v117-ground-order} now proves that the complete
operator at this window has a simple even ground, with
$0<\mu_0<3.644\cdot10^{-38}$ and $\mu_1>10^{-36}$.
Corollary~\ref{cor:v117-ground-overlap} bounds the ordinary sine angle
of the exact normalized $P_{64}p_3$ by $0.01931$. The stronger
$10^{-34}$ threshold for all other even states is essential here.
The earlier $N=64$ finite-matrix angle bound and this complete-operator
bound are different statements.

The remaining G2 spectral task is a complete lower bound tending to
zero along unbounded windows, or sufficient quantitative source-to-ground
overlap on such a family. Proposition~\ref{prop:v121-cofinal-rh}
shows why this is RH-strength: even a cofinal sequence of complete
lower bounds $-o(1)$ implies nonnegativity of every compact test.
A finite list of certified windows cannot replace that estimate.
Proposition~\ref{prop:v118-gap-free} removes an unnecessary positive-gap
requirement from one route. The already proved rank-one source residual
reduces the target to asymptotic nonnegativity of its entire orthogonal
complement. A negative compact test would persist in that complement;
the required arithmetic sign is still unproved. Growing source blocks
also need an independent uniform operator-residual bound.
Proposition~\ref{prop:v121-complement-cancellation} rules out the
claim that substantial cancellation is confined to the single source:
an exactly source-orthogonal finite test has non-prime energy greater
than $0.4$ but positive total energy below $3\cdot10^{-31}$.
This does not disprove complementary positivity; it prevents using
basis or random-vector samples as a uniform dominance argument.

The unsigned-prime shortcut is now excluded by
Propositions~\ref{prop:v119-prime-essential}--\ref{prop:v119-prime-scale}:
every finite-dimensional source or Fourier-head removal retains the
full prime norm, asymptotic to $\lambda$. The rank-two pole term cannot
make their combined ordinary operator norm smaller either. The next
estimate must therefore use the archimedean energy, directional inverse
weights or signed correlations; enlarging a finite source block alone
cannot make that unweighted arithmetic norm small. This is an
obstruction to a sufficient norm bound, not to full Weil positivity.

Proposition~\ref{prop:v120-lambda4-tail} now proves a signed lower
bound on the entire $\lambda=4$ tail $|n|>16$, retaining its exact
archimedean diagonal. Both complex parity sectors are included.
The effective head $F-B^*T^{-1}B$ has $33$ dimensions.
Proposition~\ref{prop:v125-even-complete} certifies its entire
$17$-dimensional even sector and hence the complete even form.
The $16$-dimensional odd sector is still a fixed-window target. Proposition~\ref{prop:v121-alpha-zero} now certifies the
raw head and specified finite-support trial forms $K_X$ positive in
both parity sectors. This gives local $\alpha_4=0$, without an
infinite-tail transfer for $K_X$. It still does not suffice for the
exact Schur sign because the nonnegative residual inverse correction
must be subtracted, as in \eqref{eq:v121-galerkin-defect}.
This leaves a separate growing-window target: a corresponding ordinary
lower error tending to zero. Proposition~\ref{prop:v120-weighted-error}
shows that a dimensionless weighted error must be multiplied by its
arithmetic scale before it can serve that purpose.

Proposition~\ref{prop:v120-structured-inverse} now reuses the signed
factorization to bound the actual inverse-tail correction, including
all remote rows through \eqref{eq:v120-structured-finite-rows}.
Proposition~\ref{prop:v120-structured-uniform} identifies the sufficient
ordinary error $\alpha_\lambda+e_\lambda^2+
 t_\lambda^2/\mu_\lambda$. Its operator norms concern the entire
growing head; individual source or column bounds do not suffice.
Proposition~\ref{prop:v120-inverse-polynomial} additionally gives
convergent upper bounds for the complete far inverse, with certified
constants at $\lambda=4$, sharpened to $20$ and $90$ by
Proposition~\ref{prop:v122-parity-inverse}. The complete far energy
now fits below the trial budget on one certified head direction,
by Proposition~\ref{prop:v122-directional-far}. This is not the full
inverse correction. Proposition~\ref{prop:v123-first-majorant-fails}
now proves that its first-polynomial scalar-head bound fails after
the complete mixed term is included. Corollary~\ref{cor:v123-updated-first-fails}
also excludes its exact positive head update with the saved initial
certificate. Proposition~\ref{prop:v124-positive-direction}
proves a complete positive value on the saved head direction.
Proposition~\ref{prop:v125-even-complete} extends this to the entire
even head through frozen energy-scaled trials and a simultaneous
complete Gram bound. The odd-head trials tested so far remain
inconclusive; they do not give a negative Weil vector. The primary
next target is the cofinal ordinary-error estimate, rather than a
longer finite list of certified windows.

Proposition~\ref{prop:v125-cutoff-cost} identifies a structural cost:
with the present scalar far majorant, its inner cutoff must grow
exponentially in $\lambda$. The fresh threshold table gives the
window dependence explicitly. A viable family must either control
the associated head dimensions, conditioning and residual Grams at
that cost, or improve the signed far estimate before constructing
the growing inner heads. Finite-only scaling diagnostics in the
research log do not provide the missing complete witnesses.

The trial-head hypothesis is not yet verified along an unbounded
window family. With $\alpha_\lambda=0$, the remaining target is
$e_\lambda\to0$ and $t_\lambda/\sqrt{\mu_\lambda}\to0$, retaining
all ordinary-energy and dimension factors. Inverse approximation
convergence alone does not establish those limits. The complete
even sign at one window closes a local obligation, not G2 or RH.

The successful one-sided bound does not require a norm contraction:
Proposition~\ref{prop:v120-norm-counterexample} certifies that the
latter fails at the very same cutoff. Nor can the whole weighted
remainder be estimated by an infinite Hilbert--Schmidt sum, by
Proposition~\ref{prop:v120-not-schatten}. Finite-column residual Grams
remain valid. The logarithmic weighted-tail asymptotic in
Proposition~\ref{prop:v120-weighted-tail-asymptotic} is strictly a
fixed-window statement and supplies no uniform joint-limit obstruction.

The continuum endpoint is now settled by
Theorem~\ref{thm:v118-eigen-endpoint-zero}: every complete eigenfunction
has continuous zero physical trace. Hence a nonzero endpoint comparison
to the repaired source is false, with exact relative mismatch one.
Corollary~\ref{cor:v118-ground-fejer} supplies a fixed-window filtered
endpoint-decay rate. It gives no sharp-cut or joint growing-window rate
for finite-compression ground vectors. Those finite endpoint estimates,
the sampler graph defect and full-strip normalization remain separate.
The next spectral work should retain these distinctions while seeking
an independent signed complementary lower bound, rather than imposing
a nonzero continuum ground trace. G2 and RH remain open.


\paragraph{What is established by the sampler correction.}
The uncentered full-scale transfer in the previous version cannot be used
under the stated nonzero-endpoint hypotheses. Centering repairs the form
transfer; an explicit shell then restores the original physical endpoint
on the same finite matrix. Ordinary norm coercivity and the reflected
top-root cross form survive. The graph identity gains an explicit defect
with unavoidable derivative growth. Strong graph convergence fails for
norm-vanishing endpoint repairs, but that fact alone does not settle the
weaker arithmetic pairing needed for the paper's objective. The new
resolvent calculation identifies a second obstruction: on a fixed test
space the rescaled endpoint converges to minus the reflected evaluation,
which is exactly the evaluation needed for the nonzero cross form.
Killing that endpoint with bounded coefficients therefore kills that
particular G2 signal.

\paragraph{Claim discipline.}
This manuscript claims neither RH nor a completed positivity transfer.
The pole-aware transfer and the earlier nonzero critical beta limit have
been removed from the established chain. The moving-phase exclusion is
restricted to its proved algebra, pending relative alias estimates.
Unspecified earlier source-quasimode, partner-blind, and certified
nonorthogonality calculations are not used as theorems. Numerical checks
below include exploratory illustrations, a rigorous finite-matrix
certificate and its exact-source identification, complete fixed-window
sign certificates, finite-trial positivity and an exact-source
cancellation certificate. These do not settle the remaining uniform
analytic obligations. A publication draft would need the listed inputs
proved or its scope narrowed to the results that are already supported.

\appendix
\section{Reproducible finite checks}\label{app:numerics}

The accompanying file \texttt{check\_sampler.py} constructs the real symmetric CCM matrix with the
separate logarithmic diagonal, using prime-power sums and numerical
quadrature only for the smooth archimedean integrals. The test is
$H(s)=\kappa(s)\zeta(s)$ with $\sigma=2$, $a=4$, $M=6$. Its global
zero-side form is exactly zero and $E=2\zeta(2)=3.289868133696453$.
No hypothetical zero location is inserted. At a fixed cutoff $N=512$:
\begin{center}
\begin{tabular}{rrrr}
\toprule
$L$ & Uncentered form & Centered form & Centered endpoint\\
\midrule
4 & 0.763119 & $1.174\,10^{-3}$ & 0.129293\\
6 & 1.036652 & $1.085\,10^{-4}$ & 0.033565\\
8 & 0.669071 & $4.000\,10^{-6}$ & 0.008042\\
10 & 1.337515 & $4.078\,10^{-7}$ & 0.001816\\
\bottomrule
\end{tabular}
\end{center}
For $L=6$, endpoint correction is exact to the displayed floating-point
precision, while its form price decreases and derivative price increases:
\begin{center}
\begin{tabular}{rrrr}
\toprule
$N=2K$ & Corrected form & Correction norm & Derivative norm\\
\midrule
128 & 1.400023 & 0.705010 & 72.5381\\
256 & 0.906192 & 0.498518 & 102.3284\\
512 & 0.492218 & 0.352505 & 144.5331\\
\bottomrule
\end{tabular}
\end{center}
These moderate cutoffs do not yet make the corrected form small. The
existence and rates in Theorem~\ref{v1:thm:repair} come from the proof, not
from extrapolating this table. Increasing the quadrature order on a
representative matrix changed entries by less than $2\,10^{-12}$.
The script uses 30-digit zeta evaluations followed by double-precision
matrix arithmetic; the checks are not interval-certified bounds.

\subsection*{A rational check of the resolvent endpoint asymptotic}
The file \texttt{check\_resolvent\_endpoints.py} tests the sign, half-period
phase, shell contribution and residue scale in
Lemma~\ref{lem:v11-resolvent-endpoint} using explicit rational functions:
\[
 H_-(s)=\frac{(\sigma+a)^M}{(\sigma-s)(s+a)^M},
 \qquad H_+(s)=(s-p)H_-(s).
\]
Here $\sigma=2$, $a=4$, $M=6$, $\mu=0.2+1.3i$ and
$p=1/2-\overline\mu$. These are model transforms, not zeta root
functions. Their Cauchy tail coefficients are $E_-=1$ and
$E_+=\sigma-p$. The centered samples, their adjoint resolvents and
the endpoint shell correction are summed explicitly in Fourier pairs.
For numerical tractability this check uses
$K=\lceil e^{L/2}L^2\rceil$, much smaller than the theorem cutoff.
It does not compute a Weil form.

Write $e_\pm=\delta V_LH_\pm$, interpreting $V_L$ here as the same
sampling and resolvent operations applied directly to these transforms.
The predicted reflected limit is
\[
 S_L:=\frac{e^{\overline\mu L/2}(1-e^{-\overline\mu L})e_-}
                  {-H_-(p)}\longrightarrow1.
\]
The output, rounded from double-precision arithmetic, is:
\begin{center}
\begin{tabular}{rrrr}
\toprule
$L$ & $K$ & $|S_L-1|$ & $|e_-/e_+|$\\
\midrule
4  & 119    & $7.347\,10^{-2}$ & 4.426\\
6  & 724    & $5.792\,10^{-3}$ & 35.261\\
8  & 3495   & $9.638\,10^{-4}$ & 282.874\\
10 & 14842  & $2.276\,10^{-4}$ & 1869.438\\
12 & 58094  & $7.847\,10^{-5}$ & 5886.940\\
14 & 214941 & $1.872\,10^{-5}$ & 24924.379\\
\bottomrule
\end{tabular}
\end{center}
This is a reproducible check of the analytic mechanism and conventions,
not an interval certificate, a verification on hypothetical zeta zeros,
or evidence for RH. The theorem and its application to genuine root
functions depend on the contour and uniformity arguments in the text.
\subsection*{A high-precision check of the normalized prolate source}
The script \texttt{check\_prolate\_source.py} constructs the even
Legendre--Galerkin matrix for
$-\partial_x((1-x^2)\partial_x)+c^2x^2$ on $[-1,1]$.
It takes the first and third even eigenvectors, corresponding to indices
$0$ and $4$, and normalizes them in $L^2(-1,1)$. The compressed Fourier
eigenvalue is recovered from the ratio of the integral to the central
value, with the factor $\sqrt{c/(2\pi)}$ retained. Thus the source
coefficients are those in \eqref{eq:v12-coefficients}, not an imposed
endpoint fit.

Let $S_\lambda$ denote the noncancellation ratio on the left of
\eqref{eq:v12-noncancellation}, and let
$Z_\lambda=|h_\lambda^\circ(0)|/|B_\lambda^+|$. With 50 even Legendre
modes and 85-decimal-digit arithmetic the rounded output is:
\begin{center}
\begin{tabular}{rrrr}
\toprule
$c$ & $b_4/(c^2b_0)$ & $S_\lambda$ & $Z_\lambda$\\
\midrule
10 & 6.40490  & 1.00164349 & $1.92330\,10^{-2}$\\
20 & 9.94110  & 1.00029407 & $4.00922\,10^{-6}$\\
30 & 10.99058 & 1.00012034 & $3.99136\,10^{-10}$\\
40 & 11.51074 & 1.00006514 & $3.10538\,10^{-14}$\\
\bottomrule
\end{tabular}
\end{center}
Repeating $c=40$ with 65 modes and 105-digit arithmetic preserves the
displayed figures. The script also reports the concentration loss divided
by Fuchs's leading term and the ratio $b_n^2/[c(1-\vartheta_n)]$.
Increased precision is important because concentration losses become
exponentially small. This is a finite spectral approximation with no
certified truncation or roundoff bound; the analytic proof, not this
table, establishes the asymptotic estimates. No zeta zero is inserted.

\subsection*{A compact-source primitive convention check}
The script \texttt{check\_bv\_radical\_passage.py} tests the powers and
sign in the primitive mechanism of Theorem~\ref{thm:v14-radical}.
Take $h(u)=u^2-\frac53u^4$ on $[-1,1]$, with zero extension. It has
the two exact moments and $h(1^-)=-2/3\ne0$. If $N=\lfloor1/u\rfloor$
and $S_j(N)=\sum_{n=1}^N n^j$, an explicit logarithmic primitive is
\[
 P(u)=\frac25u^{5/2}S_2(N)-\frac{10}{27}u^{9/2}S_4(N)
       +\frac4{135}\zeta(\tfrac12,N+1),
\]
where the last function is the Hurwitz zeta function. Away from a cut,
$uP'(u)=\mathcal E(h)(u)$; the jump cancellation in this expression
agrees with continuity of the primitive. At 70-digit working precision,
testing $1/u=N+\theta$ for $N=10,100,1000,10000$ and
$\theta=0.2,0.5,0.8$ gives maximum derivative discrepancy less than
$3.3\,10^{-69}$. The sampled values of $P(u)/u^{3/2}$ remain bounded.
This is a toy-source convention check, not an interval certificate or
a numerical proof of the uniform primitive bound, radicality, G1, or RH.
Those conclusions use the analytic arguments in the text.


\subsection*{Finite-core metric and projection checks}
The script \texttt{check\_section19\_metric.py} checks the exact budget
\eqref{eq:v15-budget}, source/shell annihilation,
\eqref{eq:v15-variance}, and the new projected graph source
\eqref{eq:v15-projected-graph}. It uses the rational source samples
\[
 z_n=\frac{(-1)^n}{\sqrt L(2+it_n)^3},\qquad t_n=2\pi n/L,
\]
with a length-two Jordan chain at $\mu=0.23+0.7i$ and a further
eigenvector at $\nu=0.37+1.4i$. These parameters are model eigenvalues,
not zeta zeros. The tests include a complex endpoint-shell correction and
the first-slot-linear convention. At $T=3$ the rounded results are
\begin{center}
\begin{tabular}{rrrr}
\toprule
$L$ & $\lambda_{\min}(J^*SJ)$ & $\|[D,S]\|$
 & $|\mathcal L_{11}|/G_{11}$\\
\midrule
8 & $2.55980\,10^{-4}$ & 0.976536 & 0.460000\\
16 & $8.40732\,10^{-4}$ & 0.997413 & 0.460000\\
32 & $8.95576\,10^{-4}$ & 1.007014 & 0.460000\\
\bottomrule
\end{tabular}
\end{center}
The nine algebra-check discrepancies are below $2.1\,10^{-15}$;
the source-polynomial remainder identity is also checked on a grid.
The final column equals $2\Re\mu$, as the exact identity predicts:
source annihilation does not make the relative Lyapunov defect small.
No Weil matrix or zeta zeros enter this model. The calculations are
floating-point convention checks, not positivity certificates or proofs
of any limiting estimate; Section 19 supplies the analytic arguments.

\subsection*{Sharp-band transform convention checks}
The script \texttt{check\_band\_cut.py} uses the rational source
$Z(t)=(2+it)^{-3}$ and $R(t)=(it-\mu)^{-1}$ with
$\mu=0.23+0.7i$. It subtracts the exact sampled source projection,
sets $T=3$, and takes $L=2\pi N/3$, $N=8,16,32,64,128,256$.
The second-difference identity in Lemma~\ref{lem:v16-band-transform}
is checked to error below $3\,10^{-15}$, and direct physical integration
at the smallest size agrees to below $10^{-14}$.
At $N=128$ the ordinary squared norm is approximately $0.0156032$,
while the transform magnitudes at $s=1$ and $s=0$ are approximately
$1.088\,10^{24}$ and $6.525\,10^{23}$, respectively. These large
values are consistent with the proved cut mechanism, not evidence about
zeta zeros or the sign of a Weil form. The experiment evaluates no Weil
matrix and supplies no certified numerical error bound. It also checks
the transform product at $v=\mu$ and $v^\sharp=-\overline\mu$ against
the edge expression in Proposition~\ref{prop:v17-pair-edge}. At $N=256$
the relative discrepancy is approximately $0.0148$ and the real pair sum
is approximately $1.414\,10^{43}$; this is a rational-model convention
check, not a zeta-zero or whole-Weil-form computation.

\subsection*{Full-space projection and holomorphic strip checks}
The file \texttt{check\_full\_projection.py} tests
Proposition~\ref{prop:v18-full-projection} with the rational model
\[
 Z(v)=\frac{(v-\mu)^2(v-\nu)}{(v+2)^8(2-v)},\qquad
 \mu=0.23+0.7i,\quad\nu=0.37+1.4i.
\]
The model root outputs $Z(v)/(v-\mu)$, $Z(v)/(v-\mu)^2$ and
$Z(v)/(v-\nu)$ are holomorphic on the closed critical strip;
no meromorphic-output assumption is hidden in this check. At
$L=8,12,16,24,32$ it uses the tractable cutoff $K=4L^2$, not the
theorem's exponential cutoff. Eight finite algebra checks, including
the complete graph source, both shell directions and the commutator
budget, have discrepancies below $4.2\,10^{-14}$.
On a 25-point grid in the strip, the largest ordinate-weighted transform
error decreases from approximately $0.02043$ at $L=8$ to
$4.121\,10^{-6}$ at $L=32$. The model first-eigenvector relative
Lyapunov defect remains $0.46=2\Re\mu$. No zeta zero or Weil matrix
enters this experiment. These are floating-point convention checks,
not uniform strip certificates or evidence for RH.

\subsection*{An arithmetic-matrix check of the Weil action}
The script \texttt{check\_weil\_action.py} uses the actual finite
arithmetic matrix from \texttt{check\_sampler.py}, including the
separate logarithmic diagonal, and the first known critical-line zero
$\rho_1\simeq1/2+14.1347251417i$. The transform is
$\kappa(s)\zeta(s)/(s-\rho_1)$ with $\sigma=2$, $a=4$, $M=6$.
The root location and derivative are evaluated at 40-digit precision,
then samples and matrices use double precision. This is a numerical
simple-root check, not a new multiplicity certificate. After the
full-space projection, it compares $WSJF$ with the exact finite
Fourier projection of
$\kappa(\rho_1)\zeta'(\rho_1)e^{(\rho_1-1/2)y}$.
At the moderate cutoff $K=256$, $N=2K=512$:
\begin{center}
\begin{tabular}{rrr}
\toprule
$L$ & Absolute action error & Relative action error\\
\midrule
4 & $7.9105\,10^{-4}$ & $5.7219\,10^{-2}$\\
6 & $2.3776\,10^{-4}$ & $1.4045\,10^{-2}$\\
8 & $4.1955\,10^{-5}$ & $2.1459\,10^{-3}$\\
10 & $1.5078\,10^{-5}$ & $6.8992\,10^{-4}$\\
\bottomrule
\end{tabular}
\end{center}
The script also checks $K=128$, the target Fourier coefficient by
independent physical quadrature, and the last matrix action with 400
additional archimedean quadrature nodes. These finite cutoffs do not
meet the theorem's exponential rule. No hypothetical off-line zero
is inserted, no interval certification is claimed, and the table is
not evidence for RH. It checks the action's sign, centered phase and
physical normalization using an actual arithmetic matrix.

\subsection*{Signed Fredholm moment check}
The accompanying \texttt{check\_fredholm\_shell.py} uses Gauss--Legendre
Nystr\"om matrices for the actual kernel $2\cos(2\pi xt)$ on $(0,a)$.
It compares the two-channel inverse, the signed $A_a$ formula and the
constant-kernel resummation, using complex test profiles at
$a=0.2,0.1,0.05,0.025$ and two quadrature orders. It separately checks
that reversing one profile reverses the mixed term. The profiles are
arbitrary interval data, not certified Sonine vectors or arithmetic
evaluators. These floating-point checks test signs, scale factors and
the remainder formula only; they do not test the arithmetic shell sign,
certify a uniform estimate, or supply evidence for RH.

\subsection*{Regularized co-Poisson adjoint check}
The script \texttt{check\_copoisson\_adjoint.py} uses
$\phi(t)=(t-1)^3(2-t)^3$ on $[1,2]$, extended by zero, and
$h=\mathcal D\phi$. Both source moments vanish exactly. This toy
source has sufficient piecewise smoothness for the cutoff estimates;
it is not asserted to be a smooth arithmetic-core vector or an evaluator.
At 45-digit precision, independent physical and adjoint integrals are
compared for $y(x)=e^{-x}$ and $M=4,8,16,32$. The weak limit is checked
against its convergent zeta-coefficient series. A separate power-response
check uses $w=0.7+1.1i$, not a zeta zero, and $M$ up to 1024. These
checks include $y(x)=e^{-(1+0.4i)x}$ to test the first-slot conjugation.
The
non-certified computations check counterterm signs and scales, not
the shell inequality or RH. Results are recorded in the supplied JSON.

\subsection*{A certified finite Weil matrix and candidate angle}
\label{app:v112-certificate}
The following finite result is certified by interval arithmetic; the
earlier numerical illustrations remain non-certified. Set
$\lambda=3$, $L=2\log3$, $N=64$ and retain all $129$ indices
$-64\le n\le64$ in the translated $[0,L]$ basis $U_n$.
The matrix $W$ is the same arithmetic Weil matrix used above, with all
prime powers $2,3,4,5,7,8,9$, the pole term and the separately computed
logarithmic diagonal. No list of zeta zeros is used.

The candidate is defined exactly, rather than by an accuracy claim:
the $65$ decimal strings in \texttt{g2\_finite\_candidate.json} are
interpreted as rational numbers $c_0,\ldots,c_{64}$, with
$c_{-n}=c_n$, and $u=c/\norm c$. The file's SHA-256 is
\begin{center}\ttfamily\small
6d326cafbe715752f89b90cd92eb1d5a68\par
ab0b78befb6afe074e797372b1a18e
\end{center}
The two lines concatenate to the identifier. The strings were generated
from a $105$-digit, $70$-even-mode approximation of the repaired source,
but the conclusions below depend only on their exact rational values.
This candidate definition alone does not certify proximity to the true
PSWF coefficients; Appendix~\ref{app:v113-source} supplies that bound.

\begin{proposition}[Computer-assisted finite result]
\label{prop:v112-finite-certificate}
For this exact $W$ and rationally specified normalized candidate,
the complement in Proposition~\ref{prop:v112-schur-angle} satisfies
\[
 C-\alpha I\succ10^{-34}I,\qquad
 \norm{(C-\alpha I)^{-1}r}<0.000216.
\]
The matrix is positive definite and has a simple even ground state, with
\begin{equation}\label{eq:v112-certified-bottom}
 3.64\,10^{-38}<\lambda_0<5.32\,10^{-38},
 \qquad \sin\theta<0.000216.
\end{equation}
\end{proposition}
\begin{proof}[Certificate and verification]
The supplied \texttt{certify\_g2\_finite.py} uses
Python-FLINT~0.9.0 and FLINT~3.6.0 ball arithmetic
\cite{PythonFlint2026} at $512$ bits. It encloses all archimedean
integrals, uses both exact reflection sectors, and constructs the
complement by the Householder transformation determined by the exact
normalized candidate. Orthogonality belongs to that underlying exact
matrix, not to arbitrary independent choices from its entrywise balls.
Interval $LDL^T$ elimination gives $64$ strictly positive pivot
enclosures for each of
\[
 C_{\rm even}-\alpha I-10^{-34}I,
 \qquad W_{\rm odd}-\alpha I-10^{-34}I.
\]
An enclosed LU solve gives $q<0.000216$ and verifies, by strict
interval comparisons,
\[
 \alpha<5.32\,10^{-38},\qquad
 \alpha-\ip{(C-\alpha I)^{-1}r}{r}>3.64\,10^{-38}.
\]
Proposition~\ref{prop:v112-schur-angle} proves the claims.

The independent program \texttt{certify\_g2\_series.py} repeats the
certificate at $768$ bits using the analytic series below, with $128$
terms and explicit tail bounds. It uses a rational, nonorthonormal
complement in the full $129$-dimensional space. For the unnormalized
rational vector $c$, let $s=c^Tc$ and choose $c_k\ne0$. Columns
$Q_j=e_j-(c_j/c_k)e_k$, $j\ne k$, span $c^\perp$ exactly. Put
\[
 G=Q^TQ,\quad H=Q^T(W-\alpha I)Q,\quad b=Q^TWc.
\]
All $128$ interval pivots of $H-10^{-34}G$ are positive.
For the enclosed solution $Hz=b$, the physical quantities are
$q^2=z^TGz/s$ and $E=b^Tz/s$. The Gram matrix $G$ is retained.
This independently verifies the same three strict inequalities.
All $8321$ parity-block entries also have overlapping enclosures in the
two independently assembled matrices. The supplied JSON files retain
the pivot and result intervals, input hash, versions and precision.
No approximate floating-point solve is used for either certificate.
\end{proof}

For reproducibility, write $\rho_\infty(y)=e^{y/2}/(2\sinh y)$,
$t_n=2\pi n/L$ and let $\psi_0,\psi_1$ denote digamma and trigamma.
The first calculation removes the apparent singularity at zero using
$\operatorname{sinc}z=\sin z/z$ and
$\sinh y/y=\operatorname{sinc}(iy)$. The second uses the following
exact formulas for the sine integral and the \emph{complete}
archimedean diagonal:
\begin{align*}
 S_n&:=\int_0^L\rho_\infty(y)\sin(t_ny)\,dy
 =-\tfrac12\Im\psi_0(z_n)
   -\sum_{k\ge0}e^{-a_kL}\frac{t_n}{a_k^2+t_n^2},\\
 A_n^\infty&:=\int_0^L\rho_\infty(y)
  \{2(1-y/L)\cos(t_ny)-2e^{-y/2}\}\,dy\\
 &\hspace{15mm}+\log(4\pi)+\gamma_{\!E}+\log\tanh(L/2)\\
 &=\log\pi-\Re\psi_0(z_n)-\frac{\Re\psi_1(z_n)}{2L}
   +\frac2L\sum_{k\ge0}e^{-a_kL}
          \frac{a_k^2-t_n^2}{(a_k^2+t_n^2)^2},
\end{align*}
where $a_k=2k+\tfrac12$ and $z_n=\tfrac14-it_n/2$.
To verify them, expand $\rho_\infty(y)=\sum_{k\ge0}e^{-a_ky}$.
Integrate the convergent combinations
$\rho_\infty\sin(t_ny)$,
$\rho_\infty(\cos(t_ny)-e^{-y/2})$, and
$y\rho_\infty\cos(t_ny)$ on $(0,\infty)$ and subtract their tails.
The identity $t_nL=2\pi n$ simplifies the endpoint terms. The odd
exponential tail cancels $\log\tanh(L/2)$, while
$\psi_0(1/2)=-\gamma_{\!E}-2\log2$ leaves $\log\pi$.
The digamma difference must not be split into two divergent sums.
After retaining $k<K$, the absolute tail bounds are
\[
 R_S\le\frac{e^{-a_KL}}{2a_K(1-e^{-2L})},\qquad
 R_A\le\frac{2e^{-a_KL}}{La_K^2(1-e^{-2L})}.
\]
These follow from $|t|/(a^2+t^2)\le1/(2a)$ and
$|a^2-t^2|/(a^2+t^2)^2\le1/a^2$. At $\lambda=3$ the weights are
the exact rationals $3^{-4k-1}$; both tails are enclosed symmetrically.

\paragraph{Scope and endpoint limitation.}
This certificate concerns one finite matrix and one explicitly frozen
candidate. It gives no lower bound on the omitted Fourier complement,
no growing-window ordering theorem. Its fixed-source approximation
prerequisite is addressed in Appendix~\ref{app:v113-source}. The preparatory non-certified source pilot gives relative
physical-endpoint errors about $1.375$ at $N=64$ and $0.1046$ at
$N=4096$ for $\lambda=3$; the latter calculation concerns source
coefficients only, not a matrix of that dimension. Small ordinary
ground-state angle therefore cannot substitute for endpoint recovery.
An endpoint-shell correction would still require its derivative and
graph terms. The full-strip transform cost and an independent uniform
arithmetic estimate remain open. Neither G2 nor RH follows from this
finite certificate.


\subsection*{Certified exact-source identification and endpoint error}
\label{app:v113-source}
The additional verifier \texttt{certify\_pswf\_source.py} upgrades the
source comparison at this one window from numerical agreement to a
computer-assisted proof. Its input \texttt{pswf\_rational\_input.json}
contains two $70$-coordinate exact decimal rational vectors and two
exact decimal rational approximate eigenvalues. Its SHA-256 is
\begin{center}\ttfamily\small
ac9dfa8c03f54a326e3ba7af59917ac435e\par
63d41e978310301e6df6e53e83076
\end{center}
The two lines concatenate. Approximate eigensolver output is used only
to select this input; its accuracy is not assumed. Python-FLINT~0.9.0 /
FLINT~3.6.0 at $768$ bits verifies the following steps.

In the normalized Legendre basis,
$Xe_\ell=A_\ell e_{\ell+1}+A_{\ell-1}e_{\ell-1}$ with
$A_\ell=(\ell+1)/\sqrt{(2\ell+1)(2\ell+3)}$ and $A_{-1}=0$.
The even Jacobi diagonal and next diagonal are
\[
 2k(2k+1)+c^2(A_{2k}^2+A_{2k-1}^2),\qquad
 c^2A_{2k}A_{2k+1},\qquad c=18\pi.
\]
The omitted diagonal lower bound is $140\cdot141=19740$.
For mode zero the two bounding matrices in
\eqref{eq:v113-inertia} both have inertia $0$ at $\xi_0-1$ and
$1$ at $\xi_0+1$. For mode four their counts are $2$ and $3$.
Every interval $LDL^T$ pivot has a certified nonzero sign.
Thus $g=1$ is a valid separation from each $\xi_j$ to all other
exact even eigenvalues. The full residuals include the first omitted
link, which dominates here. The following rounded numbers are strict
upper bounds, not error estimates inferred from precision settings:
\begin{center}
\begin{tabular}{@{}lrr@{}}
\toprule
Quantity & Mode $0$ & Mode $4$\\
\midrule
Full residual $\epsilon$ & $5.913\,10^{-56}$ & $4.796\,10^{-52}$\\
Angular $L^2$ error & $8.362\,10^{-56}$ & $6.783\,10^{-52}$\\
Angular uniform error & $2.723\,10^{-52}$ & $2.509\,10^{-48}$\\
\bottomrule
\end{tabular}
\end{center}
The approximate physical central values exceed $1.187$ and $0.716$,
respectively, by interval comparison with their errors, fixing the true
sign convention $h_{j,3}(0)>0$.
The error budget gives
\[
 \eta<1.725\,10^{-48},\quad
 |\widehat h^\circ(0)|<2.611\,10^{-39},\quad
 U<3.368\,10^{-37}.
\]
The exact moment repair is bounded by $\norm\psi_\infty\le129$;
no uncertified evaluation of its defining integral ratio is inserted.

The polynomial Fourier coefficients are also evaluated with enclosures.
If $\widehat h^\circ(v)=\sum_r d_r(v/\lambda)^{2r}$ and
$t_n=2\pi n/L$, the translated coefficients of its inversion average
are exactly
\begin{equation}\label{eq:v113-mellin-polynomial}
 \widehat c_n=\frac1{\sqrt L}\Re\sum_r d_r
 \frac{\displaystyle\sqrt\lambda\sum_{m=1}^{9}m^{-1/2-it_n}
       -\displaystyle\lambda^{-4r-1/2}\sum_{m=1}^{9}m^{2r}}
      {2r+1/2+it_n}.
\end{equation}
To check all scale factors, integrate each monomial on
$-\log\lambda<x<\log(\lambda/m)$ with the sole co-Poisson factor
$e^{x/2}$. Translation from the centered interval contributes
$(-1)^n$, which cancels $e^{\pm it_n\log\lambda}$ at the integration
limits. Taking the real part performs inversion averaging. At $m=9$
the interval has zero length, so its two terms cancel exactly; this
integral endpoint convention differs from the strict trace in
\eqref{eq:v113-endpoint-budget}.

The interval polynomial coefficient vector differs from the frozen
rational vector by less than $6.422\,10^{-40}$ in ordinary norm.
Adding \eqref{eq:v113-projection-budget} gives a bound below
$2.595\,10^{-36}$, in particular the slightly looser strict threshold
$2.60\,10^{-36}$ in Proposition~\ref{prop:v113-source-certificate}.
The normalized projector distance is below $4\,10^{-36}$; no division
by the tiny Weil spectral gap is used in this final source replacement.
The saved interval sine bound is less than $0.000215893$.

Direct Legendre evaluation at $1$ and $m/9$, $1\le m\le8$, together
with \eqref{eq:v113-endpoint-budget}, encloses $B_3$.
Physical evaluation of the projection is
$L^{-1/2}(c_0+2\sum_{n=1}^{64}c_n)$, and its propagated error is
at most $\sqrt{129/L}$ times the coefficient error. Strict interval
comparisons verify \eqref{eq:v113-endpoint-failure}. Thus the endpoint
mismatch is now proved for the exact repaired source, not just observed
for a numerical approximation.

The JSON output records the input hashes, inertia counts, full residual
and pointwise bounds, coefficient error, physical endpoint and relative
error intervals. The earlier frozen-candidate certificate remains a
separate dependency. These two source modes use the canonical regular
angular realization; no closed arithmetic operator assumption is needed
for this fixed-source computation. This source certificate alone supplies
no omitted \emph{Weil} Fourier-complement bound, growing-window sign,
endpoint-shell graph estimate or full-strip normalization. G2 remains open.

\subsection*{Scale check for the explicit omitted-tail constant}
The new operator-core and polynomial-endpoint arguments are analytic
proofs, not extrapolations from numerical data. A small additional
interval calculation evaluates the explicit constant in
\eqref{eq:v114-tail-constant} at $\lambda=3$. Using 256-bit Arb
quadrature of the regularized zero-mode archimedean integral gives
\[
 17.69849<C_3^{\rm tail}<17.69851,\qquad
 \Gamma_{3,50\,000\,000}>1.0797.
\]
The supplied script \texttt{g2\_deep\_next/check\_tail\_constant.py}
and its output record the complete intervals. No matrix of that size
was assembled: only the scalar constant was enclosed. This bound
illustrates the cost of the unsigned arithmetic estimate. It does not
certify the intervening low modes, their Schur correction, or positivity
of the whole window. In particular the earlier $N=64$ certificate is
not extended to the continuum by this calculation.

\subsection*{Sharper tail certificate and ordinary residual scope}
The script \texttt{g2\_signed\_tail/check\_sharp\_tail\_constants.py}
evaluates the exact elementary constants of
Proposition~\ref{prop:v115-sharp-tail} with 256-bit Arb intervals.
Its output gives
\[
 \Gamma^{\rm all}_{3,256}>0.0148329,\qquad
 \Gamma^{\rm all}_{3,512}>0.7085077,\qquad
 \Gamma^{\rm even}_{3,256}>1.5867887.
\]
The interval arithmetic certifies these scalar evaluations of analytic
bounds. That v1.15 calculation did not certify a head matrix or its
inverse-action correction; the complete certificate below now does so
at the same physical window. The ordinary residual theorem is an analytic
asymptotic estimate with no claimed effective threshold at $\lambda=3$.
It concerns the same exact source and literal Fourier projection as
the endpoint theorem. An independent internal adverse review checked
the norm normalization, low-output duality, high-output matrix bound,
graph tail, all scale factors and the remaining lack of ground overlap.

\subsection*{Complete fixed-window certificate and its reproduction}
Proposition~\ref{prop:v116-window-positive} is certified by
\texttt{g2\_schur\_directional/certify\_infinite\_schur.py}, using
the coefficient assembly in \texttt{assembly.py} and the scalar check
\texttt{certify\_weighted\_prime\_tail.py}. The bundle contains both
exact dyadic inverse-action witnesses, all positive LDL pivot enclosures,
and 768-bit and 896-bit replay reports. Its reproduction guide gives
the complete commands and analytic dependencies. The frozen witnesses
are unchanged between the two precision runs.

All finite output rows through $J=4096$ are enclosed, and the moment
inequality bounds every row beyond $J$. Thus the result is not an
extrapolation from finite eigenvalues. The finite Fourier entries through
$N=64$ were also checked against the earlier independent interval
integration: all 8,321 even and odd entries overlap their reference
enclosures. This consistency check supplements, but does not replace,
the analytic coefficient identities and series-remainder proof.

The unshifted result certifies complete fixed-window positivity.
The shifted extension below also establishes ground-state ordering.
Neither certificate alone supplies an endpoint comparison or a bound
for unbounded windows. The separate boundary theorem now determines
the continuum eigenfunction trace to be zero. Internal adversarial review checked the inverse order, dyadic
freezing residual, parity factors, remote Gram and form-domain passage;
this is not external peer review or proof-assistant verification.
Exploratory failures of weaker sufficient bounds are preserved in the
research log and are not presented as negative Weil-form results.

\subsection*{Shifted inertia and complete ground-overlap reproduction}
Proposition~\ref{prop:v117-ground-order} uses
\texttt{g2\_ground\_order/certify\_shifted\_inertia.py} with
$\texttt{shift}=10^{-34}$ in the even sector and $10^{-36}$ in the
odd sector. The bundle retains the signed pivot reports at 768 and
896 bits. It reuses the unchanged exact witnesses from
\texttt{g2\_schur\_directional}; their identity is checked by hash.
The script shifts both the actual logarithmic diagonal and all inverse
comparison weights, recomputes the middle residual and retains the
full remote moment bound. A negative lower-matrix pivot is not itself
a negative direction of the complete operator: the proved trial quotient
supplies that direction for the shifted even operator.

The separate \texttt{certify\_trial\_reuse.py} checks the exact rational
trial's full-form Rayleigh quotient, constructs the finite-supported
correction and imports the proved exact-source error from
Proposition~\ref{prop:v113-source-certificate}. Its corrected Rayleigh
enclosure is contained in
\[
 (3.64339965699233062,3.64339965699233064)\cdot10^{-38}.
\]
The separate scalar script also verifies the conservative constants
displayed in the theorem.
The complete ground-overlap conclusion uses the independently proved
even-sector spectral threshold. It does not reuse the finite matrix's
gap as if it were a gap for the complete operator.
All reproducibility commands and unsuccessful exploratory tests are
recorded separately in the guide and research log.

\subsection*{Analytic endpoint and complement checks}
The proof of Theorem~\ref{thm:v118-eigen-endpoint-zero} is analytic.
It does not use a numerical endpoint extrapolation. Its three gates are
the exact physical kernel identity with scalar $-\log(2\pi)$, a common
closed-form core, and the simultaneous $L^2/L^\infty$ resolvent argument
that makes the boundary theorem applicable. The archived independent
reviews check each gate, the retained pole signs and the distinction
between a restricted logarithmic Laplacian and a spectral logarithm.
The script \texttt{check\_log\_kernel\_normalization.py} independently
compares the original and decomposed constant-mode archimedean diagonal
at $\lambda=1.2,2,3,10$, with 90-decimal working precision. Agreement is
a diagnostic only; the scalar integral in the proof establishes the
normalization exactly.

The Fej\'er bound retains the physical $L^{-1/2}$ Fourier factor and
uses a positive kernel of total mass one. It does not replace the
paper's literal Fourier cutoff in any earlier formula.
Proposition~\ref{prop:v118-graph-trace} explicitly shows why graph
convergence does not transfer the physical trace.
The complement estimate in Proposition~\ref{prop:v118-gap-free}
is an exact Hilbert-space inequality with a sharp two-dimensional
example. Its use of the rank-one source invokes the proved ordinary
operator residual; its arithmetic complementary-sign hypothesis remains
open. These checks close the continuum endpoint question and improve
the formulation of the growing-window target, without closing G2.

\subsection*{Prime norm recurrence and growing-window scale}
Propositions~\ref{prop:v119-prime-essential}--\ref{prop:v119-prime-scale}
are analytic. The archived proof and adverse review distinguish a
finite-codimension norm obstruction from the still-open signed estimate.
The script \texttt{check\_prime\_norm.py} checks the positive trial
$\phi=e^{-x/2}+e^{-(L-x)/2}$ using
\[
 \|\phi\|_2^2=2(1-\lambda^{-2})+2L/\lambda,
 \qquad
 C_\phi(s)=2e^{-s/2}(1-e^{-(L-s)})
             +2(L-s)\lambda^{-1}\cosh(s/2).
\]
Thus its prime Rayleigh quotient is
$2\sum_m w_m C_\phi(\log m)/\|\phi\|_2^2$.
Independent piecewise physical quadrature at $\lambda=2,3$ agrees
with the autocorrelation formula to below $10^{-50}$ at 60-decimal
precision. At $\lambda=300$ the trial quotient divided by $\lambda$
is approximately $1.026889$; this is illustrative, not a norm certificate
or an asymptotic proof.

The same script tests phase recurrence at $\lambda=2$. Modulation by
Fourier index $381074$ has normalized $P_{64}$ head norm about
$9.52\cdot10^{-13}$, while its prime Rayleigh quotient differs from
the unmodulated value by about $4.7\cdot10^{-14}$.
These floating-point checks illustrate the analytic mechanism only.
They give no negative complete-Weil direction, because the
archimedean logarithmic energy grows with the modulation frequency.

\subsection*{A signed energy certificate at $\lambda=4$}
The directory \texttt{g2\_weighted\_signed} in the cumulative bundle
contains the proof, internal adverse review, exact dyadic witnesses
and interval reports for Proposition~\ref{prop:v120-lambda4-tail}.
\texttt{assembly\_general.py} encloses the actual digamma/trigamma
series with explicit exponential remainders. The separate pole diagonal
is retained. \texttt{certify\_weighted\_tail.py} evaluates the
verified-solve lower matrix, including every row up to $4096$ and the
infinite moment remainder, and verifies every Gershgorin row after an
exact dyadic congruence. The saved solve and congruence witnesses are
reused at higher precision; their origin is not a proof assumption.

\texttt{certify\_weighted\_counterexample.py} verifies the finite
dyadic positive direction disproving the two-sided norm contraction.
\texttt{weighted\_pilot.py} is only a floating-point exploration:
its small negative values at larger windows are unresolved numerical
roundoff, not negative Weil directions. The analytic arguments for
weighted compactness, non-Schatten behavior, fixed-window norm decay
and ordinary-error scaling do not rely on these computations.
The primary-source comparison with \cite{Zhu2026Windows} is contextual;
no external certificate or conjectured decay law is imported.

\subsection*{Structured low-mode inverse bounds and their scope}
The low-mode investigation retains the exact $\lambda=4$ matrix,
$E_{16}$ and the previously certified complete tail. The scripts
\texttt{certify\_low\_schur.py} and
\texttt{structured\_low\_schur.py} evaluate the coarse and structured
inverse bounds, respectively. All original arithmetic diagonals,
normalized parity factors and infinite residual moment bounds remain.
The saved tail witnesses are used as exact dyadics; their hashes and
previously certified congruence margins are checked before reuse.
The outer residual belongs to $\mathsf W_4$, whereas the reused inner
factors belong to $\mathsf W_4-10^{-8}D$. Inverse order justifies
that distinction.

These experiments have not certified the low Schur sign. Failed
lower enclosures and their exact diagnostic witnesses are kept in the
research log, not asserted as negative Weil directions. The script \texttt{certify\_far\_inverse.py} verifies the
strict far-inverse constants at 320 and 384 bits; the proof includes
the upper archimedean diagonal bound and the exact pole scale.
The analytic
results of Propositions~\ref{prop:v120-structured-inverse} and
\ref{prop:v120-structured-uniform} are proved directly above and
do not rely on the success of a numerical positivity gate.

\subsection*{Simultaneous even certificate and window scaling}
\label{app:v125-simultaneous}
The cumulative bundle adds \texttt{g2\_simultaneous}. Its main
\texttt{simultaneous\_trial.py} verifier reconstructs the frozen
$17$-column even trial and evaluates all complete residual Grams;
\texttt{complete\_margin.py} certifies the exact rational generalized
margin. Both use the same witness at $768$ and $896$ bits.
The optional \texttt{joint\_remote\_gate.py} retains one common
remote Gram and its cross block. The complete inner tail certificate
and its original exact witnesses are retained unchanged.

The odd trials and the unscaled even trial are retained as
inconclusive experiments in the log. They are not sign certificates.
All trial heads are checked for full rank; midpoint proposals never
replace outward evaluation in a proof gate. The recorded $LDL^*$
pivots are coordinate pivots, not ordinary spectral gaps.

\texttt{window\_scaling\_probe.py} freshly evaluates the analytic
far-cut thresholds at $320$ and $384$ bits. Its separate finite
Galerkin study uses windows $3,4,5,6,8$, heads through $\lambda^2$
and fixed retained/test cuts $256/384$. It records the finite
Schur correction, dimensionless margins and conditioning. Modes
beyond $384$ are absent from that diagnostic; it is not an execution
of the complete-tail pipeline at the larger windows. Every finite
tail used in its elimination must pass its own positive interval
gate. Complete window-specific inner certificates remain necessary
before interpreting such measurements as bounds for G2.

\subsection*{Complete verification of a new finite trial}
\label{app:v124-posterior}
The cumulative bundle adds \texttt{g2\_posterior\_trial}, containing
\nolinkurl{posterior_trial.py}, its imported helper scripts, the
saved remote-Gram reports, the frozen dyadic corrections and the
independent internal proof/code review. The main certificate uses
\nolinkurl{posterior_witness_M4096_s16.json.gz}; its $4080$ dyadic
pairs define the correction on indices $17$ through $4096$.
The original head vector and support-$256$ solve are unchanged.

The proposal uses $16$ steps of a finite preconditioned conjugate
gradient calculation at $256$ bits. This is only a way of choosing
coefficients. The separate verification routine reconstructs the
exact frozen vector, checks its dimensions and witness hashes,
and evaluates the unshifted complete residual at $768$ or $896$
bits. Its final gate includes the full moment remainder beyond
$J=65536$ and the old complete inner metric. It also checks nine
rows by independent dense assembly and checks the expanded energy
identity. No convergence theorem for the proposing iteration is
needed. A separate four-step trial also passes with lower bound
$2.9028\cdot10^{-20}$; it is not asserted to be the old degree-$4$
inverse polynomial. All trial histories are retained in the log.

\subsection*{Complete mixed-correlation certificate and head refinement}
\label{app:v123-mixed}
The cumulative bundle retains the exact outer and inner witnesses
from the preceding revision. The new files in
\texttt{g2\_mixed\_complete} are:
\begin{itemize}[nosep]
\item \nolinkurl{mixed_refinement.py}, for retained mixed vectors,
complete first-polynomial upper enclosures, and explicitly separate
higher-degree finite-prefix diagnostics;
\item \nolinkurl{mixed_directional_check.py}, for the complete
directional lower certificate of
Proposition~\ref{prop:v123-first-majorant-fails};
\item \nolinkurl{mixed_directional_witness_J65536.json}, specifying
the exact $496$-component probe by integer numerators over $2^{160}$;
\item the $768$- and $896$-bit interval reports and the independent
internal adverse review of the mixed, monotone-head and probe bounds.
\end{itemize}
The inner witness SHA-256 is checked against the saved positive-tail
report. The prefix inverse and its head pairings use the actual
unshifted residual and consistently shifted inner factors. The
certified probe norm is retained rather than rounded to one.
Both remote input and remote first-polynomial output enter the
error estimate. Replaying the analytic and interval inequality is
the sign certificate; a small floating residual is not.

The higher-polynomial experiments are retained in the research log.
Their positive finite-prefix margins do not certify the corresponding
complete corrections and are not used as propositions in this paper.
The strengthened head bound is proved analytically; its simultaneous
application to the full $\lambda=4$ Schur matrix remains open.

\subsection*{Parity inverse refinement and its directional gate}
\label{app:v122-inverse}
The cumulative bundle's \texttt{g2\_inverse\_refinement} directory
contains \texttt{refine\_inverse.py}, both parity-constant interval
reports, both final directional-gate reports and the internal adverse
review. The constants replay at $320$ and $384$ bits; the directional
gate replays at $768$ and $896$ bits. The mathematical arguments of
Propositions~\ref{prop:v122-prefix-refinement} and
\ref{prop:v122-refined-structured} are independent of a numerical
positivity gate.

The exact direction uses the $17$ numerators and denominator $2^{160}$
from the archived \texttt{g2\_low\_schur} directory:
\begin{itemize}[nosep]
\item \nolinkurl{structured_ceiling_counterwitness.json} defines $v$;
\item \nolinkurl{low_witness_even_M256_b768.json.gz} defines the dyadic $X$.
\end{itemize}
These archived files define the inputs, not a new rounded eigenvector.
The actual matrix assembly retains its analytic archimedean-series
remainder and correct logarithmic diagonal.

For consecutive positive indices, four Arb polynomial products
evaluate the Toeplitz divided difference and the Hankel matrix action.
The first uses the convolution kernel $1/j$ for $j\ne0$, with zero
central coefficient, applied to $z_n$ and $b_nz_n$. The second
uses reversed input coefficients and the kernel $1/(n+m)$.
Thus the output is exactly
\[
 ((U_Jz)_n)=(d_n-ca_n)z_n+
 \sum_{m\ne n}\frac{b_n-b_m}{n-m}z_m+
 \sum_m\frac{b_n+b_m}{n+m}z_m
\]
in the even sector, including the opposite-sign contribution on its
parity diagonal. The polynomial products retain outward ball errors;
a dense $25$-row evaluation independently checks their indexing.
Only the computational method changes. The tail beyond $J=65536$
is still included by the complete moment bound and
\eqref{eq:v122-remote-H}. Earlier inconclusive cutoffs are retained
in the research log. No full mixed-term, matrix-head or G2 sign
certificate is asserted by this directional test.

\subsection*{Finite trial positivity and exact-source cancellation audit}
\label{app:v121-audit}
The cumulative bundle's \texttt{g2\_schur\_cancellation} directory
contains the exact witnesses, interval outputs and internal adverse
review for Propositions~\ref{prop:v121-alpha-zero} and
\ref{prop:v121-complement-cancellation}. These use the same actual
matrix assembly as the earlier certificates, with all physical scale
and parity factors retained. No externally reported F-68 or F-69
certificate is imported as a premise.

\texttt{certify\_finite\_K.py} reconstructs $K_X$ from its definition
for $\lambda=4$, split $16$ and solve support through $256$.
The files \texttt{finite\_K\_X\_even.json.gz} and
\texttt{finite\_K\_X\_odd.json.gz} contain the frozen exact dyadics.
Every interval $LDL^*$ pivot is positive in both the $768$- and
$896$-bit replays. Conservative lower bounds for the smallest pivot
are listed below; they are \emph{not} eigenvalue lower bounds.
\begin{center}
\begin{tabular}{@{}lll@{}}
\toprule Sector & Trial & Smallest pivot exceeds\\\midrule
Even & $X=0$ & $1.4513\cdot10^{-13}$\\
Odd & $X=0$ & $8.4162\cdot10^{-13}$\\
Even & Frozen solve & $1.2207\cdot10^{-58}$\\
Odd & Frozen solve & $1.3872\cdot10^{-56}$\\\bottomrule
\end{tabular}
\end{center}
This certifies $K_X$, not $S_\infty$ or the omitted inverse-residual
subtraction. The local $\alpha=0$ conclusion therefore does not
contradict previously failed lower Schur enclosures.

\nolinkurl{check_cancellation_subspace.py} uses the exact columns in
\nolinkurl{cancellation_two_mode_witness.json}, each with denominator
$2^{300}$. The $2\times2$ Gram positivity, nonzero constructed vector,
actual quadratic values and all matrix row-sum bounds are verified by
outward ball arithmetic. The source candidate's hash is checked
against the archived \texttt{pswf\_source\_certificate.json}; the
existing exact-source certificate is an explicit dependency.
Projection to the exact source complement is justified analytically,
as in the proof, rather than by a rounded orthogonality residual.

For clarity, the approximate intermediate energies are
\begin{center}
\begin{tabular}{@{}lrr@{}}
\toprule Contribution & Source candidate & Orthogonal trial $w$\\\midrule
Pure archimedean & $-1.4324413624$ & $-0.4195793594$\\
Pole & $+1.4877438175$ & $+0.8369031184$\\
Signed prime & $-0.0553024551$ & $-0.4173237590$\\
Archimedean plus pole & $+0.0553024551$ & $+0.4173237590$\\
Complete Weil & $5.3123817\cdot10^{-38}$ & $2.5844134\cdot10^{-31}$\\
\bottomrule
\end{tabular}
\end{center}
The displayed rounded components cannot resolve the tiny totals by
subtraction; the separate full-precision interval evaluations establish
the strict inequalities. Both precision replays use identical witness
hashes. The prior source identification and the finite-support equality
then supply the exact-source statement in the main text.

These are internal computer-assisted proofs with reproducible inputs,
not external peer verification. The witness proposals are numerical;
the stated sign gates and perturbation margins are interval-certified.
No growing-window convergence or new G2 sign conclusion follows.


\section{Integration ledger for the complete previous manuscript}\label{app:ledger}
The table accounts for all 72 numbered theorems, lemmas, propositions and
corollaries in the supplied 89-page r11 source. Numbers in the first column
are historical identifiers; theorem numbering in this version is new.
A retained local implication does not certify an unproved upstream input.
The supplied audits were treated as arguments to check, not as authority.

All original subject areas are represented: Sections 1--19 are revised in
place, the original Section 20 is rebuilt with a proved compact-source passage
and restored older matrix/weighted estimates, and the original Section 21
is replaced by the consistent sampler and the rebuilt downstream analysis.
The original Sections 22--23 are replaced by the criterion and goals in
Section~\ref{sec:goals}. Invalid proofs are withdrawn instead of being
preserved as established results in an appendix.

\small
\begin{longtable}{@{}p{0.08\textwidth}p{0.85\textwidth}@{}}
\toprule Old & Disposition in v1.25 \\
\midrule\endfirsthead
\toprule Old & Disposition in v1.25 (continued) \\
\midrule\endhead
\bottomrule\endfoot
4.1 & \textbf{Pass.} Three-piece arithmetic source decomposition. Support covariance and a partition of unity identify the closure of the three-source sum; no closed-range theorem is asserted. \\[4pt]
4.2 & \textbf{Pass.} Arithmetic projection disappears in two adjoint channels. The three orthogonality channels remove the arithmetic projections exactly. \\[4pt]
6.1 & \textbf{Pass.} Finite-cutoff span positivity would imply RH. Full-span positivity is sufficient but stronger than the diagonal criterion; the research target is corrected to each evaluator. \\[4pt]
12.1 & \textbf{Repaired.} Exact multiplicity-$m$ leakage factorization. Off-line hypothesis restored; compact division and finite-jet density prove the missing kernel inclusion. Quantitative persistence still uses the evaluator assumption. \\[4pt]
13.1 & \textbf{Pass.} No generic domination between cross-scale and fixed-space defects. The finite-dimensional examples establish that generic projection geometry gives no defect domination. \\[4pt]
14.1 & \textbf{Repaired.} Exact generator range defect. Compact division and finite-jet density supply a proof. The range is the algebraic source-core power before closure. \\[4pt]
15.1 & \textbf{Pass.} Exact zero-incidence criterion. Pole removability proves zero incidence under the stated ambient Sonine division framework. The skew-extension sign is corrected nearby. \\[4pt]
17.1 & \textbf{Pass.} Exact diagonal contractions force orthogonality. Dense unimodular eigenvectors of a contraction force a unitary and phase orthogonality. \\[4pt]
17.2 & \textbf{Conditional.} Finite-mode asymptotic calibration. Retained under the explicit quantitative evaluator assumption. \\[4pt]
18.1 & \textbf{Conditional.} Exact off-line spectrum of the fixed arithmetic operator. Retained under an explicit closed Hilbert realization, domain, graph-core and spectral-completeness assumption; the construction is not supplied. \\[4pt]
19.1 & \textbf{Retained; finite-core alternative added.} The Riesz-window implication retains its realization hypothesis. Theorem~\ref{thm:v15-finite-metric} bypasses it on algebraic root spaces. The full-space projection now preserves coercivity and Weil transfer while annihilating source and shell; a small signed defect is still missing, and a fixed scalar Weil shift fails. \\[4pt]
20.1 & \textbf{Pass.} Endpoint-normalized radial bound. The fixed-order radial estimate is retained with its cited asymptotic input. \\[4pt]
20.2 & \textbf{Pass.} Away-boundary and exterior-energy bounds. Away-boundary and exterior-energy estimates are retained for each fixed radial order. \\[4pt]
20.3 & \textbf{Pass.} Uniform far field with controlled $c$-dependence. The radial ODE and integrable far-field perturbation argument are retained. \\[4pt]
20.4 & \textbf{Pass.} Endpoint-normalized radial lattice bound. The lattice summation argument and resonance control are retained. \\[4pt]
20.5 & \textbf{Pass for the explicit repaired source.} Fixed-order co-Poisson cross-tail. The normalized source, repaired tail, inversion-endpoint comparison and correlation estimates are proved. Theorem~\ref{thm:v14-radical} now supplies the global radical/semilocal-domain passage; no source assumption remains. \\[4pt]
20.6 & \textbf{Pass.} Quantitative endpoint-jump bound. The normalized source-tail bounds and the PNT remainder prove the rate; both translated-tail terms are retained. \\[4pt]
20.7 & \textbf{Pass.} Inversion symmetrization preserves the G1 envelope. Inversion preserves the chosen test class and transfers the proved envelope. \\[4pt]
20.8 & \textbf{Pass for Sobolev pairings.} Theorem~\ref{thm:v14-radical} closes the remaining domain input. This does not supply a growing-block operator-norm estimate. \\[4pt]
20.9 & \textbf{Pass.} Arithmetic BV structure and absence of a periodic boundary jump. The finite arithmetic sum has BV regularity and matching periodic endpoint representatives. \\[4pt]
20.10 & \textbf{Pass.} Fixed-cutoff boundary-preserving Galerkin approximation. Fixed-cutoff BV Fourier/form approximation is retained, with the logarithmic form weight and endpoint convergence distinguished. \\[4pt]
20.11 & \textbf{Pass.} Quantitative and cofinal finite-section transfer of G1. The finite-band approximation transfers the proved continuum estimate; its eventual threshold need not be economical. \\[4pt]
20.12 & \textbf{Pass.} Zero-independent high--low decay of the Weil matrix. The off-diagonal matrix estimate is retained and supplemented by the separately restored diagonal bound. \\[4pt]
20.13 & \textbf{Pass.} Explicit BV rate for the fixed-band Galerkin defect. BV coefficients and the off-diagonal estimate give the summable high-low tail. \\[4pt]
20.14 & \textbf{Pass.} Quantitative endpoint recovery by de la Vall\'ee--Poussin filtering. The near/far endpoint convolution estimate is retained. \\[4pt]
20.15 & \textbf{Pass.} Filtered BV transfer with explicit boundary control. The filtered transfer uses the now-proved continuum G1; zero local Lipschitz constant is handled without division. \\[4pt]
20.16 & \textbf{Pass.} High-frequency boundary repair has unavoidable square-root derivative cost. Weighted Cauchy--Schwarz gives the derivative cost; a general trace obstruction is added for the new sampler. \\[4pt]
20.17 & \textbf{Pass in the fixed-band sense.} Constructive zero-independent finite G1 closure. The exact-boundary construction transfers the now-proved continuum estimate with the stated deterministic cutoffs. No economical cutoff or growing-block uniformity is inferred. \\[4pt]
20.18 & \textbf{Repaired.} Form-Schur lower bound. The Schur error parameters are nonnegative. Rayleigh calibration and operator-domain overlap are separated from weak G1. \\[4pt]
21.1 & \textbf{Retained.} One-sided off-line Weil isotropy. Global one-sided isotropy is proved in the corrected convention within the centered transfer theorem. \\[4pt]
21.2 & \textbf{Replaced.} Off-line semilocal form transfer. The full-scale raw transfer is withdrawn. Centered strip transfer retains jumps; the explicit endpoint repair preserves its form limit. \\[4pt]
21.3 & \textbf{Rebuilt.} G2b under the off-line contradiction hypothesis. Restricted coercivity follows for the corrected sampler. A negative compact-test calibration is optional and needs its own form approximation. \\[4pt]
21.4 & \textbf{Rebuilt.} Restricted scalar-shift coercivity is automatic after form transfer. Positive scalar-shift coercivity follows from corrected form and ordinary Gram transfer; it is not independent positivity. \\[4pt]
21.5 & \textbf{Repaired.} Smooth-test continuity of the semilocal Weil pairing. The logarithmic diagonal is restored. Absolute matrix-series estimates and the weighted convolution proof replace the false uniform-diagonal bound. \\[4pt]
21.6 & \textbf{Rebuilt.} The proved G1 estimate applies to the full Hardy sampler. The restored weighted theorem gives weak G1 for centered Hardy samples using the now-proved source passage; the even endpoint correction is invisible to the zeroth pairing by parity. \\[4pt]
21.7 & \textbf{Scoped.} Finite G1 jet hierarchy on the full Hardy sampler. Sufficient smoothing controls each fixed finite derivative family. No growing-order uniformity or automatic parity preservation for all corrected jets is asserted. \\[4pt]
21.8 & \textbf{Retained.} Fixed residual-jet Gramians cannot supply G2 coercivity. A fixed finite Gramian of vanishing normalized residual pairings also vanishes; the fixed-family G1 input is now proved. \\[4pt]
21.9 & \textbf{Rebuilt.} Off-line coercivity forces a boundary-weighted sampler residual. The exact Lyapunov residual now includes the corrected graph term. Its coercive lower bound is retained. \\[4pt]
21.10 & \textbf{Rebuilt.} Exact sampler--prolate compatibility decomposition. The finite candidate decomposition uses the total source e z plus the graph defect. \\[4pt]
21.11 & \textbf{Rebuilt.} Commuting-correction cancellation and the scalar-shift no-go. Commuting and scalar cancellation remain exact with the total source. The old raw-transfer asymptotic is removed. \\[4pt]
21.12 & \textbf{Rebuilt.} Exact noncommuting G2 excess. The noncommuting excess includes the graph defect and still requires an independent upper estimate. \\[4pt]
21.13 & \textbf{Retained.} $W^2$ benchmark: positivity moves the gap to norm transfer. The positive square and rank-four commutator are exact; form smallness does not imply norm smallness. \\[4pt]
21.14 & \textbf{Repaired.} Reflected-partner $W^2$ coercivity at arbitrary multiplicity. The reflected top-root lower bound is proved for the corrected sampler; conjugation is on the second factor. This is not full-block coercivity. \\[4pt]
21.15 & \textbf{Rebuilt.} A boundary-annihilating G2 test isolates the exact sampler--prolate mismatch. Boundary annihilation is an exact finite construction when its denominator is nonzero. The proposed bounded endpoint ratio is now ruled out by Corollary~\ref{cor:v11-ratio}; bounded tests from a fixed input space lose the reflected signal by Theorem~\ref{thm:v11-fixed-tests}. \\[4pt]
21.16 & \textbf{Withdrawn.} Pole-aware radical transfer for the fixed Hardy sampler. The pole-aware transfer proof is insufficient. A replacement must bound paired errors over all zeros, including slow tails and actual cut boundaries. \\[4pt]
21.17 & \textbf{Conditional only.} The surviving simple-zero obstruction is a boundary commutator channel. The finite commutator decomposition survives; no nonzero beta limit is claimed. \\[4pt]
21.18 & \textbf{Rebuilt.} Multiplicity-robust reflected-partner block obstruction. The corrected top-Jordan relation contains an extra graph term. The cross-form lower bound survives. The original endpoint-annihilating test cannot retain it with bounded norm; a different test construction and graph estimates are required. \\[4pt]
21.20 & \textbf{Scoped.} Rank--two boundary traces exhaust the reflected G2 signal. Rank-two trace algebra survives, but it does not exhaust the signal without controlling the new graph pairing. \\[4pt]
21.21 & \textbf{Withdrawn.} Critical $L^{-1/2}$ scale of the surviving boundary trace. The critical nonzero beta limit is removed from the established chain. One former premise, a bounded ratio for the corrected resolvent, is disproved; the other graph and form obligations remain explicit. \\[4pt]
21.22 & \textbf{Retained.} Exact arithmetic realization of the surviving CCM boundary channel. The exact finite arithmetic identity is retained with physical scaling, integer index domain, and the first-slot convention. \\[4pt]
21.23 & \textbf{Conditional only.} G2 is an explicit finite prime/pole/archimedean forcing statement. The finite observable is explicit. Its asserted nonzero asymptotic is not established for the corrected test family. \\[4pt]
21.24 & \textbf{Scoped.} The present absolute PNT remainder is quantitatively insufficient for the G2 boundary channel. The absolute PNT estimate is retained as an upper-bound limitation, not a lower bound or impossibility theorem. \\[4pt]
21.26 & \textbf{Retained.} Hardy phase shifts preserve the global G2 signal while damping the Hardy endpoint. Global translation, Cauchy damping, and reflected cross-form invariance are retained in the fixed Fourier convention. \\[4pt]
21.27 & \textbf{Withdrawn.} A translated Hardy kernel recovers an absolute PNT gain at the continuum-kernel level. The periodic upper bound omitted a reflected image. The missing image and an explicit exponent comparison are now displayed. \\[4pt]
21.28 & \textbf{Scoped.} Reflected-pole alias obstruction to the moving Hardy phase. Threshold algebra is retained. A no-go conclusion requires relative alias remainder estimates that are not supplied. \\[4pt]
21.30 & \textbf{Retained.} Exact continuum normal form of the boundary-annihilated G2 test. The continuum meromorphic normal form and its residue are retained; it is not asserted to be the corrected finite endpoint-ratio limit. \\[4pt]
21.31 & \textbf{Repaired.} The reflected pole is forced by any cross-retaining holomorphic correction. Residue persistence is retained and second-slot scaling uses the conjugated scalar. \\[4pt]
21.32 & \textbf{Repaired.} Bounded Hardy filtering cannot suppress the reflected pole while retaining G2. The modulus inequality is proved; global admissibility and a nonzero endpoint factor are required for form interpretations. \\[4pt]
21.33 & \textbf{Scoped.} Absolute-PNT rescue by pole damping lies outside the audited G1 regime. Exponential norm cost describes this envelope-based estimate, not a necessary lower bound on actual transfer error. \\[4pt]
21.35 & \textbf{Retained.} Projective boundary-quotient invariance of the G1/G2 mismatch. Projective invariance is exact for a fixed endpoint-zero test. Its asymptotic consequence assumes both normalized G1 pairings vanish. \\[4pt]
21.36 & \textbf{Scoped.} The prolate core may be replaced by a pure boundary shell in the lower obstruction. A shell can replace the candidate in the exact paired identity; no shared nonzero G2 limit is imported. \\[4pt]
21.37 & \textbf{Scoped.} Moving Hardy evaluation aligns the graph functional. The moving Cauchy scalar converges to the graph functional on a fixed root space; moving-parameter uniform transfer remains open. \\[4pt]
21.38 & \textbf{Retained.} Candidate-adapted graph sampler: exact vector matching. Coefficientwise exact graph matching and scalar endpoint error are proved. Injectivity, norm, form and G1 compatibility are additional obligations. \\[4pt]
21.39 & \textbf{Retained.} Exact-radical factorization removes the zeta quotient but exposes an endpoint phase. Global Mellin divisibility and the endpoint/bulk integration-by-parts identity remain, with the Fourier-sign conversion explicit. \\[4pt]
21.40 & \textbf{Scoped.} Finite boundary-jet expansion of the exact matcher. The fixed-cutoff high-frequency jet expansion remains. Subtraction through zero frequency requires a prescription. \\[4pt]
21.42 & \textbf{Retained.} Exact periodic-resolvent formula for the adapted endpoint. The periodic ODE and parameter derivatives give the exact normalized Laplace-jet formula. \\[4pt]
21.43 & \textbf{Retained.} Boundary-perfect high-frequency matching has vanishing bulk mass. The shell exact matcher has asymptotically correct endpoints but vanishing ordinary mass, so it is not coercive. \\[4pt]
21.45 & \textbf{Rebuilt.} Parity projection gives a bounded noncommuting G2 benchmark. Parity coercivity uses only ordinary norm transfer and survives the corrected sampler. Its Lyapunov identity includes the graph defect. \\[4pt]
21.48 & \textbf{Repaired.} The exact-boundary shell is compatibility-efficient in the coercive pairing. Correct diagonal and row-sum bounds give shell weak estimates. They do not remove the commutator term in the graph pairing. \\[4pt]
21.49 & \textbf{Scoped.} One shell cutoff neutralizes both the G1 repair and its G2 pairing cost. Common-cutoff shell smallness is retained as an eventual-threshold argument; it controls auxiliary shell terms only. \\[4pt]
21.50 & \textbf{Replaced.} A joint exact-boundary G1/pairing diagonal. An explicit normalized common-cutoff selection lemma replaces ambiguous cofinality claims. It proves no missing core graph estimate. \\[4pt]
21.51 & \textbf{Scoped.} Scalar-shift positivity cannot remove the final defect. The reverse-triangle obstruction is retained only with its explicit full-residual and derivative-norm premises; weak G1 does not supply them. \\[4pt]
\end{longtable}
\normalsize
\paragraph{Earlier-version material.}
The separate correct diagonal bound, weighted-sampler convolution estimate,
arithmetic jump train, jump/continuous split, and the limited near-Slepian
jump cutoff were recovered from the older 65-page source. The 32-page v10
was checked for useful earlier structure. Its unsupported exact jet
Gram asymptotic, incomplete far-field comparison, and unseparated endpoint
jumps were not restored. The more careful one-sided Gram bound, radial ODE
argument, and endpoint-jump separation of the later source are retained.

\paragraph{Audit reconciliation.}
The supplied G1/G2 audit overstated the finite lower mechanism. The supplied
Claude audit checked useful local algebra but did not settle the analytic
transfers affected by the cut and diagonal errors. This revision adopts
locally sound calculations, corrects their conventions and scope, and
replaces the central transfer rather than treating either audit as a
certification of the whole paper.

\paragraph{Authorship and assistance.}
AI tools assisted with drafting, exploratory calculations, and adversarial
review. They are not mathematical authorities or sources of certification.
Responsibility for the manuscript's claims and any eventual publication
remains with its author. The reproducible finite checks are supplied separately. Exploratory
computations are distinguished from the rigorous finite interval
matrix and exact-source certificates in Appendix~\ref{app:numerics}.
The new endpoint and operator arguments received internal adversarial
review, including growing derivative orders, arithmetic threshold
resonance, and identification of the physical form closure. This is
not external referee verification.

\begin{thebibliography}{99}
\small
\setlength{\itemsep}{2pt}
\bibitem{Burnol2001}
J.-F. Burnol,
\emph{Sur certains espaces de Hilbert de fonctions entieres, lies a la transformation de Fourier et aux fonctions L de Dirichlet et de Riemann},
arXiv:math/0105120v1 (2001).

\bibitem{Burnol2002Sonine}
J.-F. Burnol,
\emph{Sur les espaces de Sonine associes par de Branges a la transformation de Fourier},
arXiv:math/0208121v1 (2002).

\bibitem{BurnolCoPoisson2004}
J.-F. Burnol,
\emph{Entrelacement de co-Poisson},
arXiv:math/0407443 (2004).

\bibitem{Burnol2004}
J.-F. Burnol,
\emph{Two complete and minimal systems associated with the zeros of the Riemann zeta function},
arXiv:math/0203120v7; J. Theorie des Nombres de Bordeaux 16 (2004), 65--94.

\bibitem{Burnol2011}
J.-F. Burnol,
\emph{On the system of the functions $\zeta(s)/(s-\rho)^k$},
arXiv:1106.4751 (2011).

\bibitem{Livsic2012}
A. Aleman, R. T. W. Martin, and W. T. Ross,
\emph{On a theorem of Livsic},
arXiv:1209.4497 (2012).

\bibitem{ConnesConsaniMoscovici2025}
A. Connes, C. Consani, and H. Moscovici,
\emph{Zeta Spectral Triples},
arXiv:2511.22755 (2025).

\bibitem{ConnesVanSuijlekom2025}
A. Connes and W. D. van Suijlekom,
\emph{Quadratic Forms, Real Zeros and Echoes of the Spectral Action},
arXiv:2511.23257 (2025).

\bibitem{Dunster2017}
T. M. Dunster,
\emph{Asymptotics of Prolate Spheroidal Wave Functions},
arXiv:1601.00699v3 (2017).

\bibitem{Fuchs1964}
W. H. J. Fuchs,
\emph{On the eigenvalues of an integral equation arising in the theory of band-limited signals},
J. Math. Anal. Appl. 9 (1964), 317--330.

\bibitem{BonamiKaroui2014}
A. Bonami and A. Karoui,
\emph{Uniform bounds of prolate spheroidal wave functions and eigenvalues decay},
C. R. Math. Acad. Sci. Paris 352 (2014), 229--234.

\bibitem{Suzuki2026}
M. Suzuki,
\emph{Weil's quadratic form via the screw function},
arXiv:2606.09096v2 (2026).

\bibitem{ConnesConsani2023}
A. Connes and C. Consani,
\emph{Spectral triples and $\zeta$-cycles},
L'Enseignement Math. 69 (2023), 93--148.

\bibitem{SlepianTransition}
S. Karnik, J. Romberg, and M. A. Davenport,
\emph{Improved bounds for the eigenvalues of prolate spheroidal wave functions and discrete prolate spheroidal sequences},
Appl. Comput. Harmon. Anal. 55 (2021), 97--128.

\bibitem{Douglas}
R. G. Douglas,
\emph{On majorization, factorization, and range inclusion of operators on Hilbert space},
Proc. Amer. Math. Soc. 17 (1966), 413--415.

\bibitem{Hardy}
G. H. Hardy, \emph{Sur les z\'eros de la fonction $\zeta(s)$ de Riemann},
C. R. Acad. Sci. Paris 158 (1914), 1012--1014.

\bibitem{Trudgian2012}
T. S. Trudgian,
\emph{An improved upper bound for the argument of the Riemann
zeta-function on the critical line II},
\href{https://arxiv.org/abs/1208.5846v2}{arXiv:1208.5846v2} (2012),
Corollary 1 and equation (2.5).

\bibitem{PythonFlint2026}
The Python-FLINT and FLINT developers,
\emph{Python-FLINT 0.9.0 documentation}, real and complex ball arithmetic,
rigorous integration and matrix solves; used with FLINT 3.6.0.
\url{https://python-flint.readthedocs.io/en/stable/arb_mat.html};
\url{https://python-flint.readthedocs.io/en/stable/acb.html}.
Accessed September 20, 2026.

\bibitem{DLMFPSWF}
NIST Digital Library of Mathematical Functions, Chapter 30,
Spheroidal Wave Functions, Sections 30.2--30.3.
\url{https://dlmf.nist.gov/30.2};
\url{https://dlmf.nist.gov/30.3}.
Regular angular equation, spectral ordering and normalization convention.

\bibitem{NISTDigamma}
NIST Digital Library of Mathematical Functions, Chapter 5,
equations 5.9.15 and 5.15.1 (Binet's digamma formula and trigamma series),
\url{https://dlmf.nist.gov/5.9.E15},
\url{https://dlmf.nist.gov/5.15.E1}.

\bibitem{Teschl2009}
G. Teschl, \emph{Mathematical Methods in Quantum Mechanics},
Graduate Studies in Mathematics 99, American Mathematical Society, 2009,
Section~4.3.
\url{https://www.mat.univie.ac.at/~gerald/ftp/book-schroe/schroe.pdf}.
\bibitem{ChenWethLog}
H.~Chen and T.~Weth,
\emph{The Dirichlet Problem for the Logarithmic Laplacian},
arXiv:1710.03416v6 (2019),
\url{https://arxiv.org/abs/1710.03416v6}.

\bibitem{HLSLogBoundary}
V.~Hern\'andez-Santamar\'ia, L.~F.~L\'opez R\'ios, and A.~Salda\~na,
\emph{Optimal boundary regularity and a Hopf-type lemma for Dirichlet
problems involving the logarithmic Laplacian},
arXiv:2401.18033v2 (2024),
\url{https://arxiv.org/abs/2401.18033v2}.

\bibitem{NISTPrimeAsymptotics}
NIST Digital Library of Mathematical Functions, Section~27.12,
Asymptotic Formulas: Primes, prime number theorem and classical error
estimates. \url{https://dlmf.nist.gov/27.12}.

\bibitem{Zhu2026Windows}
X. Zhu, \emph{Weil positivity in compact windows: a finite reduction,
certified two-sided bounds, and a Landau--Widom decay law},
arXiv:2608.24827v2 (September 2, 2026).
\url{https://arxiv.org/abs/2608.24827v2}.

\end{thebibliography}

\end{document}
